% 来自2026-10-21统一复核稿；题号、图号和解答保留。
\clearpage
\begingroup
\fancyhead[R]{\small 模块四：稳恒电流与电阻网络}
\clubpenalty=10000\widowpenalty=10000\displaywidowpenalty=10000
\chapter{模块四：稳恒电流与电阻网络}\label{m4:ux6a21ux5757ux56dbux7a33ux6052ux7535ux6d41ux4e0eux7535ux963bux7f51ux7edc}

连续导电介质的场方程可积分为电压、电流和电阻等端口关系。有限与无穷电阻网络沿用相同的电荷守恒；矩阵只是基尔霍夫方程的紧凑写法，第一次阅读可先用节点和支路方程计算。

在网络方法之前，先讨论不同介质的载流机制、半导体器件和仪表内阻。第一轮可读到基尔霍夫方程；矩阵、谱分解和无穷网络留作选讲。

{\def\LTcaptype{none} % do not increment counter
\begin{longtable}[]{@{}
  >{\raggedright\arraybackslash}p{(\linewidth - 4\tabcolsep) * \real{0.3333}}
  >{\raggedright\arraybackslash}p{(\linewidth - 4\tabcolsep) * \real{0.3333}}
  >{\raggedright\arraybackslash}p{(\linewidth - 4\tabcolsep) * \real{0.3333}}@{}}
\toprule\noalign{}
\begin{minipage}[b]{\linewidth}\raggedright
层级
\end{minipage} & \begin{minipage}[b]{\linewidth}\raggedright
内容
\end{minipage} & \begin{minipage}[b]{\linewidth}\raggedright
阅读定位
\end{minipage} \\
\midrule\noalign{}
\endhead
\bottomrule\noalign{}
\endlastfoot
基础必讲 &
电流密度、德鲁德直流图像、液体与气体及真空中的电流、半导体器件、电学仪表、补偿法、稳恒电流边界、基尔霍夫定律、等效电源、星形---三角形变换、等势点与自相似网络
& 建立从局域输运到可测端口和宏观电路的常规复赛方法 \\
拔高选讲 &
非线性器件分段模型、网络拓扑、福斯特迹定理、多端响应矩阵、谱分解、变分原理与离散格林函数
& 统一处理器件状态切换、高对称、无穷及多端复杂网络 \\
\end{longtable}
}

\par\noindent\begin{minipage}{\linewidth}
\begin{center}
\includegraphics[width=\linewidth,height=.40\textheight,keepaspectratio]{\detokenize{assets/module-04/theory/tikz-v11-module4-method-map.pdf}}
\par\smallskip{\small 图示：模块四从连续电流场过渡到有限与无穷电阻网络}
\end{center}
\end{minipage}\par

\section{专题二十七：微观电流与德鲁德模型（基础必讲与拔高选讲）}\label{m4:ux4e13ux9898ux4e8cux5341ux4e03ux5faeux89c2ux7535ux6d41ux4e0eux5fb7ux9c81ux5fb7ux6a21ux578bux57faux7840ux5fc5ux8bb2ux4e0eux62d4ux9ad8ux9009ux8bb2}

\begin{quote}
阅读准备：一阶弛豫方程只含一个记忆时间 方程 \(\dot v+v/\tau=a\)
的通解等于稳态解加指数衰减的暂态解。给定一个初值即可确定唯一过程；\(\tau\)
表示初始偏差缩小到原来的 \(e^{-1}\) 所需时间。
\end{quote}

\par\addvspace{.4\baselineskip}
\noindent\begin{minipage}{\linewidth}

\textbf{例 M4-001：一阶弛豫的建立过程}

一个一维弛豫变量 \(v(t)\) 在 \(t\ge0\) 满足 \(dv/dt+v/\tau=a\)，其中
\(a>0\) 为常量驱动，\(\tau>0\) 为弛豫时间，初值为
\(v(0)=0\)。求完整建立过程，并分别在 \(t\ll\tau\)、\(t\gg\tau\)
的范围给出主阶行为，说明稳态值与弛豫时间的含义。

\end{minipage}\par

\par\noindent\begin{minipage}{\linewidth}
\textbf{解答}

\begin{center}
\includegraphics[width=\linewidth,height=.40\textheight,keepaspectratio]{\detokenize{assets/module-04/examples/M4-001.pdf}}
\par\smallskip{\small 图 M4-001：一阶弛豫：初始切线与稳态}
\end{center}
\end{minipage}\par

图 M4-001
中横轴按弛豫时间归一化，蓝线从零初值出发；水平虚线是稳态值。积分因子给出的指数项控制两者差值，短时展开对应曲线起点的切线。

乘以积分因子 \(e^{t/\tau}\) 后，方程化为
\(d(ve^{t/\tau})/dt=ae^{t/\tau}\)。从0积分到 \(t\) 并用零初值，得到
\(v=a\tau(1-e^{-t/\tau})\)。图中的水平渐近线对应
\(a\tau\)；初始切线斜率为 \(a\)，在 \(t=\tau\) 时已达到稳态的
\(1-e^{-1}\)。

当 \(t\ll\tau\) 时，展开 \(e^{-t/\tau}=1-t/\tau+O((t/\tau)^2)\)，得到

\[
v(t)=at-\frac{at^2}{2\tau}+\cdots\simeq at.
\]

因此图中起点附近近似沿初始切线运动，弛豫项还来不及积累。当 \(t\gg\tau\)
时，

\[
v(t)\simeq a\tau,\qquad
v_{\infty}-v(t)=a\tau e^{-t/\tau}.
\]

稳态值 \(v_{\infty}=a\tau\) 是驱动与弛豫平衡的结果。每经过一个
\(\tau\)，与稳态的差值缩小为原来的 \(e^{-1}\)；\(\tau\)
并不表示到达稳态所需的某个截然截止的时刻。

宏观电路中的 \(I\)、\(U\)、\(R\)
来自导体内局域的载流子运动。德鲁德模型用一个动量弛豫时间描述散射对平均定向动量的破坏，从而在经典近似下建立漂移速度、电流密度、电导率和焦耳热之间的联系。

\subsection{一、微观参量与方向约定（基础必讲）}\label{m4:ux4e00ux5faeux89c2ux53c2ux91cfux4e0eux65b9ux5411ux7ea6ux5b9aux57faux7840ux5fc5ux8bb2}

设某一类载流子的电荷为 \(q\)、有效质量为 \(m\)、体数密度为
\(n\)、平均漂移速度为 \(\mathbf v_d\)，则：

\begin{itemize}
\tightlist
\item
  \(n\) 的单位为 \(\mathrm{m^{-3}}\)；
\item
  \(\mathbf v_d\) 是速度分布的定向平均，不是单个粒子的实际速度；
\item
  \(\tau_m\) 是\textbf{动量弛豫时间}，表示平均定向动量衰减的时间尺度；
\item
  \(\mathbf J\) 是约定电流密度，总电流为
\end{itemize}

\[
I=\iint_S\mathbf J\cdot d\mathbf S.
\]

金属中主要载流子为电子，此时 \(q=-e\)。电子的 \(\mathbf v_d\) 与
\(\mathbf E\) 反向，但 \(\mathbf J\) 与 \(\mathbf E\)
同向。无规则运动速度远大于 \(v_d\)，却因方向杂乱而不贡献净电流。

\par\noindent\begin{minipage}{\linewidth}
\begin{center}
\includegraphics[width=\linewidth,height=.40\textheight,keepaspectratio]{\detokenize{assets/module-04/original/tikz-m4-topic27-v8-01.png}}
\par\smallskip{\small 图示：导体中无规则运动、电子平均漂移和约定电流的方向关系}
\end{center}
\end{minipage}\par

\subsection{二、电流密度与漂移速度（基础必讲）}\label{m4:ux4e8cux7535ux6d41ux5bc6ux5ea6ux4e0eux6f02ux79fbux901fux5ea6ux57faux7840ux5fc5ux8bb2}

取法向为 \(\mathbf n\) 的微小面元 \(\Delta S\)。在 \(\Delta t\)
内，载流子的定向运动扫过有向体积

\[
(\mathbf v_d\cdot\mathbf n)\Delta S\,\Delta t.
\]

因而穿过面元的代数电荷为

\[
\Delta Q=nq(\mathbf v_d\cdot\mathbf n)\Delta S\,\Delta t.
\]

与 \(\Delta Q=(\mathbf J\cdot\mathbf n)\Delta S\,\Delta t\) 比较，得

\[
\boxed{\mathbf J=nq\mathbf v_d}.
\]

对电子 \(q=-e\)，所以

\[
\boxed{\mathbf J=-ne\mathbf v_d},
\qquad
j=nev_d.
\]

\par\noindent\begin{minipage}{\linewidth}
\begin{center}
\includegraphics[width=\linewidth,height=.40\textheight,keepaspectratio]{\detokenize{assets/module-04/theory/review-tikz-m4-topic27-v8-02.pdf}}
\par\smallskip{\small 图示：由载流子在短时间内扫过的有向体积推得电流密度}
\end{center}
\end{minipage}\par

\par\addvspace{.4\baselineskip}
\noindent\begin{minipage}{\linewidth}

\textbf{例 M4-002：电子漂移与约定电流}

均匀金属中每单位体积有 \(n>0\) 个传导电子，每个电子电荷为
\(-e\)，\(e>0\)。选取沿导线的 \(+x\) 方向，已知平均漂移速度为
\(\mathbf v_d=-v_d\mathbf e_x\)，\(v_d>0\)。求约定电流密度矢量，标明它与电子漂移的方向关系；不把电子无规则热运动的速率当作漂移速率。

\end{minipage}\par

\par\noindent\begin{minipage}{\linewidth}
\textbf{解答}

\begin{center}
\includegraphics[width=\linewidth,height=.40\textheight,keepaspectratio]{\detokenize{assets/module-04/examples/M4-002.pdf}}
\par\smallskip{\small 图 M4-002：负电子漂移与正电流方向}
\end{center}
\end{minipage}\par

图 M4-002 的电子漂移箭头向左，约定电流箭头向右。将两个带负号的量 \(q\)
与 \(\mathbf v_d\) 相乘后，电流密度沿图示
\(+x\)，其大小还要乘电子数密度。

若电子以 \(\mathbf v_d=-v_d\mathbf e_x\) 漂移，则
\(\mathbf J=nev_d\mathbf e_x\)。负电荷向左运动，约定电流向右。

\subsection{三、动量弛豫方程与直流电导率（基础必讲）}\label{m4:ux4e09ux52a8ux91cfux5f1bux8c6bux65b9ux7a0bux4e0eux76f4ux6d41ux7535ux5bfcux7387ux57faux7840ux5fc5ux8bb2}

德鲁德模型不跟踪每一次碰撞，而对系统平均漂移速度使用弛豫时间近似：

\[
\boxed{
m\frac{d\mathbf v_d}{dt}
=q\mathbf E-\frac{m}{\tau_m}\mathbf v_d
}.
\]

式中 \(-m\mathbf v_d/\tau_m\)
表示散射使\textbf{平均定向动量}衰减，不表示每个电子在碰撞后都回到静止。

\par\noindent\begin{minipage}{\linewidth}
\begin{center}
\includegraphics[width=\linewidth,height=.40\textheight,keepaspectratio]{\detokenize{assets/module-04/theory/review-tikz-m4-topic27-v8-03.pdf}}
\par\smallskip{\small 图示：电场输入定向动量，散射使平均定向动量衰减}
\end{center}
\end{minipage}\par

在不随时间变化的均匀电场中，稳态满足 \(d\mathbf v_d/dt=0\)，故

\[
\mathbf v_d=\frac{q\tau_m}{m}\mathbf E.
\]

代入 \(\mathbf J=nq\mathbf v_d\)，得

\[
\boxed{
\mathbf J=\sigma_0\mathbf E
},
\qquad
\boxed{
\sigma_0=\frac{nq^2\tau_m}{m}
}.
\]

电子的 \(q=-e\) 在 \(q^2\) 中消去，所以各向同性普通导体的 \(\mathbf J\)
与 \(\mathbf E\) 同向。直流电阻率为

\[
\boxed{
\rho=\frac{1}{\sigma_0}
=\frac{m}{nq^2\tau_m}
}.
\]

\par\addvspace{.4\baselineskip}
\noindent\begin{minipage}{\linewidth}

\textbf{例 M4-003：散射增强与电阻率}

某均匀金属在弱直流场中可用单类载流子的德鲁德模型描述，体数密度为
\(n\)、有效质量为 \(m\)、电荷为 \(q\ne0\)、动量弛豫时间为
\(\tau_m\)。在温度等因素改变散射强度的比较中，近似保持 \(n,m\)
不变。利用稳态漂移关系判断散射增强、\(\tau_m\) 缩短时，电导率
\(\sigma_0\) 和电阻率 \(\rho\) 如何变化。

\end{minipage}\par

\par\noindent\begin{minipage}{\linewidth}
\textbf{解答}

\begin{center}
\includegraphics[width=\linewidth,height=.40\textheight,keepaspectratio]{\detokenize{assets/module-04/examples/M4-003.pdf}}
\par\smallskip{\small 图 M4-003：弛豫时间控制直流电导}
\end{center}
\end{minipage}\par

图 M4-003
的曲线说明弛豫时间是漂移建立的时间尺度；右侧两个材料关系分别正比、反比于同一个
\(\tau_m\)。比较散射前后时固定
\(n,m\)，不能把载流子浓度变化混入这个判断。

稳态时 \(q\mathbf E=m\mathbf v_d/\tau_m\)，所以
\(\mathbf J=nq\mathbf v_d=(nq^2\tau_m/m)\mathbf E\)。由此
\(\sigma_0=nq^2\tau_m/m\)，\(\rho=m/(nq^2\tau_m)\)；保持 \(n,m\)
不变后，弛豫时间缩短使电导率降低、电阻率增大。

\subsection{四、直流建立与交流响应（拔高选讲）}\label{m4:ux56dbux76f4ux6d41ux5efaux7acbux4e0eux4ea4ux6d41ux54cdux5e94ux62d4ux9ad8ux9009ux8bb2}

\subsubsection{1.
恒定电场突然加上}\label{m4:ux6052ux5b9aux7535ux573aux7a81ux7136ux52a0ux4e0a}

设 \(t=0\) 时加上恒定场 \(\mathbf E\)。动量弛豫方程的通解为

\[
\mathbf v_d(t)
=\mathbf v_d(0)e^{-t/\tau_m}
+\frac{q\tau_m}{m}\mathbf E
\left(1-e^{-t/\tau_m}\right).
\]

若 \(\mathbf v_d(0)=\mathbf0\)，则

\[
\boxed{
\mathbf J(t)=\sigma_0\mathbf E
\left(1-e^{-t/\tau_m}\right)
}.
\]

当 \(t=\tau_m\) 时，局域电流密度达到直流稳态值的
\(1-e^{-1}\approx63.2\%\)；当 \(t\approx5\tau_m\) 时已达到约
\(99.3\%\)。

\subsubsection{2. 简谐电场}\label{m4:ux7b80ux8c10ux7535ux573a}

取

\[
\mathbf E(t)=\operatorname{Re}
\left(\mathbf E_0e^{-i\omega t}\right),
\]

令
\(\mathbf J(t)=\operatorname{Re}(\mathbf J_0e^{-i\omega t})\)，代入弛豫方程，得复电导率

\[
\boxed{
\sigma(\omega)
=\frac{\sigma_0}{1-i\omega\tau_m}
}.
\]

其实部与虚部为

\[
\operatorname{Re}\sigma
=\frac{\sigma_0}{1+\omega^2\tau_m^2},
\qquad
\operatorname{Im}\sigma
=\frac{\sigma_0\omega\tau_m}{1+\omega^2\tau_m^2}.
\]

当 \(\omega\tau_m\ll1\) 时，\(\mathbf J\) 近似与 \(\mathbf E\)
同相，回到直流欧姆定律；当 \(\omega\tau_m\gg1\)
时，载流子惯性不可忽略，电流幅值减小并与电场出现显著相位差。

\par\noindent\begin{minipage}{\linewidth}
\begin{center}
\includegraphics[width=\linewidth,height=.40\textheight,keepaspectratio]{\detokenize{assets/module-04/theory/review-tikz-m4-topic27-v8-04.pdf}}
\par\smallskip{\small 图示：恒定场下电流以动量弛豫时间建立，简谐场下电导率具有幅值和相位响应}
\end{center}
\end{minipage}\par

\subsection{五、能量传递与焦耳热（基础必讲）}\label{m4:ux4e94ux80fdux91cfux4f20ux9012ux4e0eux7126ux8033ux70edux57faux7840ux5fc5ux8bb2}

单位体积内，电场对载流子做功的速率为

\[
p= n(q\mathbf E)\cdot\mathbf v_d
=\boxed{\mathbf J\cdot\mathbf E}.
\]

对平均漂移运动方程两边点乘 \(n\mathbf v_d\)，得

\[
\frac{d}{dt}
\left(\frac12nmv_d^2\right)
=\mathbf J\cdot\mathbf E
-\frac{nmv_d^2}{\tau_m}.
\]

这个能量方程表明：电场输入的功率一部分暂时增加定向运动能，另一部分通过散射转化为晶格、杂质等自由度的无规则能量。直流稳态时定向运动能不再增加，因而

\[
\boxed{
p=\mathbf J\cdot\mathbf E
=\sigma_0E^2
=\frac{j^2}{\sigma_0}
}.
\]

对简谐场，周期平均耗散功率密度为

\[
\boxed{
\langle p\rangle
=\frac12\operatorname{Re}\sigma(\omega)
|\mathbf E_0|^2
}.
\]

\par\noindent\begin{minipage}{\linewidth}
\begin{center}
\includegraphics[width=\linewidth,height=.40\textheight,keepaspectratio]{\detokenize{assets/module-04/original/tikz-m4-topic27-v8-05.png}}
\par\smallskip{\small 图示：电场输入的能量在定向运动能与散射耗散之间分配}
\end{center}
\end{minipage}\par

\subsection{六、德鲁德模型的物理图像与使用条件（基础必讲与拔高选讲）}\label{m4:ux516dux5fb7ux9c81ux5fb7ux6a21ux578bux7684ux7269ux7406ux56feux50cfux4e0eux4f7fux7528ux6761ux4ef6ux57faux7840ux5fc5ux8bb2ux4e0eux62d4ux9ad8ux9009ux8bb2}

\begin{enumerate}
\def\labelenumi{\arabic{enumi}.}
\tightlist
\item
  \textbf{漂移慢不等于信号传播慢。} \(v_d\)
  是载流子的平均定向速度；电路中的电磁扰动由场在导线及其周围空间传播，传播时间不是导线长度除以
  \(v_d\)。
\item
  \textbf{\(\tau_m\) 不是一切电路过程的时间常数。}
  它描述局域平均动量弛豫；完整电路的建立还受场传播、几何分布参数及
  \(RC\)、\(L/R\) 等时间尺度控制。
\item
  \textbf{金属电子不是经典理想气体。}
  德鲁德模型可给出电导率的结构，但实际金属电子遵守费米---狄拉克统计；微观速度尺度应以费米速率理解，不应把经典能量均分定理当成金属电子的定量描述。
\item
  \textbf{局域欧姆关系有条件。} \(\mathbf J=\sigma\mathbf E\)
  默认线性、局域、近似各向同性且温度等状态已给定；强场、强磁场、尺寸接近平均自由程或材料明显各向异性时，必须修改。
\end{enumerate}

\subsection{七、练习与解答}\label{m4:ux4e03ux7ec3ux4e60ux4e0eux89e3ux7b54}

\par\addvspace{.4\baselineskip}
\noindent\begin{minipage}{\linewidth}

\subsubsection{作业一：铜导线的漂移速度与动量弛豫时间}\label{m4:ux4f5cux4e1aux4e00ux94dcux5bfcux7ebfux7684ux6f02ux79fbux901fux5ea6ux4e0eux52a8ux91cfux5f1bux8c6bux65f6ux95f4}

\textbf{题号：M4-004}

\textbf{层级：} 基础　　\textbf{方法：} 载流子输运、德鲁德模型

导线为均匀圆直导线，电流在截面上均匀；取约定电流沿
\(+x\)，铜中传导电子沿相反方向漂移。忽略温升，并以给定电子质量代入经典德鲁德模型。

截面积为 \(S=1.0\times10^{-6}\ \mathrm{m^2}\) 的铜导线中通有
\(I=1.0\ \mathrm A\) 的稳恒电流。已知铜的质量密度
\(\rho_m=8.9\times10^3\ \mathrm{kg/m^3}\)、摩尔质量
\(M=64\times10^{-3}\ \mathrm{kg/mol}\)、电阻率
\(\rho=1.7\times10^{-8}\ \Omega\cdot\mathrm m\)。假设每个铜原子贡献一个传导电子。取

\[
e=1.6\times10^{-19}\ \mathrm C,
\quad
m_e=9.1\times10^{-31}\ \mathrm{kg},
\quad
N_A=6.02\times10^{23}\ \mathrm{mol^{-1}}.
\]

求：

\begin{enumerate}
\def\labelenumi{(\arabic{enumi})}
\item
  传导电子的漂移速率 \(v_d\)；
\item
  动量弛豫时间 \(\tau_m\)。
\end{enumerate}

\end{minipage}\par

\par\noindent\begin{minipage}{\linewidth}
\textbf{解答}

\begin{center}
\includegraphics[width=\linewidth,height=.40\textheight,keepaspectratio]{\detokenize{assets/module-04/examples/M4-004.pdf}}
\par\smallskip{\small 图 M4-004：铜导线：从原子计数到漂移}
\end{center}
\end{minipage}\par

图 M4-004
把电子漂移与测得的干路电流分开。先由质量密度和摩尔质量得到单位体积电子数，再由穿过截面
\(S\) 的电荷流求 \(v_d\)；电阻率提供的是动量弛豫时间的另一条独立关系。

每个原子贡献一个传导电子，故

\[
n=\frac{\rho_m}{M}N_A
=\frac{8.9\times10^3}{64\times10^{-3}}
\times6.02\times10^{23}
\approx8.4\times10^{28}\ \mathrm{m^{-3}}.
\]

由 \(I=ne v_dS\) 得

\[
v_d=\frac{I}{neS}
=\frac{1.0}{
(8.4\times10^{28})(1.6\times10^{-19})(1.0\times10^{-6})}
\approx7.4\times10^{-5}\ \mathrm{m/s}.
\]

电子的漂移方向与约定电流方向相反。再由

\[
\rho=\frac{m_e}{ne^2\tau_m}
\]

得

\[
\tau_m
=\frac{m_e}{ne^2\rho}
=\frac{9.1\times10^{-31}}
{(8.4\times10^{28})(1.6\times10^{-19})^2
(1.7\times10^{-8})}
\approx2.5\times10^{-14}\ \mathrm s.
\]

因而

\[
\boxed{v_d\approx7.4\times10^{-5}\ \mathrm{m/s}},
\qquad
\boxed{\tau_m\approx2.5\times10^{-14}\ \mathrm s}.
\]

两个量级并不矛盾：\(v_d\)
是巨大无规则速度背景上的极小定向偏置，\(\tau_m\)
是这个定向偏置的局域弛豫时间。

\subsection{八、附录一：局域欧姆定律与宏观电阻}\label{m4:ux516bux9644ux5f55ux4e00ux5c40ux57dfux6b27ux59c6ux5b9aux5f8bux4e0eux5b8fux89c2ux7535ux963b}

局域关系 \(\mathbf J=\sigma\mathbf E\)
经过截面积分和沿程积分，可得到均匀导体的 \(U=IR\) 及
\(R=L/(\sigma S)\)。下面在德鲁德模型条件下写出这一推导，并区分材料参数
\(\sigma\) 与几何参数 \(L,S\)。

\subsubsection{一、局域欧姆定律与端口伏安关系}\label{m4:ux4e00ux5c40ux57dfux6b27ux59c6ux5b9aux5f8bux4e0eux7aefux53e3ux4f0fux5b89ux5173ux7cfb}

\textbf{1. 两种表述}

\begin{itemize}
\tightlist
\item
  \textbf{宏观表述}：\(U = I R\)（描述整体电路元件的两端电压与通过电流的关系）
\item
  \textbf{微观表述}：\(\mathbf{J} = \sigma \mathbf{E}\)（描述导体内局域点的电流密度矢量与电场强度矢量的关系）
\end{itemize}

\par\noindent\begin{minipage}{\linewidth}
\begin{center}
\includegraphics[width=\linewidth,height=.40\textheight,keepaspectratio]{\detokenize{assets/module-04/theory/review-tikz-m4-topic27-v8-07.pdf}}
\par\smallskip{\small 图示：宏观欧姆定律来自导体内局域电流密度与电场关系的空间累加}
\end{center}
\end{minipage}\par

\par\noindent\begin{minipage}{\linewidth}
\begin{center}
\includegraphics[width=\linewidth,height=.40\textheight,keepaspectratio]{\detokenize{assets/module-04/theory/review-tikz-m4-topic27-v8-02.pdf}}
\par\smallskip{\small 图示：统计单位时间穿过截面的电子数时，需区分电子漂移方向与约定电流方向}
\end{center}
\end{minipage}\par

图中 \(\mathbf v_d\) 是电子的平均漂移速度，方向与常规电流密度
\(\mathbf J\) 相反；若只写大小，才有
\(j=ne v_d\)。竞赛推导中一旦从标量转到矢量，必须恢复
\(\mathbf J=-ne\mathbf v_d\)，否则后续与 \(\mathbf E\)
的方向关系会被写反。

\textbf{2. 推导}

考察一根长度为 \(L\)、横截面积为 \(S\)
的均匀各向同性圆柱体导体，两端施加恒定电压
\(U\)。忽略边缘效应，导体内建立均匀的轴向稳恒电场
\(\mathbf{E}\)，电场线平行于导体轴线。

由电势差与电场强度的积分关系：

\[
U = \int_{0}^{L} \mathbf{E} \cdot d\mathbf{l} = E L
\]

由电流与电流密度的积分关系（电流均匀分布）：

\[
I = \iint_{S} \mathbf{J} \cdot d\mathbf{S} = j S
\]

将 \(E = \frac{U}{L}\) 与 \(j = \frac{I}{S}\) 代入微观欧姆定律的标量形式
\(j = \sigma E\) 中，可得：

\[
\frac{I}{S} = \sigma \frac{U}{L}
\]

对上式进行代数重排，分离出电压 \(U\)：

\[
U = I \left( \frac{L}{\sigma S} \right)
\]

定义 \(R=L/(\sigma S)\)，上式即为 \(U=IR\)。其中 \(\sigma\)
来自局域载流子响应，\(L/S\) 来自导体几何；二者共同决定端口电阻。

\subsubsection{二、材料参数与导体几何}\label{m4:ux4e8cux6750ux6599ux53c2ux6570ux4e0eux5bfcux4f53ux51e0ux4f55}

\textbf{1. 两种表述}

\begin{itemize}
\tightlist
\item
  \textbf{宏观表述}：\(R = \rho \frac{L}{S}\)
\item
  \textbf{德鲁德模型}：\(\rho = \frac{m_e}{n e^2 \tau_m}\)
\end{itemize}

\textbf{2. 电阻的两个因数}

由 \(R = \frac{1}{\sigma} \frac{L}{S}\) 可直接分出几何因数和材料因数：

\begin{itemize}
\item
  \textbf{几何因数}
  \(\left( \frac{L}{S} \right)\)：由导体长度和横截面积决定。
\item
  \textbf{材料因数}（电阻率
  \(\rho\)）：在德鲁德模型中，\(\rho = \frac{1}{\sigma}\) 可写成

  \[
  R = \left( \frac{m_e}{n e^2 \tau_m} \right) \frac{L}{S}
  \]
\end{itemize}

\textbf{3. 材料与温度的影响}

\begin{itemize}
\tightlist
\item
  \textbf{材料差异}：不同金属在几何尺寸相同时电阻不同，可用有效质量、载流子密度与动量弛豫时间的差异来描述。
\item
  \textbf{温度效应}：对普通纯金属的一定温区，\(n\)
  和有效质量变化较小，升温主要通过增强声子散射使 \(\tau_m\) 缩短，从而使
  \(\rho\)
  增大。``主要''不等于``唯一''；半导体、载流子浓度随温度显著改变的材料以及超导转变都不能沿用这一简化论证。
\end{itemize}

\par\noindent\begin{minipage}{\linewidth}
\begin{center}
\includegraphics[width=\linewidth,height=.40\textheight,keepaspectratio]{\detokenize{assets/module-04/original/tikz-m4-topic27-v8-08.png}}
\par\smallskip{\small 图示：均匀导体的电阻可拆成几何因数与载流子、散射因数；普通金属升温主要通过缩短弛豫时间增大电阻}
\end{center}
\end{minipage}\par

这张图用于防止把 \(R=\rho L/S\) 误读为单一的``材料公式''：\(L/S\)
是几何因素，\(n\) 与 \(\tau\)
是材料和散射因素。对普通金属而言，升温通常先改变散射率；但半导体、超导体等体系的温度响应不能由这一条链条直接外推。

\par\noindent\begin{minipage}{\linewidth}
\begin{center}
\includegraphics[width=\linewidth,height=.40\textheight,keepaspectratio]{\detokenize{assets/module-04/theory/review-tikz-m4-topic27-v8-03.pdf}}
\par\smallskip{\small 图示：电场持续加速电子而散射不断削弱有序漂移，二者平衡给出电导率}
\end{center}
\end{minipage}\par

稳态并不是电子停止运动，而是热运动叠加在一个不再增长的平均漂移速度上。图中的阻尼项是弛豫时间近似；它解释了常规金属的欧姆响应，却不应被误读为单个电子在每次碰撞后都严格回到静止。

\subsubsection{三、
竞赛进阶：德鲁德模型的局限与索末菲量子修正}\label{m4:ux4e09-ux7adeux8d5bux8fdbux9636ux5fb7ux9c81ux5fb7ux6a21ux578bux7684ux5c40ux9650ux4e0eux7d22ux672bux83f2ux91cfux5b50ux4feeux6b63}

经典德鲁德模型把电子当作经典稀薄气体，因而无法正确说明金属电子比热、传热与散射相空间等量子性质。索末菲模型保留弛豫时间的输运结构，但将载流子改为服从费米---狄拉克统计的电子。

在普通金属中，输运主要由费米面附近可改变占据状态的电子贡献，其典型速度尺度为费米速率
\(v_F\)，因而平均自由程应写成

\[
\lambda=v_F\tau_m.
\]

在高于德拜温度的典型高温区，声子散射率常近似正比于 \(T\)，从而导致
\(\tau_m\propto T^{-1}\)、\(\rho\propto T\)。低温下声子散射的幂律、杂质残余电阻与电子间散射会改变温度依赖，不能把
\(\rho\propto T\) 外推到所有温区。

\subsection{九、使用条件}\label{m4:ux4e5dux4f7fux7528ux6761ux4ef6}

\begin{enumerate}
\def\labelenumi{\arabic{enumi}.}
\tightlist
\item
  德鲁德模型中的 \(\tau_m\)
  是平均动量弛豫时间，不是每个电子两次碰撞之间完全相同的时间间隔。
\item
  \(\mathbf J=nq\mathbf v_d\)
  中必须保留载流子电荷的符号；金属电子的漂移速度与约定电流方向相反。
\item
  \(\mathbf J=\sigma\mathbf E\)
  是线性、局域、各向同性近似；强场、霍尔效应、各向异性或非局域输运需要更一般的本构关系。
\item
  直流电导率 \(\sigma_0=nq^2\tau_m/m\) 不能在 \(\omega\tau_m\gtrsim1\)
  时直接替代复电导率。
\item
  动量弛豫时间、电荷弛豫时间和宏观电路建立时间是不同尺度，不能相互混用。
\end{enumerate}

\section{专题二十八：液体、气体与真空中的电流（基础必讲与拔高选讲）}\label{m4:ux4e13ux9898ux4e8cux5341ux516bux6db2ux4f53ux6c14ux4f53ux4e0eux771fux7a7aux4e2dux7684ux7535ux6d41ux57faux7840ux5fc5ux8bb2ux4e0eux62d4ux9ad8ux9009ux8bb2}

金属中的电流主要由电子漂移形成；电解质、气体和真空中的载流子、产生机制与伏安关系并不相同。以下用电流密度表示各类载流子的贡献，并讨论电解、气体放电和电子束运动中可以定量计算的模型。

\subsection{一、不同介质中的载流子与电流密度}\label{m4:ux4e00ux4e0dux540cux4ecbux8d28ux4e2dux7684ux8f7dux6d41ux5b50ux4e0eux7535ux6d41ux5bc6ux5ea6}

若第 \(s\) 类载流子的电荷为 \(q_s\)、数密度为 \(n_s\)、平均漂移速度为
\(\mathbf v_s\)，总电流密度为

\[
\boxed{\mathbf J=\sum_s n_sq_s\mathbf v_s}.
\]

负载流子的漂移方向与常规电流方向相反，但 \(q_s\mathbf v_s\)
仍沿常规电流方向。只有在线性、局域、各向同性且载流子密度近似不变时，才可进一步写成
\(\mathbf J=\sigma\mathbf E\)。

\par\noindent\begin{minipage}{\linewidth}
\begin{center}
\includegraphics[width=\linewidth,height=.40\textheight,keepaspectratio]{\detokenize{assets/module-04/theory/tikz-m4-topic28-v9-01-carriers.pdf}}
\par\smallskip{\small 图示：金属、电解质、气体和真空中的载流子与常规电流方向}
\end{center}
\end{minipage}\par

\par\addvspace{.4\baselineskip}
\noindent\begin{minipage}{\linewidth}

\textbf{例 M4-005：两类离子共同贡献电导}

一均匀电解质含两种完全解离的一价离子，正、负离子电荷分别为
\(+e,-e\)，两者数密度均为 \(n\)，迁移率分别为正常数
\(\mu_+,\mu_-\)。在弱匀强电场 \(\mathbf E\)
下，仅考虑迁移漂移，忽略浓度梯度引起的扩散及对流。写出两种漂移速度和各自电流密度，求总电流密度及电导率，说明异号离子为何对约定电流同向贡献。

\end{minipage}\par

\par\noindent\begin{minipage}{\linewidth}
\textbf{解答}

\begin{center}
\includegraphics[width=\linewidth,height=.40\textheight,keepaspectratio]{\detokenize{assets/module-04/examples/M4-005.pdf}}
\par\smallskip{\small 图 M4-005：两类离子的电流同向相加}
\end{center}
\end{minipage}\par

图 M4-005
的正离子沿场漂移、负离子逆场漂移；但每类电流都等于数密度乘电荷再乘速度。因此下排的两个负号相消，右侧应相加两项电导。

正离子速度为 \(\mu_+\mathbf E\)，负离子速度为
\(-\mu_-\mathbf E\)。将各自电荷与速度相乘，负离子一项出现两个负号，因此两项均沿外场，总电导率为
\(ne(\mu_++\mu_-)\)。

\subsection{二、电解质中的电流与法拉第电解定律}\label{m4:ux4e8cux7535ux89e3ux8d28ux4e2dux7684ux7535ux6d41ux4e0eux6cd5ux62c9ux7b2cux7535ux89e3ux5b9aux5f8b}

设离子价数为 \(z\)，每个离子携带电荷量 \(ze\)。通过电解池的总电荷为

\[
Q=\int I(t)\,dt.
\]

每沉积 \(1\ \mathrm{mol}\) 物质需要电荷 \(zF\)，其中

\[
F=N_Ae\simeq9.65\times10^4\ \mathrm{C/mol}
\]

为法拉第常量。因此摩尔质量为 \(M\) 的物质在电流效率 \(\eta\) 下沉积质量

\[
\boxed{m=\eta\frac{M}{zF}Q}.
\]

这里 \(M\) 必须使用 \(\mathrm{kg/mol}\)
或与结果质量单位相容的单位。电极反应决定 \(z\)，不能只看化学式下标。

\par\noindent\begin{minipage}{\linewidth}
\begin{center}
\includegraphics[width=\linewidth,height=.40\textheight,keepaspectratio]{\detokenize{assets/module-04/theory/tikz-m4-topic28-v9-02-electrolysis.pdf}}
\par\smallskip{\small 图示：电解池中的离子迁移、电子交换与沉积质量记账}
\end{center}
\end{minipage}\par

\par\addvspace{.4\baselineskip}
\noindent\begin{minipage}{\linewidth}

\textbf{例 M4-006：银离子的电解沉积}

电解装置以银离子 \(\mathrm{Ag}^+\)
在阴极还原为金属银。流过电解液的总电荷为
\(Q=965\ \mathrm C\)，银的摩尔质量为
\(M=0.108\ \mathrm{kg/mol}\)，电流效率 \(\eta=1\)。取法拉第常量
\(F=9.65\times10^4\ \mathrm{C/mol}\)，不考虑其他阴极反应。求沉积银的质量，写明一价离子对应的电子转移数。

\end{minipage}\par

\par\noindent\begin{minipage}{\linewidth}
\textbf{解答}

\begin{center}
\includegraphics[width=\linewidth,height=.40\textheight,keepaspectratio]{\detokenize{assets/module-04/examples/M4-006.pdf}}
\par\smallskip{\small 图 M4-006：电解沉积：电荷转化为物质量}
\end{center}
\end{minipage}\par

图 M4-006
左电极为阴极，银离子向它迁移并获得一个电子。反应中的一电子计数使每摩尔沉积银对应电荷
\(F\)，再乘摩尔质量和电流效率即得沉积质量。

阴极反应为 \(\mathrm{Ag}^++e^-\to\mathrm{Ag}\)，每摩尔银需要电荷
\(F\)。因此 \(m=\eta MQ/F=0.00108\ \mathrm{kg}=1.08\ \mathrm g\)。

\subsection{三、气体放电的产生、维持与工作点}\label{m4:ux4e09ux6c14ux4f53ux653eux7535ux7684ux4ea7ux751fux7ef4ux6301ux4e0eux5de5ux4f5cux70b9}

常态气体中自由载流子很少。外部辐射、碰撞电离或阴极发射先产生少量电子和离子；电场加速电子后，若碰撞电离能够补偿复合与扩散损失，放电可以自持。

气体放电的伏安特性一般不是一条直线。竞赛中常把它分为非自持区、击穿过渡区和自持放电区。击穿电压
\(U_b\) 是触发条件；放电形成后维持电压可能降到
\(U_h<U_b\)。外接串联电阻用来限制电流，实际工作点由

\[
I=\frac{U_s-U_g(I)}{R}
\]

与放电管伏安曲线的交点决定。

\par\noindent\begin{minipage}{\linewidth}
\begin{center}
\includegraphics[width=\linewidth,height=.40\textheight,keepaspectratio]{\detokenize{assets/module-04/theory/tikz-m4-topic28-v9-03-gas-discharge.pdf}}
\par\smallskip{\small 图示：气体放电伏安曲线、击穿电压、维持电压与负载线}
\end{center}
\end{minipage}\par

\par\addvspace{.4\baselineskip}
\noindent\begin{minipage}{\linewidth}

\textbf{例 M4-007：放电管的稳态限流}

一只已处于稳定导通状态的放电管，与 \(R=20\ \mathrm{k}\Omega\)
电阻串联接到 \(U_s=400\ \mathrm V\) 理想直流电源。用恒定维持电压
\(U_h=200\ \mathrm V\)
近似导通管的伏安特性，不讨论点火瞬态。求回路稳态电流，指出限流电阻两端的电压，并解释为什么这里应使用维持电压而非击穿电压。

\end{minipage}\par

\par\noindent\begin{minipage}{\linewidth}
\textbf{解答}

\begin{center}
\includegraphics[width=\linewidth,height=.40\textheight,keepaspectratio]{\detokenize{assets/module-04/examples/M4-007.pdf}}
\par\smallskip{\small 图 M4-007：已导通放电管的限流}
\end{center}
\end{minipage}\par

图 M4-007
的电源经串联电阻向已导通放电管供电。沿回路把电源升压减去电阻、放电管压降，电阻上的剩余电压才用于求电流。

沿图示电源送电方向写KVL，\(U_s-IR-U_h=0\)。电阻承担
\(200\ \mathrm V\)，故
\(I=(400-200)/(20\times10^3)=0.010\ \mathrm A\)。本题预先限定已导通，不能拿击穿门槛替代稳态管压降。

\subsection{四、真空中的电子束}\label{m4:ux56dbux771fux7a7aux4e2dux7684ux7535ux5b50ux675f}

真空中没有足够多的自由载流子，必须由热发射、光电发射或场发射等方式从阴极提供电子。电子从静止经过加速电压
\(U_a>0\) 后，忽略相对论效应和空间电荷，有

\[
eU_a=\frac12mv^2,
\qquad
\boxed{v=\sqrt{\frac{2eU_a}{m}}}.
\]

\par\noindent\begin{minipage}{\linewidth}
\begin{center}
\includegraphics[width=\linewidth,height=.40\textheight,keepaspectratio]{\detokenize{assets/module-04/theory/tikz-m4-topic28-v9-04-vacuum-beam.pdf}}
\par\smallskip{\small 图示：电子由阴极发射、经加速电压形成准直电子束}
\end{center}
\end{minipage}\par

若 \(eU_a\) 与电子静能 \(mc^2\)
相比不再很小，应改用相对论能量关系。若束流较强，电子间排斥产生的空间电荷也会改变电势分布。

\subsection{五、电子束的电偏转与磁偏转}\label{m4:ux4e94ux7535ux5b50ux675fux7684ux7535ux504fux8f6cux4e0eux78c1ux504fux8f6c}

电子以速度 \(v\) 沿 \(x\) 方向进入长度为 \(\ell\)、板间距为 \(d\)
的平行板。偏转电压为 \(U_y\)，忽略边缘场，则竖直加速度大小为

\[
a_y=\frac{eU_y}{md}.
\]

若屏到偏转板末端的距离为 \(L\)，小偏转近似下屏上位移大小为

电子在板间运动时间为
\(t_1=\ell/v\)，板内偏移和离开偏转板时的竖直速度分别为

\[
y_1=\frac12a_yt_1^2,
\qquad
v_y=a_yt_1.
\]

电子离开板后再以竖直速度 \(v_y\) 漂移 \(t_2=L/v\)，因而

\[
y=y_1+v_yt_2
=\frac{eU_y}{mdv^2}\left(\frac{\ell^2}{2}+\ell L\right).
\]

再用 \(v^2=2eU_a/m\) 消去电子的电荷和质量，得到

\[
\boxed{y=\frac{U_y\ell(L+\ell/2)}{2U_ad}}.
\]

\par\noindent\begin{minipage}{\linewidth}
\begin{center}
\includegraphics[width=\linewidth,height=.40\textheight,keepaspectratio]{\detokenize{assets/module-04/theory/tikz-m4-topic28-v9-05-beam-deflection.pdf}}
\par\smallskip{\small 图示：电子束经过加速区、电偏转板和漂移区后到达荧光屏}
\end{center}
\end{minipage}\par

匀强磁场垂直于速度时只改变方向。电子轨道半径满足

\[
\boxed{R=\frac{mv}{eB}}.
\]

电偏转与磁偏转的符号都必须由电子带负电这一事实单独判断，不能把公式中的
\(e\) 同时当作正电荷和带符号电荷。

\subsection{六、示波器的时间展开与电压标定}\label{m4:ux516dux793aux6ce2ux5668ux7684ux65f6ux95f4ux5c55ux5f00ux4e0eux7535ux538bux6807ux5b9a}

示波器把待测电压送入竖直偏转系统，把近似线性的扫描电压送入水平偏转系统。稳定波形要求扫描周期与信号周期保持确定比例；触发电路的作用是让每次扫描从信号的相同相位开始。

\par\noindent\begin{minipage}{\linewidth}
\begin{center}
\includegraphics[width=\linewidth,height=.40\textheight,keepaspectratio]{\detokenize{assets/module-04/theory/tikz-m4-topic28-v9-06-oscilloscope.pdf}}
\par\smallskip{\small 图示：示波器由竖直输入、水平扫描和触发同步形成稳定波形}
\end{center}
\end{minipage}\par

若竖直灵敏度为 \(k_y\ \mathrm{V/div}\)、峰值占 \(n_y\) 格，则峰值为
\(U_m=k_yn_y\)。若水平时基为 \(k_t\ \mathrm{s/div}\)、一个周期占 \(n_t\)
格，则

\[
T=k_tn_t,
\qquad
f=\frac1T.
\]

\subsection{七、使用条件与常见错误}\label{m4:ux4e03ux4f7fux7528ux6761ux4ef6ux4e0eux5e38ux89c1ux9519ux8bef}

\begin{enumerate}
\def\labelenumi{\arabic{enumi}.}
\tightlist
\item
  法拉第电解定律统计通过界面的总电荷，不要求电流恒定。
\item
  气体放电的击穿电压与维持电压不是同一个量，负载电阻决定工作点和安全电流。
\item
  电子束公式默认忽略相对论、空间电荷、碰撞和边缘场。
\item
  示波器屏上的几何高度必须乘以档位系数，探头衰减倍数也要计入。
\item
  电解质、气体和半导体不一定满足常数电阻模型，使用叠加原理前必须检查线性。
\end{enumerate}

\subsection{八、练习与解答}\label{m4:ux516bux7ec3ux4e60ux4e0eux89e3ux7b54}

\par\addvspace{.4\baselineskip}
\noindent\begin{minipage}{\linewidth}

\subsubsection{作业一：铜的电解沉积}\label{m4:ux4f5cux4e1aux4e00ux94dcux7684ux7535ux89e3ux6c89ux79ef}

\textbf{题号：M4-008}

\textbf{层级：} 基础　　\textbf{方法：} 电解定律、载流子输运

电流在给定时段保持不变，\(\eta\)
定义为实际用于铜沉积的电荷占总电荷的比例。取法拉第常量
\(F=9.65\times10^4\ \mathrm{C/mol}\)，铜离子每沉积一个原子需要两个电子。

硫酸铜溶液中以 \(I=2.00\ \mathrm A\) 的恒定电流电解
\(30.0\ \mathrm{min}\)。铜离子为二价，铜的摩尔质量
\(M=63.5\ \mathrm{g/mol}\)，电流效率为 \(90\%\)。求阴极沉积的铜质量。

\end{minipage}\par

\par\noindent\begin{minipage}{\linewidth}
\textbf{解答}

\begin{center}
\includegraphics[width=\linewidth,height=.40\textheight,keepaspectratio]{\detokenize{assets/module-04/examples/M4-008.pdf}}
\par\smallskip{\small 图 M4-008：铜的电解沉积}
\end{center}
\end{minipage}\par

图 M4-008
中阴极附近的离子还原反应给出每个铜原子所需的电子数。先把电流与持续时间相乘得到总电荷，再取有效沉积部分，最后按反应计量换为铜的质量。

通过电解池的电荷为

\[
Q=It=2.00\times30.0\times60=3.60\times10^3\ \mathrm C.
\]

二价铜离子每沉积 \(1\ \mathrm{mol}\) 需要 \(2F\) 电荷，因此

\[
m=\eta\frac{M}{2F}Q
=0.90\times\frac{63.5}{2\times9.65\times10^4}\times3.60\times10^3\ \mathrm g.
\]

故

\[
\boxed{m=1.07\ \mathrm g}.
\]

\par\addvspace{.4\baselineskip}
\noindent\begin{minipage}{\linewidth}

\subsubsection{作业二：两类离子的电导贡献}\label{m4:ux4f5cux4e1aux4e8cux4e24ux7c7bux79bbux5b50ux7684ux7535ux5bfcux8d21ux732e}

\textbf{题号：M4-009}

\textbf{层级：} 基础　　\textbf{方法：} 载流子输运、电导叠加

两种离子完全解离，溶液除这两类外无其他载流子；弱场下迁移率为正，速度分别沿电场与逆电场，忽略扩散和对流。取
\(e>0\)，先由电中性求负离子数密度。

某电解质中正离子带电 \(+2e\)、数密度为 \(n\)、迁移率为
\(\mu_+\)；负离子带电 \(-e\)、迁移率为
\(\mu_-\)。溶液宏观电中性。求弱电场中的电流密度和电导率。

\end{minipage}\par

\par\noindent\begin{minipage}{\linewidth}
\textbf{解答}

\begin{center}
\includegraphics[width=\linewidth,height=.40\textheight,keepaspectratio]{\detokenize{assets/module-04/examples/M4-009.pdf}}
\par\smallskip{\small 图 M4-009：两类离子的电导贡献}
\end{center}
\end{minipage}\par

图 M4-009
将两种离子的速度方向和电流方向分开。宏观电中性首先固定负离子数密度，随后将各自的数密度、电荷及漂移速度相乘，两类电流同向相加。

设负离子数密度为 \(n_-\)。电中性条件给出

\[
2en-en_-=0,
\qquad n_-=2n.
\]

正离子的漂移速度为 \(\mu_+\mathbf E\)，负离子的漂移速度为
\(-\mu_-\mathbf E\)。于是

\[
\mathbf J
=n(2e)\mu_+\mathbf E+(2n)(-e)(-\mu_-\mathbf E).
\]

整理得

\[
\boxed{\mathbf J=2ne(\mu_++\mu_-)\mathbf E},
\qquad
\boxed{\sigma=2ne(\mu_++\mu_-)}.
\]

\par\addvspace{.4\baselineskip}
\noindent\begin{minipage}{\linewidth}

\subsubsection{作业三：放电管的触发与限流}\label{m4:ux4f5cux4e1aux4e09ux653eux7535ux7ba1ux7684ux89e6ux53d1ux4e0eux9650ux6d41}

\textbf{题号：M4-010}

\textbf{层级：} 基础　　\textbf{方法：} 气体放电、分段电路

电源、放电管及限流电阻构成唯一串联回路，导线电阻忽略；击穿前管电流近似为零，题中给定两种电压分别用于截止判定和导通后的稳态模型。

直流电源电压为 \(U_s=400\ \mathrm V\)，经串联电阻
\(R=20.0\ \mathrm{k}\Omega\) 接一只放电管。放电管击穿电压为
\(U_b=300\ \mathrm V\)，导通后的维持电压近似为
\(U_h=200\ \mathrm V\)。判断能否击穿，并求导通后的电流和电阻功率。

\end{minipage}\par

\par\noindent\begin{minipage}{\linewidth}
\textbf{解答}

\begin{center}
\includegraphics[width=\linewidth,height=.40\textheight,keepaspectratio]{\detokenize{assets/module-04/examples/M4-010.pdf}}
\par\smallskip{\small 图 M4-010：放电管的触发与限流}
\end{center}
\end{minipage}\par

图 M4-010
以电源和串联电阻的负载线区分击穿前与稳定导通后。先用截止时管压判断能否触发，再用维持电压与负载线的交点求电流和电阻功率。

击穿前电流极小，电阻压降近似为零，管端电压约为
\(400\ \mathrm V>U_b\)，因此可以击穿。

导通后管压降降为 \(U_h\)，工作电流为

\[
I=\frac{U_s-U_h}{R}
=\frac{400-200}{20.0\times10^3}
=10.0\ \mathrm{mA}.
\]

限流电阻功率为

\[
P_R=I^2R=(10.0\times10^{-3})^2\times20.0\times10^3=2.00\ \mathrm W.
\]

故

\[
\boxed{I=10.0\ \mathrm{mA}},
\qquad
\boxed{P_R=2.00\ \mathrm W}.
\]

\par\addvspace{.4\baselineskip}
\noindent\begin{minipage}{\linewidth}

\subsubsection{作业四：电子束的电偏转}\label{m4:ux4f5cux4e1aux56dbux7535ux5b50ux675fux7684ux7535ux504fux8f6c}

\textbf{题号：M4-011}

\textbf{层级：} 基础　　\textbf{方法：} 洛伦兹力、带电粒子运动

电子电荷为 \(-e\)、质量为
\(m_e\)，采用非相对论近似，不受磁场及重力影响。电子沿偏转板中线、平行极板进入，入射横向速度为零；\(L\)
从板末端量到屏，求偏移的大小，方向由负电子与场向关系判断。

电子由静止经 \(U_a=2.00\ \mathrm{kV}\) 加速后，进入长度
\(\ell=4.00\ \mathrm{cm}\)、板间距 \(d=1.00\ \mathrm{cm}\)
的偏转板。板间电压为 \(U_y=20.0\ \mathrm V\)，偏转板末端到屏的距离为
\(L=20.0\ \mathrm{cm}\)。忽略边缘场，求屏上偏移量。

\end{minipage}\par

\par\noindent\begin{minipage}{\linewidth}
\textbf{解答}

\begin{center}
\includegraphics[width=\linewidth,height=.40\textheight,keepaspectratio]{\detokenize{assets/module-04/examples/M4-011.pdf}}
\par\smallskip{\small 图 M4-011：电子束的电偏转}
\end{center}
\end{minipage}\par

图 M4-011
从偏转板入口、出口再追踪到屏。板内电子受与电场反向的恒定横向力，出板后保留横向速度作匀速偏移；屏距从板末端起算，所以两段位移要相加。

利用已经推得的偏转公式

\[
y=\frac{U_y\ell(L+\ell/2)}{2U_ad},
\]

统一换成米后代入：

\[
y=\frac{20.0\times0.0400\times(0.200+0.0200)}
{2\times2000\times0.0100}
=4.40\times10^{-3}\ \mathrm m.
\]

所以偏移量大小为

\[
\boxed{y=4.40\ \mathrm{mm}}.
\]

电子的实际偏转方向与电场方向相反。

\par\addvspace{.4\baselineskip}
\noindent\begin{minipage}{\linewidth}

\subsubsection{作业五：由示波器读出信号参数}\label{m4:ux4f5cux4e1aux4e94ux7531ux793aux6ce2ux5668ux8bfbux51faux4fe1ux53f7ux53c2ux6570}

\textbf{题号：M4-012}

\textbf{层级：} 基础　　\textbf{方法：} 电学测量、示波器读图

所给竖直灵敏度是示波器输入端的刻度，尚未自动补偿探头的10倍衰减；波形关于中心线对称、无直流偏置。水平时基不因探头倍率改变。

示波器竖直灵敏度为 \(2.00\ \mathrm{V/div}\)，探头处于 \(10\times\)
档；屏上正弦波从中心线到波峰为 \(3.00\) 格。水平时基为
\(0.500\ \mathrm{ms/div}\)，一个周期占 \(4.00\)
格。求被测电压的峰值、有效值和频率。

\end{minipage}\par

\par\noindent\begin{minipage}{\linewidth}
\textbf{解答}

\begin{center}
\includegraphics[width=\linewidth,height=.40\textheight,keepaspectratio]{\detokenize{assets/module-04/examples/M4-012.pdf}}
\par\smallskip{\small 图 M4-012：由示波器读出信号参数}
\end{center}
\end{minipage}\par

图 M4-012
的竖直三格表示中心线到波峰，水平四格表示一个周期。竖直刻度还要乘探头衰减倍率，时间轴按原时基读取，最后只对正弦波用峰值除以根号二。

示波器输入端对应的峰值为

\[
U_{m,\mathrm{in}}=2.00\times3.00=6.00\ \mathrm V.
\]

计入 \(10\times\) 探头后，被测点峰值为

\[
U_m=60.0\ \mathrm V,
\qquad
U_{\mathrm{rms}}=\frac{U_m}{\sqrt2}=42.4\ \mathrm V.
\]

周期与频率为

\[
T=0.500\times4.00=2.00\ \mathrm{ms},
\qquad
f=\frac1T=500\ \mathrm{Hz}.
\]

因此

\[
\boxed{U_m=60.0\ \mathrm V},
\qquad
\boxed{U_{\mathrm{rms}}=42.4\ \mathrm V},
\qquad
\boxed{f=500\ \mathrm{Hz}}.
\]

\section{专题二十九：半导体导电、PN
结与基本器件（基础必讲与拔高选讲）}\label{m4:ux4e13ux9898ux4e8cux5341ux4e5dux534aux5bfcux4f53ux5bfcux7535pn-ux7ed3ux4e0eux57faux672cux5668ux4ef6ux57faux7840ux5fc5ux8bb2ux4e0eux62d4ux9ad8ux9009ux8bb2}

半导体的载流子浓度会随温度、掺杂、光照和电场显著变化，不能只用``电阻介于导体和绝缘体之间''来概括。器件题应先判断载流子和结的状态，再选择线性或分段线性模型。

\subsection{一、电子、空穴与本征半导体}\label{m4:ux4e00ux7535ux5b50ux7a7aux7a74ux4e0eux672cux5f81ux534aux5bfcux4f53}

价带中缺少一个电子时，可把其余电子的集体运动等效为一个带正电的空穴。若电子和空穴数密度分别为
\(n,p\)，迁移率分别为 \(\mu_n,\mu_p\)，弱场电导率为

\[
\boxed{\sigma=e(n\mu_n+p\mu_p)}.
\]

本征半导体中电子和空穴成对产生，故 \(n=p=n_i\)。温度升高通常使 \(n_i\)
急剧增加，因此半导体电阻率常随温度下降；这与普通金属中散射增强导致电阻率上升的机制不同。

\par\noindent\begin{minipage}{\linewidth}
\begin{center}
\includegraphics[width=\linewidth,height=.40\textheight,keepaspectratio]{\detokenize{assets/module-04/theory/tikz-m4-topic29-v9-01-electron-hole.pdf}}
\par\smallskip{\small 图示：价带电子跃迁后形成导带电子与价带空穴}
\end{center}
\end{minipage}\par

\par\addvspace{.4\baselineskip}
\noindent\begin{minipage}{\linewidth}

\textbf{例 M4-013：本征载流子增多时的电导率}

一块均匀本征半导体中电子、空穴浓度均为
\(n_i=10^{16}\ \mathrm{m^{-3}}\)，对应迁移率为
\(\mu_n,\mu_p\)。考虑温度变化使平衡本征浓度增至原来的十倍，并近似认为两种迁移率不变。在弱电场线性输运范围内，比较变化前后的电导率，明确电子与空穴的电导贡献如何相加。

\end{minipage}\par

\par\noindent\begin{minipage}{\linewidth}
\textbf{解答}

\begin{center}
\includegraphics[width=\linewidth,height=.40\textheight,keepaspectratio]{\detokenize{assets/module-04/examples/M4-013.pdf}}
\par\smallskip{\small 图 M4-013：本征载流子浓度决定电导比较}
\end{center}
\end{minipage}\par

图 M4-013
以两类相反电荷的定向运动表示双载流子贡献。对本征半导体，应将它们解释为电子与空穴；二者浓度同时增大，而迁移率按题设固定，所以总电导按同一倍数变化。

本征条件使 \(n=p=n_i\)，故
\(\sigma=e n_i(\mu_n+\mu_p)\)。图中两类载流子浓度同时乘10，迁移率固定，电导率也乘10；这个比较只在题设忽略迁移率温度变化的近似下成立。

\subsection{二、施主、受主与多数载流子}\label{m4:ux4e8cux65bdux4e3bux53d7ux4e3bux4e0eux591aux6570ux8f7dux6d41ux5b50}

施主杂质容易提供导带电子，形成 n
型半导体；受主杂质容易俘获价带电子并留下空穴，形成 p
型半导体。在非简并、热平衡且完全电离的近似下，有

\[
np=n_i^2.
\]

若 \(N_D\gg N_A,n_i\)，则 n 型材料中

\[
n\simeq N_D-N_A,
\qquad
p\simeq\frac{n_i^2}{n}.
\]

\par\noindent\begin{minipage}{\linewidth}
\begin{center}
\includegraphics[width=\linewidth,height=.40\textheight,keepaspectratio]{\detokenize{assets/module-04/theory/tikz-m4-topic29-v9-02-doping.pdf}}
\par\smallskip{\small 图示：施主和受主掺杂分别形成 n 型与 p 型多数载流子}
\end{center}
\end{minipage}\par

掺杂并不使材料带有宏观净电荷：施主失去电子后留下固定正离子，受主得到电子后留下固定负离子，整体仍近似电中性。

\subsection{三、PN
结、耗尽层与偏置方向}\label{m4:ux4e09pn-ux7ed3ux8017ux5c3dux5c42ux4e0eux504fux7f6eux65b9ux5411}

p 区与 n
区接触后，多数载流子跨结扩散并复合，在结附近留下不能自由移动的离化杂质，形成耗尽层。耗尽层中的内建电场从
n 区固定正电荷指向 p 区固定负电荷，它阻碍多数载流子继续扩散。

\par\noindent\begin{minipage}{\linewidth}
\begin{center}
\includegraphics[width=\linewidth,height=.40\textheight,keepaspectratio]{\detokenize{assets/module-04/theory/tikz-m4-topic29-v9-03-pn-junction.pdf}}
\par\smallskip{\small 图示：PN 结耗尽层中的固定离子、内建电场与扩散方向}
\end{center}
\end{minipage}\par

\begin{itemize}
\tightlist
\item
  正向偏置降低势垒并缩窄耗尽层，电流迅速增大。
\item
  反向偏置提高势垒并加宽耗尽层，只留下很小的反向电流；达到击穿区后不能再用开路模型。
\end{itemize}

理想指数模型可写成

\[
I=I_s\left(e^{U_D/(n_{\mathrm{id}}U_T)}-1\right),
\qquad
U_T=\frac{k_BT}{e},
\]

其中 \(n_{\mathrm{id}}\) 是理想因子，\(U_T\)
是热电压。但多数竞赛电路题会明确采用理想二极管模型或恒压降模型。

\subsection{四、二极管的三种常用模型}\label{m4:ux56dbux4e8cux6781ux7ba1ux7684ux4e09ux79cdux5e38ux7528ux6a21ux578b}

\par\noindent\begin{minipage}{\linewidth}
\begin{center}
\includegraphics[width=\linewidth,height=.40\textheight,keepaspectratio]{\detokenize{assets/module-04/theory/tikz-m4-topic29-v9-04-diode-models.pdf}}
\par\smallskip{\small 图示：理想、恒压降与指数三种二极管模型的适用层次}
\end{center}
\end{minipage}\par

\begin{enumerate}
\def\labelenumi{\arabic{enumi}.}
\tightlist
\item
  \textbf{理想模型}：正向导通时 \(U_D=0\)，反向截止时 \(I_D=0\)。
\item
  \textbf{恒压降模型}：导通时 \(U_D\simeq U_F\)，硅二极管常取 \(0.6\) 至
  \(0.7\ \mathrm V\)。
\item
  \textbf{指数模型}：需要研究小信号电阻、温度或较宽电流范围时使用。
\end{enumerate}

分段模型的标准步骤是：先假设每只二极管的状态，在线性电路中求解，再检查导通电流和端电压是否与假设一致。若不一致，必须更换状态重新求解。

\par\addvspace{.4\baselineskip}
\noindent\begin{minipage}{\linewidth}

\textbf{例 M4-014：硅二极管的正向与反向}

一个 \(5.00\ \mathrm V\) 理想直流电源、\(1.00\ \mathrm{k}\Omega\)
电阻和一只硅二极管构成串联回路。二极管用正向导通压降
\(0.700\ \mathrm V\)、反向完全截止的分段模型描述，不考虑反向击穿。分别讨论阳极朝向电源正端的正接与二极管反接两种接法，求回路电流并检验所假设工作状态。

\end{minipage}\par

\par\noindent\begin{minipage}{\linewidth}
\textbf{解答}

\begin{center}
\includegraphics[width=\linewidth,height=.40\textheight,keepaspectratio]{\detokenize{assets/module-04/examples/M4-014.pdf}}
\par\smallskip{\small 图 M4-014：先假设导通，再检查极性}
\end{center}
\end{minipage}\par

图 M4-014
中串联电阻承担电源电压减去正向压降后的部分。正接试算电流必须与阳极到阴极的方向一致；反接时用开路模型，不能沿用正向压降公式。

正接的导通试算给
\(I=(5.00-0.700)/(1.00\times10^3)=4.30\ \mathrm{mA}>0\)，与从阳极到阴极的方向一致。反接时取
\(I=0\)，电阻无压降，二极管承受反向电压，在题设模型内截止。

\subsection{五、三极管的开关与电流放大模型}\label{m4:ux4e94ux4e09ux6781ux7ba1ux7684ux5f00ux5173ux4e0eux7535ux6d41ux653eux5927ux6a21ux578b}

以 NPN
管为例，发射结正向偏置、集电结反向偏置时处于放大区。在简单直流模型中

\[
I_C\simeq\beta I_B,
\qquad
I_E=I_B+I_C.
\]

\par\noindent\begin{minipage}{\linewidth}
\begin{center}
\includegraphics[width=\linewidth,height=.40\textheight,keepaspectratio]{\detokenize{assets/module-04/theory/tikz-m4-topic29-v9-05-transistor-regions.pdf}}
\par\smallskip{\small 图示：NPN 三极管的截止、放大和饱和三个工作区}
\end{center}
\end{minipage}\par

若基极驱动过小，管子截止；若 \(\beta I_B\)
要求的集电流超过外电路能够提供的最大值，管子进入饱和区，此时不能继续使用
\(I_C=\beta I_B\)。常取
\(U_{BE}\simeq0.7\ \mathrm V\)、\(U_{CE,\mathrm{sat}}\simeq0.2\ \mathrm V\)
作分段估算。

三极管的``电流放大''不意味着凭空产生能量。输出功率来自集电极电源，基极小信号只控制能量释放。

\subsection{六、半导体模型的使用条件与常见错误}\label{m4:ux516dux534aux5bfcux4f53ux6a21ux578bux7684ux4f7fux7528ux6761ux4ef6ux4e0eux5e38ux89c1ux9519ux8bef}

\begin{enumerate}
\def\labelenumi{\arabic{enumi}.}
\tightlist
\item
  电子和空穴的常规电流贡献同向，不能因电子带负电就把两项相减。
\item
  n 型、p 型描述多数载流子，不表示材料带宏观净电荷。
\item
  二极管是非线性器件，整个电路一般不能直接使用叠加原理。
\item
  恒压降不是材料常数；电流、温度和器件类型改变时会变化。
\item
  \(I_C=\beta I_B\) 只适用于放大区，截止和饱和区必须改用相应约束。
\end{enumerate}

\subsection{七、练习与解答}\label{m4:ux4e03ux7ec3ux4e60ux4e0eux89e3ux7b54-1}

\par\addvspace{.4\baselineskip}
\noindent\begin{minipage}{\linewidth}

\subsubsection{作业一：本征半导体的电阻率}\label{m4:ux4f5cux4e1aux4e00ux672cux5f81ux534aux5bfcux4f53ux7684ux7535ux963bux7387}

\textbf{题号：M4-015}

\textbf{层级：} 基础　　\textbf{方法：} 半导体器件、载流子输运

材料处于均匀本征热平衡，电子与空穴都是参与导电的载流子；采用弱场线性输运，迁移率按正值给出。取基本电荷
\(e=1.60\times10^{-19}\ \mathrm C\)。

某本征半导体中
\(n=p=1.00\times10^{16}\ \mathrm{m^{-3}}\)，电子与空穴迁移率分别为
\(\mu_n=0.140\ \mathrm{m^2/(V\,s)}\)、\(\mu_p=0.0500\ \mathrm{m^2/(V\,s)}\)。求电导率和电阻率。

\end{minipage}\par

\par\noindent\begin{minipage}{\linewidth}
\textbf{解答}

\begin{center}
\includegraphics[width=\linewidth,height=.40\textheight,keepaspectratio]{\detokenize{assets/module-04/examples/M4-015.pdf}}
\par\smallskip{\small 图 M4-015：本征半导体的电阻率}
\end{center}
\end{minipage}\par

图 M4-015
把电子与空穴两项电导分别列出。两者电荷、漂移方向虽相反，约定电流仍同向，因此求和后再取倒数得到电阻率，不能把迁移率相减。

由

\[
\sigma=e(n\mu_n+p\mu_p)
\]

得到

\[
\sigma
=1.60\times10^{-19}\times10^{16}\times(0.140+0.0500)
=3.04\times10^{-4}\ \mathrm{S/m}.
\]

因此

\[
\boxed{\rho=\frac1\sigma=3.29\times10^3\ \Omega\,\mathrm m}.
\]

\par\addvspace{.4\baselineskip}
\noindent\begin{minipage}{\linewidth}

\subsubsection{作业二：补偿掺杂后的多数载流子}\label{m4:ux4f5cux4e1aux4e8cux8865ux507fux63baux6742ux540eux7684ux591aux6570ux8f7dux6d41ux5b50}

\textbf{题号：M4-016}

\textbf{层级：} 基础　　\textbf{方法：} 半导体器件、载流子计数

样品均匀且处于热平衡，施主受主全部电离，采用非简并近似及质量作用关系
\(np=n_i^2\)；正的净施主浓度与自由电子浓度需由电中性同时确定。

某硅样品的施主和受主浓度分别为
\(N_D=5.00\times10^{21}\ \mathrm{m^{-3}}\)、\(N_A=2.00\times10^{21}\ \mathrm{m^{-3}}\)，本征载流子浓度为
\(n_i=1.00\times10^{16}\ \mathrm{m^{-3}}\)。假设完全电离，求电子和空穴浓度，并判断材料类型。

\end{minipage}\par

\par\noindent\begin{minipage}{\linewidth}
\textbf{解答}

\begin{center}
\includegraphics[width=\linewidth,height=.40\textheight,keepaspectratio]{\detokenize{assets/module-04/examples/M4-016.pdf}}
\par\smallskip{\small 图 M4-016：补偿掺杂后的多数载流子}
\end{center}
\end{minipage}\par

图 M4-016
先比较施主、受主的净浓度，再用电子与空穴的平衡关系补足少数载流子。电中性计入已电离的固定杂质电荷，而质量作用关系限制可动电子、空穴的乘积。

施主多于受主，材料为 n 型。多数电子浓度近似为

\[
n\simeq N_D-N_A=3.00\times10^{21}\ \mathrm{m^{-3}}.
\]

由质量作用关系 \(np=n_i^2\)，

\[
p=\frac{n_i^2}{n}
=\frac{10^{32}}{3.00\times10^{21}}
=3.33\times10^{10}\ \mathrm{m^{-3}}.
\]

故

\[
\boxed{n=3.00\times10^{21}\ \mathrm{m^{-3}}},
\qquad
\boxed{p=3.33\times10^{10}\ \mathrm{m^{-3}}}.
\]

\par\addvspace{.4\baselineskip}
\noindent\begin{minipage}{\linewidth}

\subsubsection{作业三：二极管串联限流}\label{m4:ux4f5cux4e1aux4e09ux4e8cux6781ux7ba1ux4e32ux8054ux9650ux6d41}

\textbf{题号：M4-017}

\textbf{层级：} 基础　　\textbf{方法：} 半导体器件、分段电路

导线电阻忽略，正接指二极管阳极经电阻接正电源、阴极接负端；反接指交换二极管两端，忽略反向漏电及击穿。功率取各元件吸收或输出的稳态值。

\(5.00\ \mathrm V\) 理想电源、\(1.00\ \mathrm{k}\Omega\)
电阻与一只硅二极管串联。采用 \(0.700\ \mathrm V\)
恒压降模型，分别求二极管正接和反接时的电流，并检查功率。

\end{minipage}\par

\par\noindent\begin{minipage}{\linewidth}
\textbf{解答}

\begin{center}
\includegraphics[width=\linewidth,height=.40\textheight,keepaspectratio]{\detokenize{assets/module-04/examples/M4-017.pdf}}
\par\smallskip{\small 图 M4-017：二极管串联限流}
\end{center}
\end{minipage}\par

图 M4-017
将导通与截止作为两种候选支路状态。正接先扣除导通压降后求电阻电流，再验证其符号；反接用截止模型，电阻与二极管耗散功率均按对应状态求。

正接时先假设导通，

\[
I=\frac{5.00-0.700}{1.00\times10^3}=4.30\ \mathrm{mA}>0,
\]

与导通假设一致。电阻和二极管功率分别为

\[
P_R=I^2R=18.5\ \mathrm{mW},
\qquad
P_D=U_DI=3.01\ \mathrm{mW}.
\]

两者之和为 \(21.5\ \mathrm{mW}=5.00I\)。

反接且未击穿时，模型给出

\[
\boxed{I\simeq0}.
\]

\par\addvspace{.4\baselineskip}
\noindent\begin{minipage}{\linewidth}

\subsubsection{作业四：二极管限幅器的分段输出}\label{m4:ux4f5cux4e1aux56dbux4e8cux6781ux7ba1ux9650ux5e45ux5668ux7684ux5206ux6bb5ux8f93ux51fa}

\textbf{题号：M4-018}

\textbf{层级：} 基础　　\textbf{方法：} 半导体器件、分段电路

电源负端与地是同一参考节点，输出电压按该节点对地定义，负载近似开路；二极管反向漏电忽略，输入范围限定在未发生反向击穿的区间。

输入电压 \(u_i\) 经 \(1.00\ \mathrm{k}\Omega\)
电阻接到输出节点。节点通过一只硅二极管接地，二极管阳极接输出、阴极接地，采用
\(0.700\ \mathrm V\) 恒压降模型；负载电阻很大。求 \(u_o(u_i)\)。

\end{minipage}\par

\par\noindent\begin{minipage}{\linewidth}
\textbf{解答}

\begin{center}
\includegraphics[width=\linewidth,height=.40\textheight,keepaspectratio]{\detokenize{assets/module-04/examples/M4-018.pdf}}
\par\smallskip{\small 图 M4-018：二极管限幅器的分段输出}
\end{center}
\end{minipage}\par

图 M4-018
左侧的二极管特性决定输出节点何时被钳位。截止时负载近似开路，输入串联电阻无压降；正向导通时输出固定在模型压降，剩余输入电压落在电阻上。

先假设二极管截止。支路无电流，电阻无压降，因此
\(u_o=u_i\)。截止要求二极管未正向达到 \(0.700\ \mathrm V\)，即
\(u_i\le0.700\ \mathrm V\)。

当 \(u_i>0.700\ \mathrm V\) 时，二极管导通并把输出近似固定在
\(0.700\ \mathrm V\)，电阻电流为
\((u_i-0.700)/1.00\ \mathrm{k}\Omega>0\)，与导通假设一致。

故

\[
\boxed{
u_o=
\begin{cases}
u_i,&u_i\le0.700\ \mathrm V,\\
0.700\ \mathrm V,&u_i>0.700\ \mathrm V.
\end{cases}}
\]

\par\noindent\begin{minipage}{\linewidth}
\begin{center}
\includegraphics[width=\linewidth,height=.40\textheight,keepaspectratio]{\detokenize{assets/module-04/process/M4-018a.pdf}}
\par\smallskip{\small 图 M4-018a：阳极接输出、阴极接地的限幅回路}
\end{center}
\end{minipage}\par

图 M4-018a
把实际连接展开：二极管阳极在输出节点侧，阴极条线朝地。输出向正方向升高才可能导通，故上限被钳位；负向输入在未击穿范围沿截止支路传递。

\par\addvspace{.4\baselineskip}
\noindent\begin{minipage}{\linewidth}

\subsubsection{作业五：判断三极管的放大与饱和}\label{m4:ux4f5cux4e1aux4e94ux5224ux65adux4e09ux6781ux7ba1ux7684ux653eux5927ux4e0eux9971ux548c}

\textbf{题号：M4-019}

\textbf{层级：} 基础　　\textbf{方法：} 半导体器件、分段电路

以地为发射极及各电压源的共同参考，忽略基极电源和集电极电源内阻。对每个给定
\(R_B\) 都应先求基极电流，再检验 \(I_C=\beta I_B\)
是否与集电极电路允许的电流相容，并给出工作区和主要电流、电压。

NPN 管发射极接地，集电极经 \(R_C=1.00\ \mathrm{k}\Omega\) 接
\(V_{CC}=12.0\ \mathrm V\)，基极经 \(R_B\) 接 \(5.00\ \mathrm V\)。取
\(U_{BE}=0.700\ \mathrm V\)、\(\beta=100\)、\(U_{CE,\mathrm{sat}}=0.200\ \mathrm V\)。分别讨论
\(R_B=100\ \mathrm{k}\Omega\) 与 \(R_B=10.0\ \mathrm{k}\Omega\)。

\end{minipage}\par

\par\noindent\begin{minipage}{\linewidth}
\textbf{解答}

\begin{center}
\includegraphics[width=\linewidth,height=.40\textheight,keepaspectratio]{\detokenize{assets/module-04/examples/M4-019.pdf}}
\par\smallskip{\small 图 M4-019：三极管的放大与饱和}
\end{center}
\end{minipage}\par

图 M4-019
将基极试算和集电极限流分开。每个给定基极电阻先决定基流，再比较放大模型所需集电流与集电极回路允许值，超出后改用饱和压降而非继续套用固定增益。

当 \(R_B=100\ \mathrm{k}\Omega\) 时，

\[
I_B=\frac{5.00-0.700}{100\times10^3}=43.0\ \mu\mathrm A,
\]

放大区预测 \(I_C=\beta I_B=4.30\ \mathrm{mA}\)。相应

\[
U_{CE}=12.0-4.30\times10^{-3}\times10^3=7.70\ \mathrm V,
\]

明显高于饱和压降，故放大区假设成立。

当 \(R_B=10.0\ \mathrm{k}\Omega\) 时，放大区预测
\(I_C=43.0\ \mathrm{mA}\)，但外电路最大只能提供

\[
I_{C,\mathrm{sat}}=\frac{12.0-0.200}{1.00\times10^3}=11.8\ \mathrm{mA}.
\]

因此管子饱和。结论为

\[
\boxed{R_B=100\ \mathrm{k}\Omega: I_C=4.30\ \mathrm{mA},\ U_{CE}=7.70\ \mathrm V},
\]

\[
\boxed{R_B=10.0\ \mathrm{k}\Omega: I_C\simeq11.8\ \mathrm{mA},\ U_{CE}\simeq0.200\ \mathrm V}.
\]

\section{专题三十：稳恒电流场与导电介质边值问题（基础必讲与拔高选讲）}\label{m4:ux4e13ux9898ux4e09ux5341ux7a33ux6052ux7535ux6d41ux573aux4e0eux5bfcux7535ux4ecbux8d28ux8fb9ux503cux95eeux9898ux57faux7840ux5fc5ux8bb2ux4e0eux62d4ux9ad8ux9009ux8bb2}

在非均匀导体中，\(\mathbf J=\sigma\mathbf E\)
还要与电荷守恒、无旋条件和边界条件联立。以下讨论导电介质分界面、表面电荷与电荷弛豫，并处理绝缘空腔、漏电和接地电极。

\subsection{一、导电介质分界面的电流边界条件}\label{m4:ux4e00ux5bfcux7535ux4ecbux8d28ux5206ux754cux9762ux7684ux7535ux6d41ux8fb9ux754cux6761ux4ef6}

在分析复杂或非均匀导体时，电流常常需要穿过两种具有不同电导率（\(\sigma_1\)
和 \(\sigma_2\)）的介质分界面。此时，电流密度 \(\mathbf{J}\) 和局域电场
\(\mathbf{E}\)
的分量一般会发生跃变。下面从电荷守恒与法拉第定律推导分界面条件。

\par\noindent\begin{minipage}{\linewidth}
\begin{center}
\includegraphics[width=\linewidth,height=.40\textheight,keepaspectratio]{\detokenize{assets/module-04/original/tikz-m4-topic30-v9-from28-v8-01.png}}
\par\smallskip{\small 图示：导电介质界面处法向电流连续来自电荷守恒，切向电场连续来自无旋条件}
\end{center}
\end{minipage}\par

\par\noindent\begin{minipage}{\linewidth}
\begin{center}
\includegraphics[width=\linewidth,height=.40\textheight,keepaspectratio]{\detokenize{assets/module-04/theory/review-tikz-m4-topic30-v9-from28-v8-02.pdf}}
\par\smallskip{\small 图示：跨界薄高斯柱必须取外法向；跨界薄回路的绕行正向由所选法向按右手定则确定}
\end{center}
\end{minipage}\par

两种积分小工具的极限过程不同：薄高斯柱保留上下底面的法向通量，而薄回路保留沿界面的两条长边环量。必须先固定法向与绕行正向再列符号，不能从``电流会连续''这一结论倒推符号。

\subsubsection{1.
法向分量的连续性（电荷守恒与恒定条件）}\label{m4:ux6cd5ux5411ux5206ux91cfux7684ux8fdeux7eedux6027ux7535ux8377ux5b88ux6052ux4e0eux6052ux5b9aux6761ux4ef6}

\textbf{1. 物理图像}

在稳恒电流状态下，导体内及分界面上的电荷分布不随时间改变。连续性方程因而给出
\(\nabla\cdot\mathbf J=0\)。

\textbf{2. 严密推导}

在介质 1 与介质 2
的分界面处，构造一个跨越界面的极薄的高斯圆柱体（柱底面积为
\(\Delta S\)，高
\(h \to 0\)）。对该闭合曲面应用电流的连续性方程积分形式：

\[
\iint_{\partial V}\mathbf J\cdot d\mathbf S=0
\]

设界面法向矢量 \(\mathbf n_{12}\) 从介质 1 指向介质 2，则有：

\[
-j_{1n} \Delta S + j_{2n} \Delta S = 0
\]

由此可得电流密度的法向分量关系：

\[
\boxed{j_{2n}-j_{1n}=0}.
\]

更一般的瞬态边界关系是

\[
\boxed{
j_{1n}-j_{2n}=\frac{\partial\sigma_s}{\partial t}
},
\]

其中 \(\sigma_s\) 是界面自由面电荷密度。稳恒时
\(\partial_t\sigma_s=0\)，才得到 \(j_{1n}=j_{2n}\)。

\textbf{物理结论}：稳恒电流场的电流密度矢量在两种介质的分界面上，其法向分量是连续的。

\subsubsection{2.
切向分量的跃变关系（电场的无旋性质）}\label{m4:ux5207ux5411ux5206ux91cfux7684ux8dc3ux53d8ux5173ux7cfbux7535ux573aux7684ux65e0ux65cbux6027ux8d28}

\textbf{1. 物理图像}

稳恒电场本质上是由电荷产生的保守场，静电力做功与路径无关，满足
\(\nabla \times \mathbf{E} = 0\)。

\textbf{2. 严密推导}

在分界面上构造一个包含部分介质 1 和部分介质 2 的微小矩形闭合环路（长为
\(\Delta l\)，宽
\(h\to0\)）。对准静态或稳恒电流场，\(\partial\mathbf B/\partial t=0\)，因而

\[
\oint_L \mathbf{E} \cdot d\mathbf{l} = 0
\]

沿环路积分可得：

\[
E_{1t} \Delta l - E_{2t} \Delta l = 0
\]

\[
\boxed{E_{1t}=E_{2t}}.
\]

根据微观欧姆定律
\(\mathbf{J} = \sigma \mathbf{E}\)，将电场切向分量替换为电流密度切向分量与电导率的比值：

\[
\boxed{
\frac{j_{1t}}{\sigma_1}
=\frac{j_{2t}}{\sigma_2}
}.
\]

\textbf{物理结论}：在 \(E_t\ne0\) 且 \(\sigma_1\ne\sigma_2\)
时，电流密度的切向分量发生跃变；若 \(E_t=0\)，两侧切向电流都为零。

\subsubsection{3.
稳恒电流的折射定律}\label{m4:ux7a33ux6052ux7535ux6d41ux7684ux6298ux5c04ux5b9aux5f8b}

结合法向分量和切向分量的关系，可以推导出电流线在分界面处的``折射''规律。

设介质 1 中的电流密度矢量 \(\mathbf{J}_1\) 与界面法线夹角为入射角
\(\theta_1\)，介质 2 中的电流密度矢量 \(\mathbf{J}_2\)
与界面法线夹角为折射角 \(\theta_2\)。

由几何关系可知：

\[
\tan \theta_1 = \frac{j_{1t}}{j_{1n}}
\]

\[
\tan \theta_2 = \frac{j_{2t}}{j_{2n}}
\]

将两式相除，并代入边界条件 \(j_{1n} = j_{2n}\) 与
\(\frac{j_{1t}}{j_{2t}} = \frac{\sigma_1}{\sigma_2}\)，可严密导出稳恒电流流线的折射定律：

\[
\boxed{
\frac{\tan\theta_1}{\tan\theta_2}
=\frac{\sigma_1}{\sigma_2}
}.
\]

\textbf{直观理解}：电流总是倾向于在电导率更高（电阻率更低）的介质中沿切向``走得更远''。当电流从低电导率介质进入高电导率介质时，流线会向远离法线的方向偏折（\(\theta_2 > \theta_1\)）。

\par\noindent\begin{minipage}{\linewidth}
\begin{center}
\includegraphics[width=\linewidth,height=.40\textheight,keepaspectratio]{\detokenize{assets/module-04/theory/review-tikz-m4-topic30-v9-from28-v8-03.pdf}}
\par\smallskip{\small 图示：稳恒电流穿过导电界面时流线折射由电导率比值决定，界面电荷负责调节电场}
\end{center}
\end{minipage}\par

\par\noindent\begin{minipage}{\linewidth}
\begin{center}
\includegraphics[width=\linewidth,height=.40\textheight,keepaspectratio]{\detokenize{assets/module-04/original/tikz-m4-topic30-v9-from28-v8-04.png}}
\par\smallskip{\small 图示：导电介质界面题从法向电流连续和切向电场连续出发，依次得到折射角与界面电荷}
\end{center}
\end{minipage}\par

这张流程图强调两个不同的积分工具分别给出两条不同的连续性条件。先固定法向，再列高斯柱和回路的符号，最后才使用
\(\mathbf J=\sigma\mathbf E\)；若直接假定整个 \(\mathbf J\)
矢量连续，会得到错误的折射关系。

\subsubsection{4.
界面真实电荷的堆积机制}\label{m4:ux754cux9762ux771fux5b9eux7535ux8377ux7684ux5806ux79efux673aux5236}

分界面处一般会堆积一层真实的宏观自由面电荷。它所产生的附加电场使两侧同时满足电流守恒和电场边界条件。

设两介质的介电常数分别为 \(\varepsilon_1\) 和
\(\varepsilon_2\)。根据高斯定理，静电场的法向边界条件为：

\[
\boxed{D_{2n}-D_{1n}=\sigma_s}.
\]

即：

\[
\varepsilon_2E_{2n}-\varepsilon_1E_{1n}=\sigma_s.
\]

其中 \(\sigma_s\) 为界面的自由面电荷密度。由于 \(E_n=j_n/\sigma\)
且稳恒时 \(j_{1n}=j_{2n}=j_n\)，代入上式可得

\[
\boxed{
\sigma_s=j_n\left(
\frac{\varepsilon_2}{\sigma_2}
-\frac{\varepsilon_1}{\sigma_1}
\right)
}.
\]

在接通电源的瞬态中，\(j_{1n}-j_{2n}=\partial_t\sigma_s\)
描述界面电荷的积累。当电荷产生的附加电场使两侧达到 \(j_{1n}=j_{2n}\)
后，\(\sigma_s\) 不再随时间变化。

若

\[
\frac{\varepsilon_2}{\sigma_2}
=\frac{\varepsilon_1}{\sigma_1},
\]

即两介质的电荷弛豫时间 \(\varepsilon/\sigma\) 相同，则纵然
\(\sigma_1\ne\sigma_2\)，界面自由面电荷仍可为零。因而``电导率不同必然堆积界面电荷''并不严格。

\subsection{二、导线弯折处的表面电荷与稳恒电流}\label{m4:ux4e8cux5bfcux7ebfux5f2fux6298ux5904ux7684ux8868ux9762ux7535ux8377ux4e0eux7a33ux6052ux7535ux6d41}

在宏观电路的常规分析中，我们往往直接假设电流能够沿任意形状的导线流动。更基本的问题是：接通电源后，局域电场怎样调整到符合弯折边界？答案不在于单个电子沿确定轨道``转弯''，而在于表面电荷从瞬态到稳恒态的重新分布。

\subsubsection{一、稳恒边值问题与表面电荷}\label{m4:ux4e00ux7a33ux6052ux8fb9ux503cux95eeux9898ux4e0eux8868ux9762ux7535ux8377}

对局域、各向同性的均匀导体，稳恒时

\[
\mathbf J=\sigma\mathbf E,
\qquad
\nabla\cdot\mathbf J=0,
\qquad
\nabla\times\mathbf E=0.
\]

若导线侧表面与绝缘介质相接，还必须满足

\[
\boxed{\mathbf J\cdot\mathbf n=0}.
\]

因此真正的稳恒电场是一个由导线几何、电极电势与表面电荷分布共同决定的边值问题。导线弯折、截面变化或材料交界时，表面电荷产生的附加场使总电场满足新的边界条件。

\par\noindent\begin{minipage}{\linewidth}
\begin{center}
\includegraphics[width=\linewidth,height=.40\textheight,keepaspectratio]{\detokenize{assets/module-04/theory/review-tikz-m4-topic30-v9-from28-v8-05.pdf}}
\par\smallskip{\small 图示：弯折导线中表面电荷逐步建立横向附加电场，使总电场沿导线切向}
\end{center}
\end{minipage}\par

\par\noindent\begin{minipage}{\linewidth}
\begin{center}
\includegraphics[width=\linewidth,height=.40\textheight,keepaspectratio]{\detokenize{assets/module-04/original/tikz-m4-topic30-v9-from28-v8-06.png}}
\par\smallskip{\small 图示：金属中电子的无规则热运动远快于平均漂移，常规电流方向与电子漂移方向相反}
\end{center}
\end{minipage}\par

弯折处的``转弯''不是每个电子沿光滑轨道被机械地拐过去，而是大量散射后的平均漂移方向由局域总电场决定。

\subsubsection{二、稳恒电流的建立过程}\label{m4:ux4e8cux7a33ux6052ux7535ux6d41ux7684ux5efaux7acbux8fc7ux7a0b}

\begin{enumerate}
\def\labelenumi{\arabic{enumi}.}
\tightlist
\item
  \textbf{电磁扰动传播。}
  接通电源后，电场与磁场的改变沿电路结构传播，不会在整条导线上同时建立。
\item
  \textbf{法向电流导致电荷积累。} 若瞬时场在弯折侧壁上产生
  \(j_n\ne0\)，表面电荷会按局域连续性方程改变。
\item
  \textbf{表面电荷重塑电场。} 新的表面电荷分布产生附加场，使侧壁的
  \(j_n\) 减小。
\item
  \textbf{达到稳恒。} 当场、电流和表面电荷都不再随时间变化时，侧壁满足
  \(\mathbf J\cdot\mathbf n=0\)，导线内部满足
  \(\nabla\cdot\mathbf J=0\)。
\end{enumerate}

\par\noindent\begin{minipage}{\linewidth}
\begin{center}
\includegraphics[width=\linewidth,height=.40\textheight,keepaspectratio]{\detokenize{assets/module-04/theory/review-tikz-m4-topic30-v9-from28-v8-07.pdf}}
\par\smallskip{\small 图示：弯折导线从接通到稳恒的过程由表面电荷产生的横向附加场完成负反馈调节}
\end{center}
\end{minipage}\par

导线弯折处的表面电荷调整局域总电场，使其满足边界条件；这里讨论的是大量载流子的平均漂移。

\subsubsection{\texorpdfstring{三、局域体电荷弛豫：\(\tau_c=\varepsilon/\sigma\)}{三、局域体电荷弛豫：\textbackslash tau\_c=\textbackslash varepsilon/\textbackslash sigma}}\label{m4:ux4e09ux5c40ux57dfux4f53ux7535ux8377ux5f1bux8c6btau_cvarepsilonsigma}

下面推导的是\textbf{均匀线性各向同性导电介质中局域自由体电荷}的衰减时间，不是任意宏观电路建立稳态的通用时间。这里还假定
\(\varepsilon\)、\(\sigma\) 在讨论区域内不随位置变化。
还要求在相关时间尺度上可使用瞬时欧姆关系
\(\mathbf J=\sigma\mathbf E\)；若载流子惯性不可忽略，则不能把直流电导率外推到任意短时。

设导电介质的电导率为 \(\sigma\)，介电常数为
\(\varepsilon\)。假设在某一极短瞬间，导体内某局域出现自由体电荷密度
\(\rho_f\)。

根据电荷守恒定律（连续性方程）：

\[
\nabla \cdot \mathbf{J} + \frac{\partial \rho_f}{\partial t} = 0
\]

根据微观欧姆定律，将 \(\mathbf{J} = \sigma \mathbf{E}\) 代入上式：

\[
\sigma (\nabla \cdot \mathbf{E}) + \frac{\partial \rho_f}{\partial t} = 0
\]

由 \(\nabla\cdot\mathbf D=\rho_f\)、\(\mathbf D=\varepsilon\mathbf E\)
以及均匀 \(\varepsilon\)，有
\(\nabla\cdot\mathbf E=\rho_f/\varepsilon\)。代入可得：

\[
\frac{\sigma}{\varepsilon} \rho_f + \frac{\partial \rho_f}{\partial t} = 0
\]

这是一个关于 \(\rho_f\) 的一阶常微分方程，对其积分求解得：

\[
\rho_f(t)=\rho_f(0)\exp\left(-\frac{t}{\tau_c}\right)
\]

\[
\boxed{\tau_c=\frac{\varepsilon}{\sigma}}.
\]

\par\noindent\begin{minipage}{\linewidth}
\begin{center}
\includegraphics[width=\linewidth,height=.40\textheight,keepaspectratio]{\detokenize{assets/module-04/theory/review-tikz-m4-topic30-v9-from28-v8-08.pdf}}
\par\smallskip{\small 图示：瞬时欧姆近似有效时，均匀介质内局域体电荷指数衰减，失去的电荷可转移到边界}
\end{center}
\end{minipage}\par

若形式上把铜的直流电导率 \(\sigma_0\simeq5.9\times10^7\,\mathrm{S/m}\)
和背景 \(\varepsilon\simeq\varepsilon_0\) 代入上式，会得到

\[
\frac{\varepsilon_0}{\sigma_0}\simeq1.5\times10^{-19}\,\mathrm s.
\]

这个数是瞬时欧姆模型的\textbf{形式外推值}，不能直接认作铜的真实超短时电荷衰减常数。前文的动量弛豫时间约为
\(2.5\times10^{-14}\,\mathrm s\)，远大于此形式值，因而短时惯性不能忽略。

保留Drude惯性时，

\[
\tau_m\frac{\partial\mathbf J}{\partial t}+\mathbf J=\sigma_0\mathbf E.
\]

对其取散度并联立连续性方程和高斯定律，得到

\[
\tau_m\ddot\rho_f+\dot\rho_f+\frac{\sigma_0}{\varepsilon}\rho_f=0.
\]

只有在慢变、过阻尼且 \(\tau_c=\varepsilon/\sigma_0\gg\tau_m\)
时，才可忽略惯性项，回到上面的一阶指数衰减。整条导线的表面电荷建立还须等待电磁扰动传播，并可受
\(RC\)、\(L/R\) 和分布参数控制；因果传播不能瞬时完成。

\par\noindent\begin{minipage}{\linewidth}
\begin{center}
\includegraphics[width=\linewidth,height=.40\textheight,keepaspectratio]{\detokenize{assets/module-04/theory/unified7-relaxation-times.pdf}}
\par\smallskip{\small 图示：局域体电荷弛豫时间、动量弛豫时间与整体电路建立时间描述不同过程}
\end{center}
\end{minipage}\par

\subsubsection{四、三个容易混淆的时间尺度}\label{m4:ux56dbux4e09ux4e2aux5bb9ux6613ux6df7ux6dc6ux7684ux65f6ux95f4ux5c3aux5ea6}

\begin{itemize}
\tightlist
\item
  \(\tau_m\)：载流子平均定向动量的弛豫时间；
\item
  \(\tau_c=\varepsilon/\sigma\)：瞬时欧姆近似有效时，均匀介质局域自由体电荷的一阶衰减时间；需另行验证响应时间尺度；
\item
  \(t_{\rm circuit}\)：有限大小电路的整体场与表面电荷建立时间，受传播延迟、电感、电容和终端条件决定。
\end{itemize}

三者数值可相差很多，不能用其中一个代替另一个。

\subsection{三、导电介质中的绝缘空腔与附加电阻}\label{m4:ux4e09ux5bfcux7535ux4ecbux8d28ux4e2dux7684ux7eddux7f18ux7a7aux8154ux4e0eux9644ux52a0ux7535ux963b}

本附录处理稳恒电流场中的绝缘缺陷。在电导率为常量且无电流源的区域内，

\[
\mathbf J=\sigma\mathbf E,
\qquad
\mathbf E=-\nabla\varphi,
\qquad
\nabla\cdot\mathbf J=0,
\]

所以

\[
\nabla^2\varphi=0.
\]

绝缘边界满足

\[
\boxed{
\mathbf J\cdot\mathbf n=0
\quad\Longleftrightarrow\quad
\frac{\partial\varphi}{\partial n}=0
}.
\]

稳恒电流场与静电场可以具有相同的边值方程，但``电流源强度''与``电荷量''、``电导率''与``介电常数''只是一种数学对应，不能混作同一物理量。

\par\addvspace{.4\baselineskip}
\noindent\begin{minipage}{\linewidth}

\subsubsection{作业一：无限导电介质中的绝缘球形空腔}\label{m4:ux4f5cux4e1aux4e00ux65e0ux9650ux5bfcux7535ux4ecbux8d28ux4e2dux7684ux7eddux7f18ux7403ux5f62ux7a7aux8154}

\textbf{题号：M4-020}

\textbf{层级：} 复赛　　\textbf{方法：} 介质边界

均匀线性各向同性导电介质充满无穷空间，其电导率为
\(\sigma>0\)，其中挖去球心在原点、半径为 \(a>0\)
的绝缘球形空腔，腔内无传导电流。稳恒时远处场趋于
\(\mathbf E_0=E_0\mathbf e_z\)，\(E_0>0\)。腔内、外均按静电无旋场处理，不设理想偶极层；腔内势要求球心有限，外部附加势在远处衰减。采用以
\(+z\)
为极轴的球坐标，求空腔内外电场及介质中的电流线方程，并在球面检验法向电流和切向电场条件。

\end{minipage}\par

\par\noindent\begin{minipage}{\linewidth}
\textbf{解答}

\begin{center}
\includegraphics[width=\linewidth,height=.40\textheight,keepaspectratio]{\detokenize{assets/module-04/examples/M4-020.pdf}}
\par\smallskip{\small 图 M4-020：绝缘球腔：法向电流为零}
\end{center}
\end{minipage}\par

图 M4-020
的球面不允许传导电流穿过，远处箭头只表示匀强背景。球外势的径向导数在
\(r=a\)
为零，同时球面切向场要与腔内解衔接；两条件共同固定偶极修正项和腔内均匀场。

取球心为原点，以 \(\mathbf E_0\) 的方向为极轴。轴对称的 \(l=1\)
阶调和解可写成

\[
\varphi_{\rm out}
=-E_0r\cos\theta+\frac{A\cos\theta}{r^2},
\qquad
\varphi_{\rm in}=-BE_0r\cos\theta+C.
\]

第一式适用于 \(r>a\)，第二式适用于 \(r<a\)。外部解在无穷远处自动满足
\(\mathbf E\to\mathbf E_0\)；内部解舍去了在原点发散的项。

空腔不导电，故球面外侧必须满足

\[
j_r(a,\theta)=\sigma E_r(a,\theta)=0.
\]

由

\[
E_r=-\frac{\partial\varphi_{\rm out}}{\partial r}
=\left(E_0+\frac{2A}{r^3}\right)\cos\theta
\]

\par\noindent\begin{minipage}{\linewidth}
\begin{center}
\includegraphics[width=\linewidth,height=.40\textheight,keepaspectratio]{\detokenize{assets/module-04/process/M4-020a.pdf}}
\par\smallskip{\small 图 M4-020a：腔内有场，球面仍无传导通量}
\end{center}
\end{minipage}\par

图 M4-020a
的绿色箭头表示腔内电场，绝缘只禁止传导电流。外部解在球面没有径向电流，切向电场却可连续地进入腔内势；这两个不同约束分别用于确定外部修正和内部幅度。

得到

\[
A=-\frac12E_0a^3.
\]

电势在界面连续，由
\(\varphi_{\rm in}(a,\theta)=\varphi_{\rm out}(a,\theta)\) 得
\(B=3/2\)；常数 \(C\) 可取为零。因此

\[
\boxed{
\varphi_{\rm out}
=-E_0\left(r+\frac{a^3}{2r^2}\right)\cos\theta
},
\]

\[
\boxed{
\varphi_{\rm in}=-\frac32E_0r\cos\theta
}.
\]

空腔外的场分量为

\[
\boxed{
E_r=E_0\left(1-\frac{a^3}{r^3}\right)\cos\theta,
\qquad
E_{\theta}=-E_0\left(1+\frac{a^3}{2r^3}\right)\sin\theta
}.
\]

空腔内的电场是匀强场：

\[
\boxed{\mathbf E_{\rm in}=\frac32\mathbf E_0}.
\]

空腔内可以有电场，但因其电导率为零，仍有 \(\mathbf J_{\rm in}=0\)。在
\(r=a\) 处，外侧 \(E_r=0\)，而切向电场连续。

对轴对称电流场定义流函数 \(\Psi\)：

\[
j_r=\frac{1}{r^2\sin\theta}\frac{\partial\Psi}{\partial\theta},
\qquad
j_{\theta}=-\frac{1}{r\sin\theta}\frac{\partial\Psi}{\partial r}.
\]

代入 \(\mathbf J=\sigma\mathbf E\) 可得

\[
\boxed{
\Psi(r,\theta)
=\frac{\sigma E_0}{2}
\left(r^2-\frac{a^3}{r}\right)\sin^2\theta
}.
\]

故空腔外的电流线由 \(\Psi=C\) 给出，即

\[
\boxed{
\left(r^2-\frac{a^3}{r}\right)\sin^2\theta=C'
}.
\]

当 \(r=a\) 时 \(\Psi=0\)，说明球面本身是一条流面，电流不会穿入空腔。

\par\addvspace{.4\baselineskip}
\noindent\begin{minipage}{\linewidth}

\subsubsection{作业二：圆柱导体中微小绝缘空腔造成的附加电阻}\label{m4:ux4f5cux4e1aux4e8cux5706ux67f1ux5bfcux4f53ux4e2dux5faeux5c0fux7eddux7f18ux7a7aux8154ux9020ux6210ux7684ux9644ux52a0ux7535ux963b}

\textbf{题号：M4-021}

\textbf{层级：} 复赛　　\textbf{方法：} 介质边界、网络方程

一根均匀各向同性导电圆柱长 \(L\)、半径 \(b\)，电导率 \(\sigma>0\)，满足
\(L\gg b\)；两端由覆盖端面的电极通以固定总电流
\(I\)，侧壁绝缘。无缺陷时远离端部的轴向场为
\(E_0\)。在远离端面和侧壁的内部挖去半径 \(a\)
的绝缘球形空腔，\(a\ll b\)，且空腔到各边界的距离都远大于
\(a\)。在局部无限介质空腔近似下求相对于完整圆柱的附加电阻主阶
\(\Delta R\)，明确比较时保持总电流不变。

\end{minipage}\par

\par\noindent\begin{minipage}{\linewidth}
\textbf{解答}

\begin{center}
\includegraphics[width=\linewidth,height=.40\textheight,keepaspectratio]{\detokenize{assets/module-04/examples/M4-021.pdf}}
\par\smallskip{\small 图 M4-021：固定总电流下空腔增加端电压}
\end{center}
\end{minipage}\par

图 M4-021
将局部小球腔与远处圆柱边界分开，侧壁和端面都远离空腔。计算先用球腔的局部偶极扰动，再把平均电导的主阶改变换算为端口电阻；比较全过程保持同一个总电流。

无缺陷时

\[
S=\pi b^2,
\qquad
I=\sigma E_0S,
\qquad
R_0=\frac{L}{\sigma S}.
\]

取球腔中心为原点，\(\theta\)
从轴向场的正方向量起。球腔壁不允许法向电流通过，因此
\(\partial_r\varphi|_{r=a}=0\)。将远处均匀场势 \(-E_0r\cos\theta\) 与按
\(r^{-2}\) 衰减的轴对称附加势相加，得到外部附加势

\[
\delta\varphi=-\frac{E_0a^3}{2r^2}\cos\theta.
\]

它等价于方向与 \(\mathbf E_0\) 相反的电流偶极子，其偶极矩为

\[
\mathbf p_j=-2\pi\sigma a^3\mathbf E_0.
\]

把一个空腔置于体积 \(V=SL\) 的代表区域内，空腔体积分数为

\[
f=\frac{4\pi a^3}{3SL}\ll1.
\]

这里还需把球腔附近的局部场换成端口测得的电阻。令 \(E_m\)
为导电基体的平均轴向场，\(E_c\) 为球腔内轴向场。球面电势连续给出腔内势
\(-3E_mr\cos\theta/2\)，所以稀释近似下
\(E_c=3E_m/2\)。孔内虽有电场，却没有传导电流。

\par\noindent\begin{minipage}{\linewidth}
\begin{center}
\includegraphics[width=\linewidth,height=.40\textheight,keepaspectratio]{\detokenize{assets/module-04/process/M4-021b.pdf}}
\par\smallskip{\small 图 M4-021b：空腔的场与导体的平均场、平均电流}
\end{center}
\end{minipage}\par

图 M4-021b 把代表体积分成导电基体和绝缘空腔。分别作体积平均，在只保留
\(f\) 一阶的精度内， \[
\overline E=(1-f)E_m+fE_c=(1+f/2)E_m,\qquad
\overline j=(1-f)\sigma E_m.
\] 端口看到的等效电导率为
\(\overline j/\overline E\)，不能只把截面积机械地减去一个空腔截面。将比值
\(\sigma(1-f)/(1+f/2)\) 按 \(f\) 展开，得到

\[
\sigma_{\rm eff}
=\sigma\left[1-\frac32f+O(f^2)\right].
\]

因此

\[
R=\frac{L}{\sigma_{\rm eff}S}
=R_0\left(1+\frac32f\right)+O(f^2),
\]

从而

\[
\boxed{
\Delta R
=\frac{2\pi a^3}{\sigma S^2}
=\frac{2a^3}{\pi\sigma b^4}
}.
\]

相应的附加电压为

\[
\Delta U=I\Delta R=\frac{2E_0a^3}{b^2}.
\]

这个结果只保留关于空腔体积分数的一阶项。若空腔靠近侧壁、靠近端面，或
\(a/b\) 不再很小，就必须重新求有限圆柱中的边值问题。

\par\addvspace{.4\baselineskip}
\noindent\begin{minipage}{\linewidth}

\subsubsection{作业三：半无限大地中的漏电与跨步电压}\label{m4:ux4f5cux4e1aux4e09ux534aux65e0ux9650ux5927ux5730ux4e2dux7684ux6f0fux7535ux4e0eux8de8ux6b65ux7535ux538b}

\textbf{题号：M4-022}

\textbf{层级：} 复赛　　\textbf{方法：} 介质边界

大地占据一半空间，视为电导率 \(\sigma>0\)
的均匀各向同性导体，地面上方为空气绝缘区。地面一处持续向大地注入稳恒电流
\(I>0\)，回流电极足够远，可取无穷远电势为零。除小接触区外，地面法向电流为零。

\begin{enumerate}
\def\labelenumi{(\arabic{enumi})}
\item
  在距注入点的空间径向距离 \(r\)
  大于接触区尺度时，求大地内的电流密度、电场及电势。
\item
  人的两脚沿地面同一径向分别位于
  \(r,r+\ell\)，\(\ell>0\)，求两点电势差，即跨步电压。
\item
  将接触区具体化为埋入地下的等势半球电极，半径为
  \(a>0\)，求它相对于远处大地的接地电阻，忽略其他电极的扰动。
\end{enumerate}

\end{minipage}\par

\par\noindent\begin{minipage}{\linewidth}
\textbf{解答}

\begin{center}
\includegraphics[width=\linewidth,height=.40\textheight,keepaspectratio]{\detokenize{assets/module-04/examples/M4-022.pdf}}
\par\smallskip{\small 图 M4-022：半空间漏电与地面电势差}
\end{center}
\end{minipage}\par

图 M4-022 的电流只能向地下半空间扩散，因此半球面积为
\(2\pi r^2\)。用同一径向场从 \(r\) 积到 \(r+\ell\)
得跨步电压，从电极半径 \(a\) 积到无穷远得接地电阻。

电流在大地中作半球面对称扩散。穿过半径为 \(r\) 的半球面的总电流为

\[
I=j(r)\,2\pi r^2,
\]

故

\[
\boxed{
j(r)=\frac{I}{2\pi r^2},
\qquad
E(r)=\frac{I}{2\pi\sigma r^2}
}.
\]

取无穷远为零电势，有

\[
\varphi(r)=\int_r^{\infty} E(s)\,ds
=\frac{I}{2\pi\sigma r}.
\]

因此有限步长 \(\ell\) 对应的跨步电压为

\[
\boxed{
U_{\rm step}(r,\ell)
=\varphi(r)-\varphi(r+\ell)
=\frac{I}{2\pi\sigma}
\left(\frac1r-\frac1{r+\ell}\right)
}.
\]

当 \(\ell\ll r\) 时，

\[
U_{\rm step}\approx\frac{I\ell}{2\pi\sigma r^2}.
\]

半球电极到无穷远的电势差为

\[
U=\varphi(a)=\frac{I}{2\pi\sigma a},
\]

所以

\[
\boxed{R_{\rm g}=\frac{1}{2\pi\sigma a}}.
\]

实际接地体的有限尺寸、土壤分层和接触电阻会破坏理想半球对称性。

\par\addvspace{.4\baselineskip}
\noindent\begin{minipage}{\linewidth}

\subsubsection{作业四：两层漏电介质中的自由电荷与极化电荷}\label{m4:ux4f5cux4e1aux56dbux4e24ux5c42ux6f0fux7535ux4ecbux8d28ux4e2dux7684ux81eaux7531ux7535ux8377ux4e0eux6781ux5316ux7535ux8377}

\textbf{题号：M4-023}

\textbf{层级：} 复赛　　\textbf{方法：} 介质边界

两个同心球形理想电极的半径为 \(r_1,r_3\)，中间的 \(r=r_2\)
是两种漏电介质的分界面，\(0<r_1<r_2<r_3\)。内层 \(r_1<r<r_2\) 的参数为
\((\varepsilon_{r1},\sigma_1)\)，外层 \(r_2<r<r_3\) 为
\((\varepsilon_{r2},\sigma_2)\)；材料均匀、线性、各向同性，绝对电容率为
\(\varepsilon_0\varepsilon_{ri}\)，\(\sigma_i>0\)。内电极接地，外电极保持
\(U>0\)，中间界面没有外接电源或面电流源，但允许稳恒建立时自由电荷堆积。取径向向外为统一法向。

\begin{enumerate}
\def\labelenumi{(\arabic{enumi})}
\item
  求由外电极流向内电极的稳恒漏电电流 \(I\)。
\item
  求 \(r_1,r_2,r_3\) 处的自由面电荷密度及介质极化面电荷；在 \(r_2\)
  处计算两侧束缚面电荷的代数和。
\item
  求两层各自开区域内的自由体电荷和束缚体电荷密度，区分界面上的集中电荷。
\end{enumerate}

\end{minipage}\par

\par\noindent\begin{minipage}{\linewidth}
\textbf{解答}

\begin{center}
\includegraphics[width=\linewidth,height=.40\textheight,keepaspectratio]{\detokenize{assets/module-04/examples/M4-023.pdf}}
\par\smallskip{\small 图 M4-023：两层漏电介质中的径向电流}
\end{center}
\end{minipage}\par

图 M4-023 的 \(r_2\) 是两层介质界面，内外两个电极分别固定零势与
\(U\)。径向电流穿过每个球面的总量相同，但场需分别除以各层电导率；面电荷则按同一外径向比较两侧电位移。

取 \(\mathbf e_r\) 径向向外。由于电流由高电势外球流向内球，

\[
\mathbf J(r)=-\frac{I}{4\pi r^2}\mathbf e_r.
\]

两层介质内分别有

\[
\mathbf E_i(r)
=-\frac{I}{4\pi\sigma_i r^2}\mathbf e_r
\qquad(i=1,2).
\]

电势差为

\[
U=\frac{I}{4\pi}
\left[
\frac1{\sigma_1}\left(\frac1{r_1}-\frac1{r_2}\right)
+\frac1{\sigma_2}\left(\frac1{r_2}-\frac1{r_3}\right)
\right],
\]

所以

\[
\boxed{
I=
\frac{4\pi U}{
\dfrac1{\sigma_1}\left(\dfrac1{r_1}-\dfrac1{r_2}\right)
+\dfrac1{\sigma_2}\left(\dfrac1{r_2}-\dfrac1{r_3}\right)}
}.
\]

\par\noindent\begin{minipage}{\linewidth}
\begin{center}
\includegraphics[width=\linewidth,height=.40\textheight,keepaspectratio]{\detokenize{assets/module-04/process/M4-023a.pdf}}
\par\smallskip{\small 图 M4-023a：中间介质界面的三个场量分开比较}
\end{center}
\end{minipage}\par

图 M4-023a 局部放大 \(r_2\)
界面，法向仍由内层指向外层，而实际电流向内。电流连续与自由面电荷允许非零并不矛盾：前者比较J，后者比较D；束缚面电荷则以相反次序比较P。

两层介质中

\[
\mathbf D_i=\varepsilon_0\varepsilon_{ri}\mathbf E_i,
\qquad
\mathbf P_i=\varepsilon_0(\varepsilon_{ri}-1)\mathbf E_i.
\]

以下各界面的法向均取径向向外。由
\(\sigma_f=\mathbf n\cdot(\mathbf D_{\rm out}-\mathbf D_{\rm in})\)，

\[
\boxed{
\sigma_{f1}
=-\frac{\varepsilon_0\varepsilon_{r1}I}
{4\pi\sigma_1r_1^2}
},
\]

\[
\boxed{
\sigma_{f2}
=\frac{\varepsilon_0I}{4\pi r_2^2}
\left(
\frac{\varepsilon_{r1}}{\sigma_1}
-\frac{\varepsilon_{r2}}{\sigma_2}
\right)
},
\]

\[
\boxed{
\sigma_{f3}
=\frac{\varepsilon_0\varepsilon_{r2}I}
{4\pi\sigma_2r_3^2}
}.
\]

以下 \(\sigma_p\) 表示同一几何界面两侧极化面电荷的代数和。由
\(\sigma_p=-\mathbf n\cdot(\mathbf P_{\rm out}-\mathbf P_{\rm in})\)，

\[
\boxed{
\sigma_{p1}
=\frac{\varepsilon_0(\varepsilon_{r1}-1)I}
{4\pi\sigma_1r_1^2}
},
\]

\[
\boxed{
\sigma_{p2}
=\frac{\varepsilon_0I}{4\pi r_2^2}
\left[
\frac{\varepsilon_{r2}-1}{\sigma_2}
-\frac{\varepsilon_{r1}-1}{\sigma_1}
\right]
},
\]

\[
\boxed{
\sigma_{p3}
=-\frac{\varepsilon_0(\varepsilon_{r2}-1)I}
{4\pi\sigma_2r_3^2}
}.
\]

在每一均匀层内部，\(\varepsilon_{ri}/\sigma_i\) 为常量，而
\(\nabla\cdot\mathbf J=0\)，因此

\[
\rho_{fi}=\nabla\cdot\mathbf D_i=0,
\qquad
\rho_{pi}=-\nabla\cdot\mathbf P_i=0
\qquad(i=1,2).
\]

电荷只需分布在导体表面与两介质的分界面上。特别地，即使
\(\sigma_1\ne\sigma_2\)，若
\(\varepsilon_{r1}/\sigma_1=\varepsilon_{r2}/\sigma_2\)，仍有
\(\sigma_{f2}=0\)。

\par\addvspace{.4\baselineskip}
\noindent\begin{minipage}{\linewidth}

\subsubsection{作业五：二维电阻片的对偶关系}\label{m4:ux4f5cux4e1aux4e94ux4e8cux7ef4ux7535ux963bux7247ux7684ux5bf9ux5076ux5173ux7cfb}

\textbf{题号：M4-024}

\textbf{层级：} 复赛　　\textbf{方法：} 介质边界、网络方程

均匀各向同性薄电阻片具有方块电阻 \(R_{\square}>0\)。其平面区域由矩形
\(0<x<a,0<y<2a\) 在右边中点 \((a,a)\)
处挖去向左凹入的半圆缺口形成，缺口半径为
\(a/2\)，电阻片仍单连通。原测量中上、下边为理想等势电极，左边及完整右侧边界（包括缺口弧）绝缘，测得两电极电阻
\(R\)。现把边界角色完全互换：左边和完整右侧边界分别作为两个电极，原上、下边绝缘。求新电阻
\(R'\) 与 \(R,R_{\square}\) 的关系；不用求出具体缺口中的完整电势。

\end{minipage}\par

\par\noindent\begin{minipage}{\linewidth}
\textbf{解答}

\begin{center}
\includegraphics[width=\linewidth,height=.40\textheight,keepaspectratio]{\detokenize{assets/module-04/examples/M4-024.pdf}}
\par\smallskip{\small 图 M4-024：电极边与绝缘边完整互换}
\end{center}
\end{minipage}\par

图 M4-024
左图上、下边为电极，左边及包含缺口弧的完整右边界绝缘；右图将两类边界整体互换。关键是把电场和片电流方向一起旋转
\(90^{\circ}\)，并保持欧姆定律。

记片电流密度为 \(\mathbf K\)（单位A/m），则
\(\mathbf K=\mathbf E/R_{\square}\)。令 \(\mathsf J\) 表示在纸面内旋转
\(90^{\circ}\)，构造 \[
\mathbf E'=R_{\square}\mathsf J\mathbf K,\qquad
\mathbf K'=\mathsf J\mathbf E/R_{\square}.
\] 它们仍满足 \(\mathbf K'=\mathbf E'/R_{\square}\)。原来的
\(\nabla\cdot\mathbf K=0\) 与 \(\nabla\times\mathbf E=0\)
在旋转后分别保证 \(\nabla\times\mathbf E'=0\) 与
\(\nabla\cdot\mathbf K'=0\)，所以新场满足相同的稳恒导电方程。

原绝缘边无电流法向分量，旋转后新电场只有法向分量，该边成为等势电极；原电极电场只有法向分量，旋转后新电流沿边界走，该边成为绝缘边。缺口弧也必须整体互换。

设原来的电压、电流大小为
\(U,I\)。沿连接两条原绝缘边的截面积分，旋转后的电压等于原电流通量乘方块电阻；新电流同理由原电压给出：
\[
U'=R_{\square} I,\qquad I'=U/R_{\square}.
\] 因此 \[
\boxed{R'=\frac{U'}{I'}=\frac{R_{\square}^2}{R}},\qquad
\boxed{RR'=R_{\square}^2}.
\]
无缺口时，原矩形的长宽比为2，\(R=2R_{\square}\)、\(R'=R_{\square}/2\)，满足乘积关系。本解只用守恒、欧姆定律和边界方向。

\par\addvspace{.4\baselineskip}
\noindent\begin{minipage}{\linewidth}

\subsubsection{作业六：绝缘平面下的埋入球形电极}\label{m4:ux4f5cux4e1aux516dux7eddux7f18ux5e73ux9762ux4e0bux7684ux57cbux5165ux7403ux5f62ux7535ux6781}

\textbf{题号：M4-025}

\textbf{层级：} 复赛　　\textbf{方法：} 介质边界

电导率为 \(\sigma>0\) 的均匀各向同性导电介质占据半空间 \(z<0\)，平面
\(z=0\) 与绝缘空气相接。半径为 \(a>0\) 的小球形理想电极球心位于
\((0,0,-h)\)，\(a\ll h\)，向介质注入稳恒电流
\(I\)。回流电极在足够远处，取无穷远零势，平面上法向电流为零。求小电极电势的主阶及平面镜像造成的首个修正，并说明把小球近似为点电流源的精度范围。

\end{minipage}\par

\par\noindent\begin{minipage}{\linewidth}
\textbf{解答}

\begin{center}
\includegraphics[width=\linewidth,height=.40\textheight,keepaspectratio]{\detokenize{assets/module-04/examples/M4-025.pdf}}
\par\smallskip{\small 图 M4-025：同号电流像源满足绝缘平面}
\end{center}
\end{minipage}\par

图 M4-025 的蓝色源只用于真实半空间以外的延拓。同号配对使平面两侧势关于
\(z\) 偶对称，法向导数为零；求小电极电势时，自身 \(1/a\) 项与距镜像
\(2h\) 的有限修正都要保留。

令导电介质占据 \(z<0\)，其表面 \(z=0\) 绝缘，故

\[
\left.\frac{\partial\varphi}{\partial z}\right|_{z=0}=0.
\]

点电流源在均匀无限介质中的电势为
\(I/(4\pi\sigma s)\)。为满足上述诺伊曼边界条件，应在镜像位置放置同号点电流源。由于
\(a\ll h\)，球形电极表面的电势可展开为

\[
\boxed{
U\simeq\frac{I}{4\pi\sigma}
\left(\frac1a+\frac1{2h}\right)
}.
\]

第一项是孤立小球电极的自电势，第二项是绝缘平面的首个镜像修正。小球表面严格等势还要求加入更高阶多极修正；这里的近似只在
\(a/h\to0\) 时成立。若只取最低阶，才有
\(U\sim I/(4\pi\sigma a)\)；在比较不同埋深时不能先删去 \(1/(2h)\) 项。

\par\addvspace{.4\baselineskip}
\noindent\begin{minipage}{\linewidth}

\subsubsection{作业七：接地球面内的偏心小电极}\label{m4:ux4f5cux4e1aux4e03ux63a5ux5730ux7403ux9762ux5185ux7684ux504fux5fc3ux5c0fux7535ux6781}

\textbf{题号：M4-026}

\textbf{层级：} 决赛　　\textbf{方法：} 电像法、介质边界

半径为 \(R>0\) 的均匀各向同性导电介质球具有电导率
\(\sigma>0\)，整个外球面由理想电极覆盖并保持零势。内部点 \(A\) 距球心
\(O\) 为 \(d\)，\(0\le d<R\)。以 \(A\) 为球心的半径 \(a\)
小球形源电极向介质注入稳恒电流 \(I\)，满足
\(a\ll R-d\)。采用点源近似求球内远离源电极表面的电流密度，并求小电极相对于外球面的接地电阻至首个有限位置修正；检验同心与接近外壁但仍保持小电极近似的情形。

\end{minipage}\par

\par\noindent\begin{minipage}{\linewidth}
\textbf{解答}

\begin{center}
\includegraphics[width=\linewidth,height=.40\textheight,keepaspectratio]{\detokenize{assets/module-04/examples/M4-026.pdf}}
\par\smallskip{\small 图 M4-026：球外负像源固定接地边界}
\end{center}
\end{minipage}\par

图 M4-026 中真实源 \(A\) 在球内，负像源 \(A^{\prime}\)
位于球外射线上。球面等势固定像源位置和强度；小电极处的自势与像势相加，除以注入电流得到相对于外球面的电阻。

外球面满足狄里克雷边界条件 \(\varphi(R,\theta)=0\)。当 \(d>0\) 时，在
\(OA\) 射线上、球外

\[
OA'=\frac{R^2}{d}
\]

处放置像电流源

\[
I'=-\frac{R}{d}I.
\]

它不是实际存在的电极，只是构造球内解的数学工具。对球内任意场点 \(P\)，记

\[
\mathbf r_1=\overrightarrow{AP},
\qquad
\mathbf r_2=\overrightarrow{A'P}.
\]

电势与电流密度为

\[
\varphi(P)=\frac{I}{4\pi\sigma}
\left(
\frac1{r_1}-\frac{R}{d}\frac1{r_2}
\right),
\]

\[
\boxed{
\mathbf J(P)=\frac{I}{4\pi}
\left(
\frac{\mathbf r_1}{r_1^3}
-\frac{R}{d}\frac{\mathbf r_2}{r_2^3}
\right)
}.
\]

\par\noindent\begin{minipage}{\linewidth}
\begin{center}
\includegraphics[width=\linewidth,height=.40\textheight,keepaspectratio]{\detokenize{assets/module-04/process/M4-026a.pdf}}
\par\smallskip{\small 图 M4-026a：小电极电势由自势与像势相加}
\end{center}
\end{minipage}\par

图 M4-026a
把源电极的半径尺度a与球心距离d分开。自身项在小球表面读取，平滑像势在小球中心读取到所保留阶次；负像势降低相对于外接地球面的端口电压。

小电极足够小时，其表面电势等于自身的 \(1/a\) 项与像源在球心 \(A\)
处的势之和。由于

\[
AA'=\frac{R^2-d^2}{d},
\]

故

\[
U_A\simeq\frac{I}{4\pi\sigma}
\left(
\frac1a-\frac{R}{R^2-d^2}
\right)
.
\]

于是

\[
\boxed{
R_g\simeq\frac{U_A}{I}
=\frac{1}{4\pi\sigma}
\left(
\frac1a-\frac{R}{R^2-d^2}
\right)
}.
\]

当 \(d\to0\) 时，结果连续地化为

\[
R_g=\frac{1}{4\pi\sigma}
\left(\frac1a-\frac1R\right),
\]

这正是同心球形电极之间的电阻。若外球面绝缘，则应满足
\(\partial\varphi/\partial n=0\)，像源构造完全不同，也不存在``相对于绝缘球面的接地电阻''。

令外壁间隙 \(h=R-d\)，在 \(a\ll h\ll R\)
时，\(R/(R^2-d^2)\sim1/(2h)\)。因此有限位置修正趋于接地平面的负像项，而
\(1/a\) 自项仍占主导，近似电阻为正。不能把公式越出 \(a\ll h\)
的范围用到小电极触壁。

\subsection{四、使用条件}\label{m4:ux56dbux4f7fux7528ux6761ux4ef6}

\begin{enumerate}
\def\labelenumi{\arabic{enumi}.}
\tightlist
\item
  \(\mathbf J=\sigma\mathbf E\)
  是局域、线性、各向同性的本构关系；强场、非欧姆材料或明显非局域输运需要另建模型。
\item
  稳恒电流区满足
  \(\nabla\cdot\mathbf J=0\)，但电流源处应保留奇点通量，不能在含源区域直接写成处处为零。
\item
  绝缘边界给出法向导数条件，接地电极给出电势条件，两者不可混用。
\item
  镜像法的像源位于求解区域之外，只服务于边界条件，不代表真实电荷或真实电流源。
\item
  小球电极与微小空腔的结果都依赖尺度分离；使用近似式时应同时检查它到边界、端面和其他缺陷的距离。
\end{enumerate}

\section{专题三十一：基尔霍夫定律与电路方程（基础必讲与拔高选讲）}\label{m4:ux4e13ux9898ux4e09ux5341ux4e00ux57faux5c14ux970dux592bux5b9aux5f8bux4e0eux7535ux8defux65b9ux7a0bux57faux7840ux5fc5ux8bb2ux4e0eux62d4ux9ad8ux9009ux8bb2}

复杂电路不必预先化成串并联组合。固定节点、支路与参考方向后，将元件伏安关系同电荷守恒、法拉第定律联立，即可得到代数方程。以下限于集总参数电路；无限网络与连续导体中的传播问题还要处理收敛性和分布参数。

\subsection{一、电路图、参考方向与支路方程（基础必讲）}\label{m4:ux4e00ux7535ux8defux56feux53c2ux8003ux65b9ux5411ux4e0eux652fux8defux65b9ux7a0bux57faux7840ux5fc5ux8bb2}

把每一段两端接在同一对节点上的元件组合视为一条支路。对从节点 \(i\)
指向节点 \(j\) 的支路，规定：

\begin{itemize}
\tightlist
\item
  支路电流 \(I_{ij}\) 的参考方向为 \(i\to j\)；
\item
  支路电压 \(U_{ij}=V_i-V_j\)；
\item
  电阻支路满足
\end{itemize}

\[
\boxed{
I_{ij}=\frac{V_i-V_j}{R_{ij}}
}.
\]

计算得到 \(I_{ij}<0\)，只表示实际电流方向为 \(j\to i\)，不需要重列方程。

若沿支路从电源负极走向正极，理想电源带来电势升高
\(\mathcal E\)；反向穿过则带来电势降低
\(\mathcal E\)。整题必须始终沿用同一组参考方向。

\par\addvspace{.4\baselineskip}
\noindent\begin{minipage}{\linewidth}

\textbf{例 M4-027：负的参考电流表示实际反向}

节点 \(A,B\) 的电势相对于同一参考分别为
\(V_A=3.0\ \mathrm V,V_B=5.0\ \mathrm V\)，两点之间只接一只
\(R=2.0\ \Omega\) 线性电阻。把支路参考电流 \(I_{AB}\) 定为由 \(A\) 流向
\(B\)。求其有符号值与实际大小、方向，说明计算结果为负时是否需要改变已建立方程的参考方向。

\end{minipage}\par

\par\noindent\begin{minipage}{\linewidth}
\textbf{解答}

\begin{center}
\includegraphics[width=\linewidth,height=.40\textheight,keepaspectratio]{\detokenize{assets/module-04/examples/M4-027.pdf}}
\par\smallskip{\small 图 M4-027：负号改变实际方向，不改变参考}
\end{center}
\end{minipage}\par

图 M4-027 的蓝箭头规定参考方向，红箭头给实际方向。电阻电流仍用
\(V_A-V_B\)
除以阻值；所得负号说明实际沿反向流动，不应再把电势差次序重复反转。

若 \(V_A=3\ \mathrm V\)、\(V_B=5\ \mathrm V\)，且 \(A\) 到 \(B\) 间接有
\(2\ \Omega\) 电阻，则按 \(A\to B\) 取正有

\[
I_{AB}=\frac{3-5}{2}\ \mathrm A=-1\ \mathrm A.
\]

实际电流为 \(1\ \mathrm A\)，方向由 \(B\) 指向 \(A\)。

\subsection{二、基尔霍夫电流定律：节点处的电荷守恒（基础必讲）}\label{m4:ux4e8cux57faux5c14ux970dux592bux7535ux6d41ux5b9aux5f8bux8282ux70b9ux5904ux7684ux7535ux8377ux5b88ux6052ux57faux7840ux5fc5ux8bb2}

取一个包围节点的小区域，规定所有支路电流均以流出该区域为正。连续性方程给出一般关系

\[
\boxed{
\sum_k I_k^{\text{出}}=-\frac{dQ_N}{dt}
}.
\]

若节点不发生可观测的净电荷积累，或所有储能效应已经用电容等集总元件表示，则

\[
\boxed{
\sum_k I_k^{\text{出}}=0
}.
\]

这就是基尔霍夫电流定律。它不是``水流不会消失''的比喻，而是电荷守恒在集总节点上的直接结果。

\par\addvspace{.4\baselineskip}
\noindent\begin{minipage}{\linewidth}

\textbf{例 M4-028：一个节点的电流守恒}

某稳恒电路节点没有电荷积累，只有三条支路与其相连。已知第一支路向节点流入
\(3.0\ \mathrm A\)，第二支路从节点流出
\(1.2\ \mathrm A\)；将第三支路的未知电流 \(I\)
规定为流出节点为正。写出一致取向的节点电流方程，求 \(I\)
并判断实际方向。

\end{minipage}\par

\par\noindent\begin{minipage}{\linewidth}
\textbf{解答}

\begin{center}
\includegraphics[width=\linewidth,height=.40\textheight,keepaspectratio]{\detokenize{assets/module-04/examples/M4-028.pdf}}
\par\smallskip{\small 图 M4-028：节点守恒先统一流入与流出}
\end{center}
\end{minipage}\par

图 M4-028
中左支路流入、右侧两支路流出。按统一流出为正，已知流入量带负号；节点不积累电荷时三项相加为零，计算结果决定未知箭头是否就是实际方向。

某节点有 \(3.0\ \mathrm A\) 流入，另有 \(1.2\ \mathrm A\) 和 \(I\)
流出，则

\[
-3.0+1.2+I=0,
\]

故 \(I=1.8\ \mathrm A\)。

\subsection{三、基尔霍夫电压定律与感应电动势（基础必讲与拔高选讲）}\label{m4:ux4e09ux57faux5c14ux970dux592bux7535ux538bux5b9aux5f8bux4e0eux611fux5e94ux7535ux52a8ux52bfux57faux7840ux5fc5ux8bb2ux4e0eux62d4ux9ad8ux9009ux8bb2}

麦克斯韦方程给出

\[
\oint_C\mathbf E\cdot d\mathbf l
=-\frac{d\Phi_B}{dt}.
\]

若回路没有被不可忽略的时变外磁通穿过，右端为零，沿闭合路径的电势变化代数和为零：

\[
\boxed{
\sum_k \Delta V_k=0
}.
\]

若回路穿过时变磁通，把已建模元件沿绕行方向的有符号电压升记为
\(\Delta V_k\)，并把外磁通引起的感应电动势单独列出，则

\[
\boxed{\sum_k\Delta V_k+\mathcal E_{\rm ind}=0},
\qquad \mathcal E_{\rm ind}=-\frac{d\Phi_B}{dt}.
\]

例如顺电流穿过电阻取 \(\Delta V_R=-IR\)，从电源负极到正极取
\(\Delta V_{\rm src}=+\mathcal E\)，故单环为
\(-IR+\mathcal E+\mathcal E_{\rm ind}=0\)。也可改用正的电阻压降写
\(IR=\mathcal E+\mathcal E_{\rm ind}\)。这里的元件电压升是支路模型中的量，并非声称时变场中存在全局单值电势，能使其沿任意回路直接望远镜相消。感应项若已包含在支路模型中，不应再重复计入。

列回路方程时可采用统一规则：

\begin{enumerate}
\def\labelenumi{\arabic{enumi}.}
\tightlist
\item
  先选定绕行方向；
\item
  沿电流方向穿过电阻，记电势降 \(-IR\)；
\item
  逆电流方向穿过电阻，记 \(+IR\)；
\item
  从电源负极到正极，记 \(+\mathcal E\)；
\item
  把单列的感应电动势一并计入，令全部有符号升压与降压之和为零。
\end{enumerate}

\par\noindent\begin{minipage}{\linewidth}
\begin{center}
\includegraphics[width=\linewidth,height=.40\textheight,keepaspectratio]{\detokenize{assets/module-04/theory/review-tikz-m4-topic31-v9-from29-v8-01.pdf}}
\par\smallskip{\small 图示：节点电流定律来自电荷守恒，回路电压定律必须正确计入感应电动势}
\end{center}
\end{minipage}\par

\par\addvspace{.4\baselineskip}
\noindent\begin{minipage}{\linewidth}

\textbf{例 M4-029：单回路的电压平衡}

理想电源电动势为 \(\mathcal E>0\)，与两只正电阻 \(R_1,R_2\)
串联构成单回路，导线电阻忽略。电流正方向取电源向外送出电流的方向，沿此方向通过电源时从负极走向正极。先在无变化磁通时列回路方程并求电流；再设回路内另有同向感应电动势
\(\mathcal E_{\rm ind}\)，说明应如何计入方程及所得电流。

\end{minipage}\par

\par\noindent\begin{minipage}{\linewidth}
\textbf{解答}

\begin{center}
\includegraphics[width=\linewidth,height=.40\textheight,keepaspectratio]{\detokenize{assets/module-04/examples/M4-029.pdf}}
\par\smallskip{\small 图 M4-029：沿同一绕向记升压和降压}
\end{center}
\end{minipage}\par

图 M4-029
的两只串联电阻合写为同一电阻段，沿约定电流绕回路时电源升压、两电阻降压。同向感应电动势作为额外升压计入，不能再把含感应的回路误写成纯静电势变化为零。

电动势 \(\mathcal E\) 与两个串联电阻 \(R_1,R_2\)
构成单回路。沿电流方向绕行：

\[
\mathcal E-IR_1-IR_2=0,
\]

所以

\[
I=\frac{\mathcal E}{R_1+R_2}.
\]

若另有同向感应电动势 \(\mathcal E_{\rm ind}\)，只需把分子改为
\(\mathcal E+\mathcal E_{\rm ind}\)。

\subsection{四、独立方程数与网络拓扑（拔高选讲）}\label{m4:ux56dbux72ecux7acbux65b9ux7a0bux6570ux4e0eux7f51ux7edcux62d3ux6251ux62d4ux9ad8ux9009ux8bb2}

设一个连通网络有 \(n\) 个节点、\(b\) 条支路。

对所有节点的 KCL
方程求和，每条支路电流恰好出现一次正号和一次负号，故总和恒为零。于是独立的节点方程只有

\[
\boxed{n-1}
\]

个。

选取一棵包含全部节点但不含回路的生成树。生成树有 \(n-1\) 条树枝，其余

\[
b-(n-1)=b-n+1
\]

条支路称为连枝。每加入一条连枝，恰好生成一个独立基本回路，所以独立回路方程数为

\[
\boxed{b-n+1}.
\]

二者之和恰为支路数：

\[
(n-1)+(b-n+1)=b.
\]

\par\noindent\begin{minipage}{\linewidth}
\begin{center}
\includegraphics[width=\linewidth,height=.40\textheight,keepaspectratio]{\detokenize{assets/module-04/original/tikz-m4-topic31-v9-from29-v8-02.png}}
\par\smallskip{\small 图示：生成树确定独立节点方程与独立回路方程的数量}
\end{center}
\end{minipage}\par

\par\addvspace{.4\baselineskip}
\noindent\begin{minipage}{\linewidth}

\textbf{例 M4-030：独立电流方程与回路方程数}

一连通电路图含 \(n=5\) 个节点、\(b=7\)
条支路，所有真实连接点均计入节点，纯几何交叉不算节点。不计特殊元件参数造成的附加退化，只按图论结构计数。求独立基尔霍夫电流方程数及独立回路数，并用一棵生成树及剩余支路解释两种计数。

\end{minipage}\par

\par\noindent\begin{minipage}{\linewidth}
\textbf{解答}

\begin{center}
\includegraphics[width=\linewidth,height=.40\textheight,keepaspectratio]{\detokenize{assets/module-04/examples/M4-030.pdf}}
\par\smallskip{\small 图 M4-030：生成树与连支给出独立方程数}
\end{center}
\end{minipage}\par

图 M4-030
画出满足节点数和支路数的一种连通拓扑。选4条无闭环树边覆盖所有节点后，每添一条剩余支路便产生一个基本回路；这个计数只依赖连接，不依赖绘图外形。

一棵生成树覆盖5个节点，需要4条树边；它对应4个独立节点约束。余下
\(7-4=3\)
条连支各自与树上路径组成基本回路，因此独立回路数为3。总节点KCL相加恒等为零，所以不能把5行全当作独立方程。

若网络有 \(c\) 个互不连通的连通分量，上述数目分别推广为 \(n-c\) 和
\(b-n+c\)。在矩阵语言中，这对应关联矩阵的秩为 \(n-c\)。

\subsection{五、节点电压法（基础必讲）}\label{m4:ux4e94ux8282ux70b9ux7535ux538bux6cd5ux57faux7840ux5fc5ux8bb2}

节点电压法以节点电势为未知量。步骤如下：

\begin{enumerate}
\def\labelenumi{\arabic{enumi}.}
\tightlist
\item
  选择一个参考节点，令其电势为零；
\item
  给其余节点标记未知电势；
\item
  把每条电阻支路电流写成 \((V_i-V_j)/R_{ij}\)；
\item
  对每个非参考节点列 KCL；
\item
  解出节点电势，再由支路方程求电流。
\end{enumerate}

令 \(g_{ij}=1/R_{ij}\)，并把从外部注入节点 \(i\) 的电流记为
\(I_i^{\text{注}}\)，则一般节点方程为

\[
\boxed{
\sum_j g_{ij}(V_i-V_j)=I_i^{\text{注}}
}.
\]

写成矩阵形式：

\[
\boxed{
\mathbf G\mathbf V=\mathbf I
}.
\]

其中 \(\mathbf G\)
的对角元是与该节点相连的电导之和，非对角元是两节点间总电导的负值。若两个非参考节点之间直接跨接理想电压源，应把它们合并为超节点，并另加电压约束式。

\par\addvspace{.4\baselineskip}
\noindent\begin{minipage}{\linewidth}

\textbf{例 M4-031：含注入电流源的节点方程}

一个未知电势节点 \(V\) 通过 \(R_1>0\) 接到已知电势为 \(\mathcal E\)
的节点，通过 \(R_2>0\) 接地；另有理想电流源从地向该节点注入有符号电流
\(I_s\)。所有电势均以同一地节点为参考。选取流出未知节点的电阻电流为正，列节点电流方程并求
\(V\)，说明电流源注入量在方程中的符号。

\end{minipage}\par

\par\noindent\begin{minipage}{\linewidth}
\textbf{解答}

\begin{center}
\includegraphics[width=\linewidth,height=.40\textheight,keepaspectratio]{\detokenize{assets/module-04/examples/M4-031.pdf}}
\par\smallskip{\small 图 M4-031：注入源进入节点方程右端}
\end{center}
\end{minipage}\par

图 M4-031
的绿色箭头向未知节点注入电流，左右电阻的参考电流则从该节点流出。用这两类方向列KCL，注入量在右端为正；\(R_2\)
的右端是与电流源共用的地参考。

节点 \(V\) 通过 \(R_1\) 接到已知电势 \(\mathcal E\)，通过 \(R_2\)
接地，并有电流源 \(I_s\) 向该节点注入电流。节点方程为

\[
\frac{V-\mathcal E}{R_1}+\frac{V}{R_2}=I_s,
\]

故

\[
\boxed{
V=
\frac{\mathcal E/R_1+I_s}
{1/R_1+1/R_2}
}.
\]

\subsection{六、网孔电流法与一般回路法（基础必讲与拔高选讲）}\label{m4:ux516dux7f51ux5b54ux7535ux6d41ux6cd5ux4e0eux4e00ux822cux56deux8defux6cd5ux57faux7840ux5fc5ux8bb2ux4e0eux62d4ux9ad8ux9009ux8bb2}

网孔电流法只直接适用于平面网络。对每个不包含其他回路的网孔任取一个网孔电流；若两个网孔共享电阻，则共享支路中的实际电流是两个网孔电流的代数差。

以两个顺时针网孔电流 \(I_1,I_2\) 为例，若两网孔分别含独有电阻总和
\(R_a,R_b\)，共享电阻为 \(R_s\)，则方程为

\[
\boxed{
(R_a+R_s)I_1-R_sI_2=\mathcal E_1
},
\]

\[
\boxed{
-R_sI_1+(R_b+R_s)I_2=\mathcal E_2
}.
\]

对角元是本网孔全部电阻之和，非对角元是共享电阻之和的负值。非平面网络没有``每个面对应一个网孔''的简单结构，此时应选择生成树对应的基本回路，或直接使用节点电压法与支路电流法。

\par\noindent\begin{minipage}{\linewidth}
\begin{center}
\includegraphics[width=\linewidth,height=.40\textheight,keepaspectratio]{\detokenize{assets/module-04/original/tikz-m4-topic31-v9-from29-v8-03.png}}
\par\smallskip{\small 图示：复杂电阻网络可按节点数、平面网孔数和电源形式选择列式方法}
\end{center}
\end{minipage}\par

\subsection{七、线性网络的叠加原理（基础必讲）}\label{m4:ux4e03ux7ebfux6027ux7f51ux7edcux7684ux53e0ux52a0ux539fux7406ux57faux7840ux5fc5ux8bb2}

若所有元件都具有线性伏安关系，则不同独立激励分别产生的节点电势和支路电流可以线性叠加。

处理某一个独立源时：

\begin{itemize}
\tightlist
\item
  独立理想电压源置零后用短路代替；
\item
  独立理想电流源置零后用开路代替；
\item
  受控源不能随意置零，必须始终保留其控制关系。
\end{itemize}

\par\noindent\begin{minipage}{\linewidth}
\begin{center}
\includegraphics[width=\linewidth,height=.40\textheight,keepaspectratio]{\detokenize{assets/module-04/theory/review-tikz-m4-topic31-v9-from29-v8-04.pdf}}
\par\smallskip{\small 图示：线性电阻网络中不同激励状态的电势差和支路电流可以逐项相加}
\end{center}
\end{minipage}\par

功率是电压与电流的乘积，一般不满足叠加原理。即使 \(I=I_1+I_2\)，通常也有

\[
I^2R\ne I_1^2R+I_2^2R.
\]

\par\addvspace{.4\baselineskip}
\noindent\begin{minipage}{\linewidth}

\subsection{八、例题：一般惠斯通桥的桥臂电流}\label{m4:ux516bux4f8bux9898ux4e00ux822cux60e0ux65afux901aux6865ux7684ux6865ux81c2ux7535ux6d41}

\textbf{题号：M4-032}

理想直流电源把上端节点 \(P\) 保持在电势 \(\mathcal E\)，下端节点 \(Q\)
接地。左支路按 \(P-A-Q\) 依次为 \(R_1,R_2\)，右支路按 \(P-B-Q\) 依次为
\(R_3,R_4\)；中间桥臂 \(A-B\) 的电阻为
\(R_5\)，全部电阻为正，除此之外没有支路。取桥臂电流正向为 \(A\to B\)。求
\(I_5\)，保留其有符号值，并给出桥平衡时应满足的电阻关系。

\end{minipage}\par

\par\noindent\begin{minipage}{\linewidth}
\textbf{解答}

\begin{center}
\includegraphics[width=\linewidth,height=.40\textheight,keepaspectratio]{\detokenize{assets/module-04/examples/M4-032.pdf}}
\par\smallskip{\small 图 M4-032：一般电桥：保留桥臂的有符号电流}
\end{center}
\end{minipage}\par

图 M4-032 逐一标出上、下端 \(P,Q\) 和中点 \(A,B\)。同一桥臂电流在节点
\(A\) 的流出项为正，在节点 \(B\)
为负；先求两个中点电势，再按图示端点次序求 \(I_5\)。

令

\[
g_k=\frac1{R_k}\qquad(k=1,\dots,5).
\]

对节点 \(A,B\) 分别列 KCL：

\[
(g_1+g_2+g_5)V_A-g_5V_B=g_1\mathcal E,
\]

\[
-g_5V_A+(g_3+g_4+g_5)V_B=g_3\mathcal E.
\]

记

\[
D=(g_1+g_2+g_5)(g_3+g_4+g_5)-g_5^2.
\]

由二元一次方程组得

\[
V_A=\frac{\mathcal E[g_1(g_3+g_4+g_5)+g_5g_3]}{D},
\]

\[
V_B=\frac{\mathcal E[g_3(g_1+g_2+g_5)+g_5g_1]}{D}.
\]

两式相减后大量交叉项抵消：

\[
V_A-V_B
=\frac{\mathcal E(g_1g_4-g_2g_3)}{D}.
\]

因此

\[
\boxed{
I_5
=g_5(V_A-V_B)
=\frac{\mathcal E g_5(g_1g_4-g_2g_3)}
{(g_1+g_2+g_5)(g_3+g_4+g_5)-g_5^2}
}.
\]

若改写成电阻形式，

\[
\boxed{
I_5=
\frac{\mathcal E(R_2R_3-R_1R_4)}
{
R_1R_2R_3+R_1R_2R_4
+R_1R_3R_4+R_2R_3R_4
+R_5(R_1R_3+R_1R_4+R_2R_3+R_2R_4)
}
}.
\]

桥路平衡条件为

\[
\boxed{
R_1R_4=R_2R_3
\quad\Longleftrightarrow\quad
I_5=0
}.
\]

这个条件与 \(R_5\) 的取值无关。若分子为负，说明实际桥臂电流由 \(B\) 流向
\(A\)。

\subsection{九、练习与解答}\label{m4:ux4e5dux7ec3ux4e60ux4e0eux89e3ux7b54}

\par\addvspace{.4\baselineskip}
\noindent\begin{minipage}{\linewidth}

\subsubsection{作业一：外方形与中点短接网络}\label{m4:ux4f5cux4e1aux4e00ux5916ux65b9ux5f62ux4e0eux4e2dux70b9ux77edux63a5ux7f51ux7edc}

\textbf{题号：M4-033}

\textbf{层级：} 基础　　\textbf{方法：} 网络方程

平面网络有九个节点：外正方形四个角 \(C_1,C_2,C_3,C_4\)，依次对应的边中点
\(M_1,M_2,M_3,M_4\)，以及中心 \(O\)（下标按模4计）。每条外边分为两段
\(C_iM_i,M_iC_{i+1}\)；中心分别与四角、四边中点相连；相邻中点
\(M_iM_{i+1}\) 也相连。以上共20条支路各接一只 \(r>0\)
电阻，除九个指定节点外的交叉不相连。现用理想导线把四角短接成端子
\(A\)，四边中点短接成端子 \(B\)；电动势 \(\mathcal E\)、内阻同为 \(r\)
的实际电源接在 \(A,B\)
之间。求外网络等效电阻及电源干路电流，说明中心和内方形支路如何处理。

\end{minipage}\par

\par\noindent\begin{minipage}{\linewidth}
\textbf{解答}

\begin{center}
\includegraphics[width=\linewidth,height=.40\textheight,keepaspectratio]{\detokenize{assets/module-04/examples/M4-033.pdf}}
\par\smallskip{\small 图 M4-033：九个真实节点经短接形成三组}
\end{center}
\end{minipage}\par

图 M4-033 只把黑点标作节点，四角同属短接组 \(A\)、四边中点同属组
\(B\)。外边框的8只电阻直接跨接两组；经过中心的另两组4只并联电阻先后串联，内方形两端同属
\(B\) 而无电流。

令 \(V_A=\varphi\)、\(V_B=0\)。中心节点 \(O\) 通过四个相同电阻接到
\(A\)，又通过四个相同电阻接到 \(B\)。由 KCL：

\[
4\frac{\varphi-V_O}{r}
+4\frac{0-V_O}{r}=0,
\]

故

\[
V_O=\frac{\varphi}{2}.
\]

\par\noindent\begin{minipage}{\linewidth}
\begin{center}
\includegraphics[width=\linewidth,height=.40\textheight,keepaspectratio]{\detokenize{assets/module-04/process/M4-033a.pdf}}
\par\smallskip{\small 图 M4-033a：短接后的三节点等效结构}
\end{center}
\end{minipage}\par

图 M4-033a
展开两类真实通路：上方保留外边框8条并联边，下方经过中心O。原内方形两端同属B，不参与端口电流；电源内阻尚未包含在这里的外网络电阻中。

内正方形四个顶点都属于节点集 \(B\)，其四条边两端等势，没有电流。

从 \(A\) 到 \(B\) 有两类通路：

\begin{itemize}
\tightlist
\item
  外边框给出 \(8\) 条并联电阻，电导为 \(8/r\)；
\item
  \(A\to O\) 的四个电阻并联为 \(r/4\)，\(O\to B\) 的四个电阻也并联为
  \(r/4\)，二者串联为 \(r/2\)，电导为 \(2/r\)。
\end{itemize}

故外网络总电导与等效电阻为

\[
G_{AB}=\frac8r+\frac2r=\frac{10}{r},
\]

\[
\boxed{R_{AB}=\frac{r}{10}}.
\]

再与电源内阻 \(r\) 串联：

\[
\boxed{
I=\frac{\mathcal E}{r+r/10}
=\frac{10\mathcal E}{11r}
}.
\]

路端电压为

\[
\varphi=IR_{AB}=\frac{\mathcal E}{11}.
\]

\subsection{十、练习与解答}\label{m4:ux5341ux7ec3ux4e60ux4e0eux89e3ux7b54}

\par\addvspace{.4\baselineskip}
\noindent\begin{minipage}{\linewidth}

\subsubsection{作业二：正七边形星形网络的支路电流}\label{m4:ux4f5cux4e1aux4e8cux6b63ux4e03ux8fb9ux5f62ux661fux5f62ux7f51ux7edcux7684ux652fux8defux7535ux6d41}

\textbf{题号：M4-034}

\textbf{层级：} 基础　　\textbf{方法：} 等效变换、网络方程

七个节点 \(A,B,C,D,E,F,G\)
按顺序排成正七边形，下标按模7计。每点分别连接到相距1及相距3的节点，共14条不同无向支路，各接一只
\(r>0\) 电阻；所有无节点标记的线条交叉均不连接。外部在 \(B,C,D\)
分别注入 \(I,2I,3I\)，在 \(G,F,E\) 分别抽出 \(I,2I,3I\)，\(I>0\)，节点
\(A\) 无注入。求真实支路 \(D-E\) 上从 \(D\) 指向 \(E\)
的有符号电流，要求同时利用连接对称性与激励的反对称性。

\end{minipage}\par

\par\noindent\begin{minipage}{\linewidth}
\textbf{解答}

\begin{center}
\includegraphics[width=\linewidth,height=.40\textheight,keepaspectratio]{\detokenize{assets/module-04/examples/M4-034.pdf}}
\par\smallskip{\small 图 M4-034：七边形网络：镜像电势反号}
\end{center}
\end{minipage}\par

图 M4-034 保留7个节点、14条支路，红边为所求 \(D-E\)。通过 \(A\) 和边
\(D-E\)
中点的对称轴交换源汇分布的正负号。镜像外部注入符号相反，因此成对电势反号；红边两端势差应为
\(2V_D\)，再除以该边真实电阻。

网络与外接电流关于通过 \(A\) 的对称轴反对称。由线性方程解的唯一性，可取

\[
V_A=0,
\]

\[
V_B=\varphi_1,\quad
V_C=\varphi_2,\quad
V_D=\varphi_3,
\]

\[
V_G=-\varphi_1,\quad
V_F=-\varphi_2,\quad
V_E=-\varphi_3.
\]

节点 \(B\) 的邻点为 \(A,C,E,F\)，故

\[
I=
\frac{\varphi_1}{r}
+\frac{\varphi_1-\varphi_2}{r}
+\frac{\varphi_1+\varphi_3}{r}
+\frac{\varphi_1+\varphi_2}{r},
\]

即

\[
4\varphi_1+\varphi_3=Ir.
\]

节点 \(C\) 的邻点为 \(B,D,F,G\)，故

\[
5\varphi_2-\varphi_3=2Ir.
\]

节点 \(D\) 的邻点为 \(C,E,G,A\)，故

\[
\varphi_1-\varphi_2+5\varphi_3=3Ir.
\]

联立得到

\[
\varphi_1=\frac{Ir}{13},
\qquad
\varphi_2=\frac{7Ir}{13},
\qquad
\varphi_3=\frac{9Ir}{13}.
\]

因此

\[
I_{DE}
=\frac{V_D-V_E}{r}
=\frac{2\varphi_3}{r},
\]

所以

\[
\boxed{
I_{DE}=\frac{18}{13}I
}.
\]

所得 \(18I/13\) 小于节点 \(D\) 的外部注入 \(3I\)。按 \(D\)
的四条真实邻边求流出电流，分别为 \(2I/13,18I/13,10I/13,9I/13\)，代数和为
\(3I\)，符合 KCL。

\subsection{十一、练习与解答}\label{m4:ux5341ux4e00ux7ec3ux4e60ux4e0eux89e3ux7b54}

\par\addvspace{.4\baselineskip}
\noindent\begin{minipage}{\linewidth}

\subsubsection{作业三：多个实际电压源并联与密耳曼定理}\label{m4:ux4f5cux4e1aux4e09ux591aux4e2aux5b9eux9645ux7535ux538bux6e90ux5e76ux8054ux4e0eux5bc6ux8033ux66fcux5b9aux7406}

\textbf{题号：M4-035}

\textbf{层级：} 基础　　\textbf{方法：} 网络方程

\(n\ge2\) 个实际电压源接在同一对端口 \(A,B\) 之间，外端暂时开路。第
\(i\) 支路由理想电动势 \(\mathcal E_i\) 与内阻 \(r_i>0\)
串联构成，\(\mathcal E_i\) 定义为该支路从 \(B\) 向 \(A\)
的电势升高，可正可负；不同支路没有额外连接。求以 \(U_{AB}=V_A-V_B\)
定向的开路等效电动势及等效内阻，说明外端开路是否意味着每个源支路都无电流。

\end{minipage}\par

\par\noindent\begin{minipage}{\linewidth}
\textbf{解答}

\begin{center}
\includegraphics[width=\linewidth,height=.40\textheight,keepaspectratio]{\detokenize{assets/module-04/examples/M4-035.pdf}}
\par\smallskip{\small 图 M4-035：并联源支路共用同一端口电压}
\end{center}
\end{minipage}\par

图 M4-035 示意三条源支路，公式适用于任意 \(n\)。所有支路承受同一
\(U=V_A-V_B\)，端口开路只要求各支路向端口的电流代数和为零；独立源置零后，各内阻才成为直接并联。

取下端为零电势，上端开路电势为 \(V\)。规定第 \(i\)
支路电流由下端流向上端为正，则

\[
I_i=\frac{\mathcal E_i-V}{r_i}.
\]

端口开路时没有外部电流，KCL 给出

\[
\sum_{i=1}^n I_i=0.
\]

于是

\[
\sum_{i=1}^n\frac{\mathcal E_i-V}{r_i}=0,
\]

从而

\[
\boxed{
\mathcal E_{\rm eq}=V
=\frac{\displaystyle\sum_{i=1}^n\mathcal E_i/r_i}
{\displaystyle\sum_{i=1}^n1/r_i}
}.
\]

求等效内阻时，把所有独立理想电压源置零，即用短路代替。各内阻并联，所以

\[
\boxed{
r_{\rm eq}
=\left(\sum_{i=1}^n\frac1{r_i}\right)^{-1}
}.
\]

\(\mathcal E_i\) 必须按统一端口方向取代数值。若某个电源极性相反，其
\(\mathcal E_i\) 应取负。不同电动势的理想电压源不能直接并联；当
\(r_i\to0\) 且电动势不一致时，模型预言环流发散，说明理想模型本身不相容。

\subsection{十二、列式方法与检查}\label{m4:ux5341ux4e8cux5217ux5f0fux65b9ux6cd5ux4e0eux68c0ux67e5}

\begin{enumerate}
\def\labelenumi{\arabic{enumi}.}
\tightlist
\item
  \textbf{先认拓扑，再列方程。}
  理想导线连接的点属于同一节点；交叉但没有连接点标记的导线不一定相连。
\item
  \textbf{节点少时优先节点电压法。} 它不要求网络平面，也便于处理电流源。
\item
  \textbf{平面网孔少时可用网孔电流法。}
  非平面网络应改用基本回路或节点法。
\item
  \textbf{KCL 的一般式允许节点储存电荷。} 使用 \(\sum I=0\)
  时，应确认储能效应已由支路元件表示。
\item
  \textbf{KVL 必须计入感应电动势。} 有时变磁通时，不能把
  \(\oint\mathbf E\cdot d\mathbf l\) 擅自写成零。
\item
  \textbf{功率检验。} 采用关联参考方向时，应满足
\end{enumerate}

\[
\boxed{
\sum_{\rm all\ branches} U_kI_k=0
}.
\]

电阻吸收功率 \(I^2R\ge0\)；其总和应等于各电源净输出功率。 7.
\textbf{极限检验。} 令某电阻趋于 \(0\) 或
\(\infty\)，结果应分别回到短路或断路情形；对称电路还应满足相应的等势与等流关系。

\par\noindent\begin{minipage}{\linewidth}
\begin{center}
\includegraphics[width=\linewidth,height=.40\textheight,keepaspectratio]{\detokenize{assets/module-04/original/tikz-m4-topic31-v9-from29-v8-09.png}}
\par\smallskip{\small 图示：复杂电路解答应依次检查拓扑、符号、守恒、功率和极限}
\end{center}
\end{minipage}\par

\section{专题三十二：电学仪表、测量电路与补偿法（基础必讲与拔高选讲）}\label{m4:ux4e13ux9898ux4e09ux5341ux4e8cux7535ux5b66ux4eeaux8868ux6d4bux91cfux7535ux8defux4e0eux8865ux507fux6cd5ux57faux7840ux5fc5ux8bb2ux4e0eux62d4ux9ad8ux9009ux8bb2}

仪表是接入原电路、具有内阻和量程的真实元件。建立测量模型时，要说明被测量、仪表对原电路的影响，以及零示法或等效关系怎样减小这种影响。

\subsection{一、表头与电流表、电压表改装}\label{m4:ux4e00ux8868ux5934ux4e0eux7535ux6d41ux8868ux7535ux538bux8868ux6539ux88c5}

设磁电式表头满偏电流为 \(I_g\)、内阻为 \(R_g\)，满偏电压为
\(U_g=I_gR_g\)。

磁电式表头本身响应电流的平均方向，以下改装公式直接适用于直流测量；交流测量还需整流或采用其他测量原理，不能把交变电流的正负半周直接送入表头后仍按同一刻度解释。

\par\noindent\begin{minipage}{\linewidth}
\begin{center}
\includegraphics[width=\linewidth,height=.40\textheight,keepaspectratio]{\detokenize{assets/module-04/theory/tikz-m4-topic32-v9-01-meter-conversion.pdf}}
\par\smallskip{\small 图示：表头并联分流电阻成为电流表，串联分压电阻成为电压表}
\end{center}
\end{minipage}\par

改装成量程 \(I\) 的电流表时，并联分流电阻 \(R_s\)，满偏时

\[
I_gR_g=(I-I_g)R_s,
\]

所以

\[
\boxed{R_s=\frac{I_gR_g}{I-I_g}}.
\]

改装成量程 \(U\) 的电压表时，串联分压电阻 \(R_m\)，

\[
U=I_g(R_g+R_m),
\]

故

\[
\boxed{R_m=\frac{U}{I_g}-R_g}.
\]

\par\addvspace{.4\baselineskip}
\noindent\begin{minipage}{\linewidth}

\textbf{例 M4-036：表头改装时的量程倍数}

磁电式表头的内阻为 \(R_g>0\)、满偏电流为
\(I_g>0\)，指针偏角与通过表头的电流成正比，允许加接理想线性电阻。对给定倍数
\(n>1\)，分别改装成满量程为 \(nI_g\) 的电流表，以及满量程为原满偏电压
\(I_gR_g\) 的 \(n\)
倍的电压表。画清分流或分压接法，求需外接的电阻，注意两种改装的连接方式不同。

\end{minipage}\par

\par\noindent\begin{minipage}{\linewidth}
\textbf{解答}

\begin{center}
\includegraphics[width=\linewidth,height=.40\textheight,keepaspectratio]{\detokenize{assets/module-04/examples/M4-036.pdf}}
\par\smallskip{\small 图 M4-036：量程扩大需要两种不同接法}
\end{center}
\end{minipage}\par

图 M4-036
的上支路是表头，下支路是电流表分流电阻，两端电压相同。改为电压表时应将附加电阻串入表头支路，不能保留图中的并联拓扑；两种公式分别来自分流和分压。

电流表在满量程时让表头仍流 \(I_g\)，分流支路流
\((n-1)I_g\)，两支路电压相等给 \(R_s=R_g/(n-1)\)。电压表仍只有表头电流
\(I_g\)，串联电阻须承担 \((n-1)I_gR_g\)，故 \(R_v=(n-1)R_g\)。

\par\noindent\begin{minipage}{\linewidth}
\begin{center}
\includegraphics[width=\linewidth,height=.40\textheight,keepaspectratio]{\detokenize{assets/module-04/process/M4-036a.pdf}}
\par\smallskip{\small 图 M4-036a：电压表串联分压的另一种拓扑}
\end{center}
\end{minipage}\par

图 M4-036a
只画电压表改装：表头与附加电阻承受同一满偏电流。和主图的并联分流比较可见，量程倍数相同也不能互换两种连接方式。

\subsection{二、仪表内阻与负载效应}\label{m4:ux4e8cux4eeaux8868ux5185ux963bux4e0eux8d1fux8f7dux6548ux5e94}

理想电流表内阻为零，理想电压表内阻为无穷大。实际电流表内阻 \(R_A>0\)
会增加所在支路电阻，实际电压表内阻 \(R_V<\infty\) 会从被测节点分流。

\par\noindent\begin{minipage}{\linewidth}
\begin{center}
\includegraphics[width=\linewidth,height=.40\textheight,keepaspectratio]{\detokenize{assets/module-04/theory/tikz-m4-topic32-v9-02-loading-effect.pdf}}
\par\smallskip{\small 图示：电流表串入产生额外压降，电压表并入产生分流}
\end{center}
\end{minipage}\par

仪表的``灵敏''与``扰动小''不是同一个概念。电压表用于高阻电路时尤其需要更大的
\(R_V\)；电流表用于低阻支路时尤其需要更小的 \(R_A\)。

\subsection{三、伏安法测电阻的两种接法}\label{m4:ux4e09ux4f0fux5b89ux6cd5ux6d4bux7535ux963bux7684ux4e24ux79cdux63a5ux6cd5}

\par\noindent\begin{minipage}{\linewidth}
\begin{center}
\includegraphics[width=\linewidth,height=.40\textheight,keepaspectratio]{\detokenize{assets/module-04/theory/tikz-m4-topic32-v9-03-voltage-current-method.pdf}}
\par\smallskip{\small 图示：伏安法两种接法分别包含电流表压降或电压表分流}
\end{center}
\end{minipage}\par

若电压表跨接在``待测电阻与电流表串联整体''两端，电流表读数就是待测支路电流，但电压表还测到电流表压降：

\[
R_{\mathrm{\text{测}}}=\frac UI=R_x+R_A.
\]

该接法适合 \(R_x\gg R_A\) 的较大电阻。

若电压表只跨接在待测电阻两端，电流表读数包含待测电阻电流和电压表电流：

\[
R_{\mathrm{\text{测}}}
=\frac U{I_x+I_V}
=\frac{R_xR_V}{R_x+R_V}.
\]

该接法适合 \(R_x\ll R_V\)
的较小电阻。不能只背``内接、外接''名称，必须看仪表实际包围了哪些元件。

\subsection{四、串联式欧姆表与非线性刻度}\label{m4:ux56dbux4e32ux8054ux5f0fux6b27ux59c6ux8868ux4e0eux975eux7ebfux6027ux523bux5ea6}

串联式欧姆表内部含电源、表头和可调电阻。把内部总串联电阻记为
\(R_0\)，外接未知电阻为 \(R_x\)，则

\[
\boxed{I=\frac{\mathcal E}{R_0+R_x}}.
\]

\par\noindent\begin{minipage}{\linewidth}
\begin{center}
\includegraphics[width=\linewidth,height=.40\textheight,keepaspectratio]{\detokenize{assets/module-04/theory/tikz-m4-topic32-v9-04-ohmmeter.pdf}}
\par\smallskip{\small 图示：串联式欧姆表的内部电源、调零电阻和反向非线性刻度}
\end{center}
\end{minipage}\par

两表笔短接时 \(R_x=0\)，调节 \(R_0\) 使表头满偏，此处标为零欧姆；开路时
\(I=0\)，标为无穷大。由于 \(I\) 与 \(R_x\)
不是线性关系，欧姆刻度不均匀，且电池电动势变化后必须重新调零。

\par\addvspace{.4\baselineskip}
\noindent\begin{minipage}{\linewidth}

\textbf{例 M4-037：欧姆表的半偏位置}

一只串联式欧姆表由理想电池、表头及内部调节电阻串联组成，总内部电阻为
\(R_0=1.50\ \mathrm{k}\Omega\)。表笔短接时已调至满偏；测量时在表笔之间接入未知电阻
\(R_x\ge0\)，指针偏角与电流成正比。求指针降至半满偏时的
\(R_x\)，解释为什么线性的电流刻度不能直接当作线性的电阻刻度。

\end{minipage}\par

\par\noindent\begin{minipage}{\linewidth}
\textbf{解答}

\begin{center}
\includegraphics[width=\linewidth,height=.40\textheight,keepaspectratio]{\detokenize{assets/module-04/examples/M4-037.pdf}}
\par\smallskip{\small 图 M4-037：欧姆刻度来自电流的反比例关系}
\end{center}
\end{minipage}\par

图 M4-037 把内部总阻 \(R_0\) 与表笔间待测阻 \(R_x\)
串联。短接时的满偏电流提供归一化基准；指针半偏意味着电流减半，故总电阻加倍，而不是把电阻数值直接减半。

设电池电动势为 \(\mathcal E\)，短接满偏时 \(I_g=\mathcal E/R_0\)。接入
\(R_x\) 后 \(I=\mathcal E/(R_0+R_x)\)，所以
\(I/I_g=R_0/(R_0+R_x)\)。令比值为1/2得
\(R_x=R_0=1.50\ \mathrm{k}\Omega\)；反比例关系导致电阻刻度不均匀。

\subsection{五、惠斯通电桥与零示法}\label{m4:ux4e94ux60e0ux65afux901aux7535ux6865ux4e0eux96f6ux793aux6cd5}

设电桥四臂依次为
\(R_1,R_2,R_3,R_x\)，电源接上下端，检流计接左右中点。平衡时检流计电流为零，两中点等势，因此

\[
\frac{R_1}{R_2}=\frac{R_3}{R_x},
\qquad
\boxed{R_x=\frac{R_2R_3}{R_1}}.
\]

\par\noindent\begin{minipage}{\linewidth}
\begin{center}
\includegraphics[width=\linewidth,height=.40\textheight,keepaspectratio]{\detokenize{assets/module-04/theory/tikz-m4-topic32-v9-05-wheatstone.pdf}}
\par\smallskip{\small 图示：惠斯通电桥平衡时两中点等势且检流计无电流}
\end{center}
\end{minipage}\par

零示法的优势是平衡时检流计不从被测支路持续取电；检流计内阻影响接近平衡时的灵敏度，却不进入理想平衡条件。滑线电桥把均匀电阻丝两段电阻比换成长度比。

\subsection{六、电势差计与补偿测量}\label{m4:ux516dux7535ux52bfux5deeux8ba1ux4e0eux8865ux507fux6d4bux91cf}

均匀电阻丝中维持稳定工作电流后，沿线电势降与长度成正比：

\[
U=k\ell.
\]

把待测电源与一段电势降反向串联，调到检流计示零时，待测支路无电流，故测得的是电动势而不是带负载端电压：

\[
\boxed{\mathcal E_x=k\ell_x}.
\]

\par\noindent\begin{minipage}{\linewidth}
\begin{center}
\includegraphics[width=\linewidth,height=.40\textheight,keepaspectratio]{\detokenize{assets/module-04/theory/tikz-m4-topic32-v9-06-potentiometer.pdf}}
\par\smallskip{\small 图示：电势差计用反向电势降补偿待测电动势并在平衡时零取流}
\end{center}
\end{minipage}\par

用标准电池和未知电池在同一工作电流下比较，

\[
\boxed{\frac{\mathcal E_x}{\mathcal E_s}=\frac{\ell_x}{\ell_s}}.
\]

若工作电流在两次测量之间变化，该比值关系失效。

\subsection{七、端口测量与黑箱辨识}\label{m4:ux4e03ux7aefux53e3ux6d4bux91cfux4e0eux9ed1ux7bb1ux8fa8ux8bc6}

两端线性有源黑箱可写成戴维南关系

\[
U=\mathcal E-Ir.
\]

只测一次端电压不能同时确定
\(\mathcal E,r\)；至少需要两个不同负载工作点，或一次开路电压加一次已知负载电压。测量方案的本质是让未知参数数目不超过独立测量方程数。

\par\noindent\begin{minipage}{\linewidth}
\begin{center}
\includegraphics[width=\linewidth,height=.40\textheight,keepaspectratio]{\detokenize{assets/module-04/theory/tikz-m4-topic32-v9-07-black-box.pdf}}
\par\smallskip{\small 图示：用开路和已知负载两个端口状态辨识黑箱电动势与内阻}
\end{center}
\end{minipage}\par

\subsection{八、使用条件与常见错误}\label{m4:ux516bux4f7fux7528ux6761ux4ef6ux4e0eux5e38ux89c1ux9519ux8bef}

\begin{enumerate}
\def\labelenumi{\arabic{enumi}.}
\tightlist
\item
  改装后的量程公式建立在表头与附加电阻同时达到相应电压或电流的条件上。
\item
  伏安法必须把仪表内阻画进等效电路，再判断系统误差方向。
\item
  欧姆表必须断开被测电路自身电源，且刻度方向与电流刻度相反。
\item
  电桥和电势差计的精确关系只在示零时成立；未平衡时不能把检流计支路删去。
\item
  电势差计工作电流必须稳定，待测电动势还必须处于可补偿量程内。
\end{enumerate}

\subsection{九、练习与解答}\label{m4:ux4e5dux7ec3ux4e60ux4e0eux89e3ux7b54-1}

\par\addvspace{.4\baselineskip}
\noindent\begin{minipage}{\linewidth}

\subsubsection{作业一：表头改装为双量程仪表}\label{m4:ux4f5cux4e1aux4e00ux8868ux5934ux6539ux88c5ux4e3aux53ccux91cfux7a0bux4eeaux8868}

\textbf{题号：M4-038}

\textbf{层级：} 基础　　\textbf{方法：} 电学测量、分流与分压

两种改装分别使用相同规格表头独立完成，表头线性、所加电阻理想，导线电阻忽略；电流表采用并联分流，电压表采用串联分压。

表头满偏电流 \(I_g=1.00\ \mathrm{mA}\)、内阻
\(R_g=100\ \Omega\)。分别求改装成 \(3.00\ \mathrm A\)
电流表所需的分流电阻，以及改装成 \(15.0\ \mathrm V\)
电压表所需的串联电阻。

\end{minipage}\par

\par\noindent\begin{minipage}{\linewidth}
\textbf{解答}

\begin{center}
\includegraphics[width=\linewidth,height=.40\textheight,keepaspectratio]{\detokenize{assets/module-04/examples/M4-038.pdf}}
\par\smallskip{\small 图 M4-038：双量程仪表改装}
\end{center}
\end{minipage}\par

图 M4-038
的并联支路承担电流表分流，串联支路承担电压表分压。两种改装都要求表头恰好满偏，分别用并联等压与串联等流确定附加电阻。

表头满偏电压为

\[
U_g=I_gR_g=0.100\ \mathrm V.
\]

电流表分流电阻为

\[
R_s=\frac{I_gR_g}{I-I_g}
=\frac{0.100}{3.00-0.00100}
=3.33\times10^{-2}\ \Omega.
\]

电压表串联电阻为

\[
R_m=\frac{15.0}{1.00\times10^{-3}}-100
=1.49\times10^4\ \Omega.
\]

故

\[
\boxed{R_s=0.0333\ \Omega},
\qquad
\boxed{R_m=14.9\ \mathrm{k}\Omega}.
\]

\par\addvspace{.4\baselineskip}
\noindent\begin{minipage}{\linewidth}

\subsubsection{作业二：比较伏安法的系统误差}\label{m4:ux4f5cux4e1aux4e8cux6bd4ux8f83ux4f0fux5b89ux6cd5ux7684ux7cfbux7edfux8befux5dee}

\textbf{题号：M4-039}

\textbf{层级：} 基础　　\textbf{方法：} 电学测量、系统误差

仪表按内阻模型处理：电流表只串入其所在支路，电压表只并入其实际跨接端。相对误差统一定义为
\((R_{\rm \text{测}}-R_x)/R_x\)，两接法分别测量，不同时接成同一回路。

待测电阻 \(R_x=100\ \Omega\)，电流表内阻
\(R_A=1.00\ \Omega\)，电压表内阻
\(R_V=10.0\ \mathrm{k}\Omega\)。分别求``电压表跨接待测电阻与电流表整体''和``电压表只跨接待测电阻''时由
\(U/I\) 得到的测量值及相对误差。

\end{minipage}\par

\par\noindent\begin{minipage}{\linewidth}
\textbf{解答}

\begin{center}
\includegraphics[width=\linewidth,height=.40\textheight,keepaspectratio]{\detokenize{assets/module-04/examples/M4-039.pdf}}
\par\smallskip{\small 图 M4-039：伏安法的系统误差}
\end{center}
\end{minipage}\par

图 M4-039
中电压表的跨接端决定它读到哪一段压降，电流表的位置决定它是否包含电压表分流。先写各自实际读数，再计算同一个比值U/I，不能先把仪表当作理想元件删除。

第一种接法中

\[
R_{\mathrm{\text{测},1}}=R_x+R_A=101\ \Omega,
\]

相对误差为 \(+1.00\%\)。

第二种接法中

\[
R_{\mathrm{\text{测},2}}=\frac{R_xR_V}{R_x+R_V}
=\frac{100\times10000}{10100}
=99.01\ \Omega,
\]

相对误差为 \(-0.990\%\)。两者数值接近，但误差来源和方向不同。

\par\addvspace{.4\baselineskip}
\noindent\begin{minipage}{\linewidth}

\subsubsection{作业三：欧姆表的非线性刻度}\label{m4:ux4f5cux4e1aux4e09ux6b27ux59c6ux8868ux7684ux975eux7ebfux6027ux523bux5ea6}

\textbf{题号：M4-040}

\textbf{层级：} 基础　　\textbf{方法：} 电学测量、非线性刻度

内部总串联电阻包括表头、电池内阻及调节电阻，视为测量期间恒定；指针偏角与电流成正比，电池电动势不随负载变化。

一只串联式欧姆表在表笔短接并调零后，内部总串联电阻为
\(R_0=1.50\ \mathrm{k}\Omega\)，满偏电流为
\(1.00\ \mathrm{mA}\)。求电池电动势、半偏和四分之一满偏位置对应的外接电阻。

\end{minipage}\par

\par\noindent\begin{minipage}{\linewidth}
\textbf{解答}

\begin{center}
\includegraphics[width=\linewidth,height=.40\textheight,keepaspectratio]{\detokenize{assets/module-04/examples/M4-040.pdf}}
\par\smallskip{\small 图 M4-040：欧姆表的非线性刻度}
\end{center}
\end{minipage}\par

图 M4-040
的横向电阻刻度由表头电流的反比例函数生成。短接满偏固定电池电动势，半偏与四分之一满偏分别对应总电阻变成两倍与四倍，随后再减去内部总阻。

短接时

\[
\mathcal E=I_gR_0=1.50\ \mathrm V.
\]

半偏时 \(I=0.500\ \mathrm{mA}\)，

\[
R_0+R_x=\frac{1.50}{0.500\times10^{-3}}=3.00\ \mathrm{k}\Omega,
\]

故 \(R_x=1.50\ \mathrm{k}\Omega\)。

四分之一满偏时 \(I=0.250\ \mathrm{mA}\)，总电阻为
\(6.00\ \mathrm{k}\Omega\)，所以

\[
\boxed{R_x=4.50\ \mathrm{k}\Omega}.
\]

\par\addvspace{.4\baselineskip}
\noindent\begin{minipage}{\linewidth}

\subsubsection{作业四：平衡电桥测未知电阻}\label{m4:ux4f5cux4e1aux56dbux5e73ux8861ux7535ux6865ux6d4bux672aux77e5ux7535ux963b}

\textbf{题号：M4-041}

\textbf{层级：} 基础　　\textbf{方法：} 电学测量、零示法

电源接桥的上、下端，检流计跨接左右中点，全部支路为线性正电阻。平衡定义为检流计电流为零，忽略电源内阻对平衡比值的影响。

惠斯通电桥四臂按左上、左下、右上、右下依次为
\(R_1=100\ \Omega\)、\(R_2=250\ \Omega\)、\(R_3=200\ \Omega\)、\(R_x\)。检流计接左右中点。求平衡时
\(R_x\)，并说明检流计内阻为何不进入结果。

\end{minipage}\par

\par\noindent\begin{minipage}{\linewidth}
\textbf{解答}

\begin{center}
\includegraphics[width=\linewidth,height=.40\textheight,keepaspectratio]{\detokenize{assets/module-04/examples/M4-041.pdf}}
\par\smallskip{\small 图 M4-041：平衡电桥测未知电阻}
\end{center}
\end{minipage}\par

图 M4-041
的检流计跨接两个分压中点。零示时这两个点等势，检流计本身的电阻不产生压降，四桥臂的比例才给出未知电阻，不能把它当成第五个参与分压的桥臂。

平衡时两中点等势，检流计电流为零。左右两支路分别只是串联分压，因此

\[
\frac{R_1}{R_2}=\frac{R_3}{R_x}.
\]

于是

\[
\boxed{R_x=\frac{R_2R_3}{R_1}
=\frac{250\times200}{100}=500\ \Omega}.
\]

平衡时检流计两端电势差为零，所以无论其有限内阻取何值，支路电流都为零；它只影响偏离平衡时的灵敏度。

\par\addvspace{.4\baselineskip}
\noindent\begin{minipage}{\linewidth}

\subsubsection{作业五：电势差计比较电动势}\label{m4:ux4f5cux4e1aux4e94ux7535ux52bfux5deeux8ba1ux6bd4ux8f83ux7535ux52a8ux52bf}

\textbf{题号：M4-042}

\textbf{层级：} 基础　　\textbf{方法：} 电学测量、零示法

工作电阻丝均匀、温度不变，两次测量中单位长度电势降相同；均在检流计零示时读取平衡长度，不从待测电池持续取电流。

同一电势差计工作电流保持不变。标准电池电动势为
\(\mathcal E_s=1.0186\ \mathrm V\)，平衡长度为
\(\ell_s=50.00\ \mathrm{cm}\)；未知电池平衡长度为
\(\ell_x=72.50\ \mathrm{cm}\)。求未知电动势。

\end{minipage}\par

\par\noindent\begin{minipage}{\linewidth}
\textbf{解答}

\begin{center}
\includegraphics[width=\linewidth,height=.40\textheight,keepaspectratio]{\detokenize{assets/module-04/examples/M4-042.pdf}}
\par\smallskip{\small 图 M4-042：补偿法比较电动势}
\end{center}
\end{minipage}\par

图 M4-042
中两次零示平衡长度沿同一均匀电阻丝量取。固定工作电流使单位长度压降相同，两个电动势之比便等于长度之比；这依赖零示时不从待测电池持续取电流。

两次平衡时工作电流相同，单位长度电势降 \(k\) 相同，因此

\[
\frac{\mathcal E_x}{\mathcal E_s}=\frac{\ell_x}{\ell_s}.
\]

代入得

\[
\boxed{\mathcal E_x
=1.0186\times\frac{72.50}{50.00}
=1.477\ \mathrm V}.
\]

\par\addvspace{.4\baselineskip}
\noindent\begin{minipage}{\linewidth}

\subsubsection{作业六：由端口测量辨识电源黑箱}\label{m4:ux4f5cux4e1aux516dux7531ux7aefux53e3ux6d4bux91cfux8fa8ux8bc6ux7535ux6e90ux9ed1ux7bb1}

\textbf{题号：M4-043}

\textbf{层级：} 基础　　\textbf{方法：} 电学测量、端口等效

两次测量采用同一均匀电阻丝和同一工作电流，电势差计均作零示测量；负载直接跨在黑箱两个外端，不将电势差计当作持续分流负载。

把一个两端有源黑箱接到电势差计。开路平衡长度为
\(\ell_0=80.0\ \mathrm{cm}\)；在黑箱端口并联 \(R=10.0\ \Omega\)
负载后，端电压的平衡长度为
\(\ell=64.0\ \mathrm{cm}\)。假设黑箱等效为理想电动势 \(\mathcal E\)
串联内阻 \(r\)，求 \(r\)。

\end{minipage}\par

\par\noindent\begin{minipage}{\linewidth}
\textbf{解答}

\begin{center}
\includegraphics[width=\linewidth,height=.40\textheight,keepaspectratio]{\detokenize{assets/module-04/examples/M4-043.pdf}}
\par\smallskip{\small 图 M4-043：端口测量辨识电源黑箱}
\end{center}
\end{minipage}\par

图 M4-043
的黑箱端口先开路、再跨接已知负载，两次用相同电势差计标度测量。开路长度对应电动势，负载长度对应路端电压；用两者比值求出内部串联电阻。

电势差计工作状态不变，所以

\[
\frac{U}{\mathcal E}=\frac{\ell}{\ell_0}=\frac{64.0}{80.0}=0.800.
\]

接负载时

\[
U=\mathcal E\frac{R}{R+r}.
\]

因此

\[
0.800=\frac{10.0}{10.0+r},
\]

解得

\[
\boxed{r=2.50\ \Omega}.
\]

\section{专题三十三：戴维南---诺顿等效、源变换与最大功率（基础必讲与拔高选讲）}\label{m4:ux4e13ux9898ux4e09ux5341ux4e09ux6234ux7ef4ux5357ux8bfaux987fux7b49ux6548ux6e90ux53d8ux6362ux4e0eux6700ux5927ux529fux7387ux57faux7840ux5fc5ux8bb2ux4e0eux62d4ux9ad8ux9009ux8bb2}

线性一端口在外部只表现为一条仿射伏安关系，因此可用戴维南或诺顿模型代替。等效限于端口伏安关系，不保留内部功率分布和各支路电流。

\subsection{一、线性一端口的仿射伏安关系（基础必讲）}\label{m4:ux4e00ux7ebfux6027ux4e00ux7aefux53e3ux7684ux4effux5c04ux4f0fux5b89ux5173ux7cfbux57faux7840ux5fc5ux8bb2}

取端口 \(a,b\)，定义

\[
U=V_a-V_b,
\]

并规定 \(I>0\) 表示电流从外部由 \(a\)
流入黑盒。在线性元件和线性受控源构成的网络中，消去内部变量以后，端口关系必为

\[
I=G_{\rm th}U+I_0.
\]

若 \(G_{\rm th}\ne0\)，改写为

\[
\boxed{U=U_{\rm oc}+I R_{\rm th}},
\qquad
R_{\rm th}=\frac1{G_{\rm th}}.
\]

这里 \(U_{\rm oc}\) 是端口开路时的电压。由于本节把 \(I\)
规定为流入黑盒，接入被动负载后通常
\(I<0\)。若改把由黑盒流向负载的电流记为 \(I_L=-I\)，则熟悉的形式为

\[
\boxed{U=U_{\rm oc}-I_LR_{\rm th}}.
\]

该端口等效为理想电压源 \(U_{\rm oc}\) 与电阻 \(R_{\rm th}\)
串联，称为戴维南等效。

\par\noindent\begin{minipage}{\linewidth}
\begin{center}
\includegraphics[width=\linewidth,height=.40\textheight,keepaspectratio]{\detokenize{assets/module-04/theory/review-tikz-m4-topic33-v9-from30-v1-01-thevenin-norton.pdf}}
\par\smallskip{\small 图示：任意线性二端黑盒可在端口处等效为戴维南源或诺顿源}
\end{center}
\end{minipage}\par

\subsection{二、等效电阻的三种求法（基础必讲）}\label{m4:ux4e8cux7b49ux6548ux7535ux963bux7684ux4e09ux79cdux6c42ux6cd5ux57faux7840ux5fc5ux8bb2}

\subsubsection{1. 独立源置零}\label{m4:ux72ecux7acbux6e90ux7f6eux96f6}

只含独立源和正电阻时，把独立理想电压源短路、独立理想电流源开路，再从端口求入端电阻，即得
\(R_{\rm th}\)。

\subsubsection{2.
端口加测试源}\label{m4:ux7aefux53e3ux52a0ux6d4bux8bd5ux6e90}

若网络含受控源，受控源必须保留。把所有独立源置零，在端口施加测试电压
\(U_t\)，求流入黑盒的测试电流 \(I_t\)：

\[
\boxed{R_{\rm th}=\frac{U_t}{I_t}}.
\]

也可施加测试电流再求端口电压。测试源法最一般，也能暴露含有源受控元件时可能出现的零电阻或负电阻。

\subsubsection{3. 开路---短路法}\label{m4:ux5f00ux8defux77edux8defux6cd5}

若一端口线性且 \(I_{\rm sc}\ne0\)，把由黑盒流向外部短路的电流定义为
\(I_{\rm sc}\)，则

\[
\boxed{R_{\rm th}=\frac{U_{\rm oc}}{I_{\rm sc}}}.
\]

这个比值依赖线性端口关系；非线性器件不能仅凭一个开路点和一个短路点确定全部伏安曲线。

\par\addvspace{.4\baselineskip}
\noindent\begin{minipage}{\linewidth}

\textbf{例 M4-044：实际电源的开路与短路}

一个实际直流电压源由理想电动势 \(\mathcal E>0\) 与内阻 \(r>0\)
串联构成，外端口上端为正端，导线电阻忽略。分别求外端开路电压
\(U_{\rm oc}\)、理想短接时向外送出的短路电流
\(I_{\rm sc}\)，由两者求戴维南电阻，并说明将独立理想电压源置零时保留下来的连接。

\end{minipage}\par

\par\noindent\begin{minipage}{\linewidth}
\textbf{解答}

\begin{center}
\includegraphics[width=\linewidth,height=.40\textheight,keepaspectratio]{\detokenize{assets/module-04/examples/M4-044.pdf}}
\par\smallskip{\small 图 M4-044：开路、短路与源置零对应三种端态}
\end{center}
\end{minipage}\par

图 M4-044
的内部串联电阻始终属于电源模型。开路时没有内阻压降，短路时电压全部落在内阻上；源置零仅短接理想电动势，不删去
\(r\)。

开路没有输出电流，内部电阻压降为零；短路时外端电压为零，电源电压全部落在
\(r\) 上。由 \(R_{\rm th}=U_{\rm oc}/I_{\rm sc}\) 得
\(r\)。求等效电阻时把理想电压源短接，内阻仍保留在两端之间。

\subsection{三、诺顿等效与源变换（基础必讲）}\label{m4:ux4e09ux8bfaux987fux7b49ux6548ux4e0eux6e90ux53d8ux6362ux57faux7840ux5fc5ux8bb2}

把戴维南关系改写为

\[
I_L=\frac{U_{\rm oc}}{R_{\rm th}}-\frac{U}{R_{\rm th}}.
\]

令

\[
I_N=\frac{U_{\rm oc}}{R_{\rm th}},
\qquad
R_N=R_{\rm th},
\]

就得到理想电流源 \(I_N\) 与电阻 \(R_N\) 并联的诺顿等效。于是

\[
\boxed{
U_{\rm th}=I_NR_N,
\qquad
R_{\rm th}=R_N
}.
\]

源变换只对外端口等效。把电压源串联电阻变成电流源并联电阻以后，源支路内部电流和各元件功率一般不再与原模型逐项相同。

\par\addvspace{.4\baselineskip}
\noindent\begin{minipage}{\linewidth}

\subsection{四、例题：桥式网络对负载的戴维南等效（基础必讲）}\label{m4:ux56dbux4f8bux9898ux6865ux5f0fux7f51ux7edcux5bf9ux8d1fux8f7dux7684ux6234ux7ef4ux5357ux7b49ux6548ux57faux7840ux5fc5ux8bb2}

\textbf{题号：M4-045}

理想电源将上节点 \(P\) 保持在电势 \(\mathcal E\)，下节点 \(Q\)
接地。左侧 \(P-A-Q\) 的电阻依次为 \(R_1,R_2\)，右侧 \(P-B-Q\) 依次为
\(R_3,R_4\)，所有电阻为正。负载 \(R_L>0\) 接在 \(A,B\) 之间，电流正向为
\(A\to B\)。求移去负载后网络在该端口的戴维南开路电压与等效电阻，进而求负载电流，检验桥平衡时的结果。

\end{minipage}\par

\par\noindent\begin{minipage}{\linewidth}
\textbf{解答}

\begin{center}
\includegraphics[width=\linewidth,height=.40\textheight,keepaspectratio]{\detokenize{assets/module-04/examples/M4-045.pdf}}
\par\smallskip{\small 图 M4-045：移去桥负载后求端口等效}
\end{center}
\end{minipage}\par

图 M4-045 的红色支路是负载端口
\(A-B\)。移去它后两侧各自分压得到开路势差；再把独立电压源置零，原上下节点短接，每侧一对电阻并联，两侧再串联。

开路时两个中点分别是分压点：

\[
V_A=\mathcal E\frac{R_2}{R_1+R_2},
\qquad
V_B=\mathcal E\frac{R_4}{R_3+R_4}.
\]

取端口电压 \(U_{AB}=V_A-V_B\)，则

\[
\boxed{
U_{\rm oc}=\mathcal E\left(
\frac{R_2}{R_1+R_2}
-\frac{R_4}{R_3+R_4}
\right)
}.
\]

把理想电压源置零后，上下端节点被短接，因此从 \(A,B\) 看入有

\[
\boxed{
R_{\rm th}=(R_1\parallel R_2)+(R_3\parallel R_4)
}.
\]

若规定负载电流由 \(A\) 流向 \(B\)，则

\[
\boxed{
I_L=\frac{U_{\rm oc}}{R_{\rm th}+R_L}
}.
\]

平衡条件 \(R_1R_4=R_2R_3\) 给出 \(U_{\rm oc}=0\)，因而任意有限 \(R_L\)
中电流均为零，与惠斯通桥结论一致。

\subsection{五、最大功率传输（基础必讲）}\label{m4:ux4e94ux6700ux5927ux529fux7387ux4f20ux8f93ux57faux7840ux5fc5ux8bb2}

对 \(R_{\rm th}>0\) 的戴维南源接正电阻负载 \(R_L\)，负载功率为

\[
P_L=I_L^2R_L
=\frac{U_{\rm oc}^2R_L}{(R_{\rm th}+R_L)^2}.
\]

求导得

\[
\frac{dP_L}{dR_L}
=\frac{U_{\rm oc}^2(R_{\rm th}-R_L)}
{(R_{\rm th}+R_L)^3}.
\]

因此

\[
\boxed{R_L=R_{\rm th}},
\qquad
\boxed{P_{L,\max}=\frac{U_{\rm oc}^2}{4R_{\rm th}}}.
\]

最大功率条件下，内阻和负载耗散相等，效率只有
\(50\%\)。它是``给定线性源、可调纯电阻负载''的最优化结论，不是一般输电系统的最佳运行条件。

\subsection{六、互易定理（拔高选讲）}\label{m4:ux516dux4e92ux6613ux5b9aux7406ux62d4ux9ad8ux9009ux8bb2}

考虑只含双向线性电阻、不含独立源和非互易元件的连通网络。选定参考节点后，约化节点方程为

\[
\boldsymbol G\boldsymbol V=\boldsymbol I,
\]

其中 \(\boldsymbol G\) 为实对称正定矩阵。把单位电流从端口对 \(a,b\)
注入，对应注入向量 \(\boldsymbol s_{ab}\)；观察端口对 \(c,d\)
的电压，得到转移电阻

\[
R_{cd,ab}
=\boldsymbol s_{cd}^{\mathsf T}
\boldsymbol G^{-1}\boldsymbol s_{ab}.
\]

由于 \(\boldsymbol G^{-1}\) 仍对称，故

\[
\boxed{R_{cd,ab}=R_{ab,cd}}.
\]

也就是说，交换激励端口与测量端口，转移电阻不变。含受控源、霍尔元件、旋转器或时变参数的网络未必满足该结论。

\subsection{七、练习与解答}\label{m4:ux4e03ux7ec3ux4e60ux4e0eux89e3ux7b54-2}

\par\addvspace{.4\baselineskip}
\noindent\begin{minipage}{\linewidth}

\subsubsection{作业一：一端口等效参数与负载功率}\label{m4:ux4f5cux4e1aux4e00ux4e00ux7aefux53e3ux7b49ux6548ux53c2ux6570ux4e0eux8d1fux8f7dux529fux7387}

\textbf{题号：M4-046}

\textbf{层级：} 基础　　\textbf{方法：} 能量法、等效变换

下面两问为两个独立的线性直流一端口，端口电压均取上端减下端；接负载时的电流以向外输出为正，而加测试源时明确采用流入黑盒为正。

\begin{enumerate}
\def\labelenumi{(\arabic{enumi})}
\item
  一个具有正等效内阻的一端口，开路电压为
  \(12\ \mathrm V\)，向外输出的短路电流为
  \(3\ \mathrm A\)。求其诺顿等效参数；接入 \(R_L=2\ \Omega\)
  后，求负载电流和功率。
\item
  另一端口含线性受控源。将独立源置零并保留受控源后，加测试电压
  \(U_t=1\ \mathrm V\)，测得流入黑盒的有符号电流为
  \(I_t=-0.2\ \mathrm A\)。求该端口
  \(R_{\rm th}\)，说明能否直接采用正电阻负载的最大功率匹配结论。测试数据只代表受控源模型的工作范围，不假设它对任意负载都稳定。
\end{enumerate}

\end{minipage}\par

\par\noindent\begin{minipage}{\linewidth}
\textbf{解答}

\begin{center}
\includegraphics[width=\linewidth,height=.40\textheight,keepaspectratio]{\detokenize{assets/module-04/examples/M4-046.pdf}}
\par\smallskip{\small 图 M4-046：一端口等效与负载功率}
\end{center}
\end{minipage}\par

图 M4-046
的开路、短路数据共同固定同一个线性端口模型。先由两读数求等效内阻，再在诺顿并联模型中接入负载；负载电流不是短路电流本身，须计入分流。

（1）正内阻端口。

\[
R_{\rm th}=\frac{12}{3}=4\ \Omega,
\qquad
I_N=3\ \mathrm A.
\]

接入 \(2\ \Omega\) 负载后

\[
I_L=\frac{12}{4+2}=2\ \mathrm A,
\qquad
\boxed{P_L=I_L^2R_L=8\ \mathrm W}.
\]

（2）含受控源的测试端口。

\[
\boxed{R_{\rm th}=\frac{U_t}{I_t}=-5\ \Omega}.
\]

该端口对小信号表现为有源负电阻。公式 \(R_L=R_{\rm th}\) 要求
\(R_{\rm th}>0\) 且
\(R_L>0\)，此处不能直接套用；还必须结合受控源的工作范围和稳定性判断允许的负载。

\par\noindent\begin{minipage}{\linewidth}
\begin{center}
\includegraphics[width=\linewidth,height=.40\textheight,keepaspectratio]{\detokenize{assets/module-04/process/M4-046a.pdf}}
\par\smallskip{\small 图 M4-046a：测试源电流为负时端口表现为负电阻}
\end{center}
\end{minipage}\par

图 M4-046a
中红箭头定义流入黑箱为正，测试结果为负表示实际电流沿反向流出。它与第一问的无源正内阻模型不同，因此不能把匹配条件解释成可接入负值的普通负载。

\subsection{八、使用条件}\label{m4:ux516bux4f7fux7528ux6761ux4ef6}

\begin{enumerate}
\def\labelenumi{\arabic{enumi}.}
\tightlist
\item
  戴维南与诺顿等效只保持指定二端口的伏安关系，不保持黑盒内部状态。
\item
  独立源置零时，受控源不能删除；含受控源时优先使用测试源。
\item
  \(R_{\rm th}=U_{\rm oc}/I_{\rm sc}\)
  依赖线性端口关系，并要求相应量有限且定义明确。
\item
  不同电压的理想电压源不能无内阻并联；不同电流的理想电流源不能无条件串联。
\item
  最大功率定理不等于最大效率定理；有源或负阻一端口还必须检查稳定性。
\item
  互易性要求网络元件和本构关系互易，不能从``网络看起来对称''代替数学条件。
\end{enumerate}

\section{专题三十四：星形---三角形变换与三端网络等效（基础必讲与拔高选讲）}\label{m4:ux4e13ux9898ux4e09ux5341ux56dbux661fux5f62ux4e09ux89d2ux5f62ux53d8ux6362ux4e0eux4e09ux7aefux7f51ux7edcux7b49ux6548ux57faux7840ux5fc5ux8bb2ux4e0eux62d4ux9ad8ux9009ux8bb2}

当电阻网络中出现非平衡电桥时，任何一只电阻都可能既不与另一只严格串联，也不与另一只严格并联。星形---三角形变换通过保持三个外部端口的伏安关系不变，把局部三端网络改写成更容易化简的拓扑。

\subsection{一、变换对象与适用条件（基础必讲）}\label{m4:ux4e00ux53d8ux6362ux5bf9ux8c61ux4e0eux9002ux7528ux6761ux4ef6ux57faux7840ux5fc5ux8bb2}

设三个外部节点为 \(1,2,3\)。

三角形网络的三条边为

\[
R_{12},\qquad R_{23},\qquad R_{31}.
\]

星形网络中，内部节点为 \(N\)，从节点 \(1,2,3\) 到 \(N\) 的支路分别为

\[
R_1,\qquad R_2,\qquad R_3.
\]

\par\noindent\begin{minipage}{\linewidth}
\begin{center}
\includegraphics[width=\linewidth,height=.40\textheight,keepaspectratio]{\detokenize{assets/module-04/theory/review-tikz-m4-topic34-v9-from31-v8-01.pdf}}
\par\smallskip{\small 图示：三角形网络与星形网络只共享三个外部节点，星形中心点必须保持为内部节点}
\end{center}
\end{minipage}\par

变换必须满足以下条件：

\begin{enumerate}
\def\labelenumi{\arabic{enumi}.}
\tightlist
\item
  子网络对外只有节点 \(1,2,3\) 三个连接端；
\item
  星形中心点 \(N\) 是内部节点，不能另接导线、元件或接地；
\item
  本专题首先讨论正电阻构成的无源线性网络；
\item
  子网络内部若含独立电源，不能只变换电阻而忽略电源；
\item
  变换后必须保留原来三个外部节点，不能按图形位置错误对应。
\end{enumerate}

``等效''不是指三只电阻逐一相等，而是指：对任意外部端口电势
\(V_1,V_2,V_3\)，两种网络给出的端口电流 \(I_1,I_2,I_3\) 完全相同。

\[
\boxed{
(V_1,V_2,V_3)\longmapsto(I_1,I_2,I_3)
\quad\text{保持不变}
}.
\]

由于网络浮置，整体电势同时增加一个常数不会改变电流，并且始终有

\[
I_1+I_2+I_3=0.
\]

\subsection{二、三端口等效的证明（拔高选讲）}\label{m4:ux4e8cux4e09ux7aefux53e3ux7b49ux6548ux7684ux8bc1ux660eux62d4ux9ad8ux9009ux8bb2}

先从星形网络出发。令

\[
g_i=\frac1{R_i},
\qquad
G=g_1+g_2+g_3.
\]

中心节点 \(N\) 不与外界相连，由 KCL：

\[
g_1(V_N-V_1)+g_2(V_N-V_2)+g_3(V_N-V_3)=0.
\]

因此

\[
\boxed{
V_N=\frac{g_1V_1+g_2V_2+g_3V_3}{G}
}.
\]

从外部节点 \(1\) 流入星形网络的电流为

\[
I_1=g_1(V_1-V_N).
\]

消去 \(V_N\)：

\[
I_1
=\frac{g_1g_2}{G}(V_1-V_2)
+\frac{g_1g_3}{G}(V_1-V_3).
\]

轮换指标可得

\[
I_2
=\frac{g_1g_2}{G}(V_2-V_1)
+\frac{g_2g_3}{G}(V_2-V_3),
\]

\[
I_3
=\frac{g_1g_3}{G}(V_3-V_1)
+\frac{g_2g_3}{G}(V_3-V_2).
\]

而三角形网络的端口电流为

\[
I_1
=\frac{V_1-V_2}{R_{12}}
+\frac{V_1-V_3}{R_{31}},
\]

其余两式由轮换得到。逐项比较便有

\[
\boxed{
\frac1{R_{12}}=\frac{g_1g_2}{G},
\quad
\frac1{R_{23}}=\frac{g_2g_3}{G},
\quad
\frac1{R_{31}}=\frac{g_3g_1}{G}
}.
\]

\par\noindent\begin{minipage}{\linewidth}
\begin{center}
\includegraphics[width=\linewidth,height=.40\textheight,keepaspectratio]{\detokenize{assets/module-04/theory/review-tikz-m4-topic34-v9-from31-v8-02.pdf}}
\par\smallskip{\small 图示：消去星形中心节点后，端口电流恰好写成三角形三条边的电流之和}
\end{center}
\end{minipage}\par

这一步证明了变换保持全部三端口伏安关系，而不只是某一种两端测量。因而把任意外部电路接在节点
\(1,2,3\) 上，外部电路都无法分辨内部采用星形还是三角形。

相同的端口伏安关系也保证端口吸收功率相同：

\[
P=V_1I_1+V_2I_2+V_3I_3.
\]

\subsection{三、星形变三角形（基础必讲）}\label{m4:ux4e09ux661fux5f62ux53d8ux4e09ux89d2ux5f62ux57faux7840ux5fc5ux8bb2}

定义

\[
Q=R_1R_2+R_2R_3+R_3R_1.
\]

由上一节

\[
R_{12}=\frac{G}{g_1g_2}.
\]

把 \(g_i=1/R_i\) 代入：

\[
R_{12}
=\frac{\dfrac1{R_1}+\dfrac1{R_2}+\dfrac1{R_3}}
{\dfrac1{R_1R_2}}
=R_1+R_2+\frac{R_1R_2}{R_3}.
\]

因此

\[
\boxed{
R_{12}=\frac{Q}{R_3}
},
\]

\[
\boxed{
R_{23}=\frac{Q}{R_1},
\qquad
R_{31}=\frac{Q}{R_2}
}.
\]

记忆时必须先看目标边：连接节点 \(1,2\) 的三角形电阻
\(R_{12}\)，分母是星形中远离这两个节点的``对臂'' \(R_3\)。

\par\addvspace{.4\baselineskip}
\noindent\begin{minipage}{\linewidth}

\textbf{例 M4-047：不等臂星形变三角形}

无源三端星形网络的中心为内部节点 \(O\)，外端为 \(1,2,3\)，三条臂
\(O-1,O-2,O-3\) 的电阻依次为
\(1\ \Omega,2\ \Omega,3\ \Omega\)。中心没有外部注入电流。求在任意外端电势激励下完全等效的三角形三条支路电阻
\(R_{12},R_{23},R_{31}\)，明确每条目标边对应哪一条星形对臂，并检查量纲。

\end{minipage}\par

\par\noindent\begin{minipage}{\linewidth}
\textbf{解答}

\begin{center}
\includegraphics[width=\linewidth,height=.40\textheight,keepaspectratio]{\detokenize{assets/module-04/examples/M4-047.pdf}}
\par\smallskip{\small 图 M4-047：对边的星形臂进入变换分母}
\end{center}
\end{minipage}\par

图 M4-047 的三条臂分别接到固定标号的外端。目标边1---2的对臂是
\(R_3\)，所以用三项乘积之和除以它；另外两边循环对应，不能把星形中心误当作第四个外端。

若

\[
R_1=1\ \Omega,\qquad
R_2=2\ \Omega,\qquad
R_3=3\ \Omega,
\]

则

\[
Q=1\times2+2\times3+3\times1=11\ \Omega^2,
\]

所以

\[
R_{12}=\frac{11}{3}\ \Omega,
\quad
R_{23}=11\ \Omega,
\quad
R_{31}=\frac{11}{2}\ \Omega.
\]

量纲检查很直接：\(Q\) 的量纲是 \(\Omega^2\)，再除以一只电阻才得到
\(\Omega\)。

\subsection{四、三角形变星形（基础必讲）}\label{m4:ux56dbux4e09ux89d2ux5f62ux53d8ux661fux5f62ux57faux7840ux5fc5ux8bb2}

令

\[
S=R_{12}+R_{23}+R_{31}.
\]

也可以令第三端开路，分别比较任意两端之间的等效电阻。以节点 \(1,2\)
为例：

\[
R_1+R_2
=R_{12}\parallel(R_{31}+R_{23})
=\frac{R_{12}(R_{31}+R_{23})}{S}.
\]

同理，

\[
R_2+R_3
=\frac{R_{23}(R_{12}+R_{31})}{S},
\]

\[
R_3+R_1
=\frac{R_{31}(R_{12}+R_{23})}{S}.
\]

取``第一式加第三式减第二式''的一半：

\[
R_1
=\frac12\left[(R_1+R_2)+(R_1+R_3)-(R_2+R_3)\right],
\]

代入上面三式并整理：

\[
\boxed{
R_1=\frac{R_{12}R_{31}}{S}
}.
\]

轮换得到

\[
\boxed{
R_2=\frac{R_{12}R_{23}}{S},
\qquad
R_3=\frac{R_{23}R_{31}}{S}
}.
\]

连接节点 \(1\) 的星形电阻 \(R_1\)，分子是三角形中夹住节点 \(1\)
的两条相邻边之积，分母是三角形三边之和。

\par\noindent\begin{minipage}{\linewidth}
\begin{center}
\includegraphics[width=\linewidth,height=.40\textheight,keepaspectratio]{\detokenize{assets/module-04/theory/review-tikz-m4-topic34-v9-from31-v8-03.pdf}}
\par\smallskip{\small 图示：三角形变星形取节点两侧相邻边的乘积，星形变三角形取两两乘积和除以对臂}
\end{center}
\end{minipage}\par

\par\addvspace{.4\baselineskip}
\noindent\begin{minipage}{\linewidth}

\textbf{例 M4-048：等边三角形变星形}

一个无源三端网络的三角形三边电阻均为
\(R_{\Delta}>0\)，求与其多端等效的星形三条臂电阻。反过来，若星形三臂均为
\(R_Y>0\)，求等效三角形每边电阻。要求外端点标号在变换前后对应，星形中心不接外部电路，并区分``等效支路阻值''和任选两端、第三端悬空时的测量电阻。

\end{minipage}\par

\par\noindent\begin{minipage}{\linewidth}
\textbf{解答}

\begin{center}
\includegraphics[width=\linewidth,height=.40\textheight,keepaspectratio]{\detokenize{assets/module-04/examples/M4-048.pdf}}
\par\smallskip{\small 图 M4-048：等臂星形与等边三角形的参数关系}
\end{center}
\end{minipage}\par

图 M4-048
把中心明确留作内部节点。三端完全等效要求保留任意外端响应，所得三角形每边是
\(3R_Y\)；做二端实验时还存在经第三端的串联通路，测量值不是单条三角形支路。

若三角形三边都为 \(R_{\Delta}\)，则

\[
R_1=R_2=R_3
=\frac{R_{\Delta}^2}{3R_{\Delta}}
=\frac{R_{\Delta}}{3}.
\]

反过来，三条星形支路都为 \(R_Y\) 时，三角形每边均为

\[
R_{\Delta}=3R_Y.
\]

这只是参数之间的等效关系，不表示某一种拓扑在所有参数相同的比较下必然``更好''或``更通畅''。

第三端悬空时，三角形的两条通路为 \(R_{\Delta}\) 与
\(2R_{\Delta}\)，故测量值为 \(2R_{\Delta}/3=2R_Y\)，与星形两臂串联相同。

\subsection{五、极限、可逆性与阻抗推广（拔高选讲）}\label{m4:ux4e94ux6781ux9650ux53efux9006ux6027ux4e0eux963bux6297ux63a8ux5e7fux62d4ux9ad8ux9009ux8bb2}

\subsubsection{1. 开路极限}\label{m4:ux5f00ux8defux6781ux9650}

若三角形中 \(R_{23}\to\infty\)，则原网络只剩两条边 \(R_{12}\) 与
\(R_{31}\) 在节点 \(1\) 相接。由变换公式：

\[
R_1\to0,
\qquad
R_2\to R_{12},
\qquad
R_3\to R_{31}.
\]

这意味着星形中心点与节点 \(1\) 合并，正好回到原拓扑。

\subsubsection{2.
正电阻保证正值}\label{m4:ux6b63ux7535ux963bux4fddux8bc1ux6b63ux503c}

若所有原电阻都严格为正，则公式中的和、积也为正，所以变换后的电阻仍严格为正。若出现负电阻或分母为零，应检查是否已经进入含源等效或理想化奇异情形。

\subsubsection{3. 可逆性}\label{m4:ux53efux9006ux6027}

把三角形变星形的结果代回星形变三角形公式，会恢复原来的三条边。实际计算若不能恢复，通常是节点下标或``对臂''对应写错。

\subsubsection{4. 交流阻抗}\label{m4:ux4ea4ux6d41ux963bux6297}

在线性正弦稳态电路中，若三支路之间没有互感耦合，可以把每个 \(R\)
替换为同一角频率下的复阻抗 \(Z\)，代数公式保持不变：

\[
Z_1=\frac{Z_{12}Z_{31}}{Z_{12}+Z_{23}+Z_{31}},
\]

\[
Z_{12}
=\frac{Z_1Z_2+Z_2Z_3+Z_3Z_1}{Z_3}.
\]

若分母为零，变换出现奇异，不能机械套用。

\par\noindent\begin{minipage}{\linewidth}
\begin{center}
\includegraphics[width=\linewidth,height=.40\textheight,keepaspectratio]{\detokenize{assets/module-04/theory/M4-R20-star-tests.pdf}}
\par\smallskip{\small 图示：对称情形、开路极限和逆变换共同检验星三角公式}
\end{center}
\end{minipage}\par

\subsection{六、星形---三角形变换的使用条件（基础必讲）}\label{m4:ux516dux661fux5f62ux4e09ux89d2ux5f62ux53d8ux6362ux7684ux4f7fux7528ux6761ux4ef6ux57faux7840ux5fc5ux8bb2}

\begin{enumerate}
\def\labelenumi{\arabic{enumi}.}
\tightlist
\item
  确认局部子网络确实只有三个外部端口；
\item
  选择能使后续出现串并联关系的变换方向；
\item
  在草图上固定节点 \(1,2,3\)，再写带下标的公式；
\item
  完成变换后重画电路，不按原图的视觉位置判断串并联；
\item
  继续化简，并用对称、量纲与极限检查结果。
\end{enumerate}

\par\noindent\begin{minipage}{\linewidth}
\begin{center}
\includegraphics[width=\linewidth,height=.40\textheight,keepaspectratio]{\detokenize{assets/module-04/theory/M4-R20-star-flow.pdf}}
\par\smallskip{\small 图示：星三角变换题应先识别三端子网，再选方向、代公式、重画拓扑并检验}
\end{center}
\end{minipage}\par

常见错误包括：

\begin{itemize}
\tightlist
\item
  星形中心点还有第四条外接支路，却强行变换；
\item
  把三角形中不相邻的两条边写进 \(R_1\) 的分子；
\item
  星形变三角形时把分母除以目标边相邻的臂，而不是对臂；
\item
  变换后保留已经被替换掉的原支路；
\item
  不重画拓扑，误判新的串联或并联关系；
\item
  子网络中含电源，却只转换电阻。
\end{itemize}

\subsection{七、练习与解答}\label{m4:ux4e03ux7ec3ux4e60ux4e0eux89e3ux7b54-3}

\par\addvspace{.4\baselineskip}
\noindent\begin{minipage}{\linewidth}

\subsubsection{作业一：非平衡电桥的等效电阻}\label{m4:ux4f5cux4e1aux4e00ux975eux5e73ux8861ux7535ux6865ux7684ux7b49ux6548ux7535ux963b}

\textbf{题号：M4-049}

\textbf{层级：} 基础　　\textbf{方法：} 等效变换、网络方程

这是无源有限网络，只有明确列出的五条电阻边，所有未列出的节点对均不连接。以
\(A\) 为流入端、\(C\)
为流出端，其余两点无外部注入，求二端等效电阻而非任一单支路阻值。

四个节点为 \(A,B,C,D\)，支路电阻为

\[
R_{AB}=1\ \Omega,\qquad
R_{AD}=2\ \Omega,\qquad
R_{BD}=1\ \Omega,
\]

\[
R_{BC}=2\ \Omega,\qquad
R_{DC}=1\ \Omega.
\]

求节点 \(A,C\) 之间的等效电阻。

\end{minipage}\par

\par\noindent\begin{minipage}{\linewidth}
\textbf{解答}

\begin{center}
\includegraphics[width=\linewidth,height=.40\textheight,keepaspectratio]{\detokenize{assets/module-04/examples/M4-049.pdf}}
\par\smallskip{\small 图 M4-049：非平衡桥保留中间真实支路}
\end{center}
\end{minipage}\par

图 M4-049 中 \(B-D\) 是真实电阻支路，不能因外形近似对称而删除。把
\(A,B,D\) 三角形变成中心 \(N\) 的星形后，\(N-B-C\) 与 \(N-D-C\)
两路并联，最后与 \(A-N\) 串联。

先检查电桥是否平衡：

\[
\frac{R_{AB}}{R_{AD}}
=\frac12,
\qquad
\frac{R_{BC}}{R_{DC}}
=2.
\]

两比值不相等，支路 \(BD\) 一般有电流，不能直接删去。

选择三角形 \(ABD\)：

\[
R_{AB}=1\ \Omega,\qquad
R_{BD}=1\ \Omega,\qquad
R_{DA}=2\ \Omega.
\]

三边之和为

\[
S=1+1+2=4\ \Omega.
\]

变换为中心点 \(N\) 的星形网络：

\[
R_A
=\frac{R_{AB}R_{AD}}S
=\frac{1\times2}{4}\ \Omega
=\frac12\ \Omega,
\]

\[
R_B
=\frac{R_{AB}R_{BD}}S
=\frac{1\times1}{4}\ \Omega
=\frac14\ \Omega,
\]

\[
R_D
=\frac{R_{AD}R_{BD}}S
=\frac{2\times1}{4}\ \Omega
=\frac12\ \Omega.
\]

变换后，从 \(N\) 到 \(C\) 有两条并联支路：

\[
R_{NBC}=R_B+R_{BC}
=\frac14+2
=\frac94\ \Omega,
\]

\par\noindent\begin{minipage}{\linewidth}
\begin{center}
\includegraphics[width=\linewidth,height=.40\textheight,keepaspectratio]{\detokenize{assets/module-04/process/M4-049a.pdf}}
\par\smallskip{\small 图 M4-049a：变换后两路并联、左臂串联}
\end{center}
\end{minipage}\par

图 M4-049a
只消去了原三角形ABD，原来的BC、DC两条外接边仍保留。新中心N到C的上、下两路先并联，A到N的电阻再串联，避免把星形中的臂误当成外端对外端边。

\[
R_{NDC}=R_D+R_{DC}
=\frac12+1
=\frac32\ \Omega.
\]

因此

\[
R_{NC}
=\frac{\dfrac94\cdot\dfrac32}
{\dfrac94+\dfrac32}
=\frac9{10}\ \Omega.
\]

最后 \(R_A\) 与 \(R_{NC}\) 串联：

\[
\boxed{
R_{AC}
=\frac12+\frac9{10}
=\frac75\ \Omega
=1.4\ \Omega
}.
\]

\textbf{节点法复核：}

取 \(V_A=1\ \mathrm V\)、\(V_C=0\)。对节点 \(B,D\) 列 KCL：

\[
\frac{V_B-1}{1}
+\frac{V_B-V_D}{1}
+\frac{V_B}{2}=0,
\]

\[
\frac{V_D-1}{2}
+\frac{V_D-V_B}{1}
+\frac{V_D}{1}=0.
\]

解得

\[
V_B=\frac47\ \mathrm V,
\qquad
V_D=\frac37\ \mathrm V.
\]

从节点 \(A\) 流入网络的总电流为

\[
I
=\frac{1-V_B}{1}
+\frac{1-V_D}{2}
=\frac57\ \mathrm A.
\]

所以

\[
R_{AC}=\frac{1\ \mathrm V}{5/7\ \mathrm A}
=\frac75\ \Omega,
\]

与星形---三角形变换结果一致。

\subsection{八、使用条件}\label{m4:ux516bux4f7fux7528ux6761ux4ef6-1}

\begin{enumerate}
\def\labelenumi{\arabic{enumi}.}
\tightlist
\item
  变换只保持外部三节点的伏安关系，不保持内部中心点的电势、内部支路电流或局部耗散逐项对应。
\item
  星形中心点若被外部使用，便不再是可以消去的内部节点。
\item
  对含独立源、受控源、互感或非线性元件的三端子网，不能直接套用纯电阻公式。
\item
  对正电阻网络，变换前后任意外接条件下的端口总功率相同，但功率在内部支路间的分配通常不同。
\item
  星形---三角形变换是局部拓扑工具；是否值得使用，应由变换后能否产生串并联或对称结构决定。
\end{enumerate}

\section{专题三十五：电路中的等势点法与对折法（基础必讲与拔高选讲）}\label{m4:ux4e13ux9898ux4e09ux5341ux4e94ux7535ux8defux4e2dux7684ux7b49ux52bfux70b9ux6cd5ux4e0eux5bf9ux6298ux6cd5ux57faux7840ux5fc5ux8bb2ux4e0eux62d4ux9ad8ux9009ux8bb2}

复杂电阻网络应先判断哪些节点必然等势。等势点法把对称性写成节点电势约束，对折法再据此降低网络维数。判断时必须同时检查网络和端口激励，不能只看几何外形。

\subsection{一、等势点的两种处理（基础必讲）}\label{m4:ux4e00ux7b49ux52bfux70b9ux7684ux4e24ux79cdux5904ux7406ux57faux7840ux5fc5ux8bb2}

若节点 \(i,j\) 的电势已经由对称性或方程证明满足

\[
V_i=V_j,
\]

则有两种常用操作。

\subsubsection{1.
删去两等势点之间的被动支路}\label{m4:ux5220ux53bbux4e24ux7b49ux52bfux70b9ux4e4bux95f4ux7684ux88abux52a8ux652fux8def}

若 \(i,j\) 之间连有电阻 \(R_{ij}>0\)，则

\[
I_{ij}=\frac{V_i-V_j}{R_{ij}}=0.
\]

因而在当前端口激励下，可删去这条支路，其他支路的电流与节点电势不变。

\subsubsection{2.
合并已证明等势的节点}\label{m4:ux5408ux5e76ux5df2ux8bc1ux660eux7b49ux52bfux7684ux8282ux70b9}

在理想节点模型中，可将 \(i,j\)
用理想导线连接并视为同一节点。新导线上的电压为零，原网络的节点电势仍是改造后网络的解；由解的唯一性，当前端口激励下的其他支路状态不变。

\par\noindent\begin{minipage}{\linewidth}
\begin{center}
\includegraphics[width=\linewidth,height=.40\textheight,keepaspectratio]{\detokenize{assets/module-04/theory/review-tikz-m4-topic35-v9-from32-v8-01.pdf}}
\par\smallskip{\small 图示：已证明等势的节点可合并，两点之间的纯电阻支路可删去}
\end{center}
\end{minipage}\par

\begin{quote}
\begin{quote}
``可合并''的前提是电势相等已被证明。两节点在电路图上看起来对称，不代表在当前端口接法下一定等势。
\end{quote}

若 \(V_i=V_j\)
只在某一指定的两端激励下成立，合并只保证该测量问题的结果不变；只有对所有允许的外部激励都有
\(V_i=V_j\)，才能声称完整的多端口伏安关系不变。
\end{quote}

\subsection{二、唯一性与对称性判据（拔高选讲）}\label{m4:ux4e8cux552fux4e00ux6027ux4e0eux5bf9ux79f0ux6027ux5224ux636eux62d4ux9ad8ux9009ux8bb2}

设连通线性电阻网络的节点电势列向量为 \(\mathbf V\)，外部注入电流为
\(\mathbf I\)，节点方程写成

\[
L\mathbf V=\mathbf I.
\]

其中 \(L\) 为加权图拉普拉斯矩阵：

\[
L_{ii}=\sum_{j\ne i}g_{ij},
\qquad
L_{ij}=-g_{ij}\quad(i\ne j),
\]

\(g_{ij}=1/R_{ij}\) 是节点 \(i,j\) 间的总电导。对任意实向量
\(\mathbf x\)，

\[
\mathbf x^{\mathsf T}L\mathbf x
=\sum_{i<j}g_{ij}(x_i-x_j)^2\ge0.
\]

因此 \(L\)
是半正定矩阵，常数向量属于其零空间。连通网络在固定一个参考电势后，约化矩阵才是正定的，节点电势解唯一。

设网络有一个节点置换 \(P\)，且它保持所有连接关系与电阻值，则

\[
PLP^{-1}=L.
\]

结论取决于激励 \(\mathbf I\) 在该置换下如何变换。

\subsubsection{1. 偶对称激励}\label{m4:ux5076ux5bf9ux79f0ux6fc0ux52b1}

若

\[
P\mathbf I=\mathbf I,
\]

则 \(P\mathbf V\) 与 \(\mathbf V\)
满足同一节点方程。在同一参考电势下，由唯一性有

\[
\boxed{P\mathbf V=\mathbf V}.
\]

所以被 \(P\)
互换的节点等势。常见情形是：电流流入端和流出端分别位于对称轴或对称面上，对称操作不交换两端。

\subsubsection{2. 奇对称激励}\label{m4:ux5947ux5bf9ux79f0ux6fc0ux52b1}

若对称操作交换电流流入端与流出端，则

\[
P\mathbf I=-\mathbf I.
\]

取两端电势中值为零，唯一性给出

\[
\boxed{P\mathbf V=-\mathbf V}.
\]

被交换的节点电势互为相反数；在对称操作下保持不动的节点满足

\[
V=-V,
\]

因而电势为零。这些节点构成``虚拟零势节点集''。

\par\noindent\begin{minipage}{\linewidth}
\begin{center}
\includegraphics[width=\linewidth,height=.40\textheight,keepaspectratio]{\detokenize{assets/module-04/theory/review-tikz-m4-topic35-v9-from32-v8-02.pdf}}
\par\smallskip{\small 图示：网络对称性只有与端口激励一起考察才能推出等势或反势}
\end{center}
\end{minipage}\par

\subsection{三、对折法与半电路法（基础必讲）}\label{m4:ux4e09ux5bf9ux6298ux6cd5ux4e0eux534aux7535ux8defux6cd5ux57faux7840ux5fc5ux8bb2}

\subsubsection{1. 偶对称对折}\label{m4:ux5076ux5bf9ux79f0ux5bf9ux6298}

对称面两侧的对应节点等势时，可将它们合并。对折后：

\begin{itemize}
\tightlist
\item
  两条互为镜像的电阻若完全重合，它们是并联；阻值都为 \(R\) 时等效为
  \(R/2\)；
\item
  一条连接两个镜像节点的电阻，两端等势，可删去；
\item
  本来位于对称面上且映射为自身的支路，只有一条，阻值不变。
\end{itemize}

\subsubsection{2.
奇对称半电路}\label{m4:ux5947ux5bf9ux79f0ux534aux7535ux8def}

对称操作交换两端时，固定节点为零势。可只计算一半网络，设正端电势为
\(U/2\)，注入电流为 \(I\)，则

\[
R_{1/2}=\frac{U/2}{I},
\]

完整两端电阻为

\[
\boxed{R_{AB}=2R_{1/2}}.
\]

若一只阻值 \(R\) 的电阻直接连接一对反势节点
\(+V,-V\)，从正半网络看，其电流为

\[
I=\frac{2V}{R}=\frac{V}{R/2},
\]

所以这条跨越奇对称面的支路在半电路中应等效为接零势节点的 \(R/2\)，而不是
\(R\)。

\par\noindent\begin{minipage}{\linewidth}
\begin{center}
\includegraphics[width=\linewidth,height=.40\textheight,keepaspectratio]{\detokenize{assets/module-04/theory/M4-R20-fold.pdf}}
\par\smallskip{\small 图示：偶对称对折后镜像电阻并联，奇对称则以零势面构造半电路}
\end{center}
\end{minipage}\par

\subsection{四、例题（基础必讲）}\label{m4:ux56dbux4f8bux9898ux57faux7840ux5fc5ux8bb2}

\par\addvspace{.4\baselineskip}
\noindent\begin{minipage}{\linewidth}

\subsubsection{对称电桥的桥臂电流}\label{m4:ux5bf9ux79f0ux7535ux6865ux7684ux6865ux81c2ux7535ux6d41}

\textbf{题号：M4-050}

四个正电阻 \(R\) 构成上、下两条从左端 \(L\) 到右端 \(K\)
的串联支路，两支路中点分别为 \(A,B\)，中点间另接桥臂电阻
\(R_b>0\)。理想电源接在 \(L,K\)，给定
\(V_L-V_K=U\)。利用保持连接和端口激励的反射对称，求桥臂电流及整个网络的端口等效电阻；再解释若把电源移到相邻节点间，为什么不能直接沿用原等势判断。

\end{minipage}\par

\par\noindent\begin{minipage}{\linewidth}
\textbf{解答}

\begin{center}
\includegraphics[width=\linewidth,height=.40\textheight,keepaspectratio]{\detokenize{assets/module-04/examples/M4-050.pdf}}
\par\smallskip{\small 图 M4-050：保持激励的反射使桥中点等势}
\end{center}
\end{minipage}\par

图 M4-050 的四条外臂相同，关于水平轴的反射保持左右端激励，只交换
\(A,B\)。唯一性使两点等势，故红色桥臂电流为零；删去该无电流支路后才可作两路串联再并联。

上、下两条支路都由两只 \(R\) 串联，中点之间连电阻
\(R_b\)，电源接在左、右端。关于水平轴反射时，网络与两端激励均不变，所以两中点等势，

\[
I_b=0.
\]

删去桥臂后，两条 \(2R\) 支路并联，故

\[
\boxed{R_{\rm eq}=R}.
\]

若改为在相邻节点之间接电源，原反射不再保持激励，两中点不再必然等势。

\par\addvspace{.4\baselineskip}
\noindent\begin{minipage}{\linewidth}

\subsubsection{奇对称面不是``无电流面''}\label{m4:ux5947ux5bf9ux79f0ux9762ux4e0dux662fux65e0ux7535ux6d41ux9762}

\textbf{题号：M4-051}

一只线性电阻 \(R>0\) 的两端分别处于电势
\(+U/2,-U/2\)，\(U>0\)。某个反射对称把这两个端点互换，并使电势变号，故对称面可取零势。把电流正向定为正势端流向负势端，求整只电阻电流，判断对称面是否是无电流面；说明若构造只含一侧的半电路，跨越该面的电阻应如何分配。

\end{minipage}\par

\par\noindent\begin{minipage}{\linewidth}
\textbf{解答}

\begin{center}
\includegraphics[width=\linewidth,height=.40\textheight,keepaspectratio]{\detokenize{assets/module-04/examples/M4-051.pdf}}
\par\smallskip{\small 图 M4-051：零势面不是无电流面}
\end{center}
\end{minipage}\par

图 M4-051 中虚线交换正负势端，属于奇对称面。电阻两端仍有完整电压
\(U\)，所以电流穿过该面；切出一半时同时把电压和跨面电阻减半，不能只把电压减半。

正向电流为
\([(+U/2)-(-U/2)]/R=U/R\)，它确实穿过零势对称面。沿电阻均匀分割到中点时，每一侧电阻为
\(R/2\)；半电路对地电压为 \(U/2\)，因此电流仍为
\((U/2)/(R/2)=U/R\)。零电势不等于零电流。

\subsection{五、等势点法和对折法的使用条件（基础必讲）}\label{m4:ux4e94ux7b49ux52bfux70b9ux6cd5ux548cux5bf9ux6298ux6cd5ux7684ux4f7fux7528ux6761ux4ef6ux57faux7840ux5fc5ux8bb2}

\begin{enumerate}
\def\labelenumi{\arabic{enumi}.}
\tightlist
\item
  先标出实际电流流入端和流出端；
\item
  寻找同时保持连接关系与电阻值的节点置换；
\item
  判断该置换是保持激励还是反转激励；
\item
  偶对称用``对应节点等势''，奇对称用``对应节点反势、固定节点零势''；
\item
  合并节点后重新统计每两节点之间的实际支路数；
\item
  用节点法、功率或极限情形复核结果。
\end{enumerate}

常见错误有：

\begin{itemize}
\tightlist
\item
  只检查几何外形，没有检查电源或电流端的接法；
\item
  把奇对称下的反势节点当成等势节点合并；
\item
  对折后不重新统计支路，漏掉并联或重复计数；
\item
  把位于对称面上的单条支路误减半；
\item
  使用半电路求出 \(R_{1/2}\) 后，忘记 \(R_{AB}=2R_{1/2}\)；
\item
  透视图中只凭线条交叉判断节点连接，而没有核对节点表或边表。
\end{itemize}

\subsection{六、综合练习一：立方体电阻网络}\label{m4:ux516dux7efcux5408ux7ec3ux4e60ux4e00ux7acbux65b9ux4f53ux7535ux963bux7f51ux7edc}

由 \(12\) 根阻值都为 \(R\) 的电阻构成立方体的 \(12\) 条棱。

\par\addvspace{.4\baselineskip}
\noindent\begin{minipage}{\linewidth}

\subsubsection{作业一：体对角线两端}\label{m4:ux4f5cux4e1aux4e00ux4f53ux5bf9ux89d2ux7ebfux4e24ux7aef}

\textbf{题号：M4-052}

\textbf{层级：} 基础　　\textbf{方法：} 网络方程

一个有限立方体电阻网络仅在12条棱上各接一只 \(R>0\)
电阻，八个顶点是全部节点，面对角线与体对角线上没有支路。取
\(A=(0,0,0),G=(1,1,1)\)
为一对体对角顶点，其余坐标为各分量取0或1。把端口接在
\(A,G\)，求二端等效电阻，写明其余六点的等势分组及各组之间的并联支路数。

\end{minipage}\par

\par\noindent\begin{minipage}{\linewidth}
\textbf{解答}

\begin{center}
\includegraphics[width=\linewidth,height=.40\textheight,keepaspectratio]{\detokenize{assets/module-04/examples/M4-052.pdf}}
\par\smallskip{\small 图 M4-052：体对角线端口把六点分成两组}
\end{center}
\end{minipage}\par

图 M4-052
的红色顶点是体对角线两端，其他六点按距端点的图距离分成三点组。合并等势点后保留的是3、6、3条真实棱边，因此各段阻值分别除以这些条数。

绕 \(AG\) 轴旋转 \(120^{\circ}\) 时，立方体网络与端口激励均不变。与
\(A\) 相邻的三个顶点等势，合并为 \(N_1\)；与 \(G\)
相邻的三个顶点等势，合并为 \(N_2\)。

\(A\) 与 \(N_1\) 之间有三只 \(R\) 并联，\(N_1\) 与 \(N_2\) 之间有六只
\(R\) 并联，\(N_2\) 与 \(G\) 之间有三只 \(R\) 并联。所以

\[
R_{AG}=\frac R3+\frac R6+\frac R3.
\]

故

\[
\boxed{R_{AG}=\frac56R}.
\]

\par\addvspace{.4\baselineskip}
\noindent\begin{minipage}{\linewidth}

\subsubsection{作业二：一条棱的两端}\label{m4:ux4f5cux4e1aux4e8cux4e00ux6761ux68f1ux7684ux4e24ux7aef}

\textbf{题号：M4-053}

\textbf{层级：} 基础　　\textbf{方法：} 网络方程

立方体电阻网络的八个节点坐标为
\((x,y,z)\in\{0,1\}^3\)，且只有相差一个坐标的顶点之间接电阻，即12条棱各为
\(R>0\)。选取相邻顶点 \(A=(0,0,0),B=(1,0,0)\)
为测量端口，其余顶点无外部注入。求
\(R_{AB}\)，在对折或等势合并时核对每两节点组之间真实支路的条数，不将透视图的交叉当作额外节点。

\end{minipage}\par

\par\noindent\begin{minipage}{\linewidth}
\textbf{解答}

\begin{center}
\includegraphics[width=\linewidth,height=.40\textheight,keepaspectratio]{\detokenize{assets/module-04/examples/M4-053.pdf}}
\par\smallskip{\small 图 M4-053：相邻端口不能沿用体对角分组}
\end{center}
\end{minipage}\par

图 M4-053 固定相邻红端点 \(A,B\)。保持这两个端点的镜面对称只允许合并
\(C,F\) 及 \(D,E\)；直连棱 \(AB\)
始终保留在外侧支路的并联位置，不能计入外围后再重复相加。

取坐标

\[
A=(0,0,0),\quad B=(1,0,0),
\]

并将其余顶点依次记为

\[
C=(1,1,0),\ D=(0,1,0),\ E=(0,0,1),\ F=(1,0,1),
\]

\[
G=(1,1,1),\ H=(0,1,1).
\]

变换 \((x,y,z)\mapsto(x,z,y)\) 保持网络与端口 \(A,B\) 不变，它交换

\[
C\leftrightarrow F,
\qquad
D\leftrightarrow E,
\]

而 \(G,H\) 不动。因此

\[
V_C=V_F,
\qquad
V_D=V_E.
\]

合并 \(C,F\) 为 \(Y\)，合并 \(D,E\) 为 \(X\)。对折后：

\[
R_{AX}=\frac R2,
\qquad
R_{YB}=\frac R2,
\]

\par\noindent\begin{minipage}{\linewidth}
\begin{center}
\includegraphics[width=\linewidth,height=.40\textheight,keepaspectratio]{\detokenize{assets/module-04/process/M4-053a.pdf}}
\par\smallskip{\small 图 M4-053a：相邻端口对折后的外围网络}
\end{center}
\end{minipage}\par

图 M4-053a
的实线画外围网络，顶端说明另有直连棱AB。X、Y之间有直接并联边以及经H、G的串联路线，先处理它们，再处理端口上的最终并联。

\[
R_{XY}^{(1)}=\frac R2,
\qquad
R_{XH}+R_{HG}+R_{GY}=\frac R2+R+\frac R2=2R.
\]

所以 \(X,Y\) 之间的等效电阻为

\[
R_{XY}=\frac R2\parallel2R=\frac25R.
\]

删去直连棱 \(AB\) 后，外围网络的等效电阻为

\[
R_{\rm out}=\frac R2+\frac25R+\frac R2=\frac75R.
\]

它与棱 \(AB\) 的电阻 \(R\) 并联，故

\[
\boxed{R_{AB}=R\parallel\frac75R=\frac7{12}R}.
\]

\subsection{七、综合练习二：方格网的等势列}\label{m4:ux4e03ux7efcux5408ux7ec3ux4e60ux4e8cux65b9ux683cux7f51ux7684ux7b49ux52bfux5217}

将 \(m\times m\) 正方格网旋转 \(45^{\circ}\) 画出，每条格边电阻均为
\(r\)，电流从一个角顶点 \(A\) 流入，从对角顶点 \(B\) 流出。

\par\addvspace{.4\baselineskip}
\noindent\begin{minipage}{\linewidth}

\subsubsection{\texorpdfstring{作业三：\(2\times2\)
方格网}{作业三：2\textbackslash times2 方格网}}\label{m4:ux4f5cux4e1aux4e092times2-ux65b9ux683cux7f51}

\textbf{题号：M4-054}

\textbf{层级：} 基础　　\textbf{方法：} 网络方程

有限二维方格网络含节点
\((i,j)\)，\(i,j=0,1,2\)，共9个节点；所有水平、竖直最近邻节点间各接一只
\(r>0\) 电阻，共12条边。边界之外没有其他节点或支路。将端口取为相对角点
\(A=(0,0),B=(2,2)\)，其余节点无外部注入。求等效电阻，说明关于对角线及端点中垂线的对称如何产生等势节点组，并核对合并后的支路数。

\end{minipage}\par

\par\noindent\begin{minipage}{\linewidth}
\textbf{解答}

\begin{center}
\includegraphics[width=\linewidth,height=.40\textheight,keepaspectratio]{\detokenize{assets/module-04/examples/M4-054.pdf}}
\par\smallskip{\small 图 M4-054：二乘二方格的对角端口}
\end{center}
\end{minipage}\par

图 M4-054 明确把端口放在相对角点
\(A,B\)，而非相邻边界点。关于端点轴的偶对称与交换端点的奇对称共同筛出等势层；相邻层间有多少条原始电阻，就并联多少条。

关于 \(AB\) 轴反射是偶对称，所以每一列中关于 \(AB\)
对称的节点等势。关于中垂线的反射交换
\(A,B\)，是奇对称，所以中间列的固定节点为零势；结合前一个对称，中间列的三个节点可合并。

从左到右合并后，相邻节点组之间的并联支路数依次为

\[
2,\quad 4,\quad 4,\quad 2.
\]

因此

\[
R_{AB}=\frac r2+\frac r4+\frac r4+\frac r2,
\]

得

\[
\boxed{R_{AB}=\frac32r}.
\]

\par\addvspace{.4\baselineskip}
\noindent\begin{minipage}{\linewidth}

\subsubsection{\texorpdfstring{作业四：\(3\times3\)
方格网}{作业四：3\textbackslash times3 方格网}}\label{m4:ux4f5cux4e1aux56db3times3-ux65b9ux683cux7f51}

\textbf{题号：M4-055}

\textbf{层级：} 基础　　\textbf{方法：} 网络方程

有限二维方格网络含节点
\((i,j)\)，\(i,j=0,1,2,3\)，共16个节点；所有水平、竖直最近邻间各接一只
\(r>0\) 电阻，共24条边，边界外没有附加连接。端口为相对角点
\(A=(0,0),B=(3,3)\)，其余节点无外部注入。求二端等效电阻，利用对角线偶对称与交换端点的奇对称构造半电路，并说明由半电路恢复全网络电阻的因子。

\end{minipage}\par

\par\noindent\begin{minipage}{\linewidth}
\textbf{解答}

\begin{center}
\includegraphics[width=\linewidth,height=.40\textheight,keepaspectratio]{\detokenize{assets/module-04/examples/M4-055.pdf}}
\par\smallskip{\small 图 M4-055：三乘三方格的半电路}
\end{center}
\end{minipage}\par

图 M4-055
的两红端点是相对角点。先沿端点轴合并镜像节点，再取到奇对称零势列的左半电路；最后必须把半电路电阻乘2，才能恢复图中完整两端的结果。

只计算从 \(A\) 到中间零势列 \(F\) 的左半网络。第二列的两节点等势，合并为
\(C\)，且

\[
R_{AC}=\frac r2.
\]

第三列的两个外侧节点合并为 \(D\)，中间节点记为 \(E\)。从 \(C\) 到零势列
\(F\) 有两条支路：

\[
R_{CDF}=\frac r2+\frac r4=\frac34r,
\]

\par\noindent\begin{minipage}{\linewidth}
\begin{center}
\includegraphics[width=\linewidth,height=.40\textheight,keepaspectratio]{\detokenize{assets/module-04/process/M4-055a.pdf}}
\par\smallskip{\small 图 M4-055a：三乘三方格约化后的左半电路}
\end{center}
\end{minipage}\par

图 M4-055a
的F代表奇对称零势列，图中只有从A到F的一半。C到F的两支路各由原边计数得到，先作并联，再与AC串联；两端电阻还要恢复另一半。

\[
R_{CEF}=\frac r2+\frac r2=r.
\]

所以

\[
R_{CF}=\frac34r\parallel r=\frac37r.
\]

左半电路的电阻为

\[
R_{AF}=\frac r2+\frac37r=\frac{13}{14}r.
\]

完整网络由奇对称的两半组成，因而

\[
\boxed{R_{AB}=2R_{AF}=\frac{13}{7}r}.
\]

\subsection{八、练习与解答}\label{m4:ux516bux7ec3ux4e60ux4e0eux89e3ux7b54-1}

\par\addvspace{.4\baselineskip}
\noindent\begin{minipage}{\linewidth}

\subsubsection{作业五：空间线框网络的边表法}\label{m4:ux4f5cux4e1aux4e94ux7a7aux95f4ux7ebfux6846ux7f51ux7edcux7684ux8fb9ux8868ux6cd5}

\textbf{题号：M4-056}

\textbf{层级：} 基础　　\textbf{方法：} 网络方程

所有未列出的节点对均无支路，带撇与不带撇的字母是不同节点；线条的几何交叉不生成额外电连接。取
\(r>0\)，从 \(A\) 注入测试电流并从 \(B\) 抽出，其余节点无外部注入。

透视图容易使人误判线条是否相交。下面的空间线框由节点

\[
A,B,C,C',D,D',E,E',F,F'
\]

构成，所有列出的支路电阻均为 \(r\)，且仅有下列 \(12\) 条边：

\[
AC,\ AC',\ CE,\ CF,\ C'E',\ C'F',
\]

\[
BD,\ BD',\ DE,\ DF,\ D'E',\ D'F'.
\]

求 \(A,B\) 之间的等效电阻。

\end{minipage}\par

\par\noindent\begin{minipage}{\linewidth}
\textbf{解答}

\begin{center}
\includegraphics[width=\linewidth,height=.40\textheight,keepaspectratio]{\detokenize{assets/module-04/examples/M4-056.pdf}}
\par\smallskip{\small 图 M4-056：边表决定连接，中央四点为零势}
\end{center}
\end{minipage}\par

图 M4-056
逐条绘出题设12条边，带撇节点没有与同名字母合并。交换左右端口的反射固定中央四个节点，使它们成为零势组；左侧两臂再由偶对称并联，形成
\(r/2+r/4\) 的半电路。

左右反射交换

\[
A\leftrightarrow B,
\qquad
C\leftrightarrow D,
\qquad
C'\leftrightarrow D',
\]

而保持 \(E,E',F,F'\)
不动。因为该反射交换流入端与流出端，中央四节点电势均为零，可合并为
\(M\)。

另一反射交换 \(C,C'\) 且保持端口不变，所以 \(V_C=V_{C'}\)，合并为
\(X\)。左半电路中，\(A\) 到 \(X\) 有两只 \(r\) 并联，\(X\) 到 \(M\)
有四只 \(r\) 并联，故

\[
R_{AM}=\frac r2+\frac r4=\frac34r.
\]

由奇对称半电路关系，

\[
\boxed{R_{AB}=2R_{AM}=\frac32r}.
\]

直接按边表也可看出：网络分成上、下两个相同的 \(3r\)
支路网，两者并联同样得 \(3r/2\)。

\subsection{九、练习与解答}\label{m4:ux4e5dux7ec3ux4e60ux4e0eux89e3ux7b54-2}

\par\addvspace{.4\baselineskip}
\noindent\begin{minipage}{\linewidth}

\subsubsection{作业六：对称多边形网络的递推}\label{m4:ux4f5cux4e1aux516dux5bf9ux79f0ux591aux8fb9ux5f62ux7f51ux7edcux7684ux9012ux63a8}

\textbf{题号：M4-057}

\textbf{层级：} 基础主问；一般闭式为选学　　\textbf{方法：} 网络方程

\(N\ge3\)
为整数，\(r>0\)，每一条对顶点支路只计一次。各节点是唯一连接点，几何中心和无节点标记的交叉不相连；测量时仅
\(P_0,P_N\) 有等量异号外部注入。

设有 \(2N\) 个节点

\[
P_0,P_1,\ldots,P_{2N-1},
\]

下标按模 \(2N\) 计算。相邻节点 \(P_kP_{k+1}\) 之间均接电阻
\(r\)，并对每个 \(k=0,1,\ldots,N-1\) 在对顶点 \(P_k,P_{k+N}\) 之间接电阻
\(r\)。(1) 用对称对折和有限递推，分别求 \(N=5\)、\(N=6\) 时 \(P_0,P_N\)
间的等效电阻，说明直接相连的直径电阻如何计入。

\begin{enumerate}
\def\labelenumi{(\arabic{enumi})}
\setcounter{enumi}{1}
\item
  求 \(N\to\infty\) 的极限。
\item
  \textbf{选学：} 求任意整数 \(N\ge3\)
  的闭式。该问可在掌握离散傅里叶模式后阅读；前两问只需对称性、串并联和递推。
\end{enumerate}

\end{minipage}\par

\par\noindent\begin{minipage}{\linewidth}
\textbf{解答}

\begin{center}
\includegraphics[width=\linewidth,height=.40\textheight,keepaspectratio]{\detokenize{assets/module-04/examples/M4-057.pdf}}
\par\smallskip{\small 图 M4-057：周界与对顶支路要分开计数}
\end{center}
\end{minipage}\par

图 M4-057 示意 \(N=5\)
的十节点网络，直径线在中央的交叉没有节点。沿端点轴对折后，周界支路成对并联；直连的
\(P_0-P_N\) 电阻始终另外保留，有限递推的末端条件还要区分奇偶。

先关于端点轴 \(P_0P_N\) 对折，镜像周界电阻变为两两并联的
\(r/2\)。再利用交换 \(P_0,P_N\)
的奇对称取半电路。为了直接表示完整网络的电阻，将半电路中的所有电阻同时乘
\(2\)：每节主干为 \(r\)，每个中间节点对零势点的旁路为 \(r/2\)。直连
\(P_0P_N\) 的直径电阻 \(r\) 最后单独并联。

\par\noindent\begin{minipage}{\linewidth}
\begin{center}
\includegraphics[width=\linewidth,height=.40\textheight,keepaspectratio]{\detokenize{assets/module-04/process/M4-057a.pdf}}
\par\smallskip{\small 图 M4-057a：有限对折的末端条件区分奇偶}
\end{center}
\end{minipage}\par

图 M4-057a 画出对折后的一次尾部替换：本级的直接支路为
\(r/2\)，另一支路串联 \(r\) 与尾负载
\(X_{k+1}\)。右侧列出的有限末端条件按 \(N\) 的奇偶选用；原网络直连
\(P_0,P_N\) 的直径电阻仍在外围计算后另行并联。

\paragraph{\texorpdfstring{1. \(N=5\)}{1. N=5}}\label{m4:n5}

最深节点的两条 \(r/2\) 零势旁路并联，故

\[
X_2=\frac r4.
\]

向外递推一节：

\[
X_1=(r+X_2)\parallel\frac r2
=\frac54r\parallel\frac r2
=\frac5{14}r.
\]

除去直连直径后，外围网络的电阻为

\[
R_{\rm out}=r+X_1=\frac{19}{14}r.
\]

所以

\[
\boxed{R_{P_0P_5}=r\parallel\frac{19}{14}r=\frac{19}{33}r}.
\]

\paragraph{\texorpdfstring{2. \(N=6\)}{2. N=6}}\label{m4:n6}

最深处的主干 \(r\) 与旁路 \(r/2\) 同时接向零势点，因而

\[
X_2=r\parallel\frac r2=\frac r3.
\]

再向外递推：

\[
X_1=(r+X_2)\parallel\frac r2
=\frac43r\parallel\frac r2
=\frac4{11}r.
\]

因此

\[
\boxed{R_{P_0P_6}=r\parallel\left(r+\frac4{11}r\right)
=\frac{15}{26}r}.
\]

\paragraph{\texorpdfstring{3.
\(N\to\infty\)}{3. N\textbackslash to\textbackslash infty}}\label{m4:ntoinfty}

设从某个中间节点向无穷远看到的等效电阻为 \(X\)，自相似性给出

\[
X=(r+X)\parallel\frac r2.
\]

化简得

\[
2X^2+2rX-r^2=0.
\]

取正根：

\[
X=\frac{\sqrt3-1}{2}r.
\]

最后与直连直径 \(r\) 并联：

\[
R_{\infty}=r\parallel(r+X),
\]

因而

\[
\boxed{R_{\infty}=\frac r{\sqrt3}}.
\]

\textbf{一般 \(N\) 的闭式}

图 M4-057 的周界平移和对顶连接还允许直接作离散傅里叶分解。以每条边电导
\(1/r\) 为尺度，周界模式 \(k=0,\ldots,2N-1\) 的无量纲本征值为

\[
\lambda_k=2-2\cos\frac{\pi k}{N}+1-(-1)^k.
\]

端口 \(P_0,P_N\) 的注入向量是 \(\mathbf e_0-\mathbf e_N\)，它在偶数
\(k\) 模式上投影为零，在奇数模式上的投影模方为 \(2/N\)。因此

\[
\frac{R_{P_0P_N}}r
=\frac1N\sum_{j=0}^{N-1}
\frac1{2-\cos[(2j+1)\pi/N]}.
\]

令 \(\rho=2-\sqrt3\)。利用 \(2=(\rho+\rho^{-1})/2\) 及收敛的泊松核展开，

\[
\frac1{2-\cos\theta}
=\frac1{\sqrt3}\left[1+2\sum_{s=1}^{\infty}\rho^s\cos(s\theta)\right].
\]

对上述 \(N\) 个奇半周期角求平均，只有 \(s\) 为 \(N\)
的倍数时保留，且对应余弦平均为 \((-1)^{s/N}\)，从而

\[
\boxed{
R_{P_0P_N}=\frac r{\sqrt3}
\frac{1-\rho^N}{1+\rho^N},\qquad \rho=2-\sqrt3
}.
\]

代入 \(N=5,6\) 分别回到 \(19r/33,15r/26\)，而 \(N\to\infty\) 时回到
\(r/\sqrt3\)。这给出了题设任意整数 \(N\ge3\)
的结果，而不只列举两个有限样例。

\subsection{十、使用条件}\label{m4:ux5341ux4f7fux7528ux6761ux4ef6}

\begin{enumerate}
\def\labelenumi{\arabic{enumi}.}
\tightlist
\item
  对称性必须同时保持网络的拓扑、元件参数和端口激励，否则不能推出等势。
\item
  偶对称给出等势节点；奇对称给出反势节点和零势固定点，两者不可混用。
\item
  合并等势节点不保留新增理想导线中的局部电流分配；在用于证明等势的指定激励下，它保持原有其他支路的状态与该次测量结果。
\item
  对折法是节点方程的降阶表示，并不是将实际电路在空间中压在一起。
\item
  对含受控源、非互易元件或参数不对称的网络，应直接检查节点方程在相应置换下是否不变。
\item
  在某一指定激励下得到的等势点合并，不能自动推广到改变端口后的另一个测量问题。
\end{enumerate}

\section{专题三十六：电流分布法与叠加原理（基础必讲与拔高选讲）}\label{m4:ux4e13ux9898ux4e09ux5341ux516dux7535ux6d41ux5206ux5e03ux6cd5ux4e0eux53e0ux52a0ux539fux7406ux57faux7840ux5fc5ux8bb2ux4e0eux62d4ux9ad8ux9009ux8bb2}

对高度对称的电阻网络，直接在两个指定端点之间通电，往往会破坏大部分对称性。电流分布法的核心是：先构造若干个更对称、且满足电流守恒的辅助状态，再利用线性叠加得到目标两端状态。

\subsection{一、叠加原理究竟叠加什么（基础必讲）}\label{m4:ux4e00ux53e0ux52a0ux539fux7406ux7a76ux7adfux53e0ux52a0ux4ec0ux4e48ux57faux7840ux5fc5ux8bb2}

对拓扑、元件参数不变的线性电路，选定参考电势后，方程可写成

\[
A\mathbf x=\mathbf b.
\]

\(\mathbf x\) 可包含节点电势和支路电流，\(\mathbf b\)
表示独立源或外部注入电流。若

\[
A\mathbf x^{(1)}=\mathbf b^{(1)},
\qquad
A\mathbf x^{(2)}=\mathbf b^{(2)},
\]

则对任意常数 \(\alpha,\beta\)，

\[
A\left(\alpha\mathbf x^{(1)}+\beta\mathbf x^{(2)}\right)
=\alpha\mathbf b^{(1)}+\beta\mathbf b^{(2)}.
\]

因而可线性叠加的是：

\begin{itemize}
\tightlist
\item
  节点电势和两点电势差；
\item
  支路电流与外部端口电流；
\item
  在线性电容、电感网络中，相应的电荷或磁通链响应。
\end{itemize}

不能直接叠加的是：

\begin{itemize}
\tightlist
\item
  功率 \(P=UI=I^2R\)，因为它是状态量的二次式；
\item
  不同端口问题的等效电阻；
\item
  非线性元件中电压、电流的一般解。
\end{itemize}

\par\noindent\begin{minipage}{\linewidth}
\begin{center}
\includegraphics[width=\linewidth,height=.40\textheight,keepaspectratio]{\detokenize{assets/module-04/original/tikz-m4-topic36-v9-from33-v8-01.png}}
\par\smallskip{\small 图示：同一线性网络中电势差与支路电流可叠加，功率和等效电阻不能如此相加}
\end{center}
\end{minipage}\par

对含多个独立源的电路，``其他源置零''的规则是：

\[
\text{理想电压源置零}\Longrightarrow\text{短路},
\]

\[
\text{理想电流源置零}\Longrightarrow\text{开路}.
\]

受控源是线性方程的一部分，分别计算独立源作用时不能擅自删去。

\par\addvspace{.4\baselineskip}
\noindent\begin{minipage}{\linewidth}

\textbf{例 M4-058：电流可叠加而功率不能直接相加}

同一个线性电阻网络在两组独立源激励下形成状态一、二，某一只电阻阻值为
\(R>0\)，按同一参考方向的支路电流分别为
\(I_1,I_2\)。把两组源按线性叠加合在一起。求该电阻在叠加态的电流、同向电压及耗散功率，判断功率能否直接相加，并指出可能出现的交叉项。

\end{minipage}\par

\par\noindent\begin{minipage}{\linewidth}
\textbf{解答}

\begin{center}
\includegraphics[width=\linewidth,height=.40\textheight,keepaspectratio]{\detokenize{assets/module-04/examples/M4-058.pdf}}
\par\smallskip{\small 图 M4-058：同一参考方向叠加后再算功率}
\end{center}
\end{minipage}\par

图 M4-058 的三根箭头共用 \(A\to B\)
的参考方向，表示两个辅助态及其叠加。电压、电流按有符号量线性相加；功率必须把红色总电流整体平方，交叉项正是不能直接相加的原因。

同一只电阻 \(R\) 在状态一、二中的支路电流分别为
\(I_1,I_2\)，则叠加态的电压为

\[
U=I_1R+I_2R=(I_1+I_2)R.
\]

但功率为

\[
P=(I_1+I_2)^2R,
\]

一般不等于 \(I_1^2R+I_2^2R\)。

\subsection{二、节点注入向量与平衡辅助态（拔高选讲）}\label{m4:ux4e8cux8282ux70b9ux6ce8ux5165ux5411ux91cfux4e0eux5e73ux8861ux8f85ux52a9ux6001ux62d4ux9ad8ux9009ux8bb2}

对连通的纯电阻网络，设节点电势为 \(\mathbf V\)，外部注入电流为
\(\mathbf I\)，则

\[
L\mathbf V=\mathbf I,
\]

\(L\) 是节点电导矩阵。稳恒状态必须满足全网络的电流守恒：

\[
\boxed{\sum_{k=1}^{N}I_k=0}.
\]

所以在有限孤立网络中，不能只在一个节点注入电流而不指明回流。最常用的高对称辅助态是：

\[
\begin{aligned}
&\text{从 }A\text{ 注入 }I,\\
&\text{从其余 }N-1\text{ 个节点各抽出 }\frac{I}{N-1}.
\end{aligned}
\]

该状态的注入向量之和恰为零，因而是合法的稳恒状态。

再构造第二个状态：

\[
\begin{aligned}
&\text{从 }B\text{ 抽出 }I,\\
&\text{向其余 }N-1\text{ 个节点各注入 }\frac{I}{N-1}.
\end{aligned}
\]

两状态叠加后：

\[
I_A^{\text{外}}=I+\frac{I}{N-1}=\frac{N}{N-1}I,
\]

\[
I_B^{\text{外}}=-I-\frac{I}{N-1}=-\frac{N}{N-1}I,
\]

其余节点的外部注入均为零。因而叠加态正是 \(A,B\) 两端测量态。

\par\noindent\begin{minipage}{\linewidth}
\begin{center}
\includegraphics[width=\linewidth,height=.40\textheight,keepaspectratio]{\detokenize{assets/module-04/theory/review-tikz-m4-topic36-v9-from33-v8-02.pdf}}
\par\smallskip{\small 图示：两个有限平衡辅助态叠加后，其余节点的外部注入两两抵消}
\end{center}
\end{minipage}\par

\subsection{三、高对称有限网络的相邻节点公式（拔高选讲）}\label{m4:ux4e09ux9ad8ux5bf9ux79f0ux6709ux9650ux7f51ux7edcux7684ux76f8ux90bbux8282ux70b9ux516cux5f0fux62d4ux9ad8ux9009ux8bb2}

设有限连通网络共有 \(N\) 个节点，每个节点都连出 \(d\) 条阻值为 \(r\)
的边。设 \(A,B\) 为相邻节点，并且满足：

\begin{enumerate}
\def\labelenumi{\arabic{enumi}.}
\tightlist
\item
  固定 \(A\) 的网络对称操作可把 \(A\) 的 \(d\) 个相邻节点任意互换；
\item
  固定 \(B\) 时也有同样的局部对称性。
\end{enumerate}

在``从 \(A\) 注入 \(I\)、其余各点均匀抽出''的辅助态中，经由 \(A\) 的
\(d\) 条边完全等价，所以

\[
I_{AB}^{(1)}=\frac Id.
\]

在``从 \(B\) 抽出 \(I\)、其余各点均匀注入''的辅助态中，沿 \(A\to B\)
的电流同样为

\[
I_{AB}^{(2)}=\frac Id.
\]

叠加后，直接支路 \(AB\) 上的电流为

\[
I_{AB}=\frac{2I}{d},
\]

所以两端电压为

\[
U_{AB}=\frac{2I}{d}r.
\]

实际端口电流为 \(NI/(N-1)\)，因而

\[
\boxed{
R_{AB}=\frac{U_{AB}}{NI/(N-1)}
=\frac{2(N-1)}{Nd}r
}.
\]

\par\noindent\begin{minipage}{\linewidth}
\begin{center}
\includegraphics[width=\linewidth,height=.40\textheight,keepaspectratio]{\detokenize{assets/module-04/theory/review-tikz-m4-topic36-v9-from33-v8-03.pdf}}
\par\smallskip{\small 图示：相邻节点公式来自辅助态的局部均分和叠加后的端口归一化}
\end{center}
\end{minipage}\par

这一公式不是对任意 \(d\)
正则网络都成立；``每点度数相同''不足以保证各相邻边电流均分，还必须有上述的相邻点置换对称性。

\par\addvspace{.4\baselineskip}
\noindent\begin{minipage}{\linewidth}

\textbf{例 M4-059：有限环网络的相邻节点电阻}

\(N\ge3\) 个节点首尾相接成单环 \(C_N\)，环上每条边接一只 \(r>0\)
电阻，除此之外没有支路。端口选任意相邻两节点，其余节点悬空。通过识别从端口出发的两条独立通路求等效电阻，再与有限边传递连通网络的相邻端口公式比较，核对节点数、边数与度数的对应。

\end{minipage}\par

\par\noindent\begin{minipage}{\linewidth}
\textbf{解答}

\begin{center}
\includegraphics[width=\linewidth,height=.40\textheight,keepaspectratio]{\detokenize{assets/module-04/examples/M4-059.pdf}}
\par\smallskip{\small 图 M4-059：环网络的相邻端口有两条通路}
\end{center}
\end{minipage}\par

图 M4-059 用六节点环示意一般
\(C_N\)，红端点之间有直接电阻与绕余下环的串联通路。中间节点均无注入，因此两路并联；只用环的几何对称而忽略另一通路会把支路值误当作测量值。

\(N\) 个阻值为 \(r\) 的电阻构成环形链 \(C_N\)，每点度数
\(d=2\)。任意一条边的两端之间，一条通路为 \(r\)，另一条通路为
\((N-1)r\)，所以

\[
R_{\text{邻}}=r\parallel(N-1)r=\frac{N-1}{N}r.
\]

通式也给出

\[
\frac{2(N-1)}{2N}r=\frac{N-1}{N}r.
\]

\subsection{四、无限网络必须通过有限网络取极限（拔高选讲）}\label{m4:ux56dbux65e0ux9650ux7f51ux7edcux5fc5ux987bux901aux8fc7ux6709ux9650ux7f51ux7edcux53d6ux6781ux9650ux62d4ux9ad8ux9009ux8bb2}

在二维无限方格网中，``从一点注入电流，从无穷远抽出''的单源电势随系统尺度对数发散。因而，不应把单源态本身当成具有有限端口电压的普通物理状态。

一种严格处理是先取 \(m\times m\)
的周期方格网：把左、右边界对应节点连接，再把上、下边界对应节点连接。得到的有限周期网络共有

\[
N=m^2
\]

个节点，每点度数

\[
d=4.
\]

对
\(m\ge5\)，每个节点的四条相邻边在网络对称下等价。上一节的公式给出周期网络相邻节点间的精确结果：

\[
R_m=\frac{2(m^2-1)}{4m^2}r
=\frac{m^2-1}{2m^2}r.
\]

令 \(m\to\infty\)，得二维无限方格网中相邻节点的等效电阻

\[
\boxed{R_{\square,\text{邻}}=\frac r2}.
\]

\par\noindent\begin{minipage}{\linewidth}
\begin{center}
\includegraphics[width=\linewidth,height=.40\textheight,keepaspectratio]{\detokenize{assets/module-04/theory/review-tikz-m4-topic36-v9-from33-v8-04.pdf}}
\par\smallskip{\small 图示：先在有限周期方格上构造平衡辅助态，再取尺度无穷大的极限}
\end{center}
\end{minipage}\par

这个证明只使用了有限、平衡的电流注入态。直观上常说的``从无穷远均匀回流''，应理解为上述有限构造的极限，而不是额外增加了一个真实的``无穷远节点''。

\subsection{五、电流分布法的计算步骤（基础必讲）}\label{m4:ux4e94ux7535ux6d41ux5206ux5e03ux6cd5ux7684ux8ba1ux7b97ux6b65ux9aa4ux57faux7840ux5fc5ux8bb2}

\begin{enumerate}
\def\labelenumi{\arabic{enumi}.}
\tightlist
\item
  确认网络在相同元件参数下是线性的；
\item
  明确目标端口和所求的是电流、电压还是等效电阻；
\item
  构造每个节点总注入为零的有限辅助态；
\item
  只对在辅助态下确实等价的支路使用``均分''；
\item
  叠加后逐点检查外部注入，确认其余节点净注入均为零；
\item
  沿真实电阻支路累加电压降，不能沿电路图上并不存在的几何连线取压；
\item
  最后用叠加后的实际端口电流归一化：
\end{enumerate}

\[
R_{AB}=\frac{U_{AB}}{I_{\text{端}}}.
\]

\par\noindent\begin{minipage}{\linewidth}
\begin{center}
\includegraphics[width=\linewidth,height=.40\textheight,keepaspectratio]{\detokenize{assets/module-04/original/tikz-m4-topic36-v9-from33-v8-05.png}}
\par\smallskip{\small 图示：电流分布法应依次检查线性、平衡注入、对称均分、叠加和端口归一化}
\end{center}
\end{minipage}\par

\subsection{六、练习与解答}\label{m4:ux516dux7ec3ux4e60ux4e0eux89e3ux7b54}

\par\addvspace{.4\baselineskip}
\noindent\begin{minipage}{\linewidth}

\subsubsection{作业一：无限方格网相邻节点}\label{m4:ux4f5cux4e1aux4e00ux65e0ux9650ux65b9ux683cux7f51ux76f8ux90bbux8282ux70b9}

\textbf{题号：M4-060}

\textbf{层级：} 复赛　　\textbf{方法：} 网络方程、周期极限

\(r>0\)，节点可用 \(\mathbb Z^2\)
标记，只连接水平、竖直最近邻。用越来越大的平衡有限周期网格确定极限，端口附近保持原有局部连接，不另规定不相容的单点向无穷远输运电流。

无限二维正方格网的每条格边电阻均为 \(r\)。求相邻节点 \(A,B\)
之间的等效电阻。

\end{minipage}\par

\par\noindent\begin{minipage}{\linewidth}
\textbf{解答}

\begin{center}
\includegraphics[width=\linewidth,height=.40\textheight,keepaspectratio]{\detokenize{assets/module-04/examples/M4-060.pdf}}
\par\smallskip{\small 图 M4-060：无限方格网的相邻节点}
\end{center}
\end{minipage}\par

图 M4-060
以相邻红端点标记源与汇，并用有限周期方格作为无限网络的定义。先在有限网中使用平衡注入与对称电流，再取尺寸极限，避免把单点向无穷远输运的势误当两点电阻。

取 \(m\times m\) 周期方格作为有限近似，\(m\ge5\)。该网络有 \(m^2\)
个节点，每点度数为 \(4\)。

在状态一中，从 \(A\) 注入 \(I\)，从其余 \(m^2-1\) 个节点各抽出

\[
\frac{I}{m^2-1}.
\]

由 \(A\) 点周围的四重对称性，\(A\) 流向 \(B\) 的电流为

\[
I_{AB}^{(1)}=\frac I4.
\]

在状态二中，从 \(B\) 抽出 \(I\)，向其余各点均匀注入 \(I/(m^2-1)\)。沿
\(A\to B\) 的电流为

\[
I_{AB}^{(2)}=\frac I4.
\]

叠加后，

\[
I_{AB}=\frac I2,
\]

因而

\[
U_{AB}=\frac I2r.
\]

但叠加后的端口电流不是 \(I\)，而是

\[
I_{\text{端}}=I+\frac{I}{m^2-1}
=\frac{m^2}{m^2-1}I.
\]

所以有限周期网络中

\[
R_m
=\frac{Ir/2}{m^2I/(m^2-1)}
=\frac{m^2-1}{2m^2}r.
\]

令 \(m\to\infty\)，得

\[
\boxed{R_{AB}=\frac r2}.
\]

若只在无限图上口头说``单点注入电流从无穷远回流''，必须补充说明这是有限平衡辅助态的极限记号。

\subsection{七、练习与解答}\label{m4:ux4e03ux7ec3ux4e60ux4e0eux89e3ux7b54-4}

\par\addvspace{.4\baselineskip}
\noindent\begin{minipage}{\linewidth}

\subsubsection{作业二：立方体的相邻顶点与面对角顶点}\label{m4:ux4f5cux4e1aux4e8cux7acbux65b9ux4f53ux7684ux76f8ux90bbux9876ux70b9ux4e0eux9762ux5bf9ux89d2ux9876ux70b9}

\textbf{题号：M4-061}

\textbf{层级：} 基础　　\textbf{方法：} 网络方程、对称性

\(r>0\)，八个顶点是全部节点，只有12条棱支路；每次测量分别向一端注入、从另一端抽出等量测试电流，其余节点均无外部注入。

立方体的 \(12\) 条棱均由阻值为 \(r\) 的电阻构成。取

\[
A=(0,0,0),\quad B=(1,0,0),\quad C=(1,1,0),
\]

且 \(A,C\) 为同一面的对角顶点。分别求 \(R_{AB}\) 与 \(R_{AC}\)。

\end{minipage}\par

\par\noindent\begin{minipage}{\linewidth}
\textbf{解答}

\begin{center}
\includegraphics[width=\linewidth,height=.40\textheight,keepaspectratio]{\detokenize{assets/module-04/examples/M4-061.pdf}}
\par\smallskip{\small 图 M4-061：改变立方体端点就改变等势分组}
\end{center}
\end{minipage}\par

图 M4-061 先以红点标出相邻端口 \(A,B\)；面对角端口 \(A,C\) 另见图
M4-061a。两次激励保持的对称操作不同，所以各自重做等势分组；不能只因网络没有改变便沿用同一个约化电路。

\paragraph{\texorpdfstring{1. 相邻顶点
\(A,B\)}{1. 相邻顶点 A,B}}\label{m4:ux76f8ux90bbux9876ux70b9-ab}

立方体共有

\[
N=8
\]

个顶点，每个顶点都与

\[
d=3
\]

个顶点相邻，且固定一个顶点时，三个相邻顶点完全等价。相邻节点公式给出

\[
R_{AB}=\frac{2(8-1)}{8\times3}r.
\]

因而

\[
\boxed{R_{AB}=\frac7{12}r}.
\]

\paragraph{\texorpdfstring{2. 面对角顶点
\(A,C\)}{2. 面对角顶点 A,C}}\label{m4:ux9762ux5bf9ux89d2ux9876ux70b9-ac}

通式只适用于有直接电阻边连接的相邻节点，\(A,C\)
之间没有直接电阻，必须沿真实路径取电压。

在状态一中，从 \(A\) 注入 \(I\)，从其余七个顶点各抽出 \(I/7\)。从 \(A\)
出发的三条棱等价，故

\[
I_{AB}^{(1)}=\frac I3.
\]

节点 \(B\) 自身向外部抽出 \(I/7\)。从 \(B\) 通往与 \(A\) 拓扑距离为
\(2\) 的两个顶点的支路等价，因而

\[
I_{BC}^{(1)}
=\frac12\left(\frac I3-\frac I7\right)
=\frac{2I}{21}.
\]

沿真实路径 \(A\to B\to C\) 计算：

\[
U_{AC}^{(1)}
=\left(\frac13+\frac{2}{21}\right)Ir
=\frac37Ir.
\]

\par\noindent\begin{minipage}{\linewidth}
\begin{center}
\includegraphics[width=\linewidth,height=.40\textheight,keepaspectratio]{\detokenize{assets/module-04/process/M4-061a.pdf}}
\par\smallskip{\small 图 M4-061a：面对角端口对应的另一组等势节点}
\end{center}
\end{minipage}\par

图 M4-061a
将测试端改为同一面的对角顶点。红点的改变也改变保持激励的对称操作，下面应按这一端口重新合并节点，而不是把相邻顶点实验中的分组直接套用。

状态二从 \(C\) 抽出 \(I\)，向其余七点各注入 \(I/7\)。由立方体的对称性，

\[
U_{AC}^{(2)}=\frac37Ir.
\]

叠加后

\[
U_{AC}=\frac67Ir,
\]

而实际端口电流为

\[
I_{\text{端}}=I+\frac I7=\frac87I.
\]

所以

\[
\boxed{R_{AC}=\frac{6Ir/7}{8I/7}=\frac34r}.
\]

\subsection{八、练习与解答}\label{m4:ux516bux7ec3ux4e60ux4e0eux89e3ux7b54-2}

\par\addvspace{.4\baselineskip}
\noindent\begin{minipage}{\linewidth}

\subsubsection{作业三：正十二面体相邻顶点}\label{m4:ux4f5cux4e1aux4e09ux6b63ux5341ux4e8cux9762ux4f53ux76f8ux90bbux9876ux70b9}

\textbf{题号：M4-062}

\textbf{层级：} 基础　　\textbf{方法：} 网络方程、对称性

仅把正十二面体的20个顶点与30条棱作为节点和支路，全部电阻为正且相同，面内无对角线连接。测试电流只在所选相邻端点流入、流出。

正十二面体的每条棱由阻值为 \(r\) 的电阻构成。求相邻顶点 \(A,B\)
之间的等效电阻。

\end{minipage}\par

\par\noindent\begin{minipage}{\linewidth}
\textbf{解答}

\begin{center}
\includegraphics[width=\linewidth,height=.40\textheight,keepaspectratio]{\detokenize{assets/module-04/examples/M4-062.pdf}}
\par\smallskip{\small 图 M4-062：正十二面体的边传递拓扑}
\end{center}
\end{minipage}\par

图 M4-062
是20节点、30边的正十二面体连接示意，红点选一条真实边的两端。每条边在网络自同构下等价，所以可把福斯特求和中的30项取为同一测量电阻；这一步依赖边等价性。

正十二面体有

\[
N=20
\]

个顶点，每个顶点连出

\[
d=3
\]

条棱，且固定任一顶点时，三个相邻顶点在旋转对称下等价。

在从 \(A\) 注入 \(I\)、其余 \(19\) 个顶点各抽出 \(I/19\) 的状态中，

\[
I_{AB}^{(1)}=\frac I3.
\]

在从 \(B\) 抽出 \(I\)、其余 \(19\) 个顶点各注入 \(I/19\) 的状态中，

\[
I_{AB}^{(2)}=\frac I3.
\]

叠加后

\[
U_{AB}=\frac{2I}{3}r,
\]

而端口电流为

\[
I_{\text{端}}=I+\frac I{19}=\frac{20}{19}I.
\]

因而

\[
\boxed{
R_{AB}=\frac{2Ir/3}{20I/19}=\frac{19}{30}r
}.
\]

也可直接使用相邻节点公式：

\[
R_{AB}=\frac{2(20-1)}{20\times3}r=\frac{19}{30}r.
\]

\subsection{九、使用条件}\label{m4:ux4e5dux4f7fux7528ux6761ux4ef6-1}

\begin{enumerate}
\def\labelenumi{\arabic{enumi}.}
\tightlist
\item
  叠加要求网络方程线性，且各辅助状态的拓扑与元件参数完全相同。
\item
  电压和电流可叠加，功率与不同端口的等效电阻不能直接叠加。
\item
  ``均分''必须由网络与当前辅助激励的共同对称性保证，不能只看节点的度数。
\item
  有限孤立网络中的注入电流之和必须为零；无限网络中的``无穷远回流''应用有限平衡网络的极限定义。
\item
  两端没有直接电阻支路时，应沿实际存在的路径累加支路电压，不能沿几何对角线直接写
  \(U=Ir\)。
\item
  相邻节点通式依赖对每个端点的相邻边置换对称；普通正则图不一定满足此条件。
\item
  对边等价的有限均匀网络，本专题的公式还可由下一专题的电阻迹定理直接得到。
\end{enumerate}

\section{专题三十七：电阻网络的福斯特迹定理（拔高选讲）}\label{m4:ux4e13ux9898ux4e09ux5341ux4e03ux7535ux963bux7f51ux7edcux7684ux798fux65afux7279ux8ff9ux5b9aux7406ux62d4ux9ad8ux9009ux8bb2}

\begin{quote}
阅读准备：节点方程可以写成一个矩阵
对每个非参考节点写基尔霍夫电流定律，可集中成
\(\mathbf L\mathbf V=\mathbf I\)。矩阵 \(\mathbf L\)
的对角元是相邻电导之和，非对角元是相连电导的负值；它记录的是网络连接关系，而不是另加的一条物理定律。
\end{quote}

\par\addvspace{.4\baselineskip}
\noindent\begin{minipage}{\linewidth}

\textbf{例 M4-063：一条支路对节点矩阵的贡献}

纯电阻网络中有一条连接节点1、2的支路，电导为 \(g>0\)，两节点电势为
\(V_1,V_2\)。电流正向取
\(1\to2\)，节点方程把外部注入量作为右端。求这条支路电流，并写出它对两节点拉普拉斯矩阵的
\(2\times2\) 子块贡献；核对对角元、非对角元和零行和。

\end{minipage}\par

\par\noindent\begin{minipage}{\linewidth}
\textbf{解答}

\begin{center}
\includegraphics[width=\linewidth,height=.40\textheight,keepaspectratio]{\detokenize{assets/module-04/examples/M4-063.pdf}}
\par\smallskip{\small 图 M4-063：一条支路同时贡献两行节点方程}
\end{center}
\end{minipage}\par

图 M4-063
的电流从节点1指向节点2。它在节点1为流出项，在节点2为反向流出项，所以同一电阻对节点矩阵贡献对称的正对角、负非对角子块，并且行和为零。

按流出节点1为正，\(I_{12}=g(V_1-V_2)\)；对节点2的流出为其相反数。因此支路对节点矩阵的贡献为
\(g\begin{pmatrix}1&-1\\-1&1\end{pmatrix}\)。行和为零，表示同时平移两个电势不会产生电流。

福斯特迹定理不逐一计算每条边两端的等效电阻，而考察所有真实边对应结果的加权和。这个总和只由网络的节点数决定；专题三十六的对称辅助态方法可作为具体网络中的对照。

\subsection{一、定理的准确表述（拔高选讲）}\label{m4:ux4e00ux5b9aux7406ux7684ux51c6ux786eux8868ux8ff0ux62d4ux9ad8ux9009ux8bb2}

设有限连通电阻网络共有 \(N\) 个节点，真实边的集合为 \(E\)。第 \(e\)
条边的电阻为 \(r_e\)，其两个端点为 \(u_e,v_e\)。保留整张网络，从端点
\(u_e,v_e\) 向外看见的等效电阻记为

\[
R_e=R_{\mathrm{eff}}(u_e,v_e).
\]

则有

\[
\boxed{\sum_{e\in E}\frac{R_e}{r_e}=N-1}.
\]

等价地，若边电导 \(c_e=1/r_e\)，则

\[
\boxed{\sum_{e\in E}c_eR_e=N-1}.
\]

这就是有限连通电阻网络的福斯特迹定理。竞赛中常用的``\(N\)
点定理''是它在所有边均为 \(r\) 时的推论：

\[
\boxed{\sum_{e\in E}R_e=(N-1)r}.
\]

必须同时看清三个对象：

\begin{enumerate}
\def\labelenumi{\arabic{enumi}.}
\tightlist
\item
  求和只遍历网络中真实存在的边，不遍历任意节点对；
\item
  \(r_e\) 是第 \(e\) 条边自身的电阻；
\item
  \(R_e\) 是保留整个网络后，这条边两端的全网等效电阻。
\end{enumerate}

\par\noindent\begin{minipage}{\linewidth}
\begin{center}
\includegraphics[width=\linewidth,height=.40\textheight,keepaspectratio]{\detokenize{assets/module-04/original/tikz-m4-foster-01-theorem-object-pdf-v11.jpg}}
\par\smallskip{\small 图示：福斯特迹定理只沿网络中真实存在的边求和}
\end{center}
\end{minipage}\par

因为第 \(e\) 条直接电阻本身也是端点之间的一条通路，所以

\[
0<R_e\le r_e,
\]

从而每个无量纲量 \(R_e/r_e\) 都在 \((0,1]\) 内。若第 \(e\)
条边是桥，则删去它后两端不再连通，故

\[
R_e=r_e.
\]

\par\addvspace{.4\baselineskip}
\noindent\begin{minipage}{\linewidth}

\textbf{例 M4-064：树网络中的福斯特求和}

某有限连通无源电阻网络是一棵含 \(N\ge2\) 个节点的树，每条边 \(e\)
的电阻为任意正数 \(r_e\)。用 \(R_e\)
表示在保持整个网络存在、其余节点无外部注入时测得的该边两端等效电阻。求各比值
\(R_e/r_e\)
及其全边求和，说明树的桥边性质如何使福斯特迹定理退化为一个拓扑计数。

\end{minipage}\par

\par\noindent\begin{minipage}{\linewidth}
\textbf{解答}

\begin{center}
\includegraphics[width=\linewidth,height=.40\textheight,keepaspectratio]{\detokenize{assets/module-04/examples/M4-064.pdf}}
\par\smallskip{\small 图 M4-064：树上的每条边都是唯一通路}
\end{center}
\end{minipage}\par

图 M4-064
的树没有闭环。以任一边两端作端口时，删除该边后的所有旁枝均不能提供替代通路，因而所测电阻就是该边本身；每一项为1，最后只计真实边的条数。

一棵含 \(N\) 个节点的树恰有 \(N-1\) 条边，每条边都是桥。因此

\[
\sum_{e\in E}\frac{R_e}{r_e}
=\sum_{e\in E}1
=N-1,
\]

定理在树上退化为``树有 \(N-1\) 条边''。

\subsection{二、每一项的物理意义（拔高选讲）}\label{m4:ux4e8cux6bcfux4e00ux9879ux7684ux7269ux7406ux610fux4e49ux62d4ux9ad8ux9009ux8bb2}

在第 \(e\) 条边的两个端点之间接入单位电流源：从 \(u_e\) 注入
\(1\,\mathrm A\)，从 \(v_e\) 抽出 \(1\,\mathrm A\)。此时端口电压就是

\[
U_e=R_e.
\]

第 \(e\) 条直接电阻上的电流为

\[
I_e^{\mathrm{\text{直}}}=\frac{U_e}{r_e}=\frac{R_e}{r_e}.
\]

因此，\(R_e/r_e\)
可以理解为：在这次单位端口实验中，总电流里有多大比例直接流过被考察的那条边。

\begin{itemize}
\tightlist
\item
  若该边是桥，全部电流都必须经过它，比例为 \(1\)；
\item
  若两端还有许多旁路，直接边只承担一部分电流，比例小于 \(1\)；
\item
  福斯特迹定理说，把每条边分别做一次这样的实验，所有``直接边电流比例''之和恒为
  \(N-1\)。
\end{itemize}

\par\noindent\begin{minipage}{\linewidth}
\begin{center}
\includegraphics[width=\linewidth,height=.40\textheight,keepaspectratio]{\detokenize{assets/module-04/theory/review-tikz-m4-topic37-v9-from34-v8-01.pdf}}
\par\smallskip{\small 图示：R\_e/r\_e 是单位端口实验中直接边承担的电流比例}
\end{center}
\end{minipage}\par

\par\addvspace{.4\baselineskip}
\noindent\begin{minipage}{\linewidth}

\textbf{例 M4-065：环网络中的福斯特求和}

有限环网络 \(C_N\) 含 \(N\ge3\) 个节点与 \(N\) 条等阻边，每条电阻为
\(r>0\)。用 \(R_e\)
表示任一真实边两端在完整网络中的等效电阻。先通过直接边与其余环路的并联求
\(R_e\)，再计算 \(\sum_eR_e/r\)，验证福斯特迹定理；不把真实支路电阻
\(r\) 误当作二端等效电阻。

\end{minipage}\par

\par\noindent\begin{minipage}{\linewidth}
\textbf{解答}

\begin{center}
\includegraphics[width=\linewidth,height=.40\textheight,keepaspectratio]{\detokenize{assets/module-04/examples/M4-065.pdf}}
\par\smallskip{\small 图 M4-065：环中的支路值与端口测量值不同}
\end{center}
\end{minipage}\par

图 M4-065
的任一红边两端还能经其余环路回流，故端口电阻小于单边阻值。所有边等价，先算一个端口再乘环边数，便得到完整的迹求和。

\(N\) 个相同电阻 \(r\) 构成环 \(C_N\)。对任一边，其两端之间有直接支路
\(r\) 和其余 \(N-1\) 条边组成的串联支路 \((N-1)r\)，所以

\[
R_e=r\parallel(N-1)r=\frac{N-1}{N}r.
\]

环上共有 \(N\) 条边，于是

\[
\sum_e\frac{R_e}{r}
=N\cdot\frac{N-1}{N}
=N-1.
\]

\subsection{三、约化矩阵证明（拔高选讲）}\label{m4:ux4e09ux7ea6ux5316ux77e9ux9635ux8bc1ux660eux62d4ux9ad8ux9009ux8bb2}

设网络有 \(M\)
条边。任意指定每条边的方向，并选定一个参考节点。删去完整关联矩阵中参考节点对应的一行，得到约化关联矩阵

\[
A\in\mathbb R^{(N-1)\times M}.
\]

第 \(e\) 列记为 \(a_e\)。边的定向只会把 \(a_e\) 整体乘以
\(-1\)，不会影响后面的二次型。

把边电导排成对角矩阵

\[
C=\operatorname{diag}(c_1,c_2,\ldots,c_M),
\]

约化节点电导矩阵为

\[
K=ACA^{\mathsf T}
=\sum_{e=1}^{M}c_ea_ea_e^{\mathsf T}.
\]

网络连通，所以约化矩阵 \(K\) 正定、可逆。

\subsubsection{1.
单条边两端的等效电阻}\label{m4:ux5355ux6761ux8fb9ux4e24ux7aefux7684ux7b49ux6548ux7535ux963b}

在第 \(e\)
条边的定向起点注入单位电流，在终点抽出单位电流。约化注入向量正是
\(a_e\)，节点电势满足

\[
K\mathbf V=a_e,
\]

因而

\[
\mathbf V=K^{-1}a_e.
\]

端点电势差为 \(a_e^{\mathsf T}\mathbf V\)，所以

\[
R_e=a_e^{\mathsf T}K^{-1}a_e.
\]

\subsubsection{2.
对所有边加权求和}\label{m4:ux5bf9ux6240ux6709ux8fb9ux52a0ux6743ux6c42ux548c}

利用标量等于自身迹以及迹的循环不变性，

\[
\begin{aligned}
\sum_{e=1}^{M}c_eR_e
&=\sum_{e=1}^{M}c_ea_e^{\mathsf T}K^{-1}a_e\\
&=\sum_{e=1}^{M}\operatorname{tr}
\!\left(K^{-1}c_ea_ea_e^{\mathsf T}\right)\\
&=\operatorname{tr}\!\left[
K^{-1}\sum_{e=1}^{M}c_ea_ea_e^{\mathsf T}
\right]\\
&=\operatorname{tr}(K^{-1}K)\\
&=\operatorname{tr}(I_{N-1})\\
&=N-1.
\end{aligned}
\]

定理得证。

\par\noindent\begin{minipage}{\linewidth}
\begin{center}
\includegraphics[width=\linewidth,height=.40\textheight,keepaspectratio]{\detokenize{assets/module-04/theory/review-tikz-m4-foster-02-trace-proof-flow.pdf}}
\par\smallskip{\small 图示：福斯特迹定理的矩阵证明}
\end{center}
\end{minipage}\par

\subsubsection{3.
与拉普拉斯广义逆的关系}\label{m4:ux4e0eux62c9ux666eux62c9ux65afux5e7fux4e49ux9006ux7684ux5173ux7cfb}

若保留全部 \(N\) 个节点，完整电导拉普拉斯矩阵记为

\[
L=BCB^{\mathsf T}.
\]

\(L\)
的零空间由全体分量相同的向量张成，因为整体电势同时平移不改变电流。用摩尔---彭若斯广义逆
\(L^+\) 可写成

\[
R_e=b_e^{\mathsf T}L^+b_e.
\]

于是

\[
\sum_ec_eR_e
=\operatorname{tr}(L^+L)
=\operatorname{rank}(L)
=N-1.
\]

约化矩阵证明与广义逆证明本质相同：前者先固定参考电势，直接删去零模；后者保留全部节点，再用广义逆投影到与常量向量正交的
\(N-1\) 维子空间。

\par\noindent\begin{minipage}{\linewidth}
\begin{center}
\includegraphics[width=\linewidth,height=.40\textheight,keepaspectratio]{\detokenize{assets/module-04/original/tikz-m4-foster-03-degree-freedom-view-pdf-v11.jpg}}
\par\smallskip{\small 图示：整体电势平移不改变电流，连通网络保留N减1个独立电势方向}
\end{center}
\end{minipage}\par

\subsection{四、竞赛中的使用方法（拔高选讲）}\label{m4:ux56dbux7adeux8d5bux4e2dux7684ux4f7fux7528ux65b9ux6cd5ux62d4ux9ad8ux9009ux8bb2}

若所有边的端点等效电阻都相同，设网络有 \(M\) 条等阻边，则

\[
MR_{\text{边}}=(N-1)r,
\]

所以

\[
\boxed{R_{\text{边}}=\frac{N-1}{M}r}.
\]

对每点度数均为 \(d\) 的边传递网络，握手定理给出

\[
M=\frac{Nd}{2},
\]

因而

\[
R_{\text{边}}
=\frac{2(N-1)}{Nd}r.
\]

这与专题三十六由平衡辅助态得到的相邻节点公式一致。

若边不完全等价，应按网络的自同构对边分类。若第 \(j\) 类有 \(m_j\)
条边，其端点等效电阻为 \(R_j\)，在等阻网络中有

\[
\boxed{\sum_jm_jR_j=(N-1)r}.
\]

这只给出一个线性约束。边分成 \(k\) 类时，通常还需另外计算 \(k-1\)
类，才能确定最后一类。

\par\noindent\begin{minipage}{\linewidth}
\begin{center}
\includegraphics[width=\linewidth,height=.40\textheight,keepaspectratio]{\detokenize{assets/module-04/theory/review-tikz-m4-foster-04-solving-workflow-pdf-v11.pdf}}
\par\smallskip{\small 图示：福斯特迹定理用于等价边和多类边网络}
\end{center}
\end{minipage}\par

\subsection{五、练习与解答}\label{m4:ux4e94ux7ec3ux4e60ux4e0eux89e3ux7b54}

\par\addvspace{.4\baselineskip}
\noindent\begin{minipage}{\linewidth}

\subsubsection{作业一：带叶边的四边形网络}\label{m4:ux4f5cux4e1aux4e00ux5e26ux53f6ux8fb9ux7684ux56dbux8fb9ux5f62ux7f51ux7edc}

\textbf{题号：M4-066}

\textbf{层级：} 决赛　　\textbf{方法：} 网络方程

\(r>0\)，节点集合恰为 \(A,B,C,D,E\)，边集合恰为
\(AB,BD,DC,CA,DE\)，没有对角线。求每一条真实边两端的测量电阻时，其他节点均不外接电流源，包括叶端悬空的情况。

节点 \(A,B,C,D\) 由四条电阻 \(AB,BD,DC,CA\) 构成一个环，节点 \(D\)
再经电阻 \(DE\) 接到叶节点 \(E\)。五条电阻均为
\(r\)。逐条计算真实边两端的等效电阻，并验证福斯特迹定理。

\end{minipage}\par

\par\noindent\begin{minipage}{\linewidth}
\textbf{解答}

\begin{center}
\includegraphics[width=\linewidth,height=.40\textheight,keepaspectratio]{\detokenize{assets/module-04/examples/M4-066.pdf}}
\par\smallskip{\small 图 M4-066：环边有回路，叶边没有替代通路}
\end{center}
\end{minipage}\par

图 M4-066 的 \(E\) 只经 \(D-E\)
接入。测环边端口时，悬空叶支路不流电；测叶边时又没有另一条通路与它并联。因此两类项不同，求和时分别计4条和1条。

先看任一环边，例如 \(AC\)。在 \(A,C\) 间测量时，节点 \(E\)
没有外部注入，且只通过 \(DE\) 与网络相连。由节点 \(E\) 的电流守恒，

\[
I_{DE}=0.
\]

所以叶支路不影响任一环边两端的等效电阻。\(A,C\)
之间有两条并联通路：直接边 \(r\)，以及路径 \(A\to B\to D\to C\)，电阻为
\(3r\)。故

\[
R_{AC}=r\parallel3r=\frac34r.
\]

对其余三条环边，同样可剪去无流叶支路；每条边都是 \(r\) 与另外三条环边
\(3r\) 并联，因此

\[
R_{AB}=R_{BD}=R_{DC}=R_{CA}=\frac34r.
\]

边 \(DE\) 是桥。从 \(D\) 到 \(E\) 的电流没有其他路径可走，所以

\[
R_{DE}=r.
\]

于是

\[
\sum_{e\in E}R_e
=4\times\frac34r+r
=4r.
\]

网络共有 \(N=5\) 个节点，定理右端为

\[
(N-1)r=4r,
\]

两者一致。

\subsection{六、练习与解答}\label{m4:ux516dux7ec3ux4e60ux4e0eux89e3ux7b54-1}

\par\addvspace{.4\baselineskip}
\noindent\begin{minipage}{\linewidth}

\subsubsection{作业二：正方体与正十二面体}\label{m4:ux4f5cux4e1aux4e8cux6b63ux65b9ux4f53ux4e0eux6b63ux5341ux4e8cux9762ux4f53}

\textbf{题号：M4-067}

\textbf{层级：} 决赛　　\textbf{方法：} 网络方程

两个网络分别测量，不互相连接。正方体含8节点12边，正十二面体含20节点30边，仅棱上有电阻，\(r>0\)；未选作端口的节点均无外部注入。

正方体和正十二面体的每条棱均由电阻 \(r\)
构成。分别求任意一条棱两端的等效电阻。

\end{minipage}\par

\par\noindent\begin{minipage}{\linewidth}
\textbf{解答}

\begin{center}
\includegraphics[width=\linewidth,height=.40\textheight,keepaspectratio]{\detokenize{assets/module-04/examples/M4-067.pdf}}
\par\smallskip{\small 图 M4-067：先确认边等价，再使用迹定理}
\end{center}
\end{minipage}\par

图 M4-067
示意正十二面体的全部连接，正方体另按12棱连接。两个网络各自具有边传递对称，所以每个求和项相同；分别代入各自节点数和边数，不能混用两者计数。

正多面体的旋转对称可把任意一条棱映到任意另一条棱，所以所有棱对应的全网等效电阻相同。

\paragraph{1. 正方体}\label{m4:ux6b63ux65b9ux4f53}

正方体有

\[
N=8,
\qquad
M=12.
\]

\par\noindent\begin{minipage}{\linewidth}
\begin{center}
\includegraphics[width=\linewidth,height=.40\textheight,keepaspectratio]{\detokenize{assets/module-04/process/M4-067a.pdf}}
\par\smallskip{\small 图 M4-067a：正方体单独使用八节点十二边计数}
\end{center}
\end{minipage}\par

图 M4-067a
给出正方体这一独立网络。12条棱在对称操作下等价，因此总和中有12个相同的端点测量值；节点计数应使用8，不能沿用正十二面体的数据。

由

\[
MR_{\text{边}}=(N-1)r,
\]

得到

\[
\boxed{R_{\text{边}}=\frac{7}{12}r}.
\]

\paragraph{2. 正十二面体}\label{m4:ux6b63ux5341ux4e8cux9762ux4f53}

正十二面体有 \(12\) 个五边形面，每条棱被两个面共有，故

\[
M=\frac{12\times5}{2}=30.
\]

由欧拉公式 \(N-M+F=2\)，

\[
N=30-12+2=20.
\]

所以

\[
\boxed{R_{\text{边}}=\frac{20-1}{30}r=\frac{19}{30}r}.
\]

这两个结果与专题三十六的平衡辅助态计算完全一致。

\subsection{七、练习与解答}\label{m4:ux4e03ux7ec3ux4e60ux4e0eux89e3ux7b54-5}

\par\addvspace{.4\baselineskip}
\noindent\begin{minipage}{\linewidth}

\subsubsection{作业三：田字格的两类边}\label{m4:ux4f5cux4e1aux4e09ux7530ux5b57ux683cux7684ux4e24ux7c7bux8fb9}

\textbf{题号：M4-068}

\textbf{层级：} 决赛　　\textbf{方法：} 网络方程

\(r>0\)，9个节点、12条最近邻边组成有限方格，边界外无连接。按图取
\(A=(0,2),B=(1,2),O=(1,1)\)；只向 \(A\) 注入、从 \(B\)
抽出等量测试电流，其他节点无注入。

一个 \(2\times2\) 方格由 \(3\times3\) 个节点组成，每一小段电阻均为
\(r\)。设 \(O\) 为中心节点，\(B\) 为上边中点，\(A\) 为左上角。求边界边
\(AB\) 两端的等效电阻。

\end{minipage}\par

\par\noindent\begin{minipage}{\linewidth}
\textbf{解答}

\begin{center}
\includegraphics[width=\linewidth,height=.40\textheight,keepaspectratio]{\detokenize{assets/module-04/examples/M4-068.pdf}}
\par\smallskip{\small 图 M4-068：田字格有八条边界边和四条中心边}
\end{center}
\end{minipage}\par

图 M4-068 的 \(A-B\) 是边界边，\(O-B\) 是中心边。虽然每条真实电阻都为
\(r\)，它们在整个网络中的端口测量值分属两类；先算中心边实验，再代入按8与4计数的福斯特关系。

网络有

\[
N=9,
\qquad
M=12.
\]

在正方形的二面体对称作用下，十二条边分成两类：

\begin{itemize}
\tightlist
\item
  八条边界边，端点等效电阻记为 \(R_{\text{外}}\)；
\item
  四条连接中心与边界中点的中心边，端点等效电阻记为 \(R_{\text{内}}\)。
\end{itemize}

福斯特迹定理给出

\[
8R_{\text{外}}+4R_{\text{内}}=8r.
\]

下面完整计算 \(R_{\text{内}}=R_{OB}\)。为避免量纲干扰，先令每条边电阻为
\(1\)，最后再乘 \(r\)。从 \(O\) 注入 \(1\,\mathrm A\)，从 \(B\) 抽出
\(1\,\mathrm A\)，并取

\[
V_B=0.
\]

关于竖直对称轴成对的节点等势。依次记：

\begin{itemize}
\tightlist
\item
  两个上角电势为 \(x\)；
\item
  左、右边界中点电势为 \(y\)；
\item
  两个下角电势为 \(z\)；
\item
  下边中点电势为 \(w\)；
\item
  中心电势为 \(V\)。
\end{itemize}

\par\noindent\begin{minipage}{\linewidth}
\begin{center}
\includegraphics[width=\linewidth,height=.40\textheight,keepaspectratio]{\detokenize{assets/module-04/process/M4-068a.pdf}}
\par\smallskip{\small 图 M4-068a：中心边实验的五个未知电势}
\end{center}
\end{minipage}\par

图 M4-068a
为了计算中心边，临时把端口换成O、B，并把每边电阻单位化。镜像节点分别具有x、y、z的相同电势，尚须保留w和V，因此下面列的是五个未知量的节点方程。

对五类未知节点分别列节点方程：

\[
2x-y=0,
\]

\[
3y-x-z-V=0,
\]

\[
2z-y-w=0,
\]

\[
3w-2z-V=0,
\]

\[
4V-2y-w=1.
\]

逐步消元。由第一式，

\[
y=2x.
\]

第二式给出

\[
z+V=5x.
\]

第三式给出

\[
w=2z-2x.
\]

代入第四式：

\[
3(2z-2x)-2z-V=0,
\]

即

\[
V=4z-6x.
\]

与 \(z+V=5x\) 联立，得到

\[
z=\frac{11}{5}x,
\qquad
V=\frac{14}{5}x,
\qquad
w=\frac{12}{5}x.
\]

最后代入中心节点方程：

\[
4V-2y-w=1,
\]

得到

\[
\frac{24}{5}x=1,
\qquad
x=\frac5{24}.
\]

因此

\[
V=\frac{7}{12}.
\]

单位电流下的端口电压就是等效电阻，所以恢复电阻量纲后

\[
R_{\text{内}}=R_{OB}=\frac7{12}r.
\]

代回总和式：

\[
8R_{\text{外}}+4\times\frac7{12}r=8r,
\]

故

\[
\boxed{R_{AB}=R_{\text{外}}=\frac{17}{24}r}.
\]

这一题体现了迹定理的典型用法：它不能单独区分多类边，但可以在算出较容易的一类后，立即反推出目标边。

\subsection{八、使用条件}\label{m4:ux516bux4f7fux7528ux6761ux4ef6-2}

\begin{enumerate}
\def\labelenumi{\arabic{enumi}.}
\tightlist
\item
  网络必须有限且连通；若有多个连通分支，应对每个分支分别应用定理。
\item
  一般网络应使用加权式 \(\sum_eR_e/r_e=N-1\)；只有所有边均为 \(r\)
  时，才能写成 \(\sum_eR_e=(N-1)r\)。
\item
  求和对象是真实边两端的全网等效电阻，不是所有节点对，也不是边自身电阻的简单总和。
\item
  ``所有边等效''必须由网络自同构保证；每点度数相同并不自动意味着边传递。
\item
  边分成多类时，迹定理只提供一个总和方程，仍需用对称性、节点电势法或其他网络变换补足方程。
\item
  含理想导线、开路支路、负电阻或受控源的网络不能不加处理地代入本节的正电阻图模型。
\item
  本定理处理有限网络的全边总和；下一专题的无穷网络必须先说明有限截断、边界条件或自相似极限。
\end{enumerate}

\section{专题三十八：无穷电阻网络与自相似性（基础必讲与拔高选讲）}\label{m4:ux4e13ux9898ux4e09ux5341ux516bux65e0ux7a77ux7535ux963bux7f51ux7edcux4e0eux81eaux76f8ux4f3cux6027ux57faux7840ux5fc5ux8bb2ux4e0eux62d4ux9ad8ux9009ux8bb2}

无穷网络不能仅凭一张``无限延伸''的电路图自动获得等效电阻。严格做法是先规定一列有限网络及其远端边界条件，再研究端口等效电阻是否存在有限极限。自相似性提供的是极限满足的固定点方程；极限是否存在、应取哪个根，还必须由有限截断或其他收敛论证决定。

\subsection{一、用有限截断定义无穷网络（基础必讲）}\label{m4:ux4e00ux7528ux6709ux9650ux622aux65adux5b9aux4e49ux65e0ux7a77ux7f51ux7edcux57faux7840ux5fc5ux8bb2}

设第 \(n\) 个有限网络的端口等效电阻为
\(R_n\)。若在明确的截断方式和远端边界条件下，

\[
\lim_{n\to\infty}R_n=R_{\infty}
\]

存在且有限，才把 \(R_{\infty}\) 定义为该无穷网络的端口等效电阻。

必须同时说明：

\begin{enumerate}
\def\labelenumi{\arabic{enumi}.}
\tightlist
\item
  端口是哪两个节点；
\item
  第 \(n\) 个有限网络包含哪些元件；
\item
  远端采用开路、短路、周期边界还是其他终端；
\item
  不同合理终端是否收敛到同一极限。
\end{enumerate}

\par\noindent\begin{minipage}{\linewidth}
\begin{center}
\includegraphics[width=\linewidth,height=.40\textheight,keepaspectratio]{\detokenize{assets/module-04/original/tikz-m4-topic38-v9-from35-v8-01.png}}
\par\smallskip{\small 图示：无穷电阻网络应由一列带明确远端边界的有限网络定义}
\end{center}
\end{minipage}\par

对被动正电阻网络，增加支路只会降低端口等效电阻，这可由瑞利单调性理解。因此，开路和短路终端常能给出上下界；若两界趋于同一值，极限便与远端终端的具体选择无关。

\subsection{二、何时可以写自相似方程（基础必讲）}\label{m4:ux4e8cux4f55ux65f6ux53efux4ee5ux5199ux81eaux76f8ux4f3cux65b9ux7a0bux57faux7840ux5fc5ux8bb2}

截去首个单元后，尾部可以用 \(R_{\infty}\) 替代，必须满足：

\begin{enumerate}
\def\labelenumi{\arabic{enumi}.}
\tightlist
\item
  尾部的两个外接节点与原端口一一对应；
\item
  尾部的连接拓扑与原网络相同；
\item
  元件参数相同，或者按一个已知比例整体缩放；
\item
  原问题没有被截断操作改变的额外接地、受控源或边界条件；
\item
  对应有限截断序列确实收敛。
\end{enumerate}

若尾部与整体完全相同，尾部等效电阻为
\(R_{\infty}\)；若所有尾部电阻都缩放为原来的 \(\lambda\)
倍，由电阻网络的齐次性，尾部等效电阻应为
\(\lambda R_{\infty}\)，不能仍写成 \(R_{\infty}\)。

\par\noindent\begin{minipage}{\linewidth}
\begin{center}
\includegraphics[width=\linewidth,height=.40\textheight,keepaspectratio]{\detokenize{assets/module-04/theory/review-tikz-m4-topic38-v9-from35-v8-02.pdf}}
\par\smallskip{\small 图示：使用自相似替换前必须同时核对端口、拓扑、参数比例和边界条件}
\end{center}
\end{minipage}\par

\subsection{三、标准无穷梯形网络（基础必讲）}\label{m4:ux4e09ux6807ux51c6ux65e0ux7a77ux68afux5f62ux7f51ux7edcux57faux7840ux5fc5ux8bb2}

每个单元由上支路串联电阻 \(R_1\) 和随后跨接上下导轨的电阻 \(R_2\)
构成，右侧继续接入完全相同的尾部。设从左端口看入的极限等效电阻为
\(R_{\infty}\)。

截去首个单元后，尾部仍为 \(R_{\infty}\)。第一个跨接电阻 \(R_2\)
与尾部并联，再与 \(R_1\) 串联，所以

\[
R_{\infty}
=R_1+\left(R_2\parallel R_{\infty}\right)
=R_1+\frac{R_2R_{\infty}}{R_2+R_{\infty}}.
\]

整理得

\[
R_{\infty}^2-R_1R_{\infty}-R_1R_2=0.
\]

两个代数根为

\[
R_{\pm}=\frac{R_1\pm\sqrt{R_1^2+4R_1R_2}}{2}.
\]

因为 \(R_1,R_2>0\)，有
\(R_-<0<R_+\)。被动正电阻有限网络的等效电阻为正，所以唯一可能的物理解是

\[
\boxed{
R_{\infty}=R_+
=\frac{R_1+\sqrt{R_1^2+4R_1R_2}}{2}
}.
\]

\par\noindent\begin{minipage}{\linewidth}
\begin{center}
\includegraphics[width=\linewidth,height=.40\textheight,keepaspectratio]{\detokenize{assets/module-04/original/tikz-m4-topic38-v9-from35-v8-03.png}}
\par\smallskip{\small 图示：标准无穷梯形网络截去首单元后得到严格对应的固定点方程}
\end{center}
\end{minipage}\par

\par\addvspace{.4\baselineskip}
\noindent\begin{minipage}{\linewidth}

\textbf{例 M4-069：等阻无穷梯形的固定点}

一条半无限单导轨梯形由相同单元重复组成：每级先在上导轨串联一只 \(r>0\)
电阻，到达下一节点后再经一只 \(r\)
跨接到理想零电阻下导轨，右侧继续相同结构；输入端之前没有额外跨接支路。用右端开路的有限截断定义输入电阻极限，求半无限等效电阻，并选择满足正电阻与串并联范围约束的固定点。

\end{minipage}\par

\par\noindent\begin{minipage}{\linewidth}
\textbf{解答}

\begin{center}
\includegraphics[width=\linewidth,height=.40\textheight,keepaspectratio]{\detokenize{assets/module-04/examples/M4-069.pdf}}
\par\smallskip{\small 图 M4-069：单导轨梯形的尾部与原端口相同}
\end{center}
\end{minipage}\par

图 M4-069
在输入端先有一只串联电阻，到下一节点才出现跨接支路。截去首级后，右尾部与输入的连接位置一致，因此并联中的尾部电阻是
\(R\)；正根还应落在 \(r\) 与 \(2r\) 之间。

设有限截断输入电阻趋于 \(R\)。首只串联电阻之后是跨接的 \(r\) 与相同尾部
\(R\) 并联，所以 \(R=r+rR/(r+R)\)，即 \(R^2-rR-r^2=0\)。正根为
\((1+\sqrt5)r/2\)，负根与无源正电阻不相容；同时 \(r<R<2r\)。

右端开路的一级截断为 \(R_1=2r\)；逐级增加单元满足
\(R_{n+1}=r+rR_n/(r+R_n)\)。此映射递增，序列从 \(2r\) 单调下降且始终大于
\(r\)，所以极限存在并满足所用正根方程。

\subsection{四、为什么有限截断确实收敛（拔高选讲）}\label{m4:ux56dbux4e3aux4ec0ux4e48ux6709ux9650ux622aux65adux786eux5b9eux6536ux655bux62d4ux9ad8ux9009ux8bb2}

定义映射

\[
f(x)=R_1+\frac{R_2x}{R_2+x}.
\]

取远端短路，令 \(x_0^{(\text{短})}=0\)；取远端开路，形式上令
\(x_0^{(\text{开})}=+\infty\)。每向左增加一个单元，都有

\[
x_{n+1}=f(x_n).
\]

第一次迭代后，

\[
x_1^{(\text{短})}=R_1,
\qquad
x_1^{(\text{开})}=R_1+R_2.
\]

映射 \(f\) 在 \([R_1,R_1+R_2]\) 上单调增加，而且

\[
0<f'(x)=\frac{R_2^2}{(R_2+x)^2}
\le\left(\frac{R_2}{R_1+R_2}\right)^2<1.
\]

因此 \(f\)
在该区间是压缩映射，只有一个固定点；开路、短路以及任何落入该区间的被动终端都会收敛到同一个
\(R_+\)。更具体地，短路序列单调递增，开路序列单调递减，并始终满足

\[
x_n^{(\text{短})}\le R_+\le x_n^{(\text{开})}.
\]

单调方向也可直接由

\[
f(x)-x
=-\frac{(x-R_+)(x-R_-)}{R_2+x}
\]

判断。由于 \(x>0>R_-\)，当 \(x<R_+\) 时有 \(f(x)>x\)，当 \(x>R_+\) 时有
\(f(x)<x\)。这就分别证明了短路截断递增、开路截断递减。

\par\noindent\begin{minipage}{\linewidth}
\begin{center}
\includegraphics[width=\linewidth,height=.40\textheight,keepaspectratio]{\detokenize{assets/module-04/theory/review-tikz-m4-topic38-v9-from35-v8-04.pdf}}
\par\smallskip{\small 图示：开路与短路有限截断从两侧收敛到同一个正固定点}
\end{center}
\end{minipage}\par

这也说明：先写 \(R=f(R)\)
再选一个看起来合理的根，只是快捷计算；有限截断收敛才是该根代表原无穷网络的依据。

\subsection{五、电压、电流与功率的几何衰减（拔高选讲）}\label{m4:ux4e94ux7535ux538bux7535ux6d41ux4e0eux529fux7387ux7684ux51e0ux4f55ux8870ux51cfux62d4ux9ad8ux9009ux8bb2}

依次把跨接节点记为 \(v_0,v_1,\ldots\)，令 \(U_n\) 为 \(v_n\)
对下导轨的电压。相邻跨接节点之间串联一个 \(R_1\)，每个跨接节点再通过
\(R_2\) 接到下导轨。在节点 \(v_{n+1}\) 处，跨接电阻 \(R_2\)
与更右侧的同构延续部分 \(R_{\infty}\) 并联，故该节点的负载等效电阻为

\[
R_2\parallel R_{\infty}=R_{\infty}-R_1.
\]

于是 \(v_n\) 到 \(v_{n+1}\) 之间正好构成 \(R_1\) 与上述负载的分压，因而

\[
\frac{U_{n+1}}{U_n}
=q
=\frac{R_{\infty}-R_1}{R_{\infty}}
=\frac{R_2}{R_2+R_{\infty}}.
\]

由于 \(0<q<1\)，

\[
\boxed{U_n=q^nU_0}.
\]

第 \(n\) 个跨接电阻中的电流为

\[
I_n^{(2)}=\frac{U_n}{R_2},
\]

其功率为

\[
P_n^{(2)}=\frac{U_n^2}{R_2}
=\frac{U_0^2}{R_2}q^{2n}.
\]

所以远端电压、电流按 \(q^n\) 衰减，单元功率按 \(q^{2n}\) 衰减。从
\(v_n\) 沿右侧延续部分看入的等效电阻为
\(R_{\infty}\)，这部分吸收的总功率为 \(U_n^2/R_{\infty}\)，也趋于零。

\par\noindent\begin{minipage}{\linewidth}
\begin{center}
\includegraphics[width=\linewidth,height=.40\textheight,keepaspectratio]{\detokenize{assets/module-04/original/tikz-m4-topic38-v9-from35-v8-05.png}}
\par\smallskip{\small 图示：标准无穷梯形网络中的节点电压和耗散功率按几何级数衰减}
\end{center}
\end{minipage}\par

\subsection{六、电流递推法：不能直接整体替换时怎么办（拔高选讲）}\label{m4:ux516dux7535ux6d41ux9012ux63a8ux6cd5ux4e0dux80fdux76f4ux63a5ux6574ux4f53ux66ffux6362ux65f6ux600eux4e48ux529eux62d4ux9ad8ux9009ux8bb2}

有些双线梯形或非对称网络，端口附近的连接方式与内部重复单元不同。此时不能把``远处看起来重复''误写成整个端口网络的同一个
\(R\)。可改用：

\begin{enumerate}
\def\labelenumi{\arabic{enumi}.}
\tightlist
\item
  对第 \(n\) 个截面定义节点电压或纵向、横向支路电流；
\item
  在内部单元列基尔霍夫方程，得到常系数差分方程；
\item
  舍去向无穷远发散的特征根；
\item
  用端口注入、抽出条件确定剩余系数；
\item
  由端口电压与端口电流求等效电阻。
\end{enumerate}

若截去首单元后的两节点端口、拓扑和参数确实与整体相同，则优先使用固定点法；否则应使用差分方程，而不是强行套二次方程。

\par\addvspace{.4\baselineskip}
\noindent\begin{minipage}{\linewidth}

\textbf{例 M4-070：双导轨梯形的电流递推}

半无限双导轨梯形的截面编号为
\(n=0,1,2,\ldots\)。每相邻截面间，上、下导轨各接一只电阻
\(a>0\)；各内部截面的上下节点间跨接一只 \(b>0\)。在差模激励下，定义
\(V_n\) 为上节点减下节点的电势差，\(I_n\) 为上导轨从第 \(n\) 列流向第
\(n+1\) 列的电流，下导轨回流反向。列局部KVL、KCL，消去 \(I_n\) 得到
\(V_n\)
的递推方程，求特征根并说明半无限远端有界条件如何选出衰减支，端口幅度不预先指定。

\end{minipage}\par

\par\noindent\begin{minipage}{\linewidth}
\textbf{解答}

\begin{center}
\includegraphics[width=\linewidth,height=.40\textheight,keepaspectratio]{\detokenize{assets/module-04/examples/M4-070.pdf}}
\par\smallskip{\small 图 M4-070：一个内部单元给出差模递推}
\end{center}
\end{minipage}\par

图 M4-070 的上下水平支路各为
\(a\)，横截面间电势差的减小要同时计入两条导轨的压降。竖直支路 \(b\)
分走的电流由节点守恒给出，消去水平电流便得到相邻三列的递推。

设上下导轨每段电阻均为 \(a\)，每个截面之间的跨接电阻均为 \(b\)。令
\(V_n\) 为第 \(n\) 个截面的上下节点电势差，\(I_n\) 为上导轨从第 \(n\)
个截面流向第 \(n+1\)
个截面的电流；由差模对称性，下导轨中有大小相等、方向相反的电流。于是

\[
V_n-V_{n+1}=2aI_n,
\qquad
I_{n-1}-I_n=\frac{V_n}{b}.
\]

消去 \(I_n\)，得到

\[
V_{n+1}-\left(2+\frac{2a}{b}\right)V_n+V_{n-1}=0.
\]

若令 \(V_n=C\lambda^n\)，特征方程为

\[
\lambda^2-\left(2+\frac{2a}{b}\right)\lambda+1=0.
\]

两根之积为 \(1\)；无穷远处电压必须有界，故舍去大于 \(1\) 的根，只保留
\(0<\lambda<1\) 的衰减模态。端口条件再决定系数
\(C\)。这正是``内部递推决定衰减，端口条件决定幅度''的标准结构。

具体地，令 \(c=1+a/b>1\)，两根为 \(\lambda_{\pm}=c\pm\sqrt{c^2-1}\)。有
\(\lambda_+>1\)、\(0<\lambda_-<1\)，有界半无限解保留后者。

\subsection{七、练习与解答}\label{m4:ux4e03ux7ec3ux4e60ux4e0eux89e3ux7b54-6}

\par\addvspace{.4\baselineskip}
\noindent\begin{minipage}{\linewidth}

\subsubsection{作业一：双线无穷梯形}\label{m4:ux4f5cux4e1aux4e00ux53ccux7ebfux65e0ux7a77ux68afux5f62}

\textbf{题号：M4-071}

\textbf{层级：} 基础　　\textbf{方法：} 网络方程

半无限双导轨梯形沿右方延伸，列号为
\(n=0,1,2,\ldots\)。每条导轨的相邻列之间各接一只 \(r>0\) 电阻，且每个
\(n\ge1\) 的上下节点之间接一只 \(r\)，输入列 \(n=0\)
没有跨接电阻。端口是第0列上下节点，远端由无外接负载的有限截断取极限定义。求输入等效电阻，说明截去首级后尾部端口与原端口的对应，并检验结果相对于必经的两只串联电阻的大小。

\end{minipage}\par

\par\noindent\begin{minipage}{\linewidth}
\textbf{解答}

\begin{center}
\includegraphics[width=\linewidth,height=.40\textheight,keepaspectratio]{\detokenize{assets/module-04/examples/M4-071.pdf}}
\par\smallskip{\small 图 M4-071：双导轨输入前没有竖直跨接}
\end{center}
\end{minipage}\par

图 M4-071 的 \(A,B\) 列没有跨接电阻，电流须先经过上、下各一只
\(r\)。到下一列后，竖直电阻才与同构尾部并联；若在输入列误加一只竖直支路，会得到另一道网络的答案。

设等效电阻为 \(R\)。通过首级到达右侧截面，上、下导轨各串联一个
\(r\)；在该截面，纵向电阻 \(r\) 与尾部 \(R\) 并联。因此

\[
R=2r+\left(r\parallel R\right)
=2r+\frac{rR}{r+R}.
\]

两边乘以 \(r+R\)：

\[
R(r+R)=2r(r+R)+rR.
\]

整理得

\[
R^2-2rR-2r^2=0.
\]

故

\[
R=(1\pm\sqrt3)r.
\]

舍去负根，得到

\[
\boxed{R=(1+\sqrt3)r}.
\]

可检验
\(R>2r\)：任何端口电流都必须先经过上、下两个串联电阻，而纵向支路与尾部并联只会使中间部分小于
\(r\)，所以还应有 \(2r<R<3r\)；所得结果满足这一范围。

一级截断为 \(R_1=3r\)，逐级扩展满足
\(R_{n+1}=2r+rR_n/(r+R_n)\)。映射递增，序列从 \(3r\) 下降且以下界 \(2r\)
为界，因此确实收敛到上述正固定点。

\subsection{八、练习与解答}\label{m4:ux516bux7ec3ux4e60ux4e0eux89e3ux7b54-3}

\par\addvspace{.4\baselineskip}
\noindent\begin{minipage}{\linewidth}

\subsubsection{作业二：按比例缩放的无穷梯形}\label{m4:ux4f5cux4e1aux4e8cux6309ux6bd4ux4f8bux7f29ux653eux7684ux65e0ux7a77ux68afux5f62}

\textbf{题号：M4-072}

\textbf{层级：} 基础　　\textbf{方法：} 网络方程

半无限单导轨梯形的下导轨为理想零电阻线。第 \(n=0,1,2,\ldots\)
级先在上导轨串联电阻 \(2^{-n}r\)，到达下一节点后经另一只 \(2^{-n}r\)
跨接到下导轨，再进入下一级，\(r>0\)；输入端没有额外跨接。以右端开路的有限截断之极限定义输入电阻。求该电阻，并解释去掉第0级后尾部阻值为何应按比例缩放，不能直接等于原网络。

\end{minipage}\par

\par\noindent\begin{minipage}{\linewidth}
\textbf{解答}

\begin{center}
\includegraphics[width=\linewidth,height=.40\textheight,keepaspectratio]{\detokenize{assets/module-04/examples/M4-072.pdf}}
\par\smallskip{\small 图 M4-072：尾部阻值缩放，端口拓扑相同}
\end{center}
\end{minipage}\par

图 M4-072 的每对串联、跨接电阻随级数乘
\(1/2\)。截去首级虽然连接形式不变，但尾部所有参数同时减半，等效电阻也减半；固定点方程必须带这个因子。

设整体等效电阻为 \(R\)。截去第零级后，尾部所有电阻都是原网络对应电阻的
\(1/2\)。由线性电阻网络的齐次性，尾部等效电阻为

\[
\frac R2,
\]

而不是 \(R\)。因此

\[
R=r+\left(r\parallel\frac R2\right)
=r+\frac{rR}{2r+R}.
\]

两边乘以 \(2r+R\)：

\[
R(2r+R)=r(2r+R)+rR.
\]

整理得

\[
R^2=2r^2.
\]

取正根，

\[
\boxed{R=\sqrt2\,r}.
\]

若误把尾部写成
\(R\)，会错误得到黄金分割结果；本题专门检验``参数比例是否保持''的判断。

有限一级截断给
\(R_1=2r\)。扩展一级并将旧尾部全部缩小一半，得到同一个递增映射
\(R_{n+1}=r+rR_n/(2r+R_n)\)；该序列单调下降、下界为
\(r\)，所以有限截断的极限确实存在。

\subsection{九、练习与解答}\label{m4:ux4e5dux7ec3ux4e60ux4e0eux89e3ux7b54-3}

\par\addvspace{.4\baselineskip}
\noindent\begin{minipage}{\linewidth}

\subsubsection{作业三：无限方格网的有限极限}\label{m4:ux4f5cux4e1aux4e09ux65e0ux9650ux65b9ux683cux7f51ux7684ux6709ux9650ux6781ux9650}

\textbf{题号：M4-073}

\textbf{层级：} 复赛　　\textbf{方法：} 网络方程

\(r>0\)，每个节点只与上下左右最近邻连接。等效电阻用平衡两点注入的有限周期正方网格取无限尺寸极限定义，求电势差，不先设无穷远单点电势为零。

二维无限正方格网的每条边电阻均为 \(r\)。求相邻节点 \(A,B\)
之间的等效电阻，并说明本题为什么不属于简单的一元自相似固定点问题。

\end{minipage}\par

\par\noindent\begin{minipage}{\linewidth}
\textbf{解答}

\begin{center}
\includegraphics[width=\linewidth,height=.40\textheight,keepaspectratio]{\detokenize{assets/module-04/examples/M4-073.pdf}}
\par\smallskip{\small 图 M4-073：无限方格由平衡有限网络取极限}
\end{center}
\end{minipage}\par

图 M4-073
红点是相邻端口，四周虚线表示网络继续延伸。先对有限周期网络建立总注入为零的源汇实验，再取尺寸极限，保持端点附近的四邻接结构；一元尾部替换不能保留这种二维连接。

取 \(m\times m\) 周期方格作为有限近似，\(m\ge5\)。网络有

\[
N=m^2
\]

个节点，每点度数
\(d=4\)，所有真实边在周期网络的对称下等价。由专题三十七的福斯特迹定理，边数为

\[
M=\frac{Nd}{2}=2m^2,
\]

任一相邻边两端的全网等效电阻为

\[
R_m=\frac{N-1}{M}r
=\frac{m^2-1}{2m^2}r.
\]

令 \(m\to\infty\)，得到

\[
\boxed{R_{AB}=\frac r2}.
\]

周期边界只在距固定端口约为 \(m\)
的位置改变连接。由狄利克雷原理或汤姆孙原理对最小能量作局部截断，可证明上述周期网络的端口电阻极限就是无限方格网的端口电阻。

这一结果也可由专题三十六的有限平衡辅助态得到。无限方格删去局部单元后，端口附近并不会化成``若干有限电阻与一个未知
\(R\)
串并联''的一端口结构；尾部仍是多端网络，因此不能直接写一个标量固定点方程。

\subsection{十、使用条件}\label{m4:ux5341ux4f7fux7528ux6761ux4ef6-1}

\begin{enumerate}
\def\labelenumi{\arabic{enumi}.}
\tightlist
\item
  自相似方程只给出极限的必要候选，必须配合有限截断、单调性、压缩映射或其他收敛论证。
\item
  截去单元后必须核对端口、拓扑、参数和边界条件；尾部整体缩放 \(\lambda\)
  倍时，其等效电阻也是 \(\lambda R\)。
\item
  固定点方程出现多个非负根时，不能只凭``取正根''决定物理解，必须由有限网络序列选择。
\item
  无穷网络的总耗散并不等于``电流流到无穷远消失''；对标准梯形网络，功率由无限多个正电阻以收敛级数吸收。
\item
  多端尾部不能总用一个标量电阻表示；必要时应使用多端导纳矩阵或差分方程。
\item
  无限二维、三维晶格通常需要有限截断、周期边界、傅里叶方法或格林函数，不能仅凭局部图形相似套用一元二次方程。
\end{enumerate}

\section{专题三十九：纯电阻网络的多端等效理论（基础必讲与拔高选讲）}\label{m4:ux4e13ux9898ux4e09ux5341ux4e5dux7eafux7535ux963bux7f51ux7edcux7684ux591aux7aefux7b49ux6548ux7406ux8bbaux57faux7840ux5fc5ux8bb2ux4e0eux62d4ux9ad8ux9009ux8bb2}

二端网络可以用一个等效电阻描述，多端网络则需要保持全部端口电势与端口电流的关系。端口响应矩阵可用于消去内部节点，消元后的矩阵还能改写成外端点之间的电阻支路。

\subsection{一、多端等效的定义（基础必讲）}\label{m4:ux4e00ux591aux7aefux7b49ux6548ux7684ux5b9aux4e49ux57faux7840ux5fc5ux8bb2}

设一个纯电阻黑盒有 \(N\) 个外端点，端点电势和注入黑盒的电流分别记为

\[
\boldsymbol v=(v_1,\ldots,v_N)^{\mathsf T},
\qquad
\boldsymbol i=(i_1,\ldots,i_N)^{\mathsf T}.
\]

约定 \(i_k>0\)
表示电流从外部流入黑盒。对不接地、内部不含独立源的网络，有

\[
\sum_{k=1}^{N}i_k=0.
\]

所有端点电势同时增加同一个常数不会改变任何支路电流，即

\[
\boldsymbol v\longrightarrow \boldsymbol v+c\boldsymbol 1
\]

不改变端口响应。因此，端口电势只有电势差有意义。

若两个 \(N\) 端网络对每个允许的 \(\boldsymbol v\) 都给出相同的
\(\boldsymbol i\)，则称它们多端等效。线性纯电阻网络必存在矩阵
\(\boldsymbol\Lambda\)，使

\[
\boxed{\boldsymbol i=\boldsymbol\Lambda\boldsymbol v}.
\]

\(\boldsymbol\Lambda\)
称为端口响应矩阵或端口导纳矩阵。两个网络多端等效，当且仅当它们的
\(\boldsymbol\Lambda\) 相同。

\par\noindent\begin{minipage}{\linewidth}
\begin{center}
\includegraphics[width=\linewidth,height=.40\textheight,keepaspectratio]{\detokenize{assets/module-04/original/tikz-m4-topic39-v9-from36-v8-01.png}}
\par\smallskip{\small 图示：多端等效要求任意端口电势激励下的全部端口电流都相同}
\end{center}
\end{minipage}\par

\par\addvspace{.4\baselineskip}
\noindent\begin{minipage}{\linewidth}

\textbf{例 M4-074：一次二端测量不足以确定三端网络}

一个线性无源三端电阻黑盒的外端为1、2、3。只进行一次测量：端3悬空，在端1、2之间测得
\(2\ \Omega\)。这一次测量是否足以确定黑盒的全部三端响应？通过比较两个与测量相容的网络，说明端3随后接入负载或指定电势时还可能出现什么差别，区分二端测量与多端完全等效。

\end{minipage}\par

\par\noindent\begin{minipage}{\linewidth}
\textbf{解答}

\begin{center}
\includegraphics[width=\linewidth,height=.40\textheight,keepaspectratio]{\detokenize{assets/module-04/examples/M4-074.pdf}}
\par\smallskip{\small 图 M4-074：同一二端读数可以来自不同三端星形}
\end{center}
\end{minipage}\par

图 M4-074 将只被一次实验约束的两条臂固定为各 \(1\,\Omega\)，第三臂记为
\(s\)。端3悬空时它不流电，所以任取两个不同的 \(s\) 都得到同一
\(R_{12}\)；换接端3后即可区分。

可取网络甲为等臂星形 \(R_1=R_2=R_3=1\ \Omega\)，网络乙为
\(R_1=R_2=1\ \Omega,R_3=2\ \Omega\)。端3悬空时，两者的1---2通路都是
\(1+1=2\ \Omega\)；改测1---3时分别为 \(2\ \Omega\) 和
\(3\ \Omega\)，所以一次测量不能确定三端响应。

\subsection{二、从节点导纳矩阵消去内部节点（拔高选讲）}\label{m4:ux4e8cux4eceux8282ux70b9ux5bfcux7eb3ux77e9ux9635ux6d88ux53bbux5185ux90e8ux8282ux70b9ux62d4ux9ad8ux9009ux8bb2}

把所有节点分为外端点集合 \(B\) 与内部节点集合
\(I\)。全网的加权拉普拉斯矩阵按这两组节点分块：

\[
\begin{pmatrix}
\boldsymbol i\\[2mm]
\boldsymbol 0
\end{pmatrix}
=
\begin{pmatrix}
\boldsymbol L_{BB}&\boldsymbol L_{BI}\\
\boldsymbol L_{IB}&\boldsymbol L_{II}
\end{pmatrix}
\begin{pmatrix}
\boldsymbol v\\[2mm]
\boldsymbol u
\end{pmatrix}.
\]

其中 \(\boldsymbol u\)
是内部节点电势；内部节点没有外部注入电流，所以右侧内部方程为

\[
\boldsymbol L_{IB}\boldsymbol v+
\boldsymbol L_{II}\boldsymbol u=\boldsymbol 0.
\]

若每个内部连通分支都与至少一个外端点相连，则 \(\boldsymbol L_{II}\)
可逆，从而

\[
\boldsymbol u=-\boldsymbol L_{II}^{-1}
\boldsymbol L_{IB}\boldsymbol v.
\]

代回外端点方程，得到

\[
\boldsymbol i=\left(
\boldsymbol L_{BB}
-\boldsymbol L_{BI}\boldsymbol L_{II}^{-1}\boldsymbol L_{IB}
\right)\boldsymbol v.
\]

因此

\[
\boxed{
\boldsymbol\Lambda
=\boldsymbol L_{BB}
-\boldsymbol L_{BI}\boldsymbol L_{II}^{-1}\boldsymbol L_{IB}
}.
\]

这就是消去内部节点的舒尔补，也称节点消元或克朗约化。它不是近似，而是在所有外端点伏安关系意义下的严格等效。

\par\noindent\begin{minipage}{\linewidth}
\begin{center}
\includegraphics[width=\linewidth,height=.40\textheight,keepaspectratio]{\detokenize{assets/module-04/theory/review-tikz-m4-topic39-v9-from36-v8-02.pdf}}
\par\smallskip{\small 图示：内部节点通过舒尔补被严格消去并转化为端口响应矩阵}
\end{center}
\end{minipage}\par

\par\addvspace{.4\baselineskip}
\noindent\begin{minipage}{\linewidth}

\textbf{例 M4-075：消去星形中心节点}

一无源星形网络只有一个内部节点 \(O\)，它通过电导 \(a_1,a_2,a_3>0\)
分别连接到外端1、2、3。中心没有外部注入，外端电势可任意指定为
\(V_1,V_2,V_3\)。用中心KCL消去
\(V_O\)，求在任意外端激励下完全等效的三角形支路电导，特别写出连接端1、2的一支，并说明不是只匹配某一次二端实验。

\end{minipage}\par

\par\noindent\begin{minipage}{\linewidth}
\textbf{解答}

\begin{center}
\includegraphics[width=\linewidth,height=.40\textheight,keepaspectratio]{\detokenize{assets/module-04/examples/M4-075.pdf}}
\par\smallskip{\small 图 M4-075：中心KCL同时确定全部外端响应}
\end{center}
\end{minipage}\par

图 M4-075 的中心 \(O\)
无外部注入，因此它的势由三条电导加权决定。将这个势代回每个端口电流，系数自然按两个外端成对出现，这才是三端完全等效的支路电导。

中心守恒给 \(V_O=(a_1V_1+a_2V_2+a_3V_3)/A\)，其中
\(A=a_1+a_2+a_3\)。端1注入为
\(a_1(V_1-V_O)=(a_1a_2/A)(V_1-V_2)+(a_1a_3/A)(V_1-V_3)\)；逐端比较便得到三条等效电导
\(g_{ij}=a_ia_j/A\)。

\subsection{三、端口响应矩阵的结构（拔高选讲）}\label{m4:ux4e09ux7aefux53e3ux54cdux5e94ux77e9ux9635ux7684ux7ed3ux6784ux62d4ux9ad8ux9009ux8bb2}

纯电阻网络的 \(\boldsymbol\Lambda\) 具有四个关键性质。

\subsubsection{1. 对称性}\label{m4:ux5bf9ux79f0ux6027}

\[
\boldsymbol\Lambda^{\mathsf T}=\boldsymbol\Lambda.
\]

这是互易性的矩阵表达，也可直接由对称矩阵的舒尔补仍对称得到。

\subsubsection{2. 零行和}\label{m4:ux96f6ux884cux548c}

\[
\boldsymbol\Lambda\boldsymbol 1=\boldsymbol 0.
\]

它表示整体平移电势不产生电流；由对称性也有
\(\boldsymbol 1^{\mathsf T}\boldsymbol\Lambda=\boldsymbol 0^{\mathsf T}\)，故端口总电流为零。

\subsubsection{3. 半正定性}\label{m4:ux534aux6b63ux5b9aux6027}

外部电源送入网络的功率为

\[
P=\boldsymbol v^{\mathsf T}\boldsymbol i
=\boldsymbol v^{\mathsf T}\boldsymbol\Lambda\boldsymbol v
=\sum_{e}g_e(\Delta v_e)^2\ge0.
\]

因此 \(\boldsymbol\Lambda\)
为半正定矩阵。若所有外端点经网络相互连通，则它只有一个零模态
\(\boldsymbol 1\)，即

\[
\operatorname{rank}\boldsymbol\Lambda=N-1.
\]

\subsubsection{4.
非对角元非正}\label{m4:ux975eux5bf9ux89d2ux5143ux975eux6b63}

对 \(i\ne j\)，有

\[
\Lambda_{ij}\le0.
\]

这是电阻网络拉普拉斯矩阵在消元后仍保持的符号结构。等价地，可以定义非负的端点间等效支路电导

\[
g_{ij}=-\Lambda_{ij}\ge0.
\]

\par\noindent\begin{minipage}{\linewidth}
\begin{center}
\includegraphics[width=\linewidth,height=.40\textheight,keepaspectratio]{\detokenize{assets/module-04/original/tikz-m4-topic39-v9-from36-v8-03.png}}
\par\smallskip{\small 图示：端口响应矩阵同时满足互易、规范不变、无源性和拉普拉斯符号结构}
\end{center}
\end{minipage}\par

这四条应联合使用。仅知道矩阵对称，并不足以保证它能够由无源正电阻实现。

\subsection{四、多端完全图等效定理（拔高选讲）}\label{m4:ux56dbux591aux7aefux5b8cux5168ux56feux7b49ux6548ux5b9aux7406ux62d4ux9ad8ux9009ux8bb2}

对每一对不同端点 \(i,j\)，在它们之间接入电导

\[
\boxed{g_{ij}=-\Lambda_{ij}}.
\]

由于 \(\boldsymbol\Lambda\) 的行和为零，完全图在端点 \(i\) 的对角元为

\[
\sum_{j\ne i}g_{ij}
=-\sum_{j\ne i}\Lambda_{ij}
=\Lambda_{ii}.
\]

非对角元则为
\(-g_{ij}=\Lambda_{ij}\)。所以这个只含外端点的加权完全图，其节点导纳矩阵恰好就是
\(\boldsymbol\Lambda\)。

因此：

\[
\boxed{
\text{任意有限无源纯电阻多端网络，都可严格约化为只含外端点的加权完全图。}
}
\]

若 \(g_{ij}=0\)，相应支路可以省略。因此实际需要的是至多

\[
\binom N2=\frac{N(N-1)}2
\]

条支路，而不是一定有这么多条非零电阻。参数计数与完全图边数相同，只能作为结果的核对；真正的存在性证明是上面的舒尔补与构造公式。

\par\noindent\begin{minipage}{\linewidth}
\begin{center}
\includegraphics[width=\linewidth,height=.40\textheight,keepaspectratio]{\detokenize{assets/module-04/original/tikz-m4-topic39-v9-from36-v8-04.png}}
\par\smallskip{\small 图示：端口响应矩阵的非对角元直接给出等效完全图的各支路电导}
\end{center}
\end{minipage}\par

\subsection{五、一般星形网络的消元公式（拔高选讲）}\label{m4:ux4e94ux4e00ux822cux661fux5f62ux7f51ux7edcux7684ux6d88ux5143ux516cux5f0fux62d4ux9ad8ux9009ux8bb2}

设内部节点 \(O\) 通过电阻 \(R_i\) 接到第 \(i\) 个外端点，记

\[
a_i=\frac1{R_i},
\qquad
A=\sum_{k=1}^{N}a_k.
\]

内部节点满足

\[
\sum_{k=1}^{N}a_k(u-v_k)=0,
\]

故

\[
u=\frac{\sum_k a_kv_k}{A}.
\]

端点 \(i\) 的注入电流为

\[
i_i=a_i(v_i-u)
=a_iv_i-\sum_j\frac{a_ia_j}{A}v_j.
\]

因此消去中心节点后，外端点 \(i,j\) 之间的完全图支路电导为

\[
\boxed{g_{ij}=\frac{a_ia_j}{A}},
\]

相应支路电阻为

\[
\boxed{R_{ij}^{(\mathrm{\text{完}})}=\frac{A}{a_ia_j}}.
\]

当 \(N=3\) 时，这一公式退化为通常的星形---三角形公式。若三条星形臂分别为
\(R_A,R_B,R_C\)，则

\[
R_{AB}^{(\Delta)}
=R_A+R_B+\frac{R_AR_B}{R_C},
\]

其余两式循环置换即可。

\par\noindent\begin{minipage}{\linewidth}
\begin{center}
\includegraphics[width=\linewidth,height=.40\textheight,keepaspectratio]{\detokenize{assets/module-04/original/tikz-m4-topic39-v9-from36-v8-05.png}}
\par\smallskip{\small 图示：消去一般星形网络的中心节点后外端点之间形成加权完全图}
\end{center}
\end{minipage}\par

\par\addvspace{.4\baselineskip}
\noindent\begin{minipage}{\linewidth}

\textbf{例 M4-076：四端等臂星形的完全图等效}

四端星形网络的四个外端分别经一只 \(r>0\) 电阻接到唯一中心
\(O\)，中心没有外部电流源。将它消元为只含四个外端的完全图
\(K_4\)，求每条等效支路电阻；再在任意两端间测量，其余两端无外部注入，求二端等效电阻，与原星形的结果核对。

\end{minipage}\par

\par\noindent\begin{minipage}{\linewidth}
\textbf{解答}

\begin{center}
\includegraphics[width=\linewidth,height=.40\textheight,keepaspectratio]{\detokenize{assets/module-04/examples/M4-076.pdf}}
\par\smallskip{\small 图 M4-076：四端星形消元为六条等值边}
\end{center}
\end{minipage}\par

图 M4-076 的四条臂都是 \(r\)，中心消元后任意一对外端都出现一条电导
\(1/(4r)\) 的边。六条完全图支路共同参与二端实验，所以每边的 \(4r\)
不等于端口测得的 \(2r\)。

四条星形臂均为 \(r\)
时，\(a_i=1/r\)、\(A=4/r\)，故任意两外端点之间的完全图支路电阻均为

\[
R_{ij}^{(\mathrm{\text{完}})}=4r.
\]

所得 \(K_4\) 网络任意两端的二端等效电阻为
\(2(4r)/4=2r\)，与原星形网络中两条臂串联的 \(2r\) 一致。

\subsection{六、支路电阻不等于二端测量值（拔高选讲）}\label{m4:ux516dux652fux8defux7535ux963bux4e0dux7b49ux4e8eux4e8cux7aefux6d4bux91cfux503cux62d4ux9ad8ux9009ux8bb2}

把多端黑盒约化为完全图后，支路
\(R_{ij}^{(\mathrm{\text{完}})}=1/g_{ij}\)
是端口模型中的参数。若只在端点 \(p,q\)
之间注入、抽出电流，而其余端点悬空，电流仍可经其他端点对应的支路绕行，因此实测二端等效电阻一般不等于完全图中的直接支路电阻。

令

\[
\boldsymbol b_{pq}=\boldsymbol e_p-\boldsymbol e_q.
\]

注入电流 \(I\boldsymbol b_{pq}\) 时，利用 \(\boldsymbol\Lambda\)
的摩尔---彭若斯逆，有

\[
\boldsymbol v=I\boldsymbol\Lambda^+\boldsymbol b_{pq}+c\boldsymbol 1.
\]

因此

\[
\boxed{
R_{pq}^{(\mathrm{\text{测}})}
=\boldsymbol b_{pq}^{\mathsf T}
\boldsymbol\Lambda^+
\boldsymbol b_{pq}
}.
\]

\par\noindent\begin{minipage}{\linewidth}
\begin{center}
\includegraphics[width=\linewidth,height=.40\textheight,keepaspectratio]{\detokenize{assets/module-04/original/tikz-m4-topic39-v9-from36-v8-06.png}}
\par\smallskip{\small 图示：多端完全图中的直接支路电阻通常不等于其余端口悬空时的二端测量值}
\end{center}
\end{minipage}\par

\par\addvspace{.4\baselineskip}
\noindent\begin{minipage}{\linewidth}

\textbf{例 M4-077：支路电阻与二端测量值}

无源三端星形的三条臂电阻均为
\(r>0\)，中心不外接电路。先求它的多端等效三角形每条真实支路阻值；再把第三外端悬空，在另外两端测量等效电阻。通过三角形中的两条通路说明两个阻值为何不同，且两种拓扑的二端测量为何一致。

\end{minipage}\par

\par\noindent\begin{minipage}{\linewidth}
\textbf{解答}

\begin{center}
\includegraphics[width=\linewidth,height=.40\textheight,keepaspectratio]{\detokenize{assets/module-04/examples/M4-077.pdf}}
\par\smallskip{\small 图 M4-077：三角形每边不是二端测量电阻}
\end{center}
\end{minipage}\par

图 M4-077 中任选两端测量，星形只经过两条臂串联。换成三角形后，直接边
\(3r\) 还要与经第三端的两条 \(3r\) 串联路并联，结果才与原星形一致。

三条星形臂都为 \(r\)，其三角形等效的每条边为
\(3r\)；但任意两端、第三端悬空时测得的是 \(2r\)。这两个量不能混用。

\subsection{七、由端口实验重建网络（拔高选讲）}\label{m4:ux4e03ux7531ux7aefux53e3ux5b9eux9a8cux91cdux5efaux7f51ux7edcux62d4ux9ad8ux9009ux8bb2}

选第 \(N\) 个端点为参考点，令 \(v_N=0\)。删去 \(\boldsymbol\Lambda\)
的第 \(N\) 行、第 \(N\) 列，得到 \((N-1)\) 阶约化矩阵
\(\widehat{\boldsymbol\Lambda}\)。对连通网络，它正定且可逆：

\[
\widehat{\boldsymbol v}
=\boldsymbol Z\widehat{\boldsymbol i},
\qquad
\boldsymbol Z=\widehat{\boldsymbol\Lambda}^{-1}.
\]

依次在端点 \(k\)
注入单位电流、由参考端点抽出，并记录全部端点电势，就得到
\(\boldsymbol Z\) 的第 \(k\) 列。测得 \(\boldsymbol Z\) 后求逆，即可恢复
\(\widehat{\boldsymbol\Lambda}\)；再利用零行和补回第 \(N\) 行、第 \(N\)
列，最后由

\[
g_{ij}=-\Lambda_{ij}
\]

得到完全图各支路。

若已经测得所有端点对之间的二端有效电阻，组成矩阵
\(\boldsymbol D=[R_{ij}^{(\mathrm{\text{测}})}]\)，令

\[
\boldsymbol J=\boldsymbol I-\frac1N\boldsymbol 1\boldsymbol 1^{\mathsf T},
\]

则

\[
\boxed{
\boldsymbol\Lambda^+
=-\frac12\boldsymbol J\boldsymbol D\boldsymbol J
}.
\]

再取一次广义逆便得到
\(\boldsymbol\Lambda\)。所以``只测一对端点''远远不够，而``测得全部端点对的有效电阻''原则上足以重建端口响应。

\par\noindent\begin{minipage}{\linewidth}
\begin{center}
\includegraphics[width=\linewidth,height=.40\textheight,keepaspectratio]{\detokenize{assets/module-04/original/tikz-m4-topic39-v9-from36-v8-07.png}}
\par\smallskip{\small 图示：通过参考端回流实验或全部端点对测量可以重建端口响应矩阵}
\end{center}
\end{minipage}\par

\subsection{八、网络拼接时响应矩阵如何变化（拔高选讲）}\label{m4:ux516bux7f51ux7edcux62fcux63a5ux65f6ux54cdux5e94ux77e9ux9635ux5982ux4f55ux53d8ux5316ux62d4ux9ad8ux9009ux8bb2}

若两个黑盒具有同一组外端点，并把对应端点直接并接，则它们承受相同的
\(\boldsymbol v\)，端口电流相加：

\[
\boxed{\boldsymbol\Lambda_{\mathrm{\text{并}}}
=\boldsymbol\Lambda_1+\boldsymbol\Lambda_2}.
\]

若只把部分端点相连，新连接点应先视为内部节点，再对总导纳矩阵做一次舒尔补。若某端点接地，则整体电势平移不再是自由度，不能继续机械使用``零行和''的未接地模型。

\par\noindent\begin{minipage}{\linewidth}
\begin{center}
\includegraphics[width=\linewidth,height=.40\textheight,keepaspectratio]{\detokenize{assets/module-04/theory/review-tikz-m4-topic39-v9-from36-v8-08.pdf}}
\par\smallskip{\small 图示：同端点黑盒并接时响应矩阵相加而部分端点拼接后还需再次消元}
\end{center}
\end{minipage}\par

\subsection{九、练习与解答}\label{m4:ux4e5dux7ec3ux4e60ux4e0eux89e3ux7b54-4}

\par\addvspace{.4\baselineskip}
\noindent\begin{minipage}{\linewidth}

\subsubsection{作业一：证明一般星形消元公式}\label{m4:ux4f5cux4e1aux4e00ux8bc1ux660eux4e00ux822cux661fux5f62ux6d88ux5143ux516cux5f0f}

\textbf{题号：M4-078}

\textbf{层级：} 决赛　　\textbf{方法：} 等效变换、网络方程

\(N\ge2\)，各
\(R_i>0\)；星形中心无外部注入，外端电势可以任意指定。结论要匹配每一个外端的电流，而不仅是某两端测量值；完全图各对外端只计一条无向支路。

一个内部节点 \(O\) 通过电阻 \(R_1,\ldots,R_N\) 分别连接到 \(N\)
个外端点。证明消去 \(O\) 后，端点 \(i,j\) 之间应接入电导

\[
g_{ij}=\frac{a_ia_j}{\sum_k a_k},
\qquad a_i=\frac1{R_i}.
\]

\end{minipage}\par

\par\noindent\begin{minipage}{\linewidth}
\textbf{解答}

\begin{center}
\includegraphics[width=\linewidth,height=.40\textheight,keepaspectratio]{\detokenize{assets/module-04/examples/M4-078.pdf}}
\par\smallskip{\small 图 M4-078：一般星形消元公式}
\end{center}
\end{minipage}\par

图 M4-078
的中心只有内部支路，没有外部注入。中心KCL确定其电势后，外端电流可改写成所有成对电势差之和；每一对的系数就是消元完全图的支路电导。

设中心节点电势为 \(u\)。由中心节点的电流守恒，

\[
\sum_k a_k(u-v_k)=0,
\]

所以

\[
u=\frac{\sum_k a_kv_k}{A},
\qquad A=\sum_k a_k.
\]

端点 \(i\) 的注入电流为

\[
i_i=a_i(v_i-u)
=\left(a_i-\frac{a_i^2}{A}\right)v_i
-\sum_{j\ne i}\frac{a_ia_j}{A}v_j.
\]

另一方面，若只在外端点之间接入电导 \(g_{ij}\)，则

\[
i_i=\sum_{j\ne i}g_{ij}(v_i-v_j).
\]

比较 \(v_j\) 的系数，得到

\[
g_{ij}=\frac{a_ia_j}{A}.
\]

并且

\[
\sum_{j\ne i}g_{ij}
=\frac{a_i(A-a_i)}{A}
=a_i-\frac{a_i^2}{A},
\]

故对角系数也完全一致，证明了任意端口电势下的全部端口电流均相同。

\subsection{十、练习与解答}\label{m4:ux5341ux7ec3ux4e60ux4e0eux89e3ux7b54-1}

\par\addvspace{.4\baselineskip}
\noindent\begin{minipage}{\linewidth}

\subsubsection{作业二：由三次二端测量重建三端网络}\label{m4:ux4f5cux4e1aux4e8cux7531ux4e09ux6b21ux4e8cux7aefux6d4bux91cfux91cdux5efaux4e09ux7aefux7f51ux7edc}

\textbf{题号：M4-079}

\textbf{层级：} 基础　　\textbf{方法：} 等效变换、网络方程

黑盒是连通、线性、无源且互易的纯电阻网络，每次测量时第三端没有外接电流。求的三角形是保持全部三端响应的等效网络，其支路电阻与给出的二端测量数值不同。

某三端黑盒在其余一端悬空时测得

\[
R_{AB}^{(\mathrm{\text{测}})}=2\ \Omega,
\quad
R_{BC}^{(\mathrm{\text{测}})}=3\ \Omega,
\quad
R_{CA}^{(\mathrm{\text{测}})}=4\ \Omega.
\]

求与该黑盒多端等效的三角形网络三条支路电阻。

\end{minipage}\par

\par\noindent\begin{minipage}{\linewidth}
\textbf{解答}

\begin{center}
\includegraphics[width=\linewidth,height=.40\textheight,keepaspectratio]{\detokenize{assets/module-04/examples/M4-079.pdf}}
\par\smallskip{\small 图 M4-079：三次测量重建三端网络}
\end{center}
\end{minipage}\par

图 M4-079
把三次二端测量分别对应到等效星形的两臂和。先解出三条星形臂，再变为三角形参数；测量时悬空的第三端并不意味着它从多端等效中消失。

先构造与三次测量相符的星形网络。若三条星形臂为
\(x_A,x_B,x_C\)，则第三端悬空时有

\[
x_A+x_B=2,
\qquad
x_B+x_C=3,
\qquad
x_C+x_A=4.
\]

解得

\[
x_A=\frac32\ \Omega,
\qquad
x_B=\frac12\ \Omega,
\qquad
x_C=\frac52\ \Omega.
\]

记

\[
S=x_Ax_B+x_Bx_C+x_Cx_A
=\frac{23}{4}\ \Omega^2.
\]

由星形---三角形公式，三角形中与星形臂 \(x_C\) 相对的 \(AB\) 支路为

\[
R_{AB}^{(\Delta)}=\frac{S}{x_C}=\frac{23}{10}\ \Omega.
\]

同理，

\[
\boxed{
R_{AB}^{(\Delta)}=\frac{23}{10}\ \Omega,
\quad
R_{BC}^{(\Delta)}=\frac{23}{6}\ \Omega,
\quad
R_{CA}^{(\Delta)}=\frac{23}{2}\ \Omega
}.
\]

检验 \(AB\) 两端、\(C\) 悬空时：

\[
R_{AB}^{(\mathrm{\text{测}})}
=R_{AB}^{(\Delta)}\parallel
\left(R_{AC}^{(\Delta)}+R_{CB}^{(\Delta)}\right)
=2\ \Omega,
\]

与题设一致。其余两对同理。

\subsection{十一、练习与解答}\label{m4:ux5341ux4e00ux7ec3ux4e60ux4e0eux89e3ux7b54-1}

\par\addvspace{.4\baselineskip}
\noindent\begin{minipage}{\linewidth}

\subsubsection{作业三：等臂多端星形网络}\label{m4:ux4f5cux4e1aux4e09ux7b49ux81c2ux591aux7aefux661fux5f62ux7f51ux7edc}

\textbf{题号：M4-080}

\textbf{层级：} 基础　　\textbf{方法：} 等效变换、网络方程

\(N\ge2\)，\(r>0\)；中心无外部电流源或接地连接。第一问要求任意外端激励的多端等效，第二问只接测试电流的两个端口，其他外端悬空。

\(N\) 个外端点分别通过一个阻值为 \(r\) 的电阻接到同一中心节点。求：

\begin{enumerate}
\def\labelenumi{\arabic{enumi}.}
\tightlist
\item
  消去中心节点后的完全图支路电阻；
\item
  任意两端之间、其余端点悬空时的二端等效电阻。
\end{enumerate}

\end{minipage}\par

\par\noindent\begin{minipage}{\linewidth}
\textbf{解答}

\begin{center}
\includegraphics[width=\linewidth,height=.40\textheight,keepaspectratio]{\detokenize{assets/module-04/examples/M4-080.pdf}}
\par\smallskip{\small 图 M4-080：等臂多端星形}
\end{center}
\end{minipage}\par

图 M4-080
区分消去中心后的完全图支路和实际二端实验。前者按任意外端电势匹配全部电流，后者将未接测试源的端点悬空，因此单条完全图支路还要和其他通路并联。

每条星形臂的电导为 \(a=1/r\)，总电导为
\(A=N/r\)。由一般星形消元公式，任意两外端点之间的完全图支路电导为

\[
g=\frac{a^2}{A}=\frac1{Nr},
\]

所以

\[
\boxed{R_{ij}^{(\mathrm{\text{完}})}=Nr}.
\]

在边电阻均为 \(Nr\) 的完全图 \(K_N\) 中，从端点 \(p\) 注入电流并由 \(q\)
抽出时，其余 \(N-2\)
个端点由交换对称性等势，连接这些等势点之间的支路无电流。标准完全图结果为

\[
R_{pq}^{(\mathrm{\text{测}})}=\frac{2}{N}(Nr)=2r.
\]

这也与原星形网络中从 \(p\) 经中心到 \(q\) 的两条电阻串联相符：

\[
\boxed{R_{pq}^{(\mathrm{\text{测}})}=2r}.
\]

再次可见：完全图的单条支路电阻 \(Nr\) 与端点对的二端测量值 \(2r\)
是不同概念。

\subsection{十二、使用条件}\label{m4:ux5341ux4e8cux4f7fux7528ux6761ux4ef6}

\begin{enumerate}
\def\labelenumi{\arabic{enumi}.}
\tightlist
\item
  多端等效要求保持全部端口伏安关系，而不是只保持某一对端点的等效电阻。
\item
  完全图可实现性的依据是舒尔补后仍具有拉普拉斯矩阵结构，不能只靠独立参数计数。
\item
  \(\binom N2\) 是非零等效支路数的上限；某些 \(g_{ij}\) 可以为零。
\item
  端口响应矩阵只描述外部行为，不唯一确定原网络的内部拓扑。
\item
  含受控源、独立源、负电阻或非互易元件时，对称性、半正定性或非对角元符号可能失效，不能直接套用纯电阻完全图结论。
\item
  接地端、公共参考端与普通悬空端的物理条件不同，建模时必须明确。
\end{enumerate}

\section{专题四十：电阻网络的谱分解、变分原理与离散格林函数（拔高选讲）}\label{m4:ux4e13ux9898ux56dbux5341ux7535ux963bux7f51ux7edcux7684ux8c31ux5206ux89e3ux53d8ux5206ux539fux7406ux4e0eux79bbux6563ux683cux6797ux51fdux6570ux62d4ux9ad8ux9009ux8bb2}

\begin{quote}
阅读准备：零模、广义逆与变分
所有节点电势同时加常量不改变任何支路电流，所以图拉普拉斯矩阵必有常量零模，不能在全空间直接求普通逆。固定参考点，或限制在总注入电流为零的子空间后，广义逆给出电势差。变分原理则把基尔霍夫方程改写为``满足约束的真实分布使耗散泛函取极值''。
\end{quote}

\par\addvspace{.4\baselineskip}
\noindent\begin{minipage}{\linewidth}

\textbf{例 M4-081：电势零模与参考值平移}

有限连通无源电阻网络中，任意支路 \(i-j\) 的电流为
\(g_{ij}(V_i-V_j)\)，\(g_{ij}>0\)，节点注入向量满足守恒。将全部节点电势同时加同一常数
\(V_0\)。判断支路电压、电流及节点注入是否改变，说明拉普拉斯矩阵中的全1向量所对应的零模代表什么，以及确定数值电势时为何需要选择一个参考。

\end{minipage}\par

\par\noindent\begin{minipage}{\linewidth}
\textbf{解答}

\begin{center}
\includegraphics[width=\linewidth,height=.40\textheight,keepaspectratio]{\detokenize{assets/module-04/examples/M4-081.pdf}}
\par\smallskip{\small 图 M4-081：共同电势平移不改变任何支路差}
\end{center}
\end{minipage}\par

图 M4-081
的每条支路只读取两个节点的电势差。将三个示意节点连同全网络一起加常量后，每一条差均不变；矩阵中的全1向量所对应的零模表达的正是这个参考值自由度。

各支路只读取 \(V_i-V_j\)，统一加常量时电压、电流均不变。矩阵形式中
\(L\mathbf1=0\)，因此
\(L(\mathbf V+V_0\mathbf1)=L\mathbf V\)，外部注入也不变。固定一个参考节点或选零和规范，只是在同一物理电流态中选唯一电势代表。

\par\noindent\begin{minipage}{\linewidth}
\begin{center}
\includegraphics[width=\linewidth,height=.40\textheight,keepaspectratio]{\detokenize{assets/module-04/theory/review-tikz-prereq-m4-topic40-v9-from37-v7-01-network-matrix.pdf}}
\par\smallskip{\small 图示：节点方程的矩阵表示、常量零模与参考电势}
\end{center}
\end{minipage}\par

对有限网络，节点电压法把问题化为线性方程；对周期网络和无限条带，直接逐点消元会迅速失去结构。谱分解的作用是找到一组彼此独立的电势模式，使高维网络分解成若干一维差分问题。变分原理则从功率最小化出发，给出等效电阻的上下界、单调性与收敛依据。

\subsection{一、加权图拉普拉斯矩阵（拔高选讲）}\label{m4:ux4e00ux52a0ux6743ux56feux62c9ux666eux62c9ux65afux77e9ux9635ux62d4ux9ad8ux9009ux8bb2}

设有限连通网络的边 \(e=(i,j)\) 电阻为 \(r_{ij}>0\)、电导为
\(c_{ij}=1/r_{ij}\)。节点 \(i\) 的外部注入电流为
\(I_i\)，取流入网络为正。KCL 给出

\[
\sum_jc_{ij}(V_i-V_j)=I_i.
\]

定义加权图拉普拉斯矩阵

\[
L_{ii}=\sum_jc_{ij},
\qquad
L_{ij}=-c_{ij}\quad(i\ne j),
\]

则节点方程写成

\[
\boxed{\boldsymbol L\boldsymbol V=\boldsymbol I}.
\]

\(\boldsymbol L\) 实对称、每行元素之和为零，并满足

\[
\boldsymbol V^{\mathsf T}\boldsymbol L\boldsymbol V
=\sum_{(i,j)\in E}c_{ij}(V_i-V_j)^2\ge0.
\]

连通网络只有一个零模态
\(\boldsymbol 1\)，对应所有节点电势同时平移。可解条件是总注入电流为零：

\[
\boldsymbol 1^{\mathsf T}\boldsymbol I=0.
\]

\par\noindent\begin{minipage}{\linewidth}
\begin{center}
\includegraphics[width=\linewidth,height=.40\textheight,keepaspectratio]{\detokenize{assets/module-04/original/tikz-m4-topic40-v9-from37-v1-01-laplacian.png}}
\par\smallskip{\small 图示：电阻网络的节点、边电导和注入电流共同组成图拉普拉斯方程}
\end{center}
\end{minipage}\par

\subsection{二、有效电阻的谱公式（拔高选讲）}\label{m4:ux4e8cux6709ux6548ux7535ux963bux7684ux8c31ux516cux5f0fux62d4ux9ad8ux9009ux8bb2}

从节点 \(a\) 注入电流 \(I\)，从节点 \(b\) 抽出同样电流：

\[
\boldsymbol I=I\boldsymbol b,
\qquad
\boldsymbol b=\boldsymbol e_a-\boldsymbol e_b.
\]

用广义逆 \(\boldsymbol L^+\) 消去零模态，有

\[
\boldsymbol V=I\boldsymbol L^+\boldsymbol b+c\boldsymbol1.
\]

端口电压为

\[
V_a-V_b
=\boldsymbol b^{\mathsf T}\boldsymbol V
=I\boldsymbol b^{\mathsf T}\boldsymbol L^+\boldsymbol b,
\]

故

\[
\boxed{
R_{ab}=\boldsymbol b^{\mathsf T}\boldsymbol L^+\boldsymbol b
}.
\]

若
\(\lambda_0=0<\lambda_1\le\cdots\le\lambda_{N-1}\)，对应规范化本征向量为
\(\boldsymbol u_k\)，则

\[
\boxed{
R_{ab}=\sum_{k=1}^{N-1}
\frac{|u_k(a)-u_k(b)|^2}{\lambda_k}
}.
\]

零模态不贡献电势差。小的非零本征值对应缓慢变化的电势模式，因分母较小，往往支配远距离等效电阻。

\par\addvspace{.4\baselineskip}
\noindent\begin{minipage}{\linewidth}

\textbf{例 M4-082：完全图的谱与等效电阻}

有限完全图 \(K_N\) 含 \(N\ge2\) 个节点，每一对不同节点间都连接一只
\(r>0\) 电阻。向节点 \(a\) 注入测试电流 \(I\)，从另一节点 \(b\)
抽出同量电流，其余节点无外部注入。写出加权拉普拉斯在零和子空间上的作用，求端口等效电阻
\(R_{ab}\)，说明全1零模怎样从谱逆中排除。

\end{minipage}\par

\par\noindent\begin{minipage}{\linewidth}
\textbf{解答}

\begin{center}
\includegraphics[width=\linewidth,height=.40\textheight,keepaspectratio]{\detokenize{assets/module-04/examples/M4-082.pdf}}
\par\smallskip{\small 图 M4-082：完全图在零和子空间只有一个非零谱值}
\end{center}
\end{minipage}\par

图 M4-082 用五节点完全图示意任意
\(K_N\)，红端点给平衡的源汇向量。所有非零模式具有同一特征值，而端口向量与全1零模正交，因而谱逆只需在零和子空间取
\(r/N\)。

完全图的加权拉普拉斯为
\(L=(N\mathsf I-\mathbf1\mathbf1^{\mathsf T})/r\)。令
\(\boldsymbol b=\boldsymbol e_a-\boldsymbol e_b\)，有
\(\mathbf1^{\mathsf T}\boldsymbol b=0\)、\(\boldsymbol b^{\mathsf T}\boldsymbol b=2\)；在零和子空间
\(L^+=r\mathsf I/N\)，所以
\(R_{ab}=\boldsymbol b^{\mathsf T}L^+\boldsymbol b=2r/N\)。

\subsection{三、狄利克雷原理：固定端口电势（拔高选讲）}\label{m4:ux4e09ux72c4ux5229ux514bux96f7ux539fux7406ux56faux5b9aux7aefux53e3ux7535ux52bfux62d4ux9ad8ux9009ux8bb2}

把端点 \(a,b\) 的电势固定为 \(1,0\)，其余节点电势任意。定义耗散泛函

\[
\mathcal P[\boldsymbol V]
=\sum_{(i,j)\in E}c_{ij}(V_i-V_j)^2.
\]

对任一内部节点 \(k\) 求偏导：

\[
\frac{\partial\mathcal P}{\partial V_k}
=2\sum_jc_{kj}(V_k-V_j).
\]

驻值条件正是内部节点的
KCL。由于泛函为正定二次型，真实电势分布给出唯一最小值。端口电压为 \(1\)
时，\(\mathcal P=I=G_{ab}\)，所以

\[
\boxed{
G_{ab}=\min_{V_a=1,V_b=0}\mathcal P[\boldsymbol V]
}.
\]

任意满足边界电势的试探分布都会给出 \(G_{ab}\) 的上界，从而给出
\(R_{ab}\) 的下界。

\subsection{四、汤姆孙原理：固定端口电流（拔高选讲）}\label{m4:ux56dbux6c64ux59c6ux5b59ux539fux7406ux56faux5b9aux7aefux53e3ux7535ux6d41ux62d4ux9ad8ux9009ux8bb2}

给每条有向边指定试探电流 \(J_e\)，要求除 \(a,b\) 外所有节点满足
KCL，并令从 \(a\) 到 \(b\) 的总流量为 \(1\)。定义

\[
\mathcal Q[\boldsymbol J]
=\sum_er_eJ_e^2.
\]

真实欧姆电流使该泛函最小：

\[
\boxed{
R_{ab}=\min_{\text{单位 }a\to b\text{ 流}}\mathcal Q[\boldsymbol J]
}.
\]

证明时把任一试探流写成真实流 \(\boldsymbol I\) 与零散度环流
\(\boldsymbol K\) 之和。交叉项为

\[
\sum_er_eI_eK_e
=\sum_e(V_i-V_j)K_e=0,
\]

最后一步由内部节点 KCL 和端点净流量为零得到。因此

\[
\mathcal Q[\boldsymbol I+\boldsymbol K]
=\mathcal Q[\boldsymbol I]+\mathcal Q[\boldsymbol K]
\ge\mathcal Q[\boldsymbol I].
\]

任意合法试探流给出 \(R_{ab}\)
的上界。它与狄利克雷原理一电势一电流，构成相互对偶的估计工具。

\subsection{五、瑞利单调性的含义（拔高选讲）}\label{m4:ux4e94ux745eux5229ux5355ux8c03ux6027ux7684ux542bux4e49ux62d4ux9ad8ux9009ux8bb2}

若保持拓扑和端点不变，只增大某些支路电阻，则端点等效电阻不会减小。由汤姆孙原理，对任一单位流，\(\sum_er_eJ_e^2\)
随 \(r_e\) 增大而不减；再对全部合法流取最小值，得到

\[
\boxed{r'_e\ge r_e\ \forall e
\quad\Longrightarrow\quad
R'_{ab}\ge R_{ab}}.
\]

等价地，增加一条正电导支路或把某些节点短接，只会使等效电阻不增；删边或把支路断开，只会使等效电阻不减。该定理能给出有限截断的上下界，但若上下界没有趋于同一极限，单调性本身不能确定无穷网络的答案。

\par\noindent\begin{minipage}{\linewidth}
\begin{center}
\includegraphics[width=\linewidth,height=.40\textheight,keepaspectratio]{\detokenize{assets/module-04/original/tikz-m4-topic40-v9-from37-v2-02-variational.png}}
\par\smallskip{\small 图示：狄利克雷原理用满足端点电势的试探电势给出电导上界，汤姆孙原理用满足节点守恒的单位试探流给出电阻上界}
\end{center}
\end{minipage}\par

\subsection{六、无限条带的横向模态分解（拔高选讲）}\label{m4:ux516dux65e0ux9650ux6761ux5e26ux7684ux6a2aux5411ux6a21ux6001ux5206ux89e3ux62d4ux9ad8ux9009ux8bb2}

设网络沿整数列 \(n\) 平移不变，每列有 \(m\)
个节点；同一列内的连接由无量纲横向拉普拉斯矩阵 \(L_{\perp}\)
描述，相邻两列对应节点之间均以电阻 \(r\) 相连。令第 \(n\) 列电势向量为
\(\boldsymbol V_n\)。无源列满足

\[
(2\boldsymbol I+L_{\perp})\boldsymbol V_n
-\boldsymbol V_{n-1}-\boldsymbol V_{n+1}=\boldsymbol0.
\]

取 \(L_{\perp}\) 的规范化本征向量

\[
L_{\perp}\boldsymbol u_q=\mu_q\boldsymbol u_q,
\]

并展开 \(\boldsymbol V_n=\sum_qv_q(n)\boldsymbol u_q\)，各模式彼此分离：

\[
(2+\mu_q)v_q(n)-v_q(n-1)-v_q(n+1)=0.
\]

对 \(\mu_q>0\)，设 \(v_q(n)\propto\rho_q^{|n|}\)，得到

\[
\rho_q+\rho_q^{-1}=2+\mu_q,
\]

取衰减根

\[
\boxed{
\rho_q=\frac{2+\mu_q-
\sqrt{\mu_q(\mu_q+4)}}{2}},
\qquad0<\rho_q<1.
\]

无限直线上的离散格林函数为

\[
\boxed{
g_{\mu_q}(n)
=\frac{\rho_q^{|n|}}
{\sqrt{\mu_q(\mu_q+4)}}
}\qquad(\mu_q>0).
\]

\(\mu_0=0\)
的均匀模式不能套用这一指数衰减式。它对应各排电势相同、总电流沿条带输运，源汇分离时产生随距离线性增长的电压降，必须单独处理。

\par\noindent\begin{minipage}{\linewidth}
\begin{center}
\includegraphics[width=\linewidth,height=.40\textheight,keepaspectratio]{\detokenize{assets/module-04/theory/review-tikz-m4-topic40-v9-from37-v1-03-strip-modes-safe-v7.pdf}}
\par\smallskip{\small 图示：无限条带先按横向本征模态分解，再分别处理轴向均匀模式和衰减模式}
\end{center}
\end{minipage}\par

\subsection{七、周期与半无限边界（拔高选讲）}\label{m4:ux4e03ux5468ux671fux4e0eux534aux65e0ux9650ux8fb9ux754cux62d4ux9ad8ux9009ux8bb2}

若横向截面是含 \(m\) 个节点的环，离散傅里叶模态为

\[
u_q(p)=\frac1{\sqrt m}e^{2\pi iqp/m},
\qquad
\mu_q=2-2\cos\frac{2\pi q}{m}.
\]

共轭的 \(q\) 与 \(m-q\) 模式必须合并，最终电势和电阻为实数。

若轴向只保留
\(n=0,1,2,\ldots\)，端面节点缺少一个邻点，边界方程不同于内部差分方程。对绝缘端面，可用偶对称镜像得到

\[
\boxed{
G_{\mu}(j,k)=g_{\mu}(j-k)+g_{\mu}(j+k+1)
}.
\]

第二项是端面镜像，不能从无限直线公式中省略。综合题中的双线条带、三排条带和半无限圆柱网络都将直接调用这些结论。

\par\addvspace{.4\baselineskip}
\noindent\begin{minipage}{\linewidth}

\subsection{八、例题：双线无限条带的模式结构}\label{m4:ux516bux4f8bux9898ux53ccux7ebfux65e0ux9650ux6761ux5e26ux7684ux6a21ux5f0fux7ed3ux6784}

\textbf{题号：M4-083}

双线电阻条带的节点为上下两排
\((s,n)\)，\(s=1,2\)，\(n\in\mathbb Z\)；所有同列竖直边及同排相邻列水平边的电阻均为
\(r>0\)。在无注入列，对列电势向量作横向本征模分解。求横向无量纲拉普拉斯的规范化本征向量及本征值，写出各模沿轴向的递推关系，求局域反相扰动的衰减比，并说明同相零模允许的轴向线性部分为何不能一概舍去。

\end{minipage}\par

\par\noindent\begin{minipage}{\linewidth}
\textbf{解答}

\begin{center}
\includegraphics[width=\linewidth,height=.40\textheight,keepaspectratio]{\detokenize{assets/module-04/examples/M4-083.pdf}}
\par\smallskip{\small 图 M4-083：双线条带：同相零模与反相衰减模}
\end{center}
\end{minipage}\par

图 M4-083 每列两个节点记为
\(V_{1,n},V_{2,n}\)。无外部注入时，上节点流向前、后和下节点的三路电流之和为零；各边均为
\(r\)，乘以 \(r\) 得 \[
3V_{1,n}-V_{1,n-1}-V_{1,n+1}-V_{2,n}=0.
\]
下节点交换下标1、2有同样的式子。将两式相加、相减，便看出两种独立变化：
\[
q_{0,n}=\frac{V_{1,n}+V_{2,n}}{\sqrt2},\qquad
q_{1,n}=\frac{V_{1,n}-V_{2,n}}{\sqrt2}.
\]
同相量表示上下电势同时升降，竖边无压降；反相量表示上下电势差，竖边参与回流。轴向递推分别为
\[
q_{0,n+1}-2q_{0,n}+q_{0,n-1}=0,\qquad
q_{1,n+1}-4q_{1,n}+q_{1,n-1}=0.
\] 第一式说明逐列电势差为常数，故 \(q_{0,n}=A+Bn\)。第二式试取
\(q_{1,n}=C\rho^n\)，得
\(\rho^2-4\rho+1=0\)；在远离局域扰动的方向保留衰减根， \[
\boxed{\rho=2-\sqrt3}.
\]
用题目要求的矩阵语言写，\(L_{\perp}=\begin{pmatrix}1&-1\\-1&1\end{pmatrix}\)。规范化本征向量为
\(\boldsymbol u_0=(1,1)^{\mathsf T}/\sqrt2\)、\(\boldsymbol u_1=(1,-1)^{\mathsf T}/\sqrt2\)，本征值为0、2；一般模式递推为
\(q_{n+1}-(2+\mu)q_n+q_{n-1}=0\)。

若有非零总电流沿条带输送，就需要恒定的逐列压降，不能只因反相量衰减便把
\(Bn\)
一并删去。若另行要求远端有界且无轴向输运，才根据这些边界条件限制该项。

\subsection{九、练习与解答}\label{m4:ux4e5dux7ec3ux4e60ux4e0eux89e3ux7b54-5}

\par\addvspace{.4\baselineskip}
\noindent\begin{minipage}{\linewidth}

\subsubsection{作业一：环形电阻网络的结构与谱公式}\label{m4:ux4f5cux4e1aux4e00ux73afux5f62ux7535ux963bux7f51ux7edcux7684ux7ed3ux6784ux4e0eux8c31ux516cux5f0f}

\textbf{题号：M4-084}

\textbf{层级：} 决赛　　\textbf{方法：} 网络方程

\(N\ge3\) 为整数，\(r>0\)，环上没有其他支路；测量端为节点0与节点
\(d\)，其余节点无外部注入。\(d\) 与 \(N-d\)
表示同一对端口的两条互补弧长。

\(N\) 个等值电阻 \(r\) 首尾相接成环。两端点沿一个方向相隔 \(d\)
条边，\(1\le d\le N-1\)。用网络结构或谱观点求两点间等效电阻，并检查
\(d\leftrightarrow N-d\) 对称性。

\end{minipage}\par

\par\noindent\begin{minipage}{\linewidth}
\textbf{解答}

\begin{center}
\includegraphics[width=\linewidth,height=.40\textheight,keepaspectratio]{\detokenize{assets/module-04/examples/M4-084.pdf}}
\par\smallskip{\small 图 M4-084：环形网络的结构与谱公式}
\end{center}
\end{minipage}\par

图 M4-084
中两端把环分成含d条边与N-d条边的互补弧。每条弧内部没有分叉，两弧先各自串联再并联，所得结构也解释了交换d与N-d的对称性。

\[
R_d=dr\parallel(N-d)r
=\frac{d(N-d)}{N}r.
\]

故

\[
\boxed{R_d=\frac{d(N-d)}N\,r}.
\]

结果在 \(d\leftrightarrow N-d\)
下不变。谱公式会把同一结果写成所有非零傅里叶模态贡献的和；串并联解在本题更短，而谱方法可直接推广到多排周期网络。

\subsection{十、使用条件}\label{m4:ux5341ux4f7fux7528ux6761ux4ef6-2}

\begin{enumerate}
\def\labelenumi{\arabic{enumi}.}
\tightlist
\item
  图拉普拉斯矩阵含零模态，不能在完整空间中直接求普通逆；应固定参考电势、使用约化矩阵或广义逆。
\item
  谱分解适合处理平移对称、多排、周期和无限网络；简单串并联题无须使用。
\item
  无限网络必须说明有限截断、边界条件和取极限方式；形式上的衰减根不能代替存在性与收敛性论证。
\item
  \(\mu=0\) 均匀模式要单独处理，不能代入只对 \(\mu>0\) 成立的格林函数。
\item
  半无限端面的缺邻点会改变边界方程，必须加入镜像项或直接求解端面条件。
\item
  狄利克雷试探电势给出电导上界，汤姆孙试探电流给出电阻上界，使用时不要把不等号方向颠倒。
\end{enumerate}

\section{模块四小结}\label{m4:ux6a21ux5757ux56dbux5c0fux7ed3}

{\def\LTcaptype{none} % do not increment counter
\begin{longtable}[]{@{}
  >{\raggedright\arraybackslash}p{(\linewidth - 4\tabcolsep) * \real{0.3333}}
  >{\raggedright\arraybackslash}p{(\linewidth - 4\tabcolsep) * \real{0.3333}}
  >{\raggedright\arraybackslash}p{(\linewidth - 4\tabcolsep) * \real{0.3333}}@{}}
\toprule\noalign{}
\begin{minipage}[b]{\linewidth}\raggedright
检验问题
\end{minipage} & \begin{minipage}[b]{\linewidth}\raggedright
结论或判断
\end{minipage} & \begin{minipage}[b]{\linewidth}\raggedright
回看位置
\end{minipage} \\
\midrule\noalign{}
\endhead
\bottomrule\noalign{}
\endlastfoot
电子漂移方向和常规电流方向是否相同？ & 相反 & 专题二十七 \\
稳恒电流界面为何可能有面电荷？ & 面电荷建立满足边界条件的电场 &
专题三十 \\
KVL 在变化磁通回路中如何写？ & 必须把感应电动势计入 & 专题三十一 \\
仪表内阻何时不能忽略？ & 负载效应与待测量同量级时 & 专题三十二 \\
合并节点前需要证明什么？ & 两点在当前激励下等势 & 专题三十五 \\
无穷网络为何先取有限极限？ & 保证电势差和能量有明确的定义 &
专题三十八至四十 \\
\end{longtable}
}

\section{模块四综合题：稳恒电流、等效电源与复杂电阻网络}\label{m4:ux6a21ux5757ux56dbux7efcux5408ux9898ux7a33ux6052ux7535ux6d41ux7b49ux6548ux7535ux6e90ux4e0eux590dux6742ux7535ux963bux7f51ux7edc}

复杂网络题要先把拓扑转化为合适的未知量。下面的综合题包括连续导体、电流折射、含源一端口、自相似网络以及有限或半无限周期网络；每题都明确端口与边界，并用量纲、简单特例或极限检查结果。

\par\noindent\begin{minipage}{\linewidth}
\begin{center}
\includegraphics[width=\linewidth,height=.40\textheight,keepaspectratio]{\detokenize{assets/module-04/theory/tikz-m4-comprehensive-v3-01-map.pdf}}
\par\smallskip{\small 图示：复杂电阻网络应先辨认拓扑与边界，再选择递推、差分或模态方法}
\end{center}
\end{minipage}\par

\par\addvspace{.4\baselineskip}
\noindent\begin{minipage}{\linewidth}

\subsection{综合题一：锥形导体的电阻与焦耳热分布}\label{m4:ux7efcux5408ux9898ux4e00ux9525ux5f62ux5bfcux4f53ux7684ux7535ux963bux4e0eux7126ux8033ux70edux5206ux5e03}

\textbf{题号：M4-085}

\textbf{层级：} 复赛　　\textbf{方法：} 导电介质、焦耳热

一段均匀各向同性导体的电导率为 \(\sigma>0\)，轴向坐标为
\(0\le x\le L\)，\(L>0\)。端面由理想电极完全覆盖，侧壁绝缘，圆截面半径缓慢变化为
\(a(x)=a_0+(a_1-a_0)x/L\)，\(a_0,a_1>0\)。采用准一维近似，认为每个截面等势、轴向电流密度近似均匀，忽略端部与局部横向修正。求两端电阻；通以固定稳恒电流
\(I\) 时，求单位轴向长度焦耳热功率，判断哪端更易发热，并检验
\(a_1\to a_0\) 的圆柱极限。

\end{minipage}\par

\par\noindent\begin{minipage}{\linewidth}
\textbf{解答}

\begin{center}
\includegraphics[width=\linewidth,height=.40\textheight,keepaspectratio]{\detokenize{assets/module-04/examples/M4-085.pdf}}
\par\smallskip{\small 图 M4-085：锥形导体的电阻与焦耳热}
\end{center}
\end{minipage}\par

图 M4-085
的圆截面半径随轴向改变。取一段薄片时电流相同、面积不同，局部电阻为dx除以电导率和截面积；同一分母也决定该处单位长度焦耳热，窄端因此更热。

截面积为 \(S(x)=\pi a(x)^2\)。稳恒时每个截面通过同一电流，故

\[
j(x)=\frac{I}{S(x)},
\qquad
E(x)=\frac{I}{\sigma S(x)}.
\]

微元电压降为 \(dU=E(x)dx\)，所以

\[
R=\frac1\sigma\int_0^L\frac{dx}{\pi a(x)^2}.
\]

\par\noindent\begin{minipage}{\linewidth}
\begin{center}
\includegraphics[width=\linewidth,height=.40\textheight,keepaspectratio]{\detokenize{assets/module-04/process/M4-085a.pdf}}
\par\smallskip{\small 图 M4-085a：相同电流下单位长度发热随面积改变}
\end{center}
\end{minipage}\par

图 M4-085a
的曲线仅以半径从a0增至2a0作示意，通式保留任意正半径。固定电流穿过每个截面时，单位长度电阻与截面积成反比，所以窄处发热较强；若两端宽窄反转，曲线方向也反转。

令 \(k=(a_1-a_0)/L\)，积分得

\[
\int_0^L\frac{dx}{(a_0+kx)^2}
=\frac{L}{a_0a_1}.
\]

因此

\[
\boxed{R=\frac{L}{\pi\sigma a_0a_1}}.
\]

单位长度的焦耳热功率为

\[
\boxed{
\frac{dP}{dx}=jE S
=\frac{I^2}{\sigma\pi a(x)^2}
}.
\]

半径较小的一端单位长度发热更强。当 \(a_1\to a_0=a\) 时，结果回到
\(R=L/(\sigma\pi a^2)\)。本题的截面均匀近似要求
\(|da/dx|\ll1\)；若锥度很大，电流线的横向分量不可忽略，应求解完整的拉普拉斯边值问题。

\par\addvspace{.4\baselineskip}
\noindent\begin{minipage}{\linewidth}

\subsection{综合题二：导电介质界面的电流折射与自由面电荷}\label{m4:ux7efcux5408ux9898ux4e8cux5bfcux7535ux4ecbux8d28ux754cux9762ux7684ux7535ux6d41ux6298ux5c04ux4e0eux81eaux7531ux9762ux7535ux8377}

\textbf{题号：M4-086}

\textbf{层级：} 复赛　　\textbf{方法：} 导电介质、边界条件

两种均匀线性各向同性漏电介质以平面相接，电导率与绝对电容率分别为
\((\sigma_1,\varepsilon_1)\)、\((\sigma_2,\varepsilon_2)\)，均为正常数。单位法向
\(\mathbf n\) 从介质1指向介质2，切向 \(\mathbf t\)
取入射电流切向分量的方向。稳恒时介质1紧邻界面的电流密度大小为
\(J_1\)，与 \(+\mathbf n\) 夹角为
\(0\le\theta_1<\pi/2\)，界面无外加面电流源且不再积累电荷，但允许已经存在自由面电荷。求介质2的电流方向角
\(\theta_2\)
及界面自由面电荷密度，说明当法向电流非零时，什么参数关系可使电导率虽突变而自由面电荷为零。

\end{minipage}\par

\par\noindent\begin{minipage}{\linewidth}
\textbf{解答}

\begin{center}
\includegraphics[width=\linewidth,height=.40\textheight,keepaspectratio]{\detokenize{assets/module-04/examples/M4-086.pdf}}
\par\smallskip{\small 图 M4-086：电流折射与自由面电荷}
\end{center}
\end{minipage}\par

图 M4-086
在同一法向、切向上分解两侧电流。法向电流由稳恒守恒连续，切向电场由无旋条件连续；自由面电荷应由两侧电位移差计算，而不能把电流连续误写成电位移连续。

稳恒电荷守恒和静电型切向条件给出

\[
j_{1n}=j_{2n}=j_n,
\qquad
E_{1t}=E_{2t}.
\]

利用 \(\mathbf J_i=\sigma_i\mathbf E_i\)：

\[
\frac{j_{1t}}{\sigma_1}
=\frac{j_{2t}}{\sigma_2}.
\]

\par\noindent\begin{minipage}{\linewidth}
\begin{center}
\includegraphics[width=\linewidth,height=.40\textheight,keepaspectratio]{\detokenize{assets/module-04/process/M4-086a.pdf}}
\par\smallskip{\small 图 M4-086a：稳恒传导比较J，静电跃变比较D}
\end{center}
\end{minipage}\par

图 M4-086a
把两种边界关系并排使用。图示方向仅说明电流折射的可能形态；通式不依赖某个指定角度，法向连续的是J，而自由面电荷由D的法向跳变决定。

因此

\[
\boxed{
\frac{\tan\theta_1}{\tan\theta_2}
=\frac{\sigma_1}{\sigma_2}
}.
\]

界面自由面电荷由电位移法向跃变决定：

\[
\sigma_s=\varepsilon_2E_{2n}-\varepsilon_1E_{1n}
=j_n\left(\frac{\varepsilon_2}{\sigma_2}
-\frac{\varepsilon_1}{\sigma_1}\right).
\]

故

\[
\boxed{
\sigma_s=j_n\left(\tau_{c2}-\tau_{c1}\right),
\qquad\tau_{ci}=\frac{\varepsilon_i}{\sigma_i}
}.
\]

当两侧电荷弛豫时间相等时，即使 \(\sigma_1\ne\sigma_2\)，也有
\(\sigma_s=0\)。这说明电流折射由电导率比决定，而界面电荷还同时依赖介电常数。

\par\addvspace{.4\baselineskip}
\noindent\begin{minipage}{\linewidth}

\subsection{综合题三：含源桥式网络的等效与最大功率}\label{m4:ux7efcux5408ux9898ux4e09ux542bux6e90ux6865ux5f0fux7f51ux7edcux7684ux7b49ux6548ux4e0eux6700ux5927ux529fux7387}

\textbf{题号：M4-087}

\textbf{层级：} 复赛　　\textbf{方法：} 能量法、等效变换、网络方程

全部 \(R_i>0\)，\(\mathcal E\) 为恒定理想源，\(R_L>0\)
可调，导线无电阻。支路只按所列 \(P-A-Q\)、\(P-B-Q\) 和负载 \(A-B\)
连接。还应说明桥平衡、开路电压为零时最优负载的退化情形。

理想电压源把节点 \(P\) 维持在电势 \(\mathcal E\)，节点 \(Q\)
接地。左支路 \(P-A-Q\) 的电阻为 \(R_1,R_2\)，右支路 \(P-B-Q\) 的电阻为
\(R_3,R_4\)。在 \(A,B\) 间接可调负载 \(R_L\)。求使负载功率最大的 \(R_L\)
和最大功率。

\end{minipage}\par

\par\noindent\begin{minipage}{\linewidth}
\textbf{解答}

\begin{center}
\includegraphics[width=\linewidth,height=.40\textheight,keepaspectratio]{\detokenize{assets/module-04/examples/M4-087.pdf}}
\par\smallskip{\small 图 M4-087：含源桥式网络与最大功率}
\end{center}
\end{minipage}\par

图 M4-087
的负载位于两个分压中点之间。先移去负载得到有符号开路电压，再将独立电源置零求端口内阻；最大功率计算应针对这个端口，而不是任意一个桥臂。

先断开负载，得到开路端口电压

\[
U_{\rm oc}
=\mathcal E\left(
\frac{R_2}{R_1+R_2}
-\frac{R_4}{R_3+R_4}
\right).
\]

把理想电压源置零，\(P,Q\) 合并为同一节点，故

\[
R_{\rm th}
=(R_1\parallel R_2)+(R_3\parallel R_4).
\]

负载功率为

\[
P_L=\frac{U_{\rm oc}^2R_L}{(R_{\rm th}+R_L)^2}.
\]

\par\noindent\begin{minipage}{\linewidth}
\begin{center}
\includegraphics[width=\linewidth,height=.40\textheight,keepaspectratio]{\detokenize{assets/module-04/process/M4-087a.pdf}}
\par\smallskip{\small 图 M4-087a：负载匹配的是端口内阻}
\end{center}
\end{minipage}\par

图 M4-087a
的横坐标是负载与端口内阻之比；峰值在1而非某条单独桥臂的阻值。归一化图要求开路电压非零；平衡桥的开路电压为零时，所有正负载的功率都为零。

因此

\[
\boxed{R_L=R_{\rm th}},
\]

\[
\boxed{
P_{L,\max}
=\frac{\mathcal E^2}{4R_{\rm th}}
\left(
\frac{R_2}{R_1+R_2}
-\frac{R_4}{R_3+R_4}
\right)^2
}.
\]

若桥平衡，括号为零，端口不能向任何被动负载输出功率；这与负载支路电流为零一致。

若桥平衡而 \(U_{\rm oc}=0\)，则任何 \(R_L>0\) 都有
\(I_L=P_L=0\)；最大功率为零，最优正负载不唯一。通常的唯一匹配点结论应在
\(U_{\rm oc}\ne0\) 时使用。

\par\addvspace{.4\baselineskip}
\noindent\begin{minipage}{\linewidth}

\subsection{综合题四：无限嵌套三角形}\label{m4:ux7efcux5408ux9898ux56dbux65e0ux9650ux5d4cux5957ux4e09ux89d2ux5f62}

\textbf{题号：M4-088}

\textbf{层级：} 复赛　　\textbf{方法：} 等效变换、网络方程

三端无源网络按嵌套三角形定义。最外层顶点为 \(A,B,C\)，三条边的总电阻各为
\(r>0\)。沿固定绕向 \(A\to B\to C\to A\)，每边先经 \(r/3\) 再经
\(2r/3\)，分点依次为
\(P,Q,S\)，这三个点成为下一层三角形顶点。下一层每条边的总电阻仍为
\(r\)，仍沿 \(P\to Q\to S\to P\)
作相同比例分割；以后逐层重复，几何缩小不意味着电阻按比例减小。只有明确的顶点和分点为节点，图中其他线条交叉不连接。用保留有限层、最内层三角形边完整存在的网络取极限。求最外层端点
\(A,B\) 间电阻，第三端 \(C\)
无外部注入，并通过有限层及单调性检查所选极限。

\end{minipage}\par

\par\noindent\begin{minipage}{\linewidth}
\textbf{解答}

\begin{center}
\includegraphics[width=\linewidth,height=.40\textheight,keepaspectratio]{\detokenize{assets/module-04/examples/M4-088.pdf}}
\par\smallskip{\small 图 M4-088：无限嵌套三角形}
\end{center}
\end{minipage}\par

图 M4-088
固定每条边的分割绕向，三个分点对应内侧三端尾网。把内网替成三角形时，应先将所求二端电阻换算为三角形单边参数，再对三个局部三角形作同一变换。

网络具有三重旋转对称性。若把整个无限网络等效为边阻均为 \(R_{\triangle}\)
的三角形，则任意相邻两端间

\[
R_{AB}=R_{\triangle}\parallel 2R_{\triangle}=\frac23R_{\triangle}.
\]

令

\[
R_{AB}=xr,
\]

便有

\[
R_{\triangle}=\frac32xr.
\]

剥去最外层后，内侧无限网络仍与整体同构，因此可以把内侧网络替换成边阻为
\(3xr/2\) 的三角形。考察任一局部三角形，其三边阻值为

\[
\frac r3,\qquad \frac{2r}{3},\qquad \frac32xr.
\]

三边之和为

\[
\Sigma=r\left(1+\frac32x\right).
\]

\par\noindent\begin{minipage}{\linewidth}
\begin{center}
\includegraphics[width=\linewidth,height=.40\textheight,keepaspectratio]{\detokenize{assets/module-04/process/M4-088a.pdf}}
\par\smallskip{\small 图 M4-088a：先把尾网换成三角形单边参数}
\end{center}
\end{minipage}\par

图 M4-088a
只取外角A附近的局部三角形，底边代表完整内网的多端等效单边。它不等于内网二端测量值xr；按三条真实边作三角形---星形变换后，才能拼回三重对称的整网。

把该三角形作 \(\triangle-Y\) 变换。星形中心分别通向图 M4-088a 中
\(S,A,P\) 的三条臂记为 \(r_1,r_2,r_3\)，得到

\[
r_1=\frac{x}{1+3x/2}r,\qquad
r_2=\frac{2/9}{1+3x/2}r,\qquad
r_3=\frac{x/2}{1+3x/2}r.
\]

三个角同时变换后，端点 \(A,B\) 各串联一个 \(r_2\)；中部是边阻均为
\(r_1+r_3\) 的三角形。因此

\[
xr=2r_2+\frac23(r_1+r_3).
\]

约去 \(r\) 并代入上式：

\[
x=\frac{4/9+x}{1+3x/2}.
\]

于是

\[
\frac32x^2=\frac49.
\]

电阻必须为正，故

\[
\boxed{R_{AB}=\frac{2\sqrt6}{9}r}.
\]

\textbf{检验}

只保留最外层三角形时，\(R_{AB}=2r/3\)；逐层加入内部支路只会减小等效电阻。所得

\[
\frac{2\sqrt6}{9}r\approx0.5443r<\frac23r
\]

符合瑞利单调性。有限 \(1,2,3\) 层的数值依次为
\(0.5556r,0.5455r,0.5444r\)，也确实收敛到上述结果。

\par\addvspace{.4\baselineskip}
\noindent\begin{minipage}{\linewidth}

\subsection{综合题五：有限方形梯形的同侧端口}\label{m4:ux7efcux5408ux9898ux4e94ux6709ux9650ux65b9ux5f62ux68afux5f62ux7684ux540cux4fa7ux7aefux53e3}

\textbf{题号：M4-089}

\textbf{层级：} 复赛　　\textbf{方法：} 等效变换、网络方程

有限双导轨梯形含 \(N\ge1\) 个方形网孔，节点为上下两排、列号
\(j=0,1,\ldots,N\)。同排相邻列间及每列上下节点间各接一只 \(r>0\)
电阻，共 \(3N+1\) 条边，左右端之外没有其他连接。端口 \(A,B\)
分别为第0列上、下节点，其余节点无外部注入。求输入等效电阻
\(R_N\)，写出首级递推及初始条件，并检验 \(N=1\) 和 \(N\to\infty\)
的结果。

\end{minipage}\par

\par\noindent\begin{minipage}{\linewidth}
\textbf{解答}

\begin{center}
\includegraphics[width=\linewidth,height=.40\textheight,keepaspectratio]{\detokenize{assets/module-04/examples/M4-089.pdf}}
\par\smallskip{\small 图 M4-089：有限梯形的同侧端口}
\end{center}
\end{minipage}\par

图 M4-089
保留有限梯形最左端的竖直电阻，它与经过一对水平电阻再进入尾部的路线并联。有限右端和左端有同样的竖直边，故递推需从单根竖直电阻的初值起算。

规定 \(N=0\) 时只剩一个竖直电阻，所以

\[
R_0=r.
\]

对 \(N\ge1\)，从端口看入有两条并联支路：一条是左端竖直电阻
\(r\)；另一条依次通过上边 \(r\)、余下 \(N-1\) 格的端口等效
\(R_{N-1}\)、下边 \(r\)。故

\[
\boxed{R_N=r\parallel(2r+R_{N-1})}
\]

即

\[
R_N=r\frac{2r+R_{N-1}}{3r+R_{N-1}}.
\]

令 \(x_N=R_N/r\)，得到 Möbius 递推

\[
x_N=\frac{x_{N-1}+2}{x_{N-1}+3},\qquad x_0=1.
\]

其两个不动点为

\[
a=\sqrt3-1,\qquad b=-\sqrt3-1.
\]

直接代数化简可得

\[
\frac{x_N-a}{x_N-b}=(7-4\sqrt3)\frac{x_{N-1}-a}{x_{N-1}-b}.
\]

迭代并还原 \(x_N\)，等价地写成

\[
\boxed{
R_N=r\frac{(2+\sqrt3)^N(1+\sqrt3)-(2-\sqrt3)^N(1-\sqrt3)}
{(2+\sqrt3)^{N+1}-(2-\sqrt3)^{N+1}}
}.
\]

\textbf{例}

当 \(N=1\) 时，一个边阻均为 \(r\) 的正方形从相邻顶点看入：

\[
R_1=r\parallel3r=\frac34r,
\]

与闭式一致。

\textbf{极限}

因为 \(0<2-\sqrt3<1\)，所以

\[
\lim_{N\to\infty}R_N=(\sqrt3-1)r.
\]

这也可由无限网络自相似方程 \(R=r\parallel(2r+R)\) 直接得到。

若将``零格''定义为仅剩一根竖直电阻，则 \(R_0=r\)。完整递推为
\(R_N=r\parallel(2r+R_{N-1})\)，适用于 \(N\ge1\)，因此 \(N=1\)
的四边环检验与同一初值一致。

\par\noindent\begin{minipage}{\linewidth}
\begin{center}
\includegraphics[width=\linewidth,height=.40\textheight,keepaspectratio]{\detokenize{assets/module-04/process/M4-089a.pdf}}
\par\smallskip{\small 图 M4-089a：一格截断检查递推初值}
\end{center}
\end{minipage}\par

图 M4-089a
的左侧竖直电阻就是端口直接支路，其余三边构成另一条串联路线。它检验有限网络的递推起点；若删去最左竖直边，就不再是这一组有限梯形。

\par\addvspace{.4\baselineskip}
\noindent\begin{minipage}{\linewidth}

\subsection{综合题六：有限方形梯形的对角端口}\label{m4:ux7efcux5408ux9898ux516dux6709ux9650ux65b9ux5f62ux68afux5f62ux7684ux5bf9ux89d2ux7aefux53e3}

\textbf{题号：M4-090}

\textbf{层级：} 复赛　　\textbf{方法：} 等效变换、网络方程

有限双导轨梯形含 \(N\ge1\) 个方形网孔，节点列号为
\(j=0,1,\ldots,N\)，每列上、下节点间及同排相邻列之间各接一只 \(r>0\)
电阻；只有这些 \(3N+1\) 条边，左右端之外无连接。测试电流 \(I\)
从第0列下节点 \(A\) 流入，从第 \(N\) 列上节点 \(B\)
流出，其余节点无注入。求
\(R_{AB}=(V_A-V_B)/I\)，说明同相与反相模态在有限两端边界下如何组合，并与单个方形的对角端口结果核对。

\end{minipage}\par

\par\noindent\begin{minipage}{\linewidth}
\textbf{解答}

\begin{center}
\includegraphics[width=\linewidth,height=.40\textheight,keepaspectratio]{\detokenize{assets/module-04/examples/M4-090.pdf}}
\par\smallskip{\small 图 M4-090：有限梯形的对角端口}
\end{center}
\end{minipage}\par

图 M4-090
的电流从左下端进入、右上端流出。任一竖截面总电流相同，因而上、下水平支路电流互补；内部KVL给非齐次递推，端点的缺边条件由两端边界单独固定。

设第 \(i\) 段下方水平支路中由左向右的电流为
\(j_i\)。任取切过上下两条水平支路的竖直截面，穿过截面的总电流恒为
\(I\)，所以上方第 \(i\) 段支路电流为 \(I-j_i\)。

在第 \(i\) 个内部竖直节点处，由电流守恒，向上的竖直支路电流为

\[
j_{i-1}-j_i.
\]

对第 \(i\) 个方格列回路电压方程，得到

\[
j_{i-1}-4j_i+j_{i+1}=-I,\qquad 1\le i\le N.
\]

用虚拟支路统一端点条件：

\[
j_0=I,\qquad j_{N+1}=0.
\]

\par\noindent\begin{minipage}{\linewidth}
\begin{center}
\includegraphics[width=\linewidth,height=.40\textheight,keepaspectratio]{\detokenize{assets/module-04/process/M4-090a.pdf}}
\par\smallskip{\small 图 M4-090a：一列截面的两条水平电流互补}
\end{center}
\end{minipage}\par

图 M4-090a
中上下参考箭头都向右，任一切面总电流为I，所以上支路是I减ji。相邻竖直支路的电流差进入一个网孔的KVL，产生非齐次项；最左、最右缺边再转换为虚拟端条件。

常数特解为 \(I/2\)。记

\[
\lambda=2+\sqrt3,\qquad \lambda^{-1}=2-\sqrt3,
\]

则通解为

\[
j_i=A\lambda^i+B\lambda^{-i}+\frac I2.
\]

代入两个边界条件可得

\[
A=-\frac I2\frac{1+\lambda^{-(N+1)}}{\lambda^{N+1}-\lambda^{-(N+1)}},\qquad
B=\frac I2\frac{1+\lambda^{N+1}}{\lambda^{N+1}-\lambda^{-(N+1)}}.
\]

从 \(A\) 沿下边到右端，再沿最后一根竖直电阻到 \(B\)：

\[
U_{AB}=r\left(\sum_{i=1}^{N}j_i+j_N\right).
\]

把差分方程改写为

\[
(j_{i-1}-j_i)-(j_i-j_{i+1})=2j_i-I,
\]

从 \(i=1\) 到 \(N\) 求和，利用 \(j_0=I,j_{N+1}=0\)，得到

\[
\sum_{i=1}^{N}j_i=\frac{(N+1)I-j_1-j_N}{2}.
\]

因此

\[
R_{AB}=\frac r2\left(N+1+\frac{j_N-j_1}{I}\right).
\]

代入通解并化简：

\[
\boxed{
R_{AB}=\frac r2\left[
N+1+\frac{2\sqrt3-(\lambda^N-\lambda^{-N})}
{\lambda^{N+1}-\lambda^{-(N+1)}}
\right]
}.
\]

\textbf{检验}

当 \(N=1\) 时是正方形对角顶点，两条 \(2r\) 支路并联，故
\(R_{AB}=r\)；闭式给出同一结果。随 \(N\) 增大，主项为
\(Nr/2\)，表示远距离输运时电流大致平均分配在上下两条导线上。

支路电流通解中的常数特解 \(I/2\) 对应上下两排平均分流的同相输运；余下的
\(\lambda^{\pm i}\)
项描述反相偏差。有限网络同时保留两支，再由左右两个端点条件固定系数，不能像半无限单端问题那样预先只留一支衰减项。

\par\addvspace{.4\baselineskip}
\noindent\begin{minipage}{\linewidth}

\subsection{综合题七：无限双线条带的错位端口}\label{m4:ux7efcux5408ux9898ux4e03ux65e0ux9650ux53ccux7ebfux6761ux5e26ux7684ux9519ux4f4dux7aefux53e3}

\textbf{题号：M4-091}

\textbf{层级：} 复赛　　\textbf{方法：} 等效变换、网络方程

双线条带向左右无限延伸，节点为上、下两排
\((s,j)\)，\(s=1,2\)、\(j\in\mathbb Z\)。每条同排相邻列边与每条同列竖直边的电阻均为
\(r>0\)，没有其他连接。测试电流 \(I\) 从下排第0列 \(A\) 流入，从上排第
\(N\) 列 \(B\) 流出，\(N\ge0\)
为整数。以平衡两端注入、两侧远端无额外输运电流的有限截断极限定义两点电阻。求
\(R_{AB}\)，保留同相模态所需的有限电势差，并检验 \(N=0\) 的同列端口。

\end{minipage}\par

\par\noindent\begin{minipage}{\linewidth}
\textbf{解答}

\begin{center}
\includegraphics[width=\linewidth,height=.40\textheight,keepaspectratio]{\detokenize{assets/module-04/examples/M4-091.pdf}}
\par\smallskip{\small 图 M4-091：无限双线条带的错位端口}
\end{center}
\end{minipage}\par

图 M4-091
将错位源汇投影为同相输运与反相局域两部分。前者保留端点间随列距增长的电势差，后者由衰减格林函数给出有限修正；两项应在同一对真实端点处组装。

把第 \(k\) 列的两排电势写成列向量。横向的两节点拉普拉斯矩阵为

\[
L_{\perp}=
\begin{pmatrix}
1&-1\\
-1&1
\end{pmatrix}.
\]

它有两个单位本征向量：

\[
\boldsymbol u_0=\frac1{\sqrt2}\binom11,\qquad \mu_0=0,
\]

\[
\boldsymbol u_1=\frac1{\sqrt2}\binom1{-1},\qquad \mu_1=2.
\]

\par\noindent\begin{minipage}{\linewidth}
\begin{center}
\includegraphics[width=\linewidth,height=.40\textheight,keepaspectratio]{\detokenize{assets/module-04/process/M4-091a.pdf}}
\par\smallskip{\small 图 M4-091a：错位源汇的输运项与局域项相加}
\end{center}
\end{minipage}\par

图 M4-091a 在下排第0列的 \(A\) 注入电流，在上排第 \(N\) 列的 \(B\)
抽出电流，左右两端都继续延伸。源汇投影后分成同相输运项与反相局域项；两项按同一对
\(A,B\) 端点相加。

对任一横向本征值 \(\mu\)，轴向格林函数 \(g_k\) 在 \(k\ne0\) 处满足

\[
(2+\mu)g_k-g_{k-1}-g_{k+1}=0.
\]

若 \(\mu>0\)，取衰减解 \(g_k=C\rho^{|k|}\)，其中

\[
\rho=\frac{2+\mu-\sqrt{\mu(\mu+4)}}{2},\qquad
C=\frac1{\sqrt{\mu(\mu+4)}}.
\]

同相模式 \(\mu_0=0\)
等价于一条沿轴向的单位电阻链；两个端点在该模式上的投影平方均为
\(1/2\)，故贡献为

\[
r\frac N2.
\]

反相模式中 \(\mu_1=2\)，所以

\[
\rho=2-\sqrt3,\qquad g_k=\frac{\rho^{|k|}}{2\sqrt3}.
\]

\(A\)、\(B\) 位于不同排，二者在反相模式上的投影异号，因此该模式贡献为

\[
r(g_0+g_N)=r\frac{1+\rho^N}{2\sqrt3}.
\]

相加得

\[
\boxed{
R_{AB}=r\left[\frac N2+\frac{1+(2-\sqrt3)^N}{2\sqrt3}\right]
}.
\]

\textbf{例与边界}

当 \(N=0\) 时，端口是同一列的上下节点：

\[
R_{AB}=\frac r{\sqrt3}.
\]

不要把无穷远电势预先规定为零。该网络是准一维网络，输运模式的电势含线性项；只有端点电势差具有确定意义。

\par\addvspace{.4\baselineskip}
\noindent\begin{minipage}{\linewidth}

\subsection{综合题八：环形梯形网络}\label{m4:ux7efcux5408ux9898ux516bux73afux5f62ux68afux5f62ux7f51ux7edc}

\textbf{题号：M4-092}

\textbf{层级：} 复赛　　\textbf{方法：} 等效变换、网络方程

环形双导轨网络有 \(N\ge3\) 列，列号按模 \(N\) 计，每列有外环节点
\(u_j\)、内环节点 \(v_j\)。三类边
\(u_ju_{j+1}\)、\(v_jv_{j+1}\)、\(u_jv_j\) 各接一只 \(r>0\) 电阻，共
\(3N\) 条边，没有其他连接。端口为同列的
\(A=u_0,B=v_0\)，其余节点无外部注入。求两端等效电阻，说明圆周周期条件如何进入横向模式求和，并检查大
\(N\) 极限。

\end{minipage}\par

\par\noindent\begin{minipage}{\linewidth}
\textbf{解答}

\begin{center}
\includegraphics[width=\linewidth,height=.40\textheight,keepaspectratio]{\detokenize{assets/module-04/examples/M4-092.pdf}}
\par\smallskip{\small 图 M4-092：环形梯形网络}
\end{center}
\end{minipage}\par

图 M4-092
的两个端点位于同一径向支路两侧，源汇向量只投影到内外环反相模式。圆周闭合要求叠加所有周期像；周期和的自项再乘端口投影因子，才是两点电阻。

端口电流在某一列为
\((I,-I)\)，与同相模式正交，所以只激发反相模式。把反相模式沿圆周的离散格林函数记为
\(G_N(k)\)，它满足周期边界

\[
G_N(k+N)=G_N(k)
\]

以及

\[
4G_N(k)-G_N(k-1)-G_N(k+1)=\delta_{k0}.
\]

令

\[
\rho=2-\sqrt3.
\]

\par\noindent\begin{minipage}{\linewidth}
\begin{center}
\includegraphics[width=\linewidth,height=.40\textheight,keepaspectratio]{\detokenize{assets/module-04/process/M4-092a.pdf}}
\par\smallskip{\small 图 M4-092a：周期闭合等于叠加相隔N列的像}
\end{center}
\end{minipage}\par

图 M4-092a
沿展开后的整数列轴标出0、正N与负N，后两点是同一周期端口的最近像。每对远像贡献按rho的N次幂继续衰减，求和后再乘反相端口投影因子2r。

无限直线上的对应格林函数为

\[
g(k)=\frac{\rho^{|k|}}{2\sqrt3}.
\]

周期格林函数可由所有周期像叠加：

\[
G_N(0)=\sum_{s=-\infty}^{\infty}g(sN)
=\frac1{2\sqrt3}\left(1+2\sum_{s=1}^{\infty}\rho^{sN}\right).
\]

求等比级数：

\[
G_N(0)=\frac1{2\sqrt3}\frac{1+\rho^N}{1-\rho^N}.
\]

端口向量在反相单位向量上的幅值为 \(\sqrt2I\)，故端口电阻为
\(2rG_N(0)\)：

\[
\boxed{
R_{AB}=\frac r{\sqrt3}\frac{1+(2-\sqrt3)^N}{1-(2-\sqrt3)^N}
}.
\]

也可写成

\[
R_{AB}=\frac r{\sqrt3}\frac{(2+\sqrt3)^N+1}{(2+\sqrt3)^N-1}.
\]

\textbf{检验}

当 \(N\to\infty\) 时，环的闭合处远离端口，得到

\[
R_{AB}\to\frac r{\sqrt3},
\]

恰好等于无限双线条带同一列上下节点间的结果。

\par\addvspace{.4\baselineskip}
\noindent\begin{minipage}{\linewidth}

\subsection{综合题九：三排无限条带}\label{m4:ux7efcux5408ux9898ux4e5dux4e09ux6392ux65e0ux9650ux6761ux5e26}

\textbf{题号：M4-093}

\textbf{层级：} 复赛　　\textbf{方法：} 等效变换、网络方程

三排电阻条带的节点为 \((s,j)\)，排号 \(s=1,2,3\)、列号
\(j\in\mathbb Z\)。每排的相邻列节点相连，每列内只有第1---2排和第2---3排相连，所有边电阻均为
\(r>0\)。端口为同一外侧排上的 \(A=(1,0),B=(1,n)\)，\(n\ge1\)。测试电流从
\(A\) 流入、从 \(B\)
流出，远端没有额外输运电流，以两侧绝缘有限截断的极限定义电阻。求
\(R_n\)，给出横向模态的贡献，并讨论随间隔增大时的主阶行为。

\end{minipage}\par

\par\noindent\begin{minipage}{\linewidth}
\textbf{解答}

\begin{center}
\includegraphics[width=\linewidth,height=.40\textheight,keepaspectratio]{\detokenize{assets/module-04/examples/M4-093.pdf}}
\par\smallskip{\small 图 M4-093：三排无限条带}
\end{center}
\end{minipage}\par

图 M4-093
的三种横向模式与三排真实连接对应。源汇都在同一外侧排，先计算它们在各模上的投影权重，再分别求轴向差分格林函数；零模给线性主项，非零模给端点修正。

第 \(k\) 列的三个电势构成向量 \(\boldsymbol V_k\)。横向拉普拉斯矩阵为

\[
L_{\perp}=
\begin{pmatrix}
1&-1&0\\
-1&2&-1\\
0&-1&1
\end{pmatrix}.
\]

它的三个规范化本征模态是

\[
\mu_0=0,\qquad
\boldsymbol u_0=\frac1{\sqrt3}(1,1,1)^{\mathsf T},
\]

\[
\mu_1=1,\qquad
\boldsymbol u_1=\frac1{\sqrt2}(1,0,-1)^{\mathsf T},
\]

\[
\mu_2=3,\qquad
\boldsymbol u_2=\frac1{\sqrt6}(1,-2,1)^{\mathsf T}.
\]

源、汇都在第一排，所以它们在三个模式上的权重平方依次为

\[
w_0^2=\frac13,\qquad w_1^2=\frac12,\qquad w_2^2=\frac16.
\]

\par\noindent\begin{minipage}{\linewidth}
\begin{center}
\includegraphics[width=\linewidth,height=.40\textheight,keepaspectratio]{\detokenize{assets/module-04/process/M4-093a.pdf}}
\par\smallskip{\small 图 M4-093a：外排源在三个横向模式上的权重}
\end{center}
\end{minipage}\par

图 M4-093a
的源向量只在第1排有非零分量，对三个规范化横向模式作投影后得到右侧权重。源汇位于同一排，所以各非零模读取的是轴向格林函数的自项减互项，而不是错排端口中的相加。

均匀模式 \(\mu_0=0\) 的贡献为

\[
r\frac n3.
\]

对 \(\mu>0\) 的模式，利用上一题已经得到的轴向格林函数

\[
g_{\mu}(k)=\frac{\rho_{\mu}^{|k|}}{\sqrt{\mu(\mu+4)}},\qquad
\rho_{\mu}=\frac{2+\mu-\sqrt{\mu(\mu+4)}}2.
\]

同一排两端的投影同号，因此该模式对两点电阻的贡献为

\[
2rw^2[g_{\mu}(0)-g_{\mu}(n)].
\]

对 \(\mu_1=1\)，

\[
\rho_1=\frac{3-\sqrt5}{2},\qquad
2w_1^2g_1(0)=\frac1{\sqrt5}.
\]

对 \(\mu_2=3\)，

\[
\rho_2=\frac{5-\sqrt{21}}{2},\qquad
2w_2^2g_2(0)=\frac1{3\sqrt{21}}.
\]

故

\[
\boxed{
\frac{R_n}{r}=\frac n3+
\frac1{\sqrt5}\left[1-\left(\frac{3-\sqrt5}{2}\right)^n\right]
+\frac1{3\sqrt{21}}\left[1-\left(\frac{5-\sqrt{21}}2\right)^n\right]
}.
\]

\textbf{检验}

\(n=0\) 时三项都为零。\(n\to\infty\)
时两个非均匀模式只留下有限的端点修正，而均匀模式给出主导项
\(nr/3\)，正对应远处电流平均分配到三排。

\begingroup\ctexset{subsubsection={beforeskip=8pt,afterskip=4pt}}\setlength{\parskip}{1pt}\setlength{\abovedisplayskip}{3pt plus 1pt minus 1pt}\setlength{\belowdisplayskip}{3pt plus 1pt minus 1pt}\setlength{\abovedisplayshortskip}{0pt plus 1pt}\setlength{\belowdisplayshortskip}{3pt plus 1pt minus 1pt}
\par\addvspace{.4\baselineskip}
\noindent\begin{minipage}{\linewidth}

\subsection{综合题十：半无限圆柱形网格}\label{m4:ux7efcux5408ux9898ux5341ux534aux65e0ux9650ux5706ux67f1ux5f62ux7f51ux683c}

\textbf{题号：M4-094}

\textbf{层级：} 复赛　　\textbf{方法：} 等效变换、网络方程

半无限圆柱网格的节点为 \((p,n)\)，\(p=0,1,\ldots,m-1\) 按模 \(m\)
计、\(n=0,1,2,\ldots\)，\(m\ge3\)。每条圆周边 \((p,n)-(p+1,n)\) 与轴向边
\((p,n)-(p,n+1)\) 均为 \(r>0\)，自由端面 \(n=0\)
没有连接到负列的支路。测试电流 \(I\) 从 \(A=(0,0)\) 流入，从不同节点
\(B=(p,n)\)
流出，其余节点无注入；\((p,n)\ne(0,0)\)。以轴向远端不另接电流源或接地极的有限截断极限定义
\(R_{p,n}=(V_A-V_B)/I\)。求真正的两点电阻，说明圆周傅里叶模、端面边界及零模分别如何处理，并检验同端面与远轴向极限。

\end{minipage}\par

\par\noindent\begin{minipage}{\linewidth}
\textbf{解答}

\begin{center}
\includegraphics[width=\linewidth,height=.40\textheight,keepaspectratio]{\detokenize{assets/module-04/examples/M4-094.pdf}}
\par\smallskip{\small 图 M4-094：半无限圆柱网格}
\end{center}
\end{minipage}\par

图 M4-094
的展开图明确圆周周期相接与轴向自由端面，只有指定的黑点是节点。圆周傅里叶分解后，端面因缺少负列邻点而使用不同的首行方程；最后必须用真实源汇对的二次型组装电阻。

\subsubsection{1.
圆周方向的离散傅里叶分解}\label{m4:ux5706ux5468ux65b9ux5411ux7684ux79bbux6563ux5085ux91ccux53f6ux5206ux89e3}

令

\[
\theta_q=\frac{2\pi q}{m},\qquad q=0,1,\ldots,m-1.
\]

圆周拉普拉斯算符的第 \(q\) 个本征值为

\[
\mu_q=2-2\cos\theta_q.
\]

在该模式中，轴向方程化为半直线上的一维方程。对于 \(q\ne0\)，定义衰减根

\[
\rho_q=
\frac{2+\mu_q-\sqrt{\mu_q(\mu_q+4)}}2,\qquad 0<\rho_q<1.
\]

\subsubsection{2.
半无限端面的边界条件}\label{m4:ux534aux65e0ux9650ux7aefux9762ux7684ux8fb9ux754cux6761ux4ef6}

内部层 \(n\ge1\) 的模式电势满足

\[
(2+\mu_q)u_n-u_{n-1}-u_{n+1}=s_n.
\]

端面节点缺少 \(n=-1\) 邻点，因此边界方程是

\[
(1+\mu_q)u_0-u_1=s_0.
\]

把半直线向左镜像延拓，并令

\[
u_{-1}=u_0,
\]

即可把端面方程写成内部方程的同一形式。于是 \(q\ne0\)
模式的半直线格林函数为

\[
G_q(j,k)=
\frac{\rho_q^{|j-k|}+\rho_q^{j+k+1}}
{\sqrt{\mu_q(\mu_q+4)}}.
\]

第一项是实际源，第二项是关于端面 \(-1/2\) 的镜像源。

\par\noindent\begin{minipage}{\linewidth}
\begin{center}
\includegraphics[width=\linewidth,height=.40\textheight,keepaspectratio]{\detokenize{assets/module-04/process/M4-094a.pdf}}
\par\smallskip{\small 图 M4-094a：半无限端面在负半列处作偶延拓}
\end{center}
\end{minipage}\par

图 M4-094a
的镜面对称位置是轴向的负半列，而不是第0列节点本身。令虚拟负一列等于第0列后，端面首行与内部差分方程一致，像源距离因而出现j+k+1。

\subsubsection{3.
组装两点电阻}\label{m4:ux7ec4ux88c5ux4e24ux70b9ux7535ux963b}

第 \(q\) 个圆周模式在 \(A=(0,0)\) 与 \(B=(p,n)\) 处的相位差为
\(e^{i\theta_qp}\)。利用

\[
R_{AB}=r\,\boldsymbol b^{\mathsf T}L^+\boldsymbol b,\qquad
\boldsymbol b=\boldsymbol e_A-\boldsymbol e_B,
\]

并将共轭模式合并取实部，非均匀模式的贡献为

\[
\frac r m\sum_{q=1}^{m-1}
\left[G_q(0,0)+G_q(n,n)-2\cos(p\theta_q)G_q(0,n)\right].
\]

由上式格林函数

\[
G_q(0,0)=\frac{1+\rho_q}{D_q},\qquad
G_q(n,n)=\frac{1+\rho_q^{2n+1}}{D_q},\qquad
G_q(0,n)=\frac{(1+\rho_q)\rho_q^n}{D_q},
\]

\par\noindent\begin{minipage}{\linewidth}
\begin{center}
\includegraphics[width=\linewidth,height=.40\textheight,keepaspectratio]{\detokenize{assets/module-04/process/M4-094b.pdf}}
\par\smallskip{\small 图 M4-094b：零圆周模式保留真实的两端电势差}
\end{center}
\end{minipage}\par

图 M4-094b 把同一轴向层的 \(m\) 个等势节点合为一点，得到从第0层到第
\(n\) 层的单链。每相邻两层间的 \(m\) 只 \(r\) 并联成
\(r/m\)，所以这条轴向链给出 \(nr/m\)；圆周边的电流在零模式下为零。

其中

\[
D_q=\sqrt{\mu_q(\mu_q+4)}.
\]

\(q=0\) 是圆周各点等势的均匀模式，等价于每一轴向步中 \(m\)
个电阻并联，因此贡献为

\[
\frac{nr}{m}.
\]

最终

\[
\boxed{
\frac{R_{p,n}}r=
\frac nm+
\frac1m\sum_{q=1}^{m-1}
\frac{
2+\rho_q+\rho_q^{2n+1}
-2(1+\rho_q)\rho_q^n\cos(p\theta_q)
}{\sqrt{\mu_q(\mu_q+4)}}
}.
\]

\subsubsection{4. 结果检验}\label{m4:ux7ed3ux679cux68c0ux9a8c}

当 \(p=0,n=0\) 时，每个求和项的分子均为零，所以 \(R_{0,0}=0\)。

由于

\[
\mu_{m-q}=\mu_q,\qquad
\cos[(m-p)\theta_q]=\cos(p\theta_q),
\]

所以

\[
R_{m-p,n}=R_{p,n},
\]

符合圆柱的反射对称性。

当 \(n\to\infty\) 时，所有 \(\rho_q^n\) 项消失，故

\[
R_{p,n}=\frac{nr}{m}+O(r).
\]

远处主项只由均匀模式决定，物理上就是轴向电流平均分配在 \(m\) 条母线中。

\subsection{方法归纳}\label{m4:ux65b9ux6cd5ux5f52ux7eb3}

\begin{enumerate}
\def\labelenumi{\arabic{enumi}.}
\tightlist
\item
  有限梯形优先寻找``剥去一格''后的端口递推；端口若在对角点，再改用支路电流差分。
\item
  无限条带先对横向有限截面求本征模态，再分别处理轴向传播；均匀模式产生线性项，非均匀模式指数衰减。
\item
  环形网络的本质是周期边界；半无限网络还必须单独处理端面缺邻点的边界方程。
\item
  无穷网络只能讨论由有限源汇对定义的电势差，不能随意把无穷远设为零，也不能把单点输运态的电压降冒充两点等效电阻。
\end{enumerate}

\backmatter
\par\endgroup
\clearpage
\endgroup
